% concatenated from root_Arxive.tex, math_commands.tex
\documentclass{article}
\usepackage[preprint]{neurips_2026}
%\usep}

%====================================================
% PACKAGES (from original header:)
%====================================================
\usepackage{microtype}
\usepackage{graphicx}
\usepackage{booktabs}
\usepackage{hyperref}

\usepackage{amsmath}
\usepackage{amssymb}
\usepackage{mathtools}
\usepackage{amsthm}
\usepackage[capitalize,noabbrev]{cleveref}

\usepackage{tikz}
  \usepackage{pgfplots}
  \usepgfplotslibrary{groupplots}
  \pgfplotsset{compat=1.17}
\usetikzlibrary{arrows.meta}
\usetikzlibrary{backgrounds}
\usepgfplotslibrary{patchplots}
\usepgfplotslibrary{fillbetween}
\pgfplotsset{%
    layers/standard/.define layer set={%
        background,axis background,axis grid,axis ticks,axis lines,axis tick labels,pre main,main,axis descriptions,axis foreground%
    }{
        grid style={/pgfplots/on layer=axis grid},%
        tick style={/pgfplots/on layer=axis ticks},%
        axis line style={/pgfplots/on layer=axis lines},%
        label style={/pgfplots/on layer=axis descriptions},%
        legend style={/pgfplots/on layer=axis descriptions},%
        title style={/pgfplots/on layer=axis descriptions},%
        colorbar style={/pgfplots/on layer=axis descriptions},%
        ticklabel style={/pgfplots/on layer=axis tick labels},%
        axis background@ style={/pgfplots/on layer=axis background},%
        3d box foreground style={/pgfplots/on layer=axis foreground},%
    },
}
\usepackage{times}

%====================================================
% CUSTOM MATH COMMANDS
%====================================================
%%%%% NEW MATH DEFINITIONS %%%%%

% \usepackage{amsmath,amsfonts,bm}

% Mark sections of captions for referring to divisions of figures
\newcommand{\figleft}{{\em (Left)}}
\newcommand{\figcenter}{{\em (Center)}}
\newcommand{\figright}{{\em (Right)}}
\newcommand{\figtop}{{\em (Top)}}
\newcommand{\figbottom}{{\em (Bottom)}}
\newcommand{\captiona}{{\em (a)}}
\newcommand{\captionb}{{\em (b)}}
\newcommand{\captionc}{{\em (c)}}
\newcommand{\captiond}{{\em (d)}}

% Highlight a newly defined term
\newcommand{\newterm}[1]{{\bf #1}}


% % Figure reference, lower-case.
% \def\figref#1{figure~\ref{#1}}
% % Figure reference, capital. For start of sentence
% \def\Figref#1{Figure~\ref{#1}}
% \def\twofigref#1#2{figures \ref{#1} and \ref{#2}}
% \def\quadfigref#1#2#3#4{figures \ref{#1}, \ref{#2}, \ref{#3} and \ref{#4}}
% % Section reference, lower-case.
% \def\secref#1{section~\ref{#1}}
% % Section reference, capital.
% \def\Secref#1{Section~\ref{#1}}
% % Reference to two sections.
% \def\twosecrefs#1#2{sections \ref{#1} and \ref{#2}}
% % Reference to three sections.
% \def\secrefs#1#2#3{sections \ref{#1}, \ref{#2} and \ref{#3}}
% % Reference to an equation, lower-case.
% \def\eqref#1{equation~\ref{#1}}
% % Reference to an equation, upper case
% \def\Eqref#1{Equation~\ref{#1}}
% % A raw reference to an equation---avoid using if possible
% \def\plaineqref#1{\ref{#1}}
% % Reference to a chapter, lower-case.
% \def\chapref#1{chapter~\ref{#1}}
% % Reference to an equation, upper case.
% \def\Chapref#1{Chapter~\ref{#1}}
% % Reference to a range of chapters
% \def\rangechapref#1#2{chapters\ref{#1}--\ref{#2}}
% % Reference to an algorithm, lower-case.
% \def\algref#1{algorithm~\ref{#1}}
% % Reference to an algorithm, upper case.
% \def\Algref#1{Algorithm~\ref{#1}}
% \def\twoalgref#1#2{algorithms \ref{#1} and \ref{#2}}
% \def\Twoalgref#1#2{Algorithms \ref{#1} and \ref{#2}}
% % Reference to a part, lower case
% \def\partref#1{part~\ref{#1}}
% % Reference to a part, upper case
% \def\Partref#1{Part~\ref{#1}}
% \def\twopartref#1#2{parts \ref{#1} and \ref{#2}}

\def\ceil#1{\lceil #1 \rceil}
\def\floor#1{\lfloor #1 \rfloor}
\def\1{\bm{1}}
\newcommand{\train}{\mathcal{D}}
\newcommand{\valid}{\mathcal{D_{\mathrm{valid}}}}
\newcommand{\test}{\mathcal{D_{\mathrm{test}}}}

\def\eps{{\epsilon}}


% Random variables
\def\reta{{\textnormal{$\eta$}}}
\def\ra{{\textnormal{a}}}
\def\rb{{\textnormal{b}}}
\def\rc{{\textnormal{c}}}
\def\rd{{\textnormal{d}}}
\def\re{{\textnormal{e}}}
\def\rf{{\textnormal{f}}}
\def\rg{{\textnormal{g}}}
\def\rh{{\textnormal{h}}}
\def\ri{{\textnormal{i}}}
\def\rj{{\textnormal{j}}}
\def\rk{{\textnormal{k}}}
\def\rl{{\textnormal{l}}}
% rm is already a command, just don't name any random variables m
\def\rn{{\textnormal{n}}}
\def\ro{{\textnormal{o}}}
\def\rp{{\textnormal{p}}}
\def\rq{{\textnormal{q}}}
\def\rr{{\textnormal{r}}}
\def\rs{{\textnormal{s}}}
\def\rt{{\textnormal{t}}}
\def\ru{{\textnormal{u}}}
\def\rv{{\textnormal{v}}}
\def\rw{{\textnormal{w}}}
\def\rx{{\textnormal{x}}}
\def\ry{{\textnormal{y}}}
\def\rz{{\textnormal{z}}}

% Random vectors
\def\rvepsilon{{\mathbf{\epsilon}}}
\def\rvtheta{{\mathbf{\theta}}}
\def\rva{{\mathbf{a}}}
\def\rvb{{\mathbf{b}}}
\def\rvc{{\mathbf{c}}}
\def\rvd{{\mathbf{d}}}
\def\rve{{\mathbf{e}}}
\def\rvf{{\mathbf{f}}}
\def\rvg{{\mathbf{g}}}
\def\rvh{{\mathbf{h}}}
\def\rvu{{\mathbf{i}}}
\def\rvj{{\mathbf{j}}}
\def\rvk{{\mathbf{k}}}
\def\rvl{{\mathbf{l}}}
\def\rvm{{\mathbf{m}}}
\def\rvn{{\mathbf{n}}}
\def\rvo{{\mathbf{o}}}
\def\rvp{{\mathbf{p}}}
\def\rvq{{\mathbf{q}}}
\def\rvr{{\mathbf{r}}}
\def\rvs{{\mathbf{s}}}
\def\rvt{{\mathbf{t}}}
\def\rvu{{\mathbf{u}}}
\def\rvv{{\mathbf{v}}}
\def\rvw{{\mathbf{w}}}
\def\rvx{{\mathbf{x}}}
\def\rvy{{\mathbf{y}}}
\def\rvz{{\mathbf{z}}}

% Elements of random vectors
\def\erva{{\textnormal{a}}}
\def\ervb{{\textnormal{b}}}
\def\ervc{{\textnormal{c}}}
\def\ervd{{\textnormal{d}}}
\def\erve{{\textnormal{e}}}
\def\ervf{{\textnormal{f}}}
\def\ervg{{\textnormal{g}}}
\def\ervh{{\textnormal{h}}}
\def\ervi{{\textnormal{i}}}
\def\ervj{{\textnormal{j}}}
\def\ervk{{\textnormal{k}}}
\def\ervl{{\textnormal{l}}}
\def\ervm{{\textnormal{m}}}
\def\ervn{{\textnormal{n}}}
\def\ervo{{\textnormal{o}}}
\def\ervp{{\textnormal{p}}}
\def\ervq{{\textnormal{q}}}
\def\ervr{{\textnormal{r}}}
\def\ervs{{\textnormal{s}}}
\def\ervt{{\textnormal{t}}}
\def\ervu{{\textnormal{u}}}
\def\ervv{{\textnormal{v}}}
\def\ervw{{\textnormal{w}}}
\def\ervx{{\textnormal{x}}}
\def\ervy{{\textnormal{y}}}
\def\ervz{{\textnormal{z}}}

% Random matrices
\def\rmA{{\mathbf{A}}}
\def\rmB{{\mathbf{B}}}
\def\rmC{{\mathbf{C}}}
\def\rmD{{\mathbf{D}}}
\def\rmE{{\mathbf{E}}}
\def\rmF{{\mathbf{F}}}
\def\rmG{{\mathbf{G}}}
\def\rmH{{\mathbf{H}}}
\def\rmI{{\mathbf{I}}}
\def\rmJ{{\mathbf{J}}}
\def\rmK{{\mathbf{K}}}
\def\rmL{{\mathbf{L}}}
\def\rmM{{\mathbf{M}}}
\def\rmN{{\mathbf{N}}}
\def\rmO{{\mathbf{O}}}
\def\rmP{{\mathbf{P}}}
\def\rmQ{{\mathbf{Q}}}
\def\rmR{{\mathbf{R}}}
\def\rmS{{\mathbf{S}}}
\def\rmT{{\mathbf{T}}}
\def\rmU{{\mathbf{U}}}
\def\rmV{{\mathbf{V}}}
\def\rmW{{\mathbf{W}}}
\def\rmX{{\mathbf{X}}}
\def\rmY{{\mathbf{Y}}}
\def\rmZ{{\mathbf{Z}}}

% Elements of random matrices
\def\ermA{{\textnormal{A}}}
\def\ermB{{\textnormal{B}}}
\def\ermC{{\textnormal{C}}}
\def\ermD{{\textnormal{D}}}
\def\ermE{{\textnormal{E}}}
\def\ermF{{\textnormal{F}}}
\def\ermG{{\textnormal{G}}}
\def\ermH{{\textnormal{H}}}
\def\ermI{{\textnormal{I}}}
\def\ermJ{{\textnormal{J}}}
\def\ermK{{\textnormal{K}}}
\def\ermL{{\textnormal{L}}}
\def\ermM{{\textnormal{M}}}
\def\ermN{{\textnormal{N}}}
\def\ermO{{\textnormal{O}}}
\def\ermP{{\textnormal{P}}}
\def\ermQ{{\textnormal{Q}}}
\def\ermR{{\textnormal{R}}}
\def\ermS{{\textnormal{S}}}
\def\ermT{{\textnormal{T}}}
\def\ermU{{\textnormal{U}}}
\def\ermV{{\textnormal{V}}}
\def\ermW{{\textnormal{W}}}
\def\ermX{{\textnormal{X}}}
\def\ermY{{\textnormal{Y}}}
\def\ermZ{{\textnormal{Z}}}

% Vectors
\def\vzero{{\bm{0}}}
\def\vone{{\bm{1}}}
\def\vmu{{\bm{\mu}}}
\def\vtheta{{\bm{\theta}}}
\def\va{{\bm{a}}}
\def\vb{{\bm{b}}}
\def\vc{{\bm{c}}}
\def\vd{{\bm{d}}}
\def\ve{{\bm{e}}}
\def\vf{{\bm{f}}}
\def\vg{{\bm{g}}}
\def\vh{{\bm{h}}}
\def\vi{{\bm{i}}}
\def\vj{{\bm{j}}}
\def\vk{{\bm{k}}}
\def\vl{{\bm{l}}}
\def\vm{{\bm{m}}}
\def\vn{{\bm{n}}}
\def\vo{{\bm{o}}}
\def\vp{{\bm{p}}}
\def\vq{{\bm{q}}}
\def\vr{{\bm{r}}}
\def\vs{{\bm{s}}}
\def\vt{{\bm{t}}}
\def\vu{{\bm{u}}}
\def\vv{{\bm{v}}}
\def\vw{{\bm{w}}}
\def\vx{{\bm{x}}}
\def\vy{{\bm{y}}}
\def\vz{{\bm{z}}}

% Elements of vectors
\def\evalpha{{\alpha}}
\def\evbeta{{\beta}}
\def\evepsilon{{\epsilon}}
\def\evlambda{{\lambda}}
\def\evomega{{\omega}}
\def\evmu{{\mu}}
\def\evpsi{{\psi}}
\def\evsigma{{\sigma}}
\def\evtheta{{\theta}}
\def\eva{{a}}
\def\evb{{b}}
\def\evc{{c}}
\def\evd{{d}}
\def\eve{{e}}
\def\evf{{f}}
\def\evg{{g}}
\def\evh{{h}}
\def\evi{{i}}
\def\evj{{j}}
\def\evk{{k}}
\def\evl{{l}}
\def\evm{{m}}
\def\evn{{n}}
\def\evo{{o}}
\def\evp{{p}}
\def\evq{{q}}
\def\evr{{r}}
\def\evs{{s}}
\def\evt{{t}}
\def\evu{{u}}
\def\evv{{v}}
\def\evw{{w}}
\def\evx{{x}}
\def\evy{{y}}
\def\evz{{z}}

% Matrix
\def\mA{{\bm{A}}}
\def\mB{{\bm{B}}}
\def\mC{{\bm{C}}}
\def\mD{{\bm{D}}}
\def\mE{{\bm{E}}}
\def\mF{{\bm{F}}}
\def\mG{{\bm{G}}}
\def\mH{{\bm{H}}}
\def\mI{{\bm{I}}}
\def\mJ{{\bm{J}}}
\def\mK{{\bm{K}}}
\def\mL{{\bm{L}}}
\def\mM{{\bm{M}}}
\def\mN{{\bm{N}}}
\def\mO{{\bm{O}}}
\def\mP{{\bm{P}}}
\def\mQ{{\bm{Q}}}
\def\mR{{\bm{R}}}
\def\mS{{\bm{S}}}
\def\mT{{\bm{T}}}
\def\mU{{\bm{U}}}
\def\mV{{\bm{V}}}
\def\mW{{\bm{W}}}
\def\mX{{\bm{X}}}
\def\mY{{\bm{Y}}}
\def\mZ{{\bm{Z}}}
\def\mBeta{{\bm{\beta}}}
\def\mPhi{{\bm{\Phi}}}
\def\mLambda{{\bm{\Lambda}}}
\def\mSigma{{\bm{\Sigma}}}

% Tensor
% \DeclareMathAlphabet{\mathsfit}{\encodingdefault}{\sfdefault}{m}{sl}
% \SetMathAlphabet{\mathsfit}{bold}{\encodingdefault}{\sfdefault}{bx}{n}
% \newcommand{\tens}[1]{\bm{\mathsfit{#1}}}
\def\tA{{\tens{A}}}
\def\tB{{\tens{B}}}
\def\tC{{\tens{C}}}
\def\tD{{\tens{D}}}
\def\tE{{\tens{E}}}
\def\tF{{\tens{F}}}
\def\tG{{\tens{G}}}
\def\tH{{\tens{H}}}
\def\tI{{\tens{I}}}
\def\tJ{{\tens{J}}}
\def\tK{{\tens{K}}}
\def\tL{{\tens{L}}}
\def\tM{{\tens{M}}}
\def\tN{{\tens{N}}}
\def\tO{{\tens{O}}}
\def\tP{{\tens{P}}}
\def\tQ{{\tens{Q}}}
\def\tR{{\tens{R}}}
\def\tS{{\tens{S}}}
\def\tT{{\tens{T}}}
\def\tU{{\tens{U}}}
\def\tV{{\tens{V}}}
\def\tW{{\tens{W}}}
\def\tX{{\tens{X}}}
\def\tY{{\tens{Y}}}
\def\tZ{{\tens{Z}}}


% Graph
\def\gA{{\mathcal{A}}}
\def\gB{{\mathcal{B}}}
\def\gC{{\mathcal{C}}}
\def\gD{{\mathcal{D}}}
\def\gE{{\mathcal{E}}}
\def\gF{{\mathcal{F}}}
\def\gG{{\mathcal{G}}}
\def\gH{{\mathcal{H}}}
\def\gI{{\mathcal{I}}}
\def\gJ{{\mathcal{J}}}
\def\gK{{\mathcal{K}}}
\def\gL{{\mathcal{L}}}
\def\gM{{\mathcal{M}}}
\def\gN{{\mathcal{N}}}
\def\gO{{\mathcal{O}}}
\def\gP{{\mathcal{P}}}
\def\gQ{{\mathcal{Q}}}
\def\gR{{\mathcal{R}}}
\def\gS{{\mathcal{S}}}
\def\gT{{\mathcal{T}}}
\def\gU{{\mathcal{U}}}
\def\gV{{\mathcal{V}}}
\def\gW{{\mathcal{W}}}
\def\gX{{\mathcal{X}}}
\def\gY{{\mathcal{Y}}}
\def\gZ{{\mathcal{Z}}}

% Sets
\def\sA{{\mathbb{A}}}
\def\sB{{\mathbb{B}}}
\def\sC{{\mathbb{C}}}
\def\sD{{\mathbb{D}}}
% Don't use a set called E, because this would be the same as our symbol
% for expectation.
\def\sF{{\mathbb{F}}}
\def\sG{{\mathbb{G}}}
\def\sH{{\mathbb{H}}}
\def\sI{{\mathbb{I}}}
\def\sJ{{\mathbb{J}}}
\def\sK{{\mathbb{K}}}
\def\sL{{\mathbb{L}}}
\def\sM{{\mathbb{M}}}
\def\sN{{\mathbb{N}}}
\def\sO{{\mathbb{O}}}
\def\sP{{\mathbb{P}}}
\def\sQ{{\mathbb{Q}}}
\def\sR{{\mathbb{R}}}
\def\sS{{\mathbb{S}}}
\def\sT{{\mathbb{T}}}
\def\sU{{\mathbb{U}}}
\def\sV{{\mathbb{V}}}
\def\sW{{\mathbb{W}}}
\def\sX{{\mathbb{X}}}
\def\sY{{\mathbb{Y}}}
\def\sZ{{\mathbb{Z}}}

% Entries of a matrix
\def\emLambda{{\Lambda}}
\def\emA{{A}}
\def\emB{{B}}
\def\emC{{C}}
\def\emD{{D}}
\def\emE{{E}}
\def\emF{{F}}
\def\emG{{G}}
\def\emH{{H}}
\def\emI{{I}}
\def\emJ{{J}}
\def\emK{{K}}
\def\emL{{L}}
\def\emM{{M}}
\def\emN{{N}}
\def\emO{{O}}
\def\emP{{P}}
\def\emQ{{Q}}
\def\emR{{R}}
\def\emS{{S}}
\def\emT{{T}}
\def\emU{{U}}
\def\emV{{V}}
\def\emW{{W}}
\def\emX{{X}}
\def\emY{{Y}}
\def\emZ{{Z}}
\def\emSigma{{\Sigma}}

% entries of a tensor
% Same font as tensor, without \bm wrapper
\newcommand{\etens}[1]{\mathsfit{#1}}
\def\etLambda{{\etens{\Lambda}}}
\def\etA{{\etens{A}}}
\def\etB{{\etens{B}}}
\def\etC{{\etens{C}}}
\def\etD{{\etens{D}}}
\def\etE{{\etens{E}}}
\def\etF{{\etens{F}}}
\def\etG{{\etens{G}}}
\def\etH{{\etens{H}}}
\def\etI{{\etens{I}}}
\def\etJ{{\etens{J}}}
\def\etK{{\etens{K}}}
\def\etL{{\etens{L}}}
\def\etM{{\etens{M}}}
\def\etN{{\etens{N}}}
\def\etO{{\etens{O}}}
\def\etP{{\etens{P}}}
\def\etQ{{\etens{Q}}}
\def\etR{{\etens{R}}}
\def\etS{{\etens{S}}}
\def\etT{{\etens{T}}}
\def\etU{{\etens{U}}}
\def\etV{{\etens{V}}}
\def\etW{{\etens{W}}}
\def\etX{{\etens{X}}}
\def\etY{{\etens{Y}}}
\def\etZ{{\etens{Z}}}

% The true underlying data generating distribution
\newcommand{\pdata}{p_{\rm{data}}}
% The empirical distribution defined by the training set
\newcommand{\ptrain}{\hat{p}_{\rm{data}}}
\newcommand{\Ptrain}{\hat{P}_{\rm{data}}}
% The model distribution
\newcommand{\pmodel}{p_{\rm{model}}}
\newcommand{\Pmodel}{P_{\rm{model}}}
\newcommand{\ptildemodel}{\tilde{p}_{\rm{model}}}
% Stochastic autoencoder distributions
\newcommand{\pencode}{p_{\rm{encoder}}}
\newcommand{\pdecode}{p_{\rm{decoder}}}
\newcommand{\precons}{p_{\rm{reconstruct}}}

\newcommand{\laplace}{\mathrm{Laplace}} % Laplace distribution

\newcommand{\E}{\mathbb{E}}
\newcommand{\Ls}{\mathcal{L}}
\newcommand{\R}{\mathbb{R}}
\newcommand{\emp}{\tilde{p}}
\newcommand{\lr}{\alpha}
\newcommand{\reg}{\lambda}
\newcommand{\rect}{\mathrm{rectifier}}
\newcommand{\softmax}{\mathrm{softmax}}
\newcommand{\sigmoid}{\sigma}
\newcommand{\softplus}{\zeta}
\newcommand{\KL}{\mathrm{KL}}
\newcommand{\Var}{\mathrm{Var}}
\newcommand{\standarderror}{\mathrm{SE}}
\newcommand{\Cov}{\mathrm{Cov}}
% Wolfram Mathworld says $L^2$ is for function spaces and $\ell^2$ is for vectors
% But then they seem to use $L^2$ for vectors throughout the site, and so does
% wikipedia.
\newcommand{\normlzero}{L^0}
\newcommand{\normlone}{L^1}
\newcommand{\normltwo}{L^2}
\newcommand{\normlp}{L^p}
\newcommand{\normmax}{L^\infty}

\newcommand{\parents}{Pa} % See usage in notation.tex. Chosen to match Daphne's book.

\DeclareMathOperator*{\argmax}{arg\,max}
\DeclareMathOperator*{\argmin}{arg\,min}

\DeclareMathOperator{\sign}{sign}
\DeclareMathOperator{\Tr}{Tr}
\let\ab\allowbreak

\renewcommand{\KL}{\mathrm{KL}}
%====================================================
% TITLE
%====================================================
\title{PAC-Bayesian Bounds for Learning Partially Observed Stochastic Linear Time-Invariant State-Space Systems with Inputs and Sub-Gaussian Noise}

%====================================================
% AUTHORS (original authors)
%====================================================
\author{
Mihaly Petreczky \\
CRIStAL, Centrale Lille, Université de Lille, France, \texttt{mihaly.petreczky@centralelille.fr} \\
\And
Mohamad Al Ahdab \\
Department of Electronic Systems, Aalborg University, Denmark, \texttt{maah@es.aau.dk} \\
\And
John Leth \\
Department of Electronic Systems, Aalborg University, Denmark, \texttt{jjl@es.aau.dk} \\
}

%====================================================
% THEOREM ENVIRONMENTS
%====================================================
%\theoremstyle{plain}
\newtheorem{theorem}{Theorem}[section]
\newtheorem{proposition}[theorem]{Proposition}
\newtheorem{lemma}[theorem]{Lemma}
\newtheorem{corollary}[theorem]{Corollary}

%\theoremstyle{definition}
\newtheorem{definition}[theorem]{Definition}
\newtheorem{assumption}[theorem]{Assumption}

%\theoremstyle{remark}
\newtheorem{remark}[theorem]{Remark}
\newtheorem{example}[theorem]{Example}

%====================================================
% USER-DEFINED COMMANDS (from original header)
%====================================================
\newcommand{\e}{\mathbf{e}}
\newcommand{\br}{\mathbf{r}}
\newcommand{\x}{\mathbf{x}}
\newcommand{\y}{\mathbf{y}}
\newcommand{\z}{\mathbf{z}}
\newcommand{\w}{\mathbf{w}}
\newcommand{\bv}{\mathbf{v}}
\newcommand{\reals}{\mathbb{R}}

\newcommand{\hyf}{\hat{\y}_{f}}

\newcommand{\bP}{\mathbf{P}}
\newcommand{\bE}{\mathbf{E}}
\newcommand{\F}{\mathcal{F}}
\newcommand{\B}{\mathcal{B}_\Theta}
\newcommand{\bnu}{\bm{\nu}}
\newcommand{\etab}{\bm{\eta}}

\newcommand{\Gp}{G_{\theta}}
\newcommand{\GpExpr}{\left (1+\|\hat{D}_\theta\|+
\frac{\hat{M}_\theta\|\hat{B}_\theta\|\|\hat{C}_\theta\|}_2{1-\hat{\gamma}_\theta} \right ) \left (\frac{\hat{M}_\theta\|\hat{C}_\theta\| \|\hat{B}_\theta\|}{(1-\hat{\gamma}_\theta)^{\frac{3}{2}}} \right )}



\usepackage{pgfplots}
\pgfplotsset{compat=1.17}
\begin{document}
\maketitle 

\begin{abstract}
%In this paper we derive a Probably Approximately Correct (PAC)-Bayesian error bound for linear time-invariant (LTI) stochastic dynamical systems with inputs.  Such bounds
%are widespread in machine learning, and they are useful for characterizing the predictive power of models learned from  
%finitely many data points.  
%The bound derived in this paper relates the expectation of prediction errors with the prediction error generated by the model on the data used for learning. In addition, we show that it can also be used to derive bounds for the parameter estimation error.
%In turn, this allows us to provide finite-sample error bounds  for the prediction error and parameret estimation error for
%a wide class of system identification algorithms. 
%Furthermore, as LTI systems are a sub-class of recurrent neural
%networks (RNNs), these error bounds could be a first step towards 
%PAC-Bayesian bounds for RNNs. 
In this paper we derive a Probably Approximately Correct (PAC)-Bayesian error bound for partially observed linear time-invariant (LTI) stochastic dynamical systems in state-space form with inputs.  Such bounds
are widespread in machine learning, and they are useful for characterizing the predictive power of models learned from  
finitely many data points.  
The bound derived in this paper relates the expectation of prediction errors with the prediction error generated by the model on the data used for learning. In addition, we show that it can also be used to derive bounds for the parameter estimation error.
In turn, this allows us to provide finite-sample error bounds  for the prediction error and parameter estimation error for
a wide class of system identification algorithms. 
Furthermore, as LTI systems are a sub-class of recurrent neural
networks (RNNs), these error bounds could be a first step towards 
PAC-Bayesian bounds for RNNs. 
\end{abstract}

\section{Introduction}
% Linear time invariant (LTI) state-space models have been widely used in control and econometric applications to model time-series and 
% have rich literature on learning (classically called identification)\cite{LjungBook}. 

% In this paper, we present PAC-Bayesian type bounds on learning LTI systems from data generated by LTI system driven by zero-mean, i.i.d., Gaussian or sub-Gaussian noise. 

% The Probably Approximately Correct (PAC)-Bayesian framework, provides theoretical guarantees (with arbitrary high probability) on the difference between learning from infinite amount of data, and learning from finite empirical data, see \cite{guedj2019primer,alquier2021userfriendly,grunwald-2012,alquier-15,nips-16,ShethK17}. 
%\textbf{Motivation}
PAC and PAC-Bayesian bounds are fundamental tools for analyzing learning algorithms. They provide generalization guarantees that relate the \emph{true error} (performance on unseen data) to the \emph{empirical error} (performance on training data), by bounding the \emph{generalization gap}, the difference between true and empirical error, in a way that is independent of the specific learning algorithm. In this framework, learning is viewed as selecting a model from a data-dependent posterior distribution over hypotheses, which is a modification of a prior based on the observed data. Different learning algorithms correspond to different choices of priors, posteriors, and selection mechanisms (e.g., sampling, averaging, or maximizing the posterior). PAC-Bayesian bounds then provide a uniform upper bound on the average generalization gap with respect to any posterior, thus offering generalization guarantees for a broad class of learning and model selection methods.

This generality makes PAC-Bayesian approaches especially useful for analyzing complex models, including neural networks, where they can yield non-vacuous generalization bounds \cite{Dziugaite2017}.
Although PAC and PAC-Bayesian theory is well developed for i.i.d.\ settings \cite{shalev2014understanding,alquier2021userfriendly,guedj2019primer}, there are comparatively few results for dynamical systems.
\par
    %Traditionally, the literature on LTI systems \cite{LjungBook} has focused on statistical consistency. More recently, several results have appeared on finite-sample bounds for learning LTI systems, but they are valid only for specific learning algorithms or for very limited sub-classes \cite{simchowitz2021statistical,Ozay2018,oymak2021revisiting}.
   \textbf{Contribution:}
    In this paper, we derive PAC-Bayesian error bounds for learning discrete-time stochastic LTI state-space systems with sub-Gaussian noise and partial observations (the full states is not observed). 
    Following standard practice in system identification, we view LTI systems as predictors that generate predictions for the current output based on past inputs and outputs.
    % Learning amounts to finding the predictor that minimizes the \emph{empirical error} on training data. However, the true quality of a learned model is determined by its performance on unseen data, quantified by the \emph{true error}. We thus focus on bounding the \emph{generalization gap}—the difference between true error and empirical error.
    Our main contribution is a PAC-Bayesian bound showing that, with high probability (at least $1-\delta$), for any posterior distribution over predictors, the average generalization gap is $O(\frac{(\ln N)^r}{\sqrt{N}})$, where $N$ is the sample size, $r=0$ for bounded noise, $r=1$ for Lipschitz loss, and $r=3/2$ for quadratic loss. The bound depends on the Kullback--Leibler (KL) divergence between prior and posterior distributions and on $\ln(1/\delta)$. We subsequently derive generalization gap and parameter estimation error bounds for concrete learning algorithms (e.g., maximum likelihood, regularized prediction-error methods (PEM), and Bayesian methods).

    \textbf{Related work:}
    PAC and PAC-Bayesian bounds for autoregressive models with bounded signals or Lipschitz loss have been studied in \cite{campi2002finite,vidyasagar2006learning,alquier2012pred,alquier2013prediction,MASSUCCI2022110532,METAKALARD2026100373}. 
    In contrast, our results address state-space models with partially observed states and  with inputs, which are more general than autoregressive models, and we  allow for unbounded (sub-Gaussian) noise, and handle quadratic loss. We also build on techniques from \cite{alquier2012pred,alquier2013prediction}. 
    %Other related works, such as \cite{}, provide PAC bounds for autoregressive systems, but do not cover the stochastic LTI state-space setting considered here.
%That is, the learning problem considered in this paper isdifferent from that of %\cite{alquier2012pred,alquier2013prediction}.
\par
In contrast to the literature on finite-sample parameter estimation bounds for specific least-squares based learning
algorithms for LTI systems, e.g., \cite{simchowitz2019learning,simchowitz2021statistical,oymak2021revisiting,lale2020logarithmic,foster2020learning,NEURIPS2018_d6288499,Ziemann1,Ziemann2,SarkarRD21}, this papers presents bounds on both the parameter estimation and the generalization gap for a wide-range of algorithms, see Remark \ref{rem:finite-sample}, Appendix \ref{app:indent} for a detailed comparison.
% \cite{simchowitz2018learning,simchowitz2019learning}
%All the cited papers propose a bound which is valid only for models generated by a specific learning algorithm, and  In particular, these bounds do not relate the generalization loss with the empirical error for arbitrary algorithms, i.e., they are not PAC(-Bayesian) bounds. This means that in contrast to the results of this paper, the bounds of the cited papers cannot be used for analyzing algorithms other than for which they were derived.
%Second, many of the cited papers
%do not derive bounds on the infinite horizon prediction error.
% More precisely, \cite{oymak2021revisiting,SarkarRD21,lale2020logarithmic,Pappas1,Simchowitz_Foster_2020} provided error bounds for the difference of the first $T$ Markov-parameters of the estimated and true system for a specific
% identification algorithm. However, in order to characterize the infinite horizon prediction error,
% we need to take $T=\infty$.
% For $T=\infty$ the cited bounds become infinite, i.e., vacuous. 
% In addition, in contrast to the present paper,  \cite{oymak2021revisiting,SarkarRD21,simchowitz2018learning} deals only with the deterministic part of the stochastic LTI, and \cite{Pappas1} deals only with the stochastic part.
\par
PAC bounds for recurrent neural networks, of which LTI state-space representations are a subclass, have been developed using VC dimension \cite{KOIRAN199863,sontag1998learning}, Rademacher complexity \cite{KOIRAN199863,sontag1998learning,pmlr-v108-chen20d}, and PAC-Bayesian approaches \cite{Dziugaite2017}. However, existing results assume noiseless models, fixed prediction horizons, i.i.d.\ training data samples, and bounded signals. In contrast, this paper considers noisy models, infinite-horizon prediction errors, a single observed time series, and unbounded signals with sub-Gaussian noise.
% Note that the error bounds of
% the cited papers
%\cite{simchowitz2019learning} ,lale2020logarithmic,NEURIPS2018_d6288499,Pappas1,SarkarRD21}
% converge to their limit at rate $O(\frac{\ln(N)}{\sqrt{N}})$, which is comparable to the rate $O(\frac{1}{\sqrt{N}})$ of this paper.
% This being said, the cited papers 
% provide bounds on the parameter estimation error, and many of them allow marginally stable systems.
%\textbf{PAC-Bayesian bounds for state-space representation.}
% In \cite{haussmann2021learning} learning of stochastic differential equations without inputs was considered and it was assumed that  several independently sampled time-series were available for learning. 
% In contrast, in this paper we deal with discrete-time systems with inputs and the learning takes place from a single time-series.
% In \cite{PACMarkov} learning of general Markov-chains was considered, but the state of the Markov-chain was assumed to be observable and no inputs were considered. The learning problem of \cite{PACMarkov} is thus different from the one of this paper.
\par
%Several recent works have addressed PAC-Bayesian bounds for LTI systems. 
In \cite{eringisRenyi}, bounds based on R\'enyi divergence were proposed, but the resulting probabilistic guarantees scale as $O(1/\delta)$ with the confidence level, rather than the standard $O(\ln(1/\delta))$, and it applies only to quadratic losses as opposed to general Lipschitz losses, and it provides parameter estimation bounds only in terms of the empirical error. 
In contrast, the bound of this paper depends on the confidence level $O(\ln(1/\delta))$ and it provides an absolute bound 
on the parameter estimation error. 
The works \cite{eringis2023pacbayesRNN,eringis2023pacbayesian} provide PAC-Bayesian generalization gap bounds for RNNs and LTI systems with bounded noise, respectively. In contrast, the present paper addresses unbounded (sub-Gaussian) noise and also derives parameter estimation error bounds. Even for bounded signals, our results refine the constants used in \cite{eringis2023pacbayesRNN,eringis2023pacbayesian} (see Remark~\ref{rem:deividas}, Appendix~\ref{app:indent}). 
PAC-Bayesian bounds for autonomous LTI systems without inputs were considered in \cite{CDC21paper} but those bounds did not converge to zero as $N \rightarrow \infty$.
% here, we extend to systems with exogenous inputs and provide tighter bounds. (see also Remark~\ref{rem:deividas}, Appendix~\ref{app:ident}).

%For the case of bounded noise, the result of this paper is a special case of \cite{eringis2023pacbayesRNN}.

\textbf{Paper Outline:}
  We start by defining the problem formulation in Section \ref{sect:Problem_formulation},
  followed by the main results in Section \ref{sec:main_results} and discussion of the bounds in Section \ref{sec:conclusions}.
  The technical proofs are presented in the appendices.
  
 % Then we discuss the PAC-Bayesian approach in Section \ref{sect:pac}, where the fundamental theorem is stated, and the use cases of PAC-Bayesian bounds are explored. The necessary assumptions are stated in Section \ref{sec:asm}, with definitions of important quantities relating to systems. We present the results for the most general case in Section \ref{sec:unbounded}, followed by narrowing the scope to bounded processes in Section \ref{sec:boundedResults}, with tighter results.  Finally, a short numerical example is presented in  Section \ref{sec:numEx}. 
%Finally, we will have the conclusion in Section %\ref{sect:concl}.

\section{Problem formulation}\label{sect:Problem_formulation}
\subsection{ Notation and terminology} 
We occasionally use $\triangleq$ to denote ''defined by''. 
Let $\F$ denote a $\sigma$-algebra on the set $\Omega$ and $\bP$ be a probability measure on $\F$. Unless otherwise stated, all probabilistic considerations will be with respect to the probability space $(\Omega,\F,\bP)$, and we let $\bE[\z]$ denote expectation of the stochastic variable $\z$.
We use bold face letters to indicate stochastic variables/processes relative to $(\Omega,\F,\bP)$. 
Each euclidean space is associated with the topology generated by the 2-norm $\|\cdot\|_2$, and the Borel $\sigma$-algebra generated by the open sets. The induced matrix 2-norm is also denoted $\|\cdot\|_2$. 
\par
Let $B_{\Theta}$ be the $\sigma$-algebra of Lebesgue-measurable subsets of $\Theta \subset \mathbb{R}^{n_\theta}$, and $m$ the Lebesgue measure. For a probability density $\rho$ on $(\Theta, B_{\Theta}, m)$ and a measurable, integrable function $g:\Theta \to \mathbb{R}$, define
\(
E_{\theta \sim \rho}[g(\theta)] \triangleq \int_{\Theta} \rho(\theta) g(\theta) dm(\theta)
\)
as the expectation of $g$ under $\rho$. For a density $\pi$ on $\Theta$, let $\mathcal{M}_{\pi}$ denote the set of all densities on $\Theta$ %absolutely continuous with respect to $\pi$.
whose probability measure  is absolutely continuous w.r.t. the probability measure of $\pi$. 
\par
Let $\mathrm{vol}(A) = \int_{A} dm(\theta)$ denote the Lebesgue measure of $A \subseteq \mathbb{R}^{n_\theta}$. The uniform distribution on a measurable set $A \in B_\Theta$ is $\rho(\theta) = \frac{\chi_A(\theta)}{\mathrm{vol}(A)}$, where $\chi_A$ is the indicator function. For $A = \Theta$, this is the uniform distribution on $\Theta$.
\par
The phrase \emph{probability over data} refers to probability under $\bP$ on $(\Omega, \F)$. Specifically, for measurable functions $\{X_i, Y_i\}_{i \in I}$ and a binary relation $\diamond$, we say $\forall i \in I: X_i \diamond Y_i$ holds with probability at least $1-\delta$ over data if the set $F = \{\omega \in \Omega : \forall i \in I,\, X_i(\omega) \diamond Y_i(\omega)\}$ satisfies $\bP(F) \ge 1-\delta$.

\subsection{\textcolor{black}{Problem formulation}}
The problem formulation follows \cite{eringis2023pacbayesRNN,eringis2023pacbayesian,eringisRenyi}.  
\par
Let us fix stationary stochastic processes $\y(t)\in \reals^{n_\y}$, and $\rvu(t)\in  \reals^{n_\rvu}$, that share a time axis $t\in\sZ$, that is, for any $t\in\sZ$, $\rvy(t):\Omega\to\sR^{n_\rvy};\omega\mapsto\rvy(t)(\omega)$, and $\rvu(t):\Omega\to\sR^{n_\rvu};\omega\mapsto\rvu(t)(\omega)$
%; hence  $y(t):\omega\to\y(t)(\omega)$, $w(t):\omega\to\w(t)(\omega)$ 
are random vectors on $(\Omega,\F,\bP)$:
$\y$ represents outputs (labels), and $\rvu$ represents inputs (features).
The goal is to learn models from the finite sample $\train=\{\rvy(t)(\omega),\rvu(t)(\omega)\}_{t=0}^{N-1}$, $\omega \in \Omega$, such that the learned model
predicts  $\rvy$  from past measurements of $(\rvy,\rvu)$ sufficiently well.

%\textcolor{red}{\textbf{Informal description}} %Predictor (hypothesis)  class and data generator:} } 
%We focus on models which are parameterised linear time-invariant (LTI) systems, 
%and which use the past and present of $\rvu$, or past of $\rvy$, to estimate the present $\rvy$. 
\textbf{Predictor class:}
To simplify notation, let us define a stochastic process $\w(t)\in \sR^{n_\w}$ by %two cases
\begin{description}
    \item $\w(t)=\begin{bmatrix} \y^T(t), & \rvu^T(t) \end{bmatrix}^T$ when $n_\w=n_\y+n_\rvu$ and both past $\y$ and $\rvu$ are used for prediction,
    \item $\w(t)=\rvu(t)$ when  $n_\w=n_\rvu$, when only past $\rvu$ is used for prediction.
\end{description}
Note that with a trivial extension of the results of this paper one can define $\w(t)$ to consist of only some of the components of $\y(t)$, extending the results to not only prediction but also filtering.

\begin{assumption}[Predictors]\label{as:parameterisation}
The class of predictors %$\mathcal{F}$ is the family 
is the class of LTI systems parametrized by $\theta \in \Theta \subseteq \mathbb{R}^{n_{\theta}} $ which generate 
predictions $\hat{\y}_\theta(t|t_0)$ of $\y(t)$ based on $\{\w(s)\}_{s=t_0}^{t}$
%tuples $\Sigma(\theta)=(\hat{A}_\theta,\hat{B}_\theta,\hat{C}_\theta,\hat{D}_\theta)$, 
%for estimating $\rvy(t)$  
\begin{equation}\label{eq:predictor}
    \begin{aligned}
        \hat{\x}(t+1)=\hat{A}_\theta\hat{\x}(t)+\hat{B}_\theta\w(t), ~ \hat{\x}(t_0)=0, \quad \hat{\y}_\theta(t|t_0)=\hat{C}_\theta\hat{\x}(t)+\hat{D}_\theta\w(t), \quad t \ge t_0,
        %\hat{\x}(0)&=E[\bar{x}(0)]=0
    \end{aligned}
\end{equation}
  where $\hat{A}_\theta$,$\hat{B}_\theta,\hat{C}_\theta,\hat{D}_\theta$ 
  %and of dimension
	are $\hat{n}\times\hat{n}$, $\hat{n}\times n_\w$,
	$n_\y\times \hat{n}$ and $n_\y\times n_\w$ matrices
	respectively, they are are continuous in $\theta$,
  and 
  $\hat{A}_\theta$ is Schur (all its eigenvalues are inside the unit disk).
  Moreover, if $\w(t)=[\y^T(t)~\rvu^T(t)]^T$, then  $\hat{D}_\theta=[0~\hat{D}_{\rvu,\theta}]$ for some $n_\y \times n_\rvu$ matrix $\hat{D}_{\rvu,\theta}$. 
  %i.e., $\hat{D}_\theta\w(t)$  depends only on $\rvu(t)$.
  \footnote{The latter assumption is necessary, as it guarantees that $\hat{D}_\theta\w(t)$ depends only on $\rvu(t)$ and hence it does not depend on $\rvy(t)$, since otherwise we would be using $\rvy(t)$ to predict $\rvy(t)$, which is not meaningful.}
    %The class of predictors %$\mathcal{F}$ 
    %is then   
	%\(
     % \mathcal{F}=\{\Sigma(\theta) \mid %\gamma(\hat{A}_\theta)<1,\;
    %\theta\in\Theta\}. 
    %\).
	%with $\gamma(\hat{A}_\theta)$ the spectral radius of $\hat{A}_\theta$.
    %i.e., the largest modulus of eigenvalues of $\hat{A}_\theta$.
    Moreover, there exist constants
     $M,C  >0$, $\gamma \in (0,1)$ such that 
   for all $\theta \in \Theta$, $\|\hat{A}_\theta^k\|_2 \le M \gamma^k$, $\max\{\|\hat{B}_\theta\|_2,\|\hat{C}_\theta\|_2,\|\hat{D}_\theta\|_2\} \le C$.
\end{assumption}
%\begin{remark}[Stable parametrisations]
%\label{rem:parameter}
Intuitively, our assumption expresses that the predictors are stable LTIs and we have a uniform upper bound on the matrices and on the spectral radius of $\hat{A}_\theta$. These assumptions are standard in system identification, see Remark \ref{rem:parameter:methods}, Appendix \ref{app:sect:Problem_formulation}.
%Using parametrizations of stable predictors is a standard practice \cite{LjungBook,CainesBook}, and  various methods exist for constructing them \cite[Section 3.2.3]{RalfPeeters}, \cite[Section 4.6]{Ribarits1}, for more details 
%\end{remark}
%, and parameters $\theta\in\Theta\subset\sR^{n_\theta}$.  Moreover, 
Note that the predictors are started at the zero initial state at inital time $t_0$. The latter represents the time instant at which the predictor was applied to generate predictions, or when learning started. Most of the time, we will set $t_0=0$ to simplify notation.
%Note that $\hat{L}_\theta$ can be set to zero, if measurements of $\rvy(t)$ are not available during implementation.

\textbf{Data generator:}
Intuitively, we would like %to assume that 
the data  $\y$ to be generated by an LTI system %in innovation form
%stochastic linear system with input in innovation form,
\begin{align}
    \begin{split}
		&\hat{\rvs}_{g}(t+1)=A_0\hat{\rvs}_{g}(t)+B_0\rvu(t)+K_0\e^s(t), \quad \y(t) = C_0\hat{\rvs}_{g}(t)+D_0\rvu(t)+\e^s(t), \quad t \in \sZ,
    \end{split}     
		\label{eq:T0sys1}
\end{align}
where $A_0,B_0,K_0,C_0,D_0$ are matrices of suitable dimension, $A_0$ and $A_0-K_0C_0$ are both Schur, and $\e^s(t)$ is a zero-mean i.i.d. sub-Gaussian process. The Schur conditions ensure stationarity of $\rvx$ and that $\e^s(t)$ is the innovation process (see \cite[Chapter 4, p.~87]{LjungBook}), i.e., the difference between $\y(t)$ and its best linear predictor based on past inputs and outputs. Thus, \eqref{eq:T0sys1} is in steady-state Kalman filter form.
Unlike predictors, the data generator \eqref{eq:T0sys1} is assumed to operate in steady state, i.e., it has run for infinite time prior to prediction. This standard assumption in system identification (e.g., \cite{LjungBook}) reflects that the physical system has been subject to noise and inputs long before learning begins. This subtle difference has significant theoretical implications.

Assuming the data is generated by a system of the same form as the predictors is a realizability assumption, as \eqref{eq:T0sys1} yields the optimal predictor.
% \eqref{eq:T0sys1}
%can be identified with the predictor 
%$f_{\mathrm{true}}$
\begin{equation}
%\label{eq:T0sys1:true}
    \begin{split}
		&\hat{\rvs}_{g}(t+1 \mid t_0)\!\!=\!\!\hat{A}_0\hat{\rvs}_{g}(t)+\hat{B}_0\w(t), ~ \hat{\rvs}_g(t_0 | t_0)=0, \quad  
    %f_{\mathrm{true}}(\{\w(s)\}_{s=0}^{t})
    \hat{\y}(t \mid t_0) \!\!=\!\! \hat{C}_0\hat{\rvs}_{g}(t \mid t_0)+\hat{D}_0\w(t), 
    \end{split}     
		% 		\text{where}&\\
		% 		&A_0=\begin{bmatrix}A_1&A_2-K_2C_3+K_1D_0C_3\\0&A_3-K_3C_3  \end{bmatrix}\\
		% 		&K_0=\begin{bmatrix}K_2-K_1D_0\\K_3\end{bmatrix}\\
		% 		&C_0=\begin{bmatrix} C_1 & C_2+D_0C_3 \end{bmatrix}\\
		% 		&D_0=(\bE[\y(t)\e_2^T(t)])^T(\bE[\e_2(t)\e_2^T(t)])^{-1}
		\label{eq:T0sys2}
\end{equation}
where $t \ge t_0$,
 $\hat{A}_0=A_0,\hat{B}_0=B_0,\hat{C}_0=C_0$, $\hat{D}_0=D_0$
if $n_w=n_u$ and 
$\hat{A}_0=A_0-K_0C_0$, $\hat{B}_0=\begin{bmatrix} K_0 & B_0-K_0D_0 \end{bmatrix}$,
$\hat{C}_0=C_0$, $\hat{D}_0=\begin{bmatrix} 0 & D_0 \end{bmatrix}$ if $n_w=n_u+n_y$. 
When $\w(s)=[\rvy^T(s),\; \rvu^T(s) ]^T$, 
the predictor \eqref{eq:T0sys2} corresponds 
to stationary Kalman-filter and it is an optimal predictor of $\y(t)$ in the mean square error sense.
%that the generalisation loss 
%$\mathcal{L}(f_{\mathrm{true}})$  is the smallest
%possible among all the predictors $f$ which arises
%via a LTI state-space representation \eqref{eq:predictor}.
%More precisely, we assume the following.
%\textcolor{red}{\textbf{Formal description}}
  Formally, instead of assuming the existence of \eqref{eq:T0sys1}, for the data generator  we use the following assumption, which is essentially equivalent to the existence of \eqref{eq:T0sys1}.
  % see Remark \ref{rem:equiv:gen} in Appendix \ref{app:sect:Problem_formulation}.}
%we will use the following assumption.
\begin{assumption}\label{as:generator}
	Let $\y(t)$ and $\mathbf{u}(t)$ be generated by an autonomous stochastic LTI system
	%\begin{subequations}\label{eq:generator}
    \begin{align}\label{eq:generator}
			\rvx(t+1) &=A_g\rvx(t)+K_g\rve_g(t), \quad
			\begin{bmatrix}\rvy(t)\\\rvu(t)\end{bmatrix}=C_g\rvx(t)+\rve_g(t)
		\end{align}
	%\end{subequations}
	with $A_g \in \mathbb{R}^{n \times n},K_g \in \mathbb{R}^{n \times n_\rve},C_g \in \mathbb{R}^{n_\rve \times n}$, for some $n>0$, and $n_\rve=n_\rvy+n_\rvu\geq2$, $A_g$ and $A_g-K_gC_g$ are Schur. Moreover, 
	%\begin{itemize}
	%\item
	%$A_g$ and $A_g-K_gC_g$ are Schur,
		$\rvx(t)$ has finite variance, $\rve_g(t)$ is independent of $\rvx(t)$ for all $t \in \mathbb{Z}$, 
	%\item
		and $\rve_g$ is square integrable, i.i.d process,
  %  \textcolor{red}{which satisfies Cramer's condition, i.e.,  
      %with the moment generating function $\Psi_{\e_g}(\lambda)=\bE[e^{\lambda \|\e_g(t)\|_2}] \le e^{4\lambda^2 \sigma^2_{\rve_g}+3\lambda\sigma_{\rve}_g}$ and 
    $\|\rve_g(t)\|_2$ is sub-Gaussian and the sub-Gaussian norm of $\|\rve_g(t)\|_2$
    \cite[Definition 2.6.4]{vershynin2026highdimensional} is denoted by $\|\rve_g\|_{\psi_2}$.
    %$\sigma_{\rve_g}$.} % such that $\bP(\|\rve_g(t)\|_2 > z) \le 2\exp(-\frac{z^2}{\sigma^2_{\rve_g}})$.   
 %is wel-defined for all $\lambda > 0$}.
        %and $\rve_g(t)$ is a zero mean Gaussian variable for all $t \in \mathbb{Z}$
%		with noise covariance $\bE[\rve_g(t)\rve_g^T(t)]=Q_e$, and $\bE[\|(\rve_g(t))\|_2^2]<\infty$, 
     %    \item
	%\end{itemize}
	% We identify the system \eqref{eq:generator} with the tuple
	% $\Sigma_{gen}\triangleq(A_g,K_g,C_g,I)$
\end{assumption}
%Note that 
If $\e_g(t)$ is sub-Gaussian, then $\|\e_g(t)\|_2$ is also sub-Gaussian, see
\cite[Lemma 1]{jin2019short}.
%\textcolor{red}{Note that if $\e_g(t)$ is sub-Gaussian, then $\|\e_g(t)\|$ is also sub-Gaussian, see
%\cite[Lemma 1]{jin2019short}.}
Under Assumption \ref{as:generator}, 
$\rvy$ and $\rvu$ are jointly stationary and square integrable processes
\cite{CainesBook}.
%\cite[Theorem 1.4, p. 80]{CainesBook}. 
%Moreover, 
%$\bE[\|(\rvy(t))\|_2^2]$ and $\bE[\|(\rvu(t))\|_2^2]<\infty$,
%$\rvy$ and $\rvu$ are essentially bounded if $\e_g$
%is essentially bounded, see
%
%Moreover,
%from classical theory of LTI systems
%it follows that $\rvy(t)$ and $\rvu(t)$, $t \in \mathbb{Z}$
%are essentially bounded if the noise 
%$\e_g(s)$ is essentially bounded for all $s \in \mathbb{Z}$
{\color{black}
%\end{color}
%\begin{remark} 
%  It follows
%   that $\e_g(t)$ is the innovation process of $[\rvy^T(t)\;\rvu^T(t)]^T$
%   and \eqref{eq:generator} is in the so called 
%   (forward) innovation form, see
%   \cite[Section 8.5.2]{Katayama:05}.
%   Intuitively, it means that $\e_g$
%   is the difference between $[\rvy^T(t)\;\rvu^T(t)]^T$ and its best linear prediction based on its own past values.
%  \end{remark} 
  %\begin{remark}[Equivalence between \eqref{eq:T0sys1} and \eqref{eq:generator}]
  %\label{rem:equiv:gen}
   %The assumption that $\y$ and $\rvu$ are generated by an autonomous linear system is in fact a mild assumption, and is fairly standard in system identification \cite[Section 9.2,Chapter 9]{Katayama:05}. In fact, under some mild conditions (no feedback from $\y$ to $\rvu$) 
   Under mild 
   conditions, Assumption \ref{as:generator} is equivalent to the existence of  an LTI system
   \eqref{eq:T0sys1}, see Remark \ref{rem:equiv_gen:proof} in Appendix \ref{app:sect:Problem_formulation}. 
   Assumption \ref{as:generator}
   is more convenient, and it is often
   used in system identification \cite[Chapter 9]{Katayama:05}.
  %\end{remark}   
}
% \todo{.}
%\textbf{Class of predictors (hypotheses):}  
%In this paper, we will be interested in the following model class, consisting of predictors realizable by LTI systems.


\textbf{Measure of prediction accuracy:} 
Following standard practice in system identification \cite{LjungBook}, we measure the predictive accuracy
using a \emph{loss function}
$\ell: \mathbb{R}^{n_y} \times \mathbb{R}^{n_y} \rightarrow [0,\infty)$
which is continuous and satisfies $\ell(y,y)=0$.
%which is assumed to be $L$-Lipschitz for some $L > 0$ in the sense that
Typical choices are 
  $\ell(y,\hat{y})=\|y-\hat{y}\|_{2}$, to which we refer to as $\ell_2$ loss,
  and 
    \( \ell(y,\hat{y})=\|y-\hat{y}\|_2^2 \) to which we refer to as the \emph{quadratic loss}. 
  %Quadratic loss functions are the typical choice in system identification. 
%\begin{itemize}
%\item
%\item 
    %(quadratic) loss. 
%\end{itemize}
%Since $\ell_2$ losses are smooth and convex functions, they are more common in the classical system identification literature. However, recently $\ell_1$
%losses have also gained popularity \cite{pillonetto2022regularized}.
%In the sequel we consider loss functions which are either quadratic or 
%Lipschitz.
\begin{assumption}[Qudaratic or Lipschitz loss functions]
\label{as:loss}
 The loss function $\ell$ is either quadratic or it is $L_{\ell}$-Lipschitz, i.e., for all $y_1,y_2,\hat{y}_1,\hat{y}_2 \in \mathbb{R}^{n_y}$,
\( |\ell(y_1,\hat{y}_1)-\ell(y_2,\hat{y}_2)| \le L_{\ell} (\|y_1-y_2\|_2+\|\hat{y}_1-\hat{y}_2\|_2). \)
\end{assumption}
%Assumption \ref{as:Lipschitz} will be  used  together with the following 
%boundedness assumption on the noise of the data generator. 
We define the \emph{true error} as the limit
\begin{align}
\mathcal{L}(\theta) &\triangleq\lim_{t_0 \rightarrow -\infty} \bE[\ell(\y(t),\hat{\rvy}_\theta(t\mid t_0))]. \label{eq:GenLoss}
\end{align}  
%We shall refer to $\mathcal{L}(\theta)$ as the generalisation loss. 
The true error captures the long-term, steady-state prediction error obtained during the deployment of the predictor on an increasing number of past inputs.
\begin{lemma}\label{l:ihp} Let Assumptions \ref{as:generator}, \ref{as:parameterisation}, and \ref{as:loss} hold.
  %\textcolor{red}{Assume that either Assumptions \ref{as:loss} holds true, 
  %or $\ell$ is the quadratic loss.  }
  %If $(\rvy,\rvu)$ is stationary 
Then the limits 
	\( \hat{\rvy}_\theta(t)=\lim_{t_0 \rightarrow -\infty} \hat{\rvy}_\theta(t\mid t_0) \)
  and $\lim_{t \rightarrow \infty} (\hat{\rvy}_{\theta}(t)-\hat{\rvy}_{\theta}(t \mid 0))=0$
	exist in the mean-square sense for all $t$ and $\theta$, the process $\hat{\rvy}_\theta(t)$ is stationary, and 
	\(	\mathcal{L}(\theta)=\bE[\ell(\y(t),\hat{\rvy}_\theta(t) )]=\lim_{t \rightarrow \infty} \bE[\ell(\y(t),\hat{\rvy}_\theta(t\mid 0))]\). 
  %Moreover,  
  %  in the mean-square sense, and 
  %  $\mathcal{L}(\theta) = \lim_{t \rightarrow \infty} \bE[\ell(\y(t),\hat{\rvy}_\theta(t\mid 0))]$
    %in the mean-square sense.}
\end{lemma}
The proof of Lemma \ref{l:ihp} is presented in Appendix \ref{app:sect:Problem_formulation}. 
Intuitively, $\hat{\rvy}_{\theta}(t)$
is the output predicted by the LTI system corresponding to
$\theta$, if the latter is started from zero at $-\infty$, and as $t$ grows, $\hat{\rvy}_{\theta}(t)$
and $\hat{\rvy}_{\theta}(t \mid 0)$ converge to each other.
%Note that, due to stationarity we can also state that 

\textbf{The Learning Problem} is to 
find a parameter $\theta_{\star}$  from sampled data $\train$ of the sequence of random variables $\{\rvy(t),\rvu(t)\}_{t=0}^{N-1}$ which minimizes the true error, i.e., $\theta_{\star}\in\arg\min \mathcal{L}(\theta)$.
%\textbf{Note:} For learning, we assume to have the training data set $\train$, i.e., a single trajectory of $(\rvy,\rvu)$, but no 
Note that we do not assume the knowledge of the matrices $A_g,K_g,C_g$ and noise process $\e_g$ when learning. %system \eqref{eq:generator} only defines the assumptions on the data generating process.

Since the true error is unknown, a standard practice is to use the \emph{empirical error} 
\begin{align}
    \hat{\mathcal{L}}_N(\theta)=\frac{1}{N}\sum_{t=0}^{N-1}\ell(\rvy(t)(\omega),\hat{\rvy}_\theta(t|0)(\omega))
\end{align}
as a proxy for the true error, and try to find a 
model $\hat{\theta}$ for which the empirical error 
$\hat{\mathcal{L}}_N(\hat{\theta})$ is small. 
%However, the classical approach above minimizes the empirical error, whereas we are interested in the predictive power of the learned model, i.e. its empirical error. 
%That is, we would like to bound the 
%\emph{generalization gap}
%$\mathcal{L}(\theta)-\hat{\Ls}_N(\theta)$.
%µTo this end, we use the so called 
%\emph{PAC-Bayesian approach:} 
%for any distribution on the models we bound the average generalization gap w.r.t. this distribution. 
%\textcolor{blue}{{\bf Alternative to the above paragraph:} 
%The classical approach above minimizes the empirical error, whereas we are interested in minimizing the 
While doing so, we would like to keep the \emph{generalization gap} $\mathcal{L}(\theta)-\hat{\Ls}_N(\theta)$, i.e., the difference between a model's performance on the training data (empirical error) and its expected performance (true error) on unseen test data, small.  The PAC-Bayesian approach provides a framework to derive such guarantees.

\subsection{Constants used in PAC-bounds }\label{sec:asm}
%By Theorem \ref{thm:initial}, in order to obtain the upper-bound $\tilde{\mathbb{B}}_N$, we only need to bound
%$\Psi(\lambda,\pi,N)$.
%To this end, we will
In order to present the main result,  we will need several constants, many of which have an interesting system-theoretic interpretation.
% We use the correspondence between the parameter values from $\Theta$ and the elements of $\cF$, and we will use densities defined on $\Theta$ instead of densities defined on $\cF$. Since $\Theta$ is a subset of an Euclidean space, densities on $\Theta$ can be defined in a classical way.  We will need to 
% % More precisely, since every  hypothesis $f \in \cF$ corresponds to some parameter
% % value $\theta \in \Theta$ such that
% % $f=f_{\theta}$, instead of probability 
% % densities on $\cF$
% % we can consider probability densities on $\Theta$.
% take expectations of
% a  function $g$ on $\cF$ w.r.t. to the probability distribution induced by a density $\rho$ on $\Theta$. To this end,  we can identify $g$ with the function $\theta \mapsto g(f_{\theta})$ on $\Theta$, and take the expectation of the latter function  w.r.t. the probability distribution induced by $\rho$.

% We will identify the system \eqref{eq:formal_predictor}
% with the tuple 
% $(\hat{A},\hat{B},\hat{C},\hat{D})$.
% For the sake of notation, throughout the paper we will use $f$, to denote $f_{\Sigma(\theta)}$, for some arbitrary $\theta\in\Theta$. 
%For every parameter $\theta\in\Theta$ we define a number of constants which will be used in the PAC-Bayesian error bound. %% $G(\theta)$ and $G_e(\theta)$. To this end, let  $f=(\hat{A},\hat{B},\hat{C},\hat{D})$ be a predictor.
	 % The constants $G(\theta)$, $G_e(\theta)$ are define as follows.
   
% The lemma abobve motivates the definition of the following constants
\begin{definition}[Convolution and $\ell_1$ norm]
  Let $\alpha=\{\alpha(k)\}_{k=0}^{\infty}$ be a sequence of matrices. 
  Define the $\ell_1$ norm $\|\alpha\|_{\ell_1}$ and 
  the \emph{weakly dependent mixing coefficient (WDMC)} $\theta_{\infty}(\alpha)$ of $\alpha$ as
  \[ \|\alpha\|_{\ell_1}=\sum_{k=0}^{\infty} \|\alpha(k)\|_2, \quad  \theta_{\infty}(\alpha)=\sum_{k=0}^{\infty} k \|\alpha(k)\|_2. \]
  If $\|\alpha\|_{\ell_1} < \infty$ and $\mathbf{r}$ is a stationary square integrable stochastic process, then define the \emph{convolution $(\alpha \star \mathbf{r})$ of $\alpha$ with $\mathbf{r}$}  as the stochastic process
  \( (\alpha \star \mathbf{r})(t)=\sum_{k=0}^{\infty} \alpha(k) \mathbf{r}(t-k) \). 
  By \cite[Theorem 1.1]{CainesBook}, $(\alpha \star \mathbf{r})$ is well-defined and stationary. 
\end{definition}
  Intuitively, the $\ell_1$ norm of a sequence is the extension of the standard norm on the space of
  absolutely summable scalar sequences to the case of matrix valued ones. The convolution 
  $(\alpha \star \mathbf{r})$ represents the response of a linear filter with impulse response $\alpha$
  to the noise $\mathbf{r}$, and $\|\alpha\|_{\ell_1}$ represents the $\ell_1$ norm of the filter.
  %For the  convolution $\alpha \star \mathbf{r}$, 
  %$\bE[\|(\alpha \star \mathbf{r})(t)\|_2] < \|\alpha\|_{\ell_1}E[\|\mathbf{r}(t)\|_2]$, and
  If  $\mathbf{r}$ is bounded, i.e., $\|\mathbf{r}(t)\|_2 < C$, then 
  the convolution  $(\alpha \star \mathbf{r})(t)$ is bounded, i.e., $\|(\alpha \star \mathbf{r})\|_2 < \|\alpha\|_{\ell_1} C$, and it is a  weakly dependent process, in the sense of \cite{alquier2012pred},
  whose mixing coefficients $\{\theta_{\infty,n}(1)\}_{n=1}^{\infty}$ 
  are bounded by $2\theta_{\infty}(\alpha)C$, \cite[Proposition 4.12]{alquier2012pred}.
   Intuitively, the smaller $\theta_{\infty}(\alpha)$ is,
  the closer $(\alpha \star \mathbf{r})(t)$ is to being i.i.d., if $\mathbf{r}$ is i.i.d. 
  If $\alpha$ is the sequence of Markov parameters of a stable LTI system, then both $\|\alpha\|_{\ell_1}$ and $\theta_{\infty}(\alpha)$ are finite and can be upper bounded in terms of bounds on the norms of the system matrices and the spectral radius of the system matrix $A$, see Remark \ref{norms:rem}.
 \begin{definition}[Auxiliary definitions]
 %$\Gp,G_e(\theta)$]
	\label{def:constants}
  Let us consider the Markov-parameters $\alpha_g$ and 
    $\alpha_{\theta}$ of the data generator \eqref{eq:generator} and predictor 
    \eqref{eq:predictor} corresponding to
    $\theta \in \Theta$ respectively:
    %$\alpha_{k,g}, \alpha_k(\theta)$ 
    %\alpha_{k,e}(\theta), \widetilde{\alpha}_{k,e}(\theta)n \widetilde{\alpha}_{k,e,t_0}(\theta)$ as follows
    \begin{equation}
    \label{alpha:lemma:alpha_def}
    \begin{aligned}
    & \alpha_{g}(k)=\left\{\begin{array}{rl}
            C_g A_g^{k-1}K_g & k > 0 \\
            I & k = 0
    \end{array}\right., \quad   \alpha_{\theta}(k)=\left\{\begin{array}{rl}
            \widehat{C}_{\theta}\widehat{A}_{\theta}^{k-1}\widehat{B}_{\theta} & k > 0 \\
           \widehat{D}_{\theta}   & k = 0
    \end{array}\right. \\
  %& \widetilde{\alpha}_{k,e}(\theta)=
  %  \sum_{l=0}^{k} \alpha_{k}(\theta)\alpha_{k-l,w}, \\
  %  & \alpha_{k,e}(\theta)=\alpha_{k,y} - \widetilde{\alpha}_{k,e}, \\
  %  &  \alpha_{k,y}=\begin{bmatrix} I_{n_{\rvy}} & 0 \end{bmatrix} \alpha{k,g}    \\
  %& \widetilde{\alpha}_{k,e,s}(\theta)=\sum_{l=0}^{\min\{k,s\}} \alpha_{k}(\theta)\alpha_{k-l,w}, \\
  %& \alpha_{k,e,s}(\theta)=\begin{bmatrix} I_{n_{\rvy}} & 0 \end{bmatrix}\alpha_{k,g} - \widetilde{\alpha}_{k,e,t_0}
\end{aligned}
\end{equation}
 Let $\bar{\alpha}(k)=\sup_{\theta \in \Theta} \|\alpha_{\theta}(k)\|_2$.
 With notation as above%\begin{color}{red}
  \begin{align*}
    %G_g=\|\alpha_g\|_{\ell_1}, ~ G_{g,1}=\theta_{\infty}(\alpha), ~
   %& G_f(\theta)=\|\alpha_{\theta}\|_1, ~ G_{f,1}(\theta)=\theta_{\infty}(\alpha_{\theta}) \\
       % & G_g \triangleq \sum_{k=0}^{\infty} \|\alpha_{k,g}\|_2, \quad 
        %G_{g,1} \triangleq \sum_{k=0}^{\infty} (k+1)\|\alpha_{k,g}\|_2 \\%1+\sum_{k=0}^\infty \|C_gA_g^{k-1}K_g\|_2\\
       %& G_f(\theta)\triangleq 
       % \sum_{k=0}^{\infty} \|\alpha_k(\theta)\|_2, ~
       % G_{f,1}(\theta)\triangleq \sum_{k=0}^{\infty} (k+1) \|\alpha_k(\theta)\|_2 \\
       & G_e(\theta) \triangleq (1+\|\alpha_{\theta}\|_{\ell_1})\|\alpha_g\|_{\ell_1}, \quad   G_{e,1}(\theta) \triangleq \theta_{\infty}(\alpha_{\theta})\|\alpha\|_{\ell_1}+\theta_{\infty}(\alpha_g)(\|\alpha_{\theta}\|_{\ell_1}+1) \\
       %& G_f \triangleq  \|\bar{\alpha}\|_{\ell_1}, ~ G_{f,1} \triangleq \theta_{\infty}(\bar{\alpha}) \\
        %\sum_{k=0}^{\infty} \sup_{\theta \in \Theta} \|\alpha_k(\theta)\|_2,  \\
        %& G_{f,1} \triangleq \sum_{k=0}^{\infty} (k+1) \sup_{\theta \in \Theta} 
        % \|\alpha_k(\theta)\|_2  \\
        & G_e \triangleq (1+\|\bar{\alpha}\|_{\ell_1})\|\alpha_g\|_{\ell_1}, \quad 
        G_{e,1} \triangleq \theta_{\infty}(\bar{\alpha})\|\alpha_g\|_{\ell_1}+\theta_{\infty}(\alpha_g)(\|\bar{\alpha}\|_{\ell_1}+1)
        %\bar{G}_{f,1}G_g+G_{g,1}\bar{G}_f \\
   \end{align*}
  %\end{color}
%% \begin{color}{red}
%%  \begin{align*}
%%        & G_g \triangleq \sum_{k=0}^{\infty} \|\alpha_{k,g}\|_2, \quad 
%%        G_{g,1} \triangleq \sum_{k=0}^{\infty} (k+1)\|\alpha_{k,g}\|_2 \\%1+\sum_{k=0}^\infty \|C_gA_g^{k-1}K_g\|_2\\
%%       & G_f(\theta)\triangleq 
%%        \sum_{k=0}^{\infty} \|\alpha_k(\theta)\|_2, ~
%%        G_{f,1}(\theta)\triangleq \sum_{k=0}^{\infty} (k+1) \|\alpha_k(\theta)\|_2 \\
%%       & G_e(\theta) \triangleq (1+\bar{G}_f(\theta))G_{g}, \\
%%       &   G_{e,1}(\theta) \triangleq \bar{G}_{f,1}(\theta)G_g+G_{g,1}\bar{G}_f(\theta) \\
%%       & G_f \triangleq 
%%        \sum_{k=0}^{\infty} \sup_{\theta \in \Theta} \|\alpha_k(\theta)\|_2,  \\
%%        & G_{f,1} \triangleq \sum_{k=0}^{\infty} (k+1) \sup_{\theta \in \Theta} 
%%         \|\alpha_k(\theta)\|_2  \\
%%        & G_e \triangleq (1+\bar{G}_f)G_{g}, \quad  G_{e,1} \triangleq \bar{G}_{f,1}G_g+G_{g,1}\bar{G}_f \\
%%   \end{align*}
  %\end{color}
\end{definition}
%\begin{color}{red}
For the definition of the upcoming bounds, we are going to use the quantities
$\|\alpha\|_{\ell_1}$, $\theta_{\infty}(\alpha)$ for 
$\alpha \in  \{\alpha_g, \bar{\alpha}\} \cup \{\alpha_\theta \mid \theta \in \Theta\}$ and
the constants $G_e(\theta), G_{e,1}(\theta), G_e, G_{e,1}$ defined above. 
Note that all these quantities are well-defined, see Lemma \ref{alpha:lemma1.5}. In addition, they are related to 
 the true outputs $\rvy(t)$ and predicted outputs $\hat{\rvy}_\theta(t)$, $\hat{\rvy}_\theta(t \mid t_0)$ as follows:
 $\begin{bmatrix}\rvy^T(t) &  \rvu^T(t)\end{bmatrix}^T=(\alpha_{g} \star \e_g)(t)$,
 $\hat{\rvy}_{\theta}(t)=(\alpha_{\theta} \star \rvw)(t)$, 
 $\hat{\rvy}_{\theta}(t \mid t_0)=(\alpha_{\theta, t-t_0} \star \rvw)(t)$.
 %their respective differences arise as the following convolutions.
 %$\alpha_g$ and $\alpha_{\theta}$
%respectively as follows. 
%If $a(k),b(k)$, $k \in \mathbb{N}$ are two sequences of matrices, then denote by
%$a \star b$ their Cauchy product, i.e., 
%\[ (a \star b)(k)=\sum_{l=0}^{k} a(l)b(k-l).  \]
%It is well-known \cite{} that if the series  
%$\sum_{k=0}^{\infty} a(k)$ and 
%$\sum_{k=0}^{\infty} b(k)$ are absolutely convergent,
%i.e., $\sum_{k=0}^{\infty} \|a(k)\|_2$ and 
%$\sum_{k=0}^{\infty} \|b(k)\|_2$ are convergent, then 
%the series $\sum_{k=0}^{\infty} (a \star b)(k)$ is also convergent, in fact it is absolutely convergent, i.e.,
%$\sum_{k=0}^{\infty} \|(a \star b)(k)\|_2$ is convergent, and 
%\[
%   \sum_{k=0}^{\infty} (a \star b)(k)=\sum_{k=0}^{\infty} (a \star b)(k)
% \] 
%It then follows that the following holds.
%\begin{lemma}
%\label{alpha:lemma2}
%\[
%\begin{aligned*}
%&  \alpha_{w}(k)= \left\{ \begin{array}{rl} \alpha_{g}(k) & \rvw = \begin{bmatrix} \rvy \rvu \end{bmatrix} \\
%            \begin{bmatrix} 0 & I_{n_{\rvu}} \end{bmatrix} \alpha_{g}(k) & \rvw=\rvu 
 %   \end{array}\right. \\
 %   & \alpha_{y}(k)=\begin{bmatrix} I_{n_{\rvy}} & 0 \end{bmatrix} \alpha_{g}(k)  \\
 % \end{aligned*}
%\]\left(\sum_{l=0}^{\min\{k,t-t_0\}} \alpha_{\theta}(k)\alpha_w(k-l)\right)
%Let $\alpha_{g,y}(k)$ and $\alpha_{g,w}(k)$ be the matrix formed by the first $n_\rvy$ and the last $n_\rvw$ 
%rows of $\alpha_g(k)$ respectively, where $n_{\rvw}$ is the number of entries of $\w(t)$. Let $c_{\theta,s}(k)=\left(\sum_{l=0}^{\min\{k,s\}} \alpha_{\theta}(k)\alpha_{g,w}(k-l)\right)$
%and $c_{\theta}(k)=\sum_{l=0}^{k} \alpha_{\theta}(k)\alpha_{g,w}(k-l)$. 
%Then 
%\begin{equation} 
% \label{alpha:lemma:eq0}
%   
%\end{equation}
%The proof of \eqref{alpha:lemma:eq0} and 
The  formal relationship between the constants is provided in 
Lemma \ref{alpha:lemma2}, Appendix \ref{app:constants}. %The intuitive  interpretation of the various constants is as follows. 
The interpretation of the various terms appearing in  Definition \ref{def:constants} is as follows
\begin{itemize}
    \item 
%\textbf{The term $G_f(\theta)$} is the $\ell_1$ norm 
%of the error system, it relates high-order
%moments of the infinite past prediction error
%and the high-order moments of the innovation 
%noise $\e_g$ of the data generator, i.e., 
%of $\begin{bmatrix} \y^T & \rvu^T \end{bmatrix}^T$
%$\bE[\|\y(t)-\hat{\y}_\theta(t)\|_2^r]  \le G_e^r(\theta) \bE[\|\e_g(t)\|_2^r]$. \\
\textbf{$\|\alpha_g\|_{\ell_1}$, $\theta_{\infty}(\alpha_g)$}
are the $\ell_1$  norm  and %of the upper bound on 
the mixing coefficients of Markov-parameters of the data
generator:  both quantities decrease with stability of the data generator,
see Remark \ref{norms:rem}.
% it relates high-order moments of $\y$ and $\rvu$
%with those of the noise.
% and it is an upper bound on the well-known $H_2$ norm
%of the data generator. 
%$\bE[\|[\y^T(t),\rvu^T(t)]^T\|_2^r]  \le \|\Sigma_{gen}\|_{\ell_1}^r\bE[\|\e_g(t)\|_2^r]$. 
%\begin{color}{red}
%\textbf{The term $\|\Si\|_{\ell_2}$}
%is the $H_2$ norm of the data generator,
%$\|\Sigma_{gen}\|_{\ell_2} \le\|\Sigma_{gen}\|_{\ell_1}$.
%\end{color}

\item \textbf{$\|\alpha_{\theta}\|_{\ell_1}$} is the $\ell_1$ norm of the predictor
 \eqref{eq:predictor}.  %it upper bounds the well-know $H_2$ norm. 
%depends only the predictor $\Sigma(\theta)$, and it gives an 
%upper bound on high-order moments
%of the difference of empirical errores 
%$V_N(\theta)-\mathcal{L}_N(\theta)$ in terms of high-order moments of $\y$ and $\rvu$, i.e,
%$\bE[|V_N(\theta)-\mathcal{L}_N(\theta)|^r\| \le \frac{2^r}{\sqrt{N}} \Gp^r \bE[\| [\y^T(t),\rvu^T(t)]^T\|^r]$.  
This term decreases with the spectral radius of 
the predictor, see Remark \ref{norms:rem}.  %the more stable the predictor, the smaller is $G_f(\theta)$.
%similarly to 
%$\|(\hat{A}_\theta,\hat{B}_\theta ,\hat{C}_\theta ,\hat{D}_\theta )%\|_{\ell_1}^2$.
\item \textbf{$\theta_{\infty}(\alpha_{\theta})$  } 
describes the mixing properties of the output generated
by the predictor, and it also bounds the sum of differences between 
the infinite and finite past prediction $\sum_{t=0}^{N-1} \|\hat{\rvy}_{\theta}(t)-\hat{\rvy}_{\theta}(t\mid 0)\|$.
Intuitively, the smaller $\theta_{\infty}(\alpha_{\theta})$ is, the closer
$\hat{\rvy}_\theta(t \mid 0)$ is to $\hat{\rvy}_{\theta}(t)$.
This quantity decreases with the spectral radius of the predictor by Remark \ref{norms:rem}.
% the more stable is the predictor, the smaller is $\theta_{\infty}(\alpha_{\theta})$.

\item \textbf{$G_{e,1}(\theta)$} characterizes
mixing properties (short memory condition \cite{alquier2013prediction})  of the prediction
error $\y(t)-\hat{\y}_\theta(t)$. Intuitively, the smaller $G_{e,1}(\theta)$ is, the more the prediction error behaves like an i.i.d process. 
%Intuitively, the smaller $G_{g,1}$ and $G_{f,1}(\theta)$ are, the closer $\rvw$ is to being i.i.d., and the closer $\hat{\rvy}_{\theta}(t)$ is to being i.i.d. , if $\rvw$ is close to i.i.d.
%This shows up on the definition of $G_{e,1}(\theta)$, which is increasing in 
%the $\ell_1$ norm and the WDMC of the Markov-parameters $\alpha_g$ and $\alpha_{\theta}$
%of the data generator and predictors respectively.
This constant is related to stability of the generator and the predictor,
the smaller the spectral radius of those systems are , the
smaller this constant is.  

\item \textbf{$\|\bar{\alpha}\|_{\ell},\theta_{\infty}(\bar{\alpha}),G_{e},G_{e,1}$} are uniform upper bounds on 
$\|\alpha_{\theta}\|_{\ell_1}$, $\theta_{\infty}(\alpha_{\theta})$,
$G_e(\theta),G_{e,1}(\theta)$ respectively. 
\end{itemize}


\section{Main result}
\label{sec:main_results}
 We start by presenting the main result of the paper, which is the following theorem.
\begin{theorem}[Main result]
\label{thm:main}
    Let Assumptions \ref{as:generator}, \ref{as:loss} and \ref{as:parameterisation} hold.
    For any a prior distribution $\pi$ over $\Theta$, confidence level  $\delta\in(0,0.5)$, integer $\epsilon \in \{-1,1\}$,  and $N \ge K_{\delta}$, 
    the inequality
    \begin{equation}
        \label{eq:initial:tildern1} 
    \begin{aligned}
         \forall \rho & \in \mathcal{M}_\pi: \epsilon E_{\theta\sim\rho} \mathcal{L}(\theta)
         \leq   \epsilon E_{\theta\sim\rho}\hat{\mathcal{L}}_N(\theta) +   \\
          & \frac{1}{\lambda}\Big(\KL(\rho\|\pi)  +\ln\dfrac{C_{\delta}}{\delta} + \frac{1}{2}\sum_{i=1}^2
            \ln E_{\theta\sim\pi} \exp\left(\mathcal{C}_i(2\lambda,\theta,N,\delta,\epsilon)\right) \Big)
          %\textcolor{red}{\left(
          %\ln E_{\theta \sim \pi} \exp(\mathcal{C}_2(2\lambda,\theta,N,\delta,\epsilon))\right)}
           %\Psi(\lambda,\pi,N,\delta,\epsilon)
    \end{aligned}
    \end{equation}
    holds with probability at least $1-2\delta$ over data, with
   % \begin{align}
   %     \forall \rho\in \mathcal{M}_\pi:\ 
   %     E_{\theta\sim\rho}\hat{\mathcal{L}}_N(\theta) 
   %     \leq 
   %     E_{\theta\sim\rho} \mathcal{L}(\theta)
   %     +        
   %     \tilde{\mathbb{B}}_N(\rho,\pi,\lambda) \label{eq:initial:tildern2} \\
    %\begin{align}       
    %where
    %\begin{align}
    %& \mathbb{B}_N(\rho,\pi,\lambda)
    %    \triangleq \frac{1}{\lambda}\left ( \KL(\rho\|\pi) + \ln\dfrac{1}{\delta} + \Psi(\lambda,\pi,N)\right ) \label{eq:initial:def:tildern} \\
    %\begin{multline}
    %      \Psi(\lambda,\pi,N,\delta,\epsilon)& \triangleq  \frac{1}{2} 
    %    \ln E_{\theta\sim\pi} \exp(\mathcal{C}(2\lambda,\theta,N,\delta,\epsilon)) \nonumber 
%%        \mathcal{C}(\theta,\lambda, N,\delta) & \triangleq \left(
%%          \left(\lambda \mathcal{C}_{1}(\theta) +  
%%          \lambda^2 \mathcal{C}_{2}(\theta) \right) \frac{1}{N}  +  \right. \nonumber \\
%%           & \left. \lambda \mathcal{C}_3(\theta)  \frac{\ln(N+1)}{N} + 
%%           \lambda^2 \mathcal{C}_4(\theta) \frac{\left(\ln(N+1)\right)^2 }{N} \right.  \nonumber  \\
%%          & \left. \left(\lambda \mathcal{C}_{5}(\theta) 
%%          %\ln\left(\frac{(N+1)}{\delta}\right) + 
%%            +\lambda^2 \mathcal{C}_{6}(\theta) \right) \frac{\ln\left(\frac{(N+1)}{\delta}\right)}{N} + 
%%          \right. \nonumber \\
%%          & \left. \lambda \mathcal{C}_{6}(\theta) \frac{\ln(N+1) \ln\left(\frac{(N+1)}{\delta}\right)}{N} +  \right. \nonumber \\
%%          & \left. \lambda^2 \mathcal{C}_{8}(\theta) \frac{\left(\ln(N+1)\right)^2 \ln\left(\frac{(N+1)}{\delta}\right)}{N} \right) \nonumber 
        %\Psi_1(\theta,\lambda,N)+ \Psi_2(\theta,\lambda,N) \right)}
        %\bE[e^{2\lambda(V_N(\theta)-\hat{\mathcal{L}}_N(\theta))}] \nonumber \\
        %+\ln E_{\theta\sim\pi}\bE[e^{2\lambda (\mathcal{L}(\theta)-V_N(\theta))}] \Big ),\label{eq:Psi:initial}
   % \end{align}
    $\KL(\rho\|\pi)=E_{\theta\sim\rho}\ln\frac{\rho(\theta)}{\pi(\theta)}$ the Kullback–Leibler divergence
    and the constants $K_{\delta}$, 
    $\mathcal{C}_i(\lambda,\theta,N,\delta,\epsilon)$, $i=1,2$ and $C_\delta$ are defined as follows:%\footnote{Note that $\mathcal{C}_i$ does not depend on $\delta$ and $\epsilon$ for bounded noise and Lipschitz loss. } 
    %, where $\mathcal{C}_i(\theta$, $i=1,2,3,4$ are defined as follows:

    \textbf{Bounded noise: }
       If %$\e_g$  is essentially bounded, i.e.,
       $\|\e_g(t)\|_{\infty} \le c_{\rve}$, then $C_{\delta}=1, K_{\delta}=2$, $\lambda > 0$, and %arbitrary and %for $i=1,2$
%    \begin{color}{red}
    \begin{equation*}
    %\label{bounded:pac}
    \begin{aligned}
         & \mathcal{C}_i(\lambda, \theta,N, \delta,\epsilon)  \triangleq 
         \mathcal{C}_{b,i}(\lambda,\theta,L(\theta),c_{\rve},N) \triangleq %\\
         %& \triangleq 
         \frac{\mathcal{C}_{m,b,i}(\lambda,\alpha_{\theta},G_e(\theta)+2G_{e,1}(\theta), L(\theta),c_{\rve})}{N} \\
         % \triangleq \\
         %&\alpha_\theta, G_e(\theta),G_{e,1}(\theta),L(\theta),c_{\rve},N)\\
         & L(\theta)   \triangleq \left\{\begin{array}{rl} 2 G_e(\theta)c_{\rve}\sqrt{n_{\rve}} & \mbox{ $\ell$ is quadratic loss} \\

              L_{\ell} &  \mbox{ $\ell$ is $L_\ell$-Lipschitz }
         \end{array}\right.   \\
         &   \mathcal{C}_{m,b,i}(\lambda,\alpha,X,L,c_{\rve}) 
        %  \footnote{Note that $\mathcal{C}_{m,b,1}$ does not depend on $X$, and $\mathcal{C}_{m,b,2}$ does not depend on $\alpha$. }
         \triangleq \left\{ \begin{array}{rl}
            \lambda Lc_{\rve}\sqrt{n_{\rve}}\theta_{\infty}(\alpha)\|\alpha_g\|_{\ell_1} & i=1 \\
             \lambda^2 \frac{L^2}{2} X (\sqrt{n_{\rve}}c_{\rve})^2 & i=2
          \end{array}\right. 
        \end{aligned}
      \end{equation*}          
          %\footnote{Note that $\mathcal{C}_{m,b,1}$ does not depend on $X$, and $\mathcal{C}_{m,b,2}$ does not depend on $\alpha$. }
      %and the functions $\mathcal{C}_{m,b,i}$, $i=1,2$ is defined as
      %\begin{equation}
      %  \label{bounded:pac:eq1}
      %  \begin{aligned}
          
           %\lambda Lc_{\rve}\sqrt{n_{\rve}}\theta_{\infty}(\alpha)\|\alpha_g\|_{\ell_1}, \quad 
           %  \mathcal{C}_{m,b,2}(\lambda,\alpha,X,L,c_{\rve}) & \triangleq 
        %& \Psi_{1}(\theta,\lambda,N) \triangleq 
        %& \mathcal{C}_{b,1}(\alpha)  \triangleq \theta_{\infty}(\alpha)\|\alpha_g\|_{\ell_1} \sqrt{n_\rve},  %\nonumber \\
        %C_{b,1}(\theta)
         %\quad  \mathcal{C}_{b,2}(\theta) \triangleq 2(X+2Y)n_e  %\nonumber \\
        %% & \mathcal{C}_i(\theta) \triangleq 0, \quad i=3,\ldots,8 \\
         %\mathcal{C}_3(\theta) \triangleq 0 \nonumber   \\
        %& C_{b,1}(\theta)=\theta_{\infty}(\alpha_{\theta})\|\alpha_g\|_{\ell_1} \sqrt{n_\rve}, 
          %& C_{b,2}(\theta) \triangleq (G_e(\theta)+2G_{e,1}(\theta))n_e.  \nonumber 
        %\left (\exp \left \{\textcolor{red}{\frac{4\lambda K^2}{N}\bar{G}_{f,1}(\theta)} \bar{G}_{f,2}(\theta)\right \}\right ) \nonumber \\ %\left ( 1-C_{1,2}(\theta)+ C_{1,2}(\theta) e^{\frac{\lambda}{N}2 C_{1,1}(\theta)}  \right ) \nonumber \\
        %& \Psi_{2}(\theta,\lambda,N) \triangleq 
        %\textcolor{red}{\frac{(2 L(\theta))^2\lambda^2 c_{\rve}^2 C_{b,2}(\theta)}{N}} \nonumber \\
    %\end{aligned}
    %\end{equation}
%  \end{color}
%    where $L(\theta)=L$ 
%    is the Lipschitz constant of the loss function $\ell$ if Assumption \ref{as:loss} holds,
%    or $L(\theta) 2G_e(\theta)c_{rve}\sqrt{n_{rve}}$ if $\ell$ is the quadratic loss. 

   %\end{color}
   \textbf{Lipschitz loss:}
     If  %instead of Assumptions \ref{as:lipschitz}--\ref{as:bounded}, 
     $\ell$ is $L_{\ell}$-Lipschitz, then $C_{\delta}=1, K_{\delta}=2$, $0 <\lambda < \frac{N}{16L_{\ell}}$, and %for $i=1,2$
     %\begin{color}{red}
     \begin{equation*}
     \begin{aligned}
         & \mathcal{C}_i(\lambda,\theta,N,\delta,\epsilon) \triangleq  \mathcal{C}_{ub,i}(\lambda,\theta,L(\theta),N)
         \triangleq 
         %& \triangleq 
         \frac{1}{2} \mathcal{C}_{b,i}\left(2\lambda,\theta,L(\theta),\frac{2\ln(N)}{\|\alpha_g\|_{\ell_1}K},N\right) +
            \bar{\mathcal{C}}_{ub}(\lambda,L(\theta),N) \\
            &  \bar{\mathcal{C}}_{ub}(\lambda,L(\theta),N)  \triangleq \frac{\lambda^2 32 L(\theta)^2 C_{u,1}}{N}+ \frac{4\lambda L(\theta) C_{u,2} \ln(N)}{N}
            =
             \frac{\lambda^2 32 L^2_{\ell} C_{u,1}}{N}+ \frac{4\lambda L_{\ell} C_{u,2} \ln(N)}{N}
           % &  \bar{\mathcal{C}}_{ub}(\lambda,L_{\ell},N)  \triangleq \frac{\lambda^2 32 L^2_{\ell} C_{u,1}}{N}+ \frac{2\lambda L_{\ell} C_{u,2} \ln(N)}{N} 
         %\\ 
        %& 
        %\lambda \frac{L\mathcal{C}_{1,b}(\alpha_{\theta})}{\|\alpha_g\|_2} \frac{\ln(N+1)}{N}+ \\
        %& \lambda^2 2(L)^2 \frac{\mathcal{C}_{b,2}(G_e(\theta),G_{e,1}(\theta))}{\|\alpha\|_g^2} 
           %\frac{\left(\ln(N+1)\right)^2}{N}
        %&\mathcal{C}_1(\theta) \triangleq 0,  \quad \mathbf{C}_i(\theta)=0, \quad i=5,6,7,8 \\
        %& \mathcal{C}_2(\theta) \triangleq 40 L^2C_{1,ub}, \quad 
        %    \mathcal{C}_3(\theta) \triangleq L C_{3,ub}(\alpha_{\theta}), \quad
        %\mathcal{C}_4(\theta)  \triangleq 2L^2 \mathcal{C}_{4,ub}(G_{e}(\theta),G_{e,1}(\theta))
        %\mathcal{C}_4(\theta)  \triangleq 2L^2 \mathcal{C}_{3,ub}(G_{e}(\theta),G_{e,1}(\theta))
     \end{aligned}
    \end{equation*}
    where the constants, $K$,$C_{u,j}$, $j=1,2$ are  defined as follows:
    \begin{equation}
      \label{lipschitz:bound1}
     \begin{aligned}
        &  C_{u,1} \triangleq \bar{\Psi}_{\e_g}(\|\bar{\alpha}\|_{\ell_1}\|\alpha_g\|_{\ell_1})+\bar{\Psi}_{\e_g}(\|\alpha_g\|_{\ell_1}),  \quad \bar{\Psi}_{\rve_g}(s)=e^{14^2s^2 \|\rve_g\|_{\psi_2}^2+3s \|\rve_g\|_{\psi_2}} \\
        & C_{u,2} \triangleq \bar{\Psi}_{\rve}(\|\bar{\alpha}\|_{\ell_1}\|\alpha_g\|_{\ell_1})
        \frac{\|\bar{\alpha}\|_{\ell_1}^2\|\alpha_g\|_{\ell_1}}{K}+
        \bar{\Psi}_{\e_g}(\|\alpha_g\|_{\ell_1}) \frac{\|\alpha_g\|_{\ell_1}}{K},  \quad
        \quad K=\min\{1,\|\bar{\alpha}\|_{\ell_1}\}
        %&  C_{u,1} \triangleq \bar{\Psi}_{\e_g}(\|\bar{\alpha}\|_{\ell_1\|\alpha_g\|_{\ell_1}})+\bar{\Psi}_{\e_g}(\|\alpha_g\|_{\ell_1}),  \quad K=\min\{1,\|\bar{\alpha}\|_{\ell_1}\} \\
        %M& C_{u,2} \triangleq \bar{\Psi}_{\rve}(\|\bar{\alpha}\|_{\ell_1}\|\alpha_g\|_{\ell_1})
        %\frac{\|\bar{\alpha}\|_{\ell_1}^2\|\alpha_g\|_{\ell_1}}{K}+
        %\bar{\Psi}_{\e_g}(\|\alpha_g\|_{\ell_1}) \frac{\|\alpha_g\|_{\ell_1}}{K},  
        %\left(C_{12}+C_{22}\right) \nonumber \\ 
      %& C_{3,ub}(\alpha) \triangleq 5C_{2,ub}
        %+2 \theta_{\infty}(\alpha)\|\alpha_g\|_{\ell_1} \sqrt{n_\rve}}{\|\alpha_g\|_{\ell_1}}  \\
   %     & C_{3,ub}(X,Y)=\left(X+2Y\right)n_e}{\|\alpha_g\|_{\ell_1}^2}  \\
   \end{aligned}
    \end{equation}
       % \Psi_1(\theta,\lambda,N) & \triangleq
        %\exp\left(K_{f,1}(\theta,\lambda,N)+K_{f,2}(\lambda,N)\right) 
       % \exp\left(\lambda L  \left(C_{11}+\frac{2 C_{b,1}(\theta)}{K}\right) \frac{\ln(N+1)}{N} + \right. \nonumber \\
       % & \left.  (L\lambda)^2C_{12}\frac{1}{N}\right)
       % \nonumber \\
         %\bar{\Psi}_{2}(\theta) & \triangleq \exp\left(K_{\infty,1}(\theta,\lambda,N)+K_{\infty,2}(\lambda,N)\right) \nonumber \\
       %  \Psi_{2}(\theta,\lambda,N) & \triangleq \exp\left((L\lambda) C_{21}\frac{\ln(N+1)}{N} + 
       %  (L\lambda)^2C_{22}\frac{1}{N} + \right.  \nonumber  \\
       %  &  \left. (L\lambda)^2 \frac{2 C_{b,2}(\theta)}{K^2} \frac{(\ln(N+1))^2}{N} \right)
         %(K_{\infty,1}(\theta,\lambda,N)+K_{\infty,2}(\lambda,N)\right) 
         %\nonumber  
 %& C_{12}=16(\Psi_{\e_g}(G_fG_g)+\Psi_{\e_g}(G_g)) \nonumber \\
% & C_{21} =\frac{3(\Psi_{\e_g}(G_fG_g)G_fG_g+\Psi_{\e_g}(G_g)G_g)}{K}  \nonumber \\
 %& C_{22}= 24(\Psi_{\e_g}(G_fG_g)+\Psi_{\e_g}(G_g)) \nonumber \\
 %& K=\min\{G_e,\|\} \nonumber  
   %  \end{aligned} 
   % \end{equation}

 \textbf{Quadratic loss:}   
   if $\ell$ is the quadratic loss, then $C_{\delta}=\frac{2}{3}$, 
   $K_{\delta}=3$,
   %$K_{\delta}=\sqrt{\frac{3}{2\delta}}$, 
   $0 < \lambda < \frac{N}{16 L_{qd}}$,
 %and %for $i=1,2$
 %using the constants and functions
 %defined in \eqref{lipschitz:bound1}, then 
  \begin{equation*}
  %\label{quadratic:bound}
  \begin{aligned}
     \mathcal{C}_i(\lambda,\theta,N,\delta,\epsilon)  \triangleq  
     %\frac{\lambda \ln(\frac{3N^2}{2\delta}) 72G_e^2\|\e_g\|_{\psi_2}^2\max\{0,\epsilon\}}{N} + 
     \frac{\lambda L^2_{qd} \max\{0,\epsilon\}}{2N} +
       \mathcal{C}_{ub,i}\left(\lambda,\theta,L_{qd},
       N\right), \quad   L_{qd}=12G_e\|\e_g\|_{\psi_2}\sqrt{\ln\left(\frac{3N^2}{2\delta}\right)} 
       %12 G_e \|\e_g\|_{\psi_2}\sqrt{\ln\left(\frac{3N^2}
       %{2\delta}\right)},
  \end{aligned}
  \end{equation*}
\end{theorem}
  For $\epsilon=1$, \eqref{eq:initial:tildern1} bounds the average generalization gap $E_{\theta \sim \rho}(\mathcal{L}(\theta) - \hat{\mathcal{L}}_N(\theta))$; 
  for $\epsilon=-1$, it bounds $E_{\theta \sim \rho}(\hat{\mathcal{L}}_N(\theta) - \mathcal{L}(\theta))$. 
  The latter is useful for deriving oracle inequalities. By applying the union bound to both cases, we obtain a bound on $\left|E_{\theta \sim \rho}(\mathcal{L}(\theta) - \hat{\mathcal{L}}_N(\theta))\right|$.
  Note that except for the quadratic loss case with potentially unbounded data, 
  $\mathcal{C}_i(\lambda,\theta,N,\delta,\epsilon)$ is independent of $\delta$ and $\epsilon$.

%  Also, the bounds for the case of unbounded data can be derived from the bounded one 
%  by choosing $c_{\rve}$ to be $O(\ln(N))$ and adding some additional terms of
%  $O(\frac{\ln (N)}{N})$. 

  \textbf{Intuition behind the proof:}
  The proof of Theorem \ref{thm:main} is in Appendix \ref{sect:pf:thm_main}. 
  The key idea is to decompose the generalization gap using an auxiliary empirical error 
    \( V_N(\theta) = \frac{1}{N}\sum_{t=0}^{N-1}\ell(\rvy(t),\hat{\rvy}_\theta(t)) \)
    where the predictor has run for infinite time (started at $t_0=-\infty$). Since $\bE[V_N(\theta)]=\mathcal{L}(\theta)$, we have 
    \(\epsilon(\mathcal{L}(\theta) - \hat{\mathcal{L}}_N(\theta)) = \epsilon (\mathcal{L}(\theta) - V_N(\theta)) + \epsilon (V_N(\theta) - \hat{\mathcal{L}}_N(\theta)) \).
    The constants $\exp(\mathcal{C}_1)$ and $\exp(\mathcal{C}_2)$ bound the moment generating functions of these two differences. The PAC-Bayesian bound then follows by standard concentration arguments \cite[Chapter 2]{alquier2021userfriendly}.
    
    \textbf{Bounded case:} $\mathcal{C}_1$ is obtained by direct calculation, and $\mathcal{C}_2$ via Rios's concentration inequality for weakly-dependent processes \cite[Prop.\ 4.2]{alquier2013prediction}, improving upon constants in \cite{eringis2023pacbayesRNN,eringis2023pacbayesian}.
    
    \textbf{Unbounded case:} We employ a truncation approach \cite{alquier2012pred}. Define truncated data $\bar{\rvy}(t), \bar{\rvu}(t)$ generated from truncated noise $\bar{\e}_g(t)=\e_g(t)\chi(\|\e_g(t)\| \le C)$ with $C=\frac{2\ln(N)}{K\|\alpha_g\|_{\ell_1}}$, where $K=\min\{1,\|\bar{\alpha}\|_{\ell_1}\}$. We bound the moment generating functions of the differences $\sup_{\theta \in \Theta} |\hat{\bar{\mathcal{L}}}_N(\theta) - \hat{\mathcal{L}}_N(\theta)|$ and $\sup_{\theta \in \Theta} |V_N(\theta) - \bar{V}_N(\theta)|$ by $\exp(\bar{\mathcal{C}}_{ub}(\lambda,L_{\ell},N))$. Combining with the bounded-case bounds yields PAC-Bayesian guarantees for unbounded data.
    
    \textbf{Quadratic case:} We show that with probability at least $1-\frac{2\delta}{3}$ over data, the prediction errors satisfy $\|\y(t)-\y_{\theta}(t)\|_2 \le C$ and $\|\y(t)-\y_{\theta}(t\mid 0)\|_2 \le C$ for $C=6G_e\|\e_g\|_{\psi_2}\sqrt{\ln(\frac{3N^2}{2\delta})}$. The empirical quadratic loss then coincides with the empirical error for the truncated Lipschitz loss $\ell_C(y,y')=\min\{\|y-y'\|_2^2, C^2\}$. The result follows by union bound and the Lipschitz case, noting that the true error difference between the quadratic loss and the truncated loss is $O(C^2/N)$.

\textbf{Interpretation and $O(\frac{\ln(N)^{r}}{N})$ rates:}
The interpretation of Theorem \ref{thm:main} is similar to other PAC-Bayesian bounds \cite{alquier2021userfriendly}:
we view learning algorithms as choices of parameters from a  data-dependent probability density on model parameters, 
  referred to as \emph{posterior}
 $\rho^{\train}$, for which 
 the right-hand side of \eqref{eq:initial:tildern1} is small for $\epsilon=1$.
One option is to choose $\rho^{\train}$ as the \emph{Gibbs posterior} 
which minimizes of the right-hand side of \eqref{eq:initial:tildern1} 
and which is given by \cite[Definition 2.1]{alquier2021userfriendly}:
%This minimizer, known as the \emph{Gibbs posterior},
%which  has an explicit 
\begin{equation*}
\begin{aligned}
\rho_{\text{Gibbs}} \triangleq  \mathrm{arg min}_{\rho \in \mathcal{M}_{\pi}}
\left(E_{\theta\sim\rho}\hat{\mathcal{L}}_N(\theta) + 
          \frac{1}{\lambda}  \KL(\rho\|\pi)  \right), \quad 
          %+ \ln\dfrac{C_{\delta}}{\delta} + \frac{1}{2} 
        %\ln E_{\theta\sim\pi} \exp(\mathcal{C}(2\lambda,\theta,N,\delta,1) \right)\right) \\
%\left( \KL(\rho\|\pi)+\ln(\frac{C_{\delta}}{\delta}) + \mathcal{C}(\delta,\epsilon)\right) \\
\rho_{\text{Gibbs}}(\theta)=\frac{\pi(\theta)\exp(-\lambda \hat{\mathcal{L}}_N(\theta))}{E_{\theta\sim\pi}\exp(-\lambda \hat{\mathcal{L}}_N(\theta))}, 
%\quad \lambda_N = \frac{\sqrt{N}}{2L(\ln(N))^r},   
%\label{eq:gibbs}
 %\nonumber
\end{aligned}
\end{equation*}
Once a posterior $\rho^{\train}$ is obtained, the learned model 
 $\theta_{\star}$
 can be chosen as follows:
 %there are several standard ways to use it for choosing a parameter
 %\begin{description}
 \par
 \textbf{Mean posterior (MEP):} 
 %we can choose $\theta_{\star}$ 
 as the mean of the posterior
 $\rho^{\train}$, i.e. $\theta_{\star}=E_{\theta \sim \rho^{\train}} \theta$,
 \par
\textbf{Single draw}:
   draw $\theta_{\star}$ randomly from $\rho^{\train}$,  and
\par
\textbf{Maximum posterior (MAP)}:
   choose $\theta_{\star}$ as %the maximum posterior estimate, i.e.,
    $\theta_{\star}=\mathrm{arg max} \rho^{\train}(\theta)$.
    %as the maximum likelhood of the posterior.
    For $\rho^{\train}=\rho^{\mathrm{Gibbs}}$,
     it minimizes the empirical error regularized using the prior, i.e.,
    $\theta_{\star}=\argmin_{\theta\in\Theta}\left ( \hat{\mathcal{L}}_N(\theta) - \frac{1}{\lambda} \ln\pi(\theta) \right )$,
   %\begin{align}
%\label{map:cost}
%    \theta_{\star} & =
%    %\argmax_{\theta\in\Theta} \rho_{\text{Gibbs}}(\theta)= 
%\end{align}

These choices represent well-known classes of system identification algorithms, such as regularized PEM, 
maximimum likelihood estimation for state-space models, and Bayesian methods, see Remarks \ref{rem:pem}, \ref{rem:fir}, and \ref{rem:ml} in Appendix \ref{app:indent}.
%\textbf{PEM methods $\leftrightarrow$ MAP}
%   By \eqref{map:cost} Regularized PEM methods \cite[Chapter 7]{LjungBook} can be translated to using
%   MAP and Gaussian prior $\pi \sim \mathcal{N}(\theta_m,P)$, 
%   in which case $-\frac{1}{\lambda} \ln(\pi(\theta))=0.5 (\theta - \theta_m)^T P^{-1} (\theta -  \theta_m)$, see Example \ref{rem:pem} for more details.

%\textbf{PEM for FIR $\leftrightarrow$ MEP}
%   For FIR systems the regularized linear least squares method \cite{pillonetto2022regularized}
%    can be  interpreted as MEP, if the prior is chosen to be Gaussian. In this case, the 
%   Gibbs posterior is also Gaussian, and MEP and MAP coincide. Example \ref{rem:fir} for more details.

%\textbf{Maximum likelihood for state-space $\leftrightarrow$ MAP}
%   For linear systems in state-space form with Gaussian noise the
%   maximum likelihood method \cite[Chapter 7]{CainesBook} can be interpreted as MAP
%   applied to the Gibbs posterior corresponding to uniform prior $\pi$, see
%   Example \ref{rem:ml}  

Theorem~\ref{thm:main} bounds the average generalization gap $E_{\theta \sim \rho} \left(\mathcal{L}(\theta) - \hat{\mathcal{L}}_N(\theta)\right)$ by a sum of two terms: the KL divergence $\frac{1}{\lambda}\KL(\rho \| \pi)$ between the posterior and the prior, and a complexity term $\mathcal{C}_i(2\lambda,\theta,N,\delta,\epsilon)$ that depends on the data and the prior. For bounded noise, $\mathcal{C}_i = O(1/N)$; for Lipschitz loss, $\mathcal{C}_i = O(\frac{\ln N)}{N})$; and for quadratic loss, $\mathcal{C}_i = O(\frac{(\ln N)^3}{N})$. By choosing $\lambda = \frac{\sqrt{N}}{L(\ln N)^r}$ with appropriately selected $L$ and $r$, and thus forcing $\frac{1}{\lambda}\KL(\rho \| \pi)$ to tend to zero as $N \to \infty$, we obtain an $O\left(\frac{(\ln N)^r}{\sqrt{N}} \right)$:
  %The result of Theorem \ref{thm:main} while general, is not easy to use. As the next step, 
  % we present a simpler version  of Theorem \ref{thm:main}, obtained by 
  %\eqref{eq:initial:tildern1}, then we obtain
 %The bound above can be  simplified to the following Catoni-like bound:
 \begin{corollary}[$O((\ln(N))^r/\sqrt{N})$ bound]
\label{thm:bounded:alt:col}
With the notation and assumptions of Theorem~\ref{thm:main}, 
\begin{equation}
\label{catoni1}
\begin{split}
  & \forall \rho \in\mathcal{M}_\pi: \epsilon E_{\theta \sim \rho } \mathcal{L}(\theta) \leq \epsilon E_{\theta\sim \rho } \hat{\mathcal{L}}_N(\theta)+  \frac{L(\ln(N))^{r}}{\sqrt{N}}\left(\KL(\rho ||\pi)+\ln\frac{C_{\delta}}{\delta} + 
  \mathcal{C}(\delta,\epsilon) \right)
\end{split}
\end{equation}
holds with probability at least $1-2\delta$ over data, 
for any prior $\pi$, $\delta \in (0,0.5)$, and for any $\epsilon \in \{-1,1\}$, and $N \ge K_{\delta}$, 
where %$C_{\delta},K_{\delta}$ are as in Theorem \ref{thm:main}, and the constants
$\mathcal{C}(\delta,\epsilon)$, $L$, $r$ are defined as follows.
\\
\textbf{Bounded noise:}
  If $\|e_g(t)\|_{\infty} \le c_{\rve}$, then %$C_{\delta}=1$, $L=1$, $r=0$ and 
      $r\triangleq 0$, $L \triangleq 1$, and 
\begin{equation*}
%\label{thm:bounded:alt:col:eq1}
\begin{split}
   & \mathcal{C}(\delta,\epsilon) \triangleq \textcolor{black}{\frac{1}{2} \sum_{i=1}^2 \mathcal{C}_{m,b,i}(2,\bar{\alpha},G_e+2G_{e,1},L_b,c_{\rve})}, \quad 
   %\left(L \|\bar{\alpha}\|_{\ell_1}\|\alpha_g\|_{\ell_1}  +  L^2 (G_e+2G_{e,1}) \right), \\
   %&  L=2G_ec_{\rve}\sqrt{n_{\rve}}
    L_b \triangleq \left\{\begin{array}{rl} 2 G_ec_{\rve}\sqrt{n_{\rve}} & \mbox{ $\ell$ is the quadratic loss} \\
              L_{\ell} & \mbox{$\ell$ is $L_\ell$-Lipschitz }
         \end{array}\right. 
\end{split}
\end{equation*}
\textbf{Lipschitz loss:}
If $\ell$ is $L_{\ell}$-Lipschitz, then %$C_{\delta}=1$, $r=1$, $L=4L_{\ell}$ and 
       $r \triangleq 1$, $L \triangleq 16L_{\ell}$, and 
\begin{equation*}
%\label{thm:unbounded:alt:col:eq1}
\begin{split}
   \mathcal{C}(\delta,& \epsilon)  \triangleq \mathcal{\bar{C}}_{ub}=\left(
        C_{u,1}+C_{u,2} %{\|\alpha_g\|_{\ell_1}^2
   +  
    \frac{1}{2} \sum_{i=1}^2 \mathcal{C}_{m,b,i}\left(\frac{1}{4},\bar{\alpha},G_e+2G_{e,1},1,\frac{1}{\|\alpha_g\|_{\ell_1}K}\right)
      \right), 
   %\\
   %\left(L \|\bar{\alpha}\|_{\ell_1}\|\alpha_g\|_{\ell_1}  +  L^2 (G_e+2G_{e,1}) \right), \\
   %&  L=2G_ec_{\rve}\sqrt{n_{\rve}}
   %& L \triangleq \left\{\begin{array}{rl} 2 G_ec_{rve}\sqrt{n_{rve}} & \mbox{ $\ell$ is the} \\
   %                                                                               &  \mbox{quadratic loss} \\
   %           L_{\ell} & \mbox{$\ell$ is $L_\ell$-Lipschitz }
   %      \end{array}\right
\end{split}
\end{equation*}
 %Assume that Assumptions  \ref{as:generator}, \ref{as:parameterisation} and Assumptions \ref{as:Lipschitz}--\ref{as:bounded} hold. Then 
%Then for any prior $\pi$ and for $\lambda_N = \sqrt{N}$, 
%In particular, for any prior 
%
%    Moreover, for quadratic loss the constant $\mathcal{C}_{b}$ is of the form
%    \[
%      \begin{aligned}
%        & \mathcal{C}_b=2G_e\theta_{\infty}(\bar{\alpha})\|\alpha_g\|_{\ell_1}c_\rve^2n_{\rve}+ 
%        %& \left. 
%          (2G_e)^2(G_e+2G_{e,1})c_{\rve}^4 n_{\rve}^2. 
%      \end{aligned} 
%    \]
\textbf{Quadaratic loss:}
If $\ell$ is quadratic, then %$C_{\delta}=3/2$, $r=3/2$ and 
$L \triangleq 192G_e \|\e_g\|_{\psi_2} \sqrt{2}\sqrt{1+\ln\left(\frac{3}{2\delta}\right)}$,  %\quad C_{\delta} \triangleq 
$r \triangleq 3/2$, % \quad K_{\delta} \triangleq \sqrt{\frac{3}{2\delta}}$
\begin{equation*}
%\label{thm:quadratic:alt:col:eq1}
\begin{split}
   \mathcal{C}(\delta,\epsilon) \triangleq & %\mathcal{\bar{C}}_{ub}+\frac{L}{128} \max\{0,\epsilon\} 
    \frac{L\max\{0,\epsilon\}}{128} +C_{u,1}+C_{u,2}+\frac{1}{2} \textcolor{black}{\sum_{i=1}^2 \mathcal{C}_{m,b,i}\left(\frac{1}{4},\bar{\alpha},G_e+2G_{e,1},1,\frac{1}{\|\alpha_g\|_{\ell_1}K}\right)} \\
   %& C_{1,ub}+C_{u,2} + 
     %\quad r=3/2
   %\\
   %\left(L \|\bar{\alpha}\|_{\ell_1}\|\alpha_g\|_{\ell_1}  +  L^2 (G_e+2G_{e,1}) \right), \\
   %&  L=2G_ec_{\rve}\sqrt{n_{\rve}}
   %& L \triangleq \left\{\begin{array}{rl} 2 G_ec_{rve}\sqrt{n_{rve}} & \mbox{ $\ell$ is the} \\
   %                                                                               &  \mbox{quadratic loss} \\
   %           L_{\ell} & \mbox{$\ell$ is $L_\ell$-Lipschitz }
   %      \end{array}\right
\end{split}
\end{equation*}
%\mathcal{C}_{ub}(\lambda,\theta,
%           2G_e\sigma_{\rve} \sqrt{\frac{3\ln(N+1)}{2\delta}},N) + \\
%Assume that Assumptions  \ref{as:generator}, \ref{as:parameterisation} and \ref{as:loss}  hold.
%Then 
%%Define 
%\[ 
%     \C_{ub}=\mathcal{C}_1+\mathcal{C}_2+\mathcal{C}_3
%\]
%Then for any prior $\pi$, $\lambda_N=\frac{\sqrt{N}}{L\ln(N+1)}$, and $N \ge 2$, 
%for  $\lambda =\frac{\sqrt{N}}{L \ln(N+1)}$,  and $N \ge 2$
% in \eqref{thm:unbounded:alt:col:eq1},
%\begin{equation}
%\label{thm:unbounded:alt:col:eq1}
%\begin{split}
%     & \Psi(\lambda_N,\pi,N) \le  \mathcal{C}_{ub}\left(\ln(N+1)\right)^2 \\
%     & \mathcal{C}_{ub}=0.5\left( \frac{5C_{2,u}+2 \theta_{\infty}(\bar{\alpha})\|\alpha_g\|_{\ell_1} \sqrt{n_\rve}}{\|\alpha_g\|_{\ell_1}} + \right. \\
 %    & \left. \frac{2 \left(G_e+2G_{e,1}\right)n_e}{\|\alpha_g\|_{\ell_1}^2} + 40 L^2C_{1,ub} \right)
 %     %& \frac{(L\lambda)^2\mathcal{C}_1}{N}+\frac{(L\lambda) \mathcal{C}_2\ln(N+1)}{N}+ 
 %     %\frac{(\lambda L)^2 \mathcal{C}_3(\ln(N+1))^2}{N}%M\end{split}      
%\end{equation}
%In particular, for any prior $\pi$
%and for any $\delta \in (0,0.5)$, $N \ge 2$, $\epsilon \in \{-1,1\}$,
%$\Phi(\lambda,\pi,N) \le 2L C_{ub} \frac{1}{N}
%\begin{equation}
%\label{catoni:unbounded1}
%\begin{split}
%   \forall \rho \in\mathcal{M}_\pi: & \epsilon E_{\theta \sim \rho } \mathcal{L}(\theta)  \leq  \epsilon E_{\theta\sim \theta } \hat{\mathcal{L}}_N(\theta)+ \\
  %& \frac{L \ln(N+1)}{\sqrt{N}}(\KL(\rho ||\pi)+\ln\frac{1}{\delta}) + \\ 
  %& \frac{L\mathcal{C}_1}{\sqrt{N}}+\frac{L\mathcal{C}_2 \ln(N+1)}{\sqrt{N}}+ \frac{L\mathcal{C}_3(\ln(N+1)}{\sqrt{N}} \\
  %& \le   E_{\theta\sim \theta } \hat{\mathcal{L}}_N(\theta)
%  & + \frac{\ln(N)}{\sqrt{N}} 2L(\KL(\rho ||\pi)+\ln\frac{1}{\delta}+C_{ub}) \\
%\end{split}
%\end{equation}
%Since the empirical error converges to the generalization loss as $N \rightarrow \infty$, we expect that for tight bounds the term . This would imply that we want  However, $\Psi(\lambda,N)$ increases with $\lambda$. Thus for $\tilde{\mathbb{B}}_N$ to converge to 0, the term $\Psi(\lambda,N)$ should converge to 0 as $N\rightarrow\infty$. This means that for \eqref{eq:initial:tildern} to be non-vacuous $\lambda$ cannot increase too fast.
\end{corollary}
The proof of Corollary~\ref{thm:bounded:alt:col} follows by substituting $\lambda=\frac{\sqrt{N}}{L(\ln N)^r}$ into Theorem~\ref{thm:main} and simplifying the resulting bounds (see Appendix~\ref{sec:proof_corollary} for the proof).
Note that while Corollary~\ref{thm:bounded:alt:col} makes the $O\left(\frac{(\ln N)^r}{\sqrt{N}}\right)$ rates explicit, Theorem~\ref{thm:main} may give a  tighter bound even for $\lambda=\frac{\sqrt{N}}{L(\ln N)^r}$ \cite{alquier2021userfriendly}.
%as  it is illustrated in Appendix \ref{sec:numerical_examples}.
%for a specific choice of $\lambda$, in general, Theorem~\ref{thm:main} allows $\lambda$ to be a free parameter, which can be optimized for tighter bounds
%for a given $N}. 


%\paragraph*{PAC-Bayesian bounds for learning and relationship with system identification}
%\label{rem:single_draw}
%Next, we discuss how 
%PAC-Bayesian bounds can be used for learning. 
%%\begin{Remark}[Application for learning]
%First, PAC-Bayesian bounds can be used to deribe uniform PAC bounds, by choosing posteriors
%which make computing the KL-divergence term easy. In turn, PAC bounds can be used to derive generalisation guarantees for standard learning algorithms, such as regularized empirical risk minimization (ERM). 
%\begin{lemma}
%  \label{lem:pac}
%   If $\hat{\theta}=\argmin_{\theta \in \Theta} \hat{\mathcal{L}}_N(\theta)+\frac{1}{\lambda}C(\theta)$, then we have the \emph{oracle inequality} with probability at least $1-2\delta$ over data,
%   \begin{equation}
%   \label{lem:pac:eq2}
%   \begin{split}
%     \mathcal{L}(\hat{\theta}) \le \min_{\theta \in \Theta} \left(\mathcal{L}(\theta)+ 2\frac{2L(\ln N))^r}{\sqrt{N}} C(\theta)\right) + 
%     2\frac{2\sigma^2L}{\sqrt{N}} +  
%     2\frac{2L(\ln N))^r}{\sqrt{N}} \left(\ln \frac{C_{\delta}}{\delta} + \mathcal{C}(\delta,\epsilon)\right)
%   \end{split}
%   \end{equation}
%  \end{lemma}
%  The proof is presented in Appendix, where the general bound from Theorem \ref{thm:initial} is used instead of Corollary \ref{thm:bounded:alt:col}, allowing to keep $\lambda$ as a free parameter.
  % We can then use standard optimization algorithms to minimize the right-hand side of \eqref{lem:pac:eq1} over $\theta \in \Theta$, leading 
  %is a slightly modified version of \cite[Example 2.1,2.2]{alquier2021userfriendly} and it is presented in Appendix. 
%  Lemma \ref{lem:pac} provides a bound for the learning algorithm which minimizes the empirical error with a regularisation term $C(\theta)$, which, in the Gaussian prior case, is  the standard ridge regularisation.
%  This idea can be used in a more general setting. 

%Theorem \ref{thm:main} and Corollary \ref{thm:bounded:alt:col} can be used for learning as follow.



%$\frac{1}{\lambda}\left ( \KL(\rho\|\pi) + \ln\dfrac{C_{\delta}}{\delta} + \Psi(\lambda,\pi,N,\delta,\epsilon)\right )$
%The term 
%\[ \tilde{\mathbb{B}}_N(\rho^{\train},\pi,\lambda)=\frac{1}{\lambda}\left ( \KL(\rho\|\pi) + \ln\dfrac{C_{\delta}}{\delta} + \Psi(\lambda,\pi,N,\delta,\epsilon)\right ) \]
 %can be viewed as a regularisation term. 
%The posterior $\rho^{\train}$ can be chosen in various ways.
% For instance, 
%if $\pi$ is a Gaussian prior $\mathcal{N}(\theta_m,P)$, then
%$\pi$ could be chosen as a Gaussian distribution $\mathcal{N}(\theta_{\train},Q)$ with mean $\theta_{\train}$ and covariance $Q$ to be optimized.
%Instead of minimizing \eqref{eq:initial:def:tildern}, 
%$\rho^{\train}$ to make the KL-divergence term easy to compute, leading to uniform PAC bounds as follows. 
%For example, if both prior $\pi$ and posterior $\rho^{\train}$ are chosen to be Gaussian distributions, then the KL-divergence has a closed form expression \cite[Example 2.1]{alquier2021userfriendly}.     
%then the posterior $\rho^{\train}$ can be chosen to be Gaussian or uniform. In this case, 
%the KL-divergence has a closed form, and an optimization problem can be formulated \cite[Example 2.2]{alquier2021userfriendly}.
%Interpreting priors or prior dependent expressions as regularisation terms was explored for LTI systems in the book \cite{pillonetto2022regularized}.
%There is an explicit expression for the data-dependent posterior which
%minimizes $E_{f\sim\rho^{\train}} \hat{\mathcal{L}}_N(f) + \tilde{\mathbb{B}}_N(\rho^{\train},\pi,\lambda)$, which is known as the 

%This minimizer, known as the \emph{Gibbs posterior}, has an explicit form \cite[Definition 2.1]{alquier2021userfriendly}:
%\begin{align}
%\rho_{\text{Gibbs}}(\theta)&\triangleq \frac{\pi(\theta)\exp(-\lambda \hat{\mathcal{L}}_N(\theta))}
%{E_{\theta\sim\pi}\exp(-\lambda \hat{\mathcal{L}}_N(\theta))}
%Z\triangleq .
%\label{eq:gibbs}
 %\nonumber
%\end{align}
 
%\end{description}
%If $\rho^{\train}$ is Gaussian, then it is well-known \cite{CainesBook,pillonetto2022regularized} that MAP and MEP are the same. 
%Moreover, if we take the Gibbs posterior, i.e., $\rho^{\train}=\rho_{\text{Gibbs}}$, then 
%\
%That is, the MAP estimate from the optimal density $\rho_{\text{Gibbs}}$ is equivalent to minimising the empirical error, with a regularisation term $-\lambda^{-1}_N\ln\pi(\theta)$.
%M

The previous results bound the average generalization gap with respect to a posterior $\rho^{\train}$. For practical learning, however, we are interested in the generalization gap for a specific choice of $\theta_{\star}$ (e.g., MEP, MAP, or a single draw from $\rho^{\train}$). Below, we provide such bounds when $\rho^{\train}$ is the Gibbs posterior with $\lambda = \frac{\sqrt{N}}{L(\ln N)^r}$, where $L, r$ are as in Corollary~\ref{thm:bounded:alt:col}. This yields a generalization gap for $\theta_{\star}$ of order $O\left(\frac{(\ln N)^{r+1}}{\sqrt{N}}\right)$.
%While this choice of $\lambda$ ensures the gap vanishes as $N \to \infty$, it may not always be optimal; see \cite[Section 2.1.4]{alquier2021userfriendly} for further discussion and Appendix~\ref{sec:proof_single_draw} for more general cases (e.g., Lemma~\ref{lem:single_draw}).
 We now introduce the necessary assumptions for these results.
\begin{assumption}
\label{as:oracle}
The following conditions are used in the sequel.

\begin{description}
    \item[\textbf{(Opt)}]
    An optimal model $\theta_{true} = \argmin_{\theta \in \Theta} \mathcal{L}(\theta)$ exists.

    \item[\textbf{(LipModel)}]
    The map $\theta \mapsto \alpha_{\theta}$ is $L_{\alpha}$-Lipschitz w.r.t.\
    to the $\|\cdot\|_{\ell_1}$ norm, i.e.,
    \[
        \|\alpha_{\theta_1}-\alpha_{\theta_2}\|_{\ell_1}
        \le L_{\alpha}\|\theta_1-\theta_2\|_2,
        \quad \forall \theta_1,\theta_2 \in \Theta.
    \]

    \item[\textbf{(Gauss)}]
    $\Theta=\mathbb{R}^{n_{\theta}}$ and $\pi$ is Gaussian
    $\mathcal{N}(\theta_m,P)$.

    \item[\textbf{(Uni)}]
    $\Theta$ is compact and $\pi$ is the uniform density on $\Theta$.

    \item[\textbf{(SConv)}]
    The true error $\mathcal{L}(\theta)$ is smooth and strongly convex
    in $\theta$ and its Hessian is bounded from below
    $m_{\Theta}I \preceq \nabla^2_{\theta}\mathcal{L}(\theta)$
    for all $\theta \in \Theta$.

    \item[\textbf{(Real)}]
    The predictor $\Sigma(\theta_{true})$ equals the predictor
    \eqref{eq:T0sys2} corresponding to the data generator.

    \item[\textbf{(PE)}]
    The spectral density $\Phi_{\w}$ of $\w$ satisfies
    $\Phi_{\w}(iz) \ge m_{\w} I$ for all $z \in [-\pi,\pi]$.
\end{description}
\end{assumption}
\begin{corollary}
\label{lem:bound:mean}
Assume the following
\\
\textbf{(1)} Assumptions \ref{as:generator}, \ref{as:parameterisation}, \ref{as:loss} and  \textbf{(Opt)} and \textbf{(LipModel)} of Assumption \ref{as:oracle}  hold, and
\\
\textbf{(2)}
either \textbf{(Gauss)} or \textbf{(Uni)} of Assumption \ref{as:oracle} holds, and \\
\textbf{(3)}
  $\theta_{\star}$ be selected from $\rho_{\mathrm{Gibbs}}$ using MEP, MAP or a single draw,
and $\rho_{\mathrm{Gibbs}}$ is computed using $\lambda=\lambda_N = \frac{\sqrt{N}}{2L(\ln(N))^r}$
with $r,L,K_{\delta}$ are as defined in Corollary \ref{thm:bounded:alt:col}.\\
\textbf{(4)}
Furthermore, if $\theta_{\star}$ is chosen using MEP, then assume that \textbf{(SConv)} of Assumption \ref{as:oracle} holds. \\
Consider the constants $\mathcal{C}_{em}(\delta,\sigma)$,
$\mathcal{C}_{or}(\delta,\sigma),r_{\text{Lip},1},r_{\text{Lip},2}$ as in Table \ref{tab:constants:sysid}, and consider the 
inequalities
\begin{align}
& \textbf{Empirical bound:} \quad    \mathcal{L}(\theta_{\star}) - \hat{\mathcal{L}}_N(\theta_{\star})  \le  \frac{L(\ln(N))^{r+1}}{\sqrt{N}} C_{em}(\theta_{\star},\delta) 
  \label{bound:emp} \\
  & \textbf{Oracle bound:} \quad \mathcal{L}(\theta_{\star})  - \mathcal{L}(\theta_{true})  \le
     \frac{L(\ln(N))^{r+1}}{\sqrt{N}} \mathcal{C}_{or}(\delta,\sigma) 
     \label{bound:oracle} \\
  & \textbf{Parameter estimation error:} \quad \|\theta_{\star}-\theta_{true}\|^2_2  \le
    \frac{L(\ln(N))^{r+1}}{\sqrt{N}m_{\Theta}}
    \mathcal{C}_{or}(\delta,\sigma), 
  \label{cor:param:est1:eq1}  \\
  & \text{\textbf{$H_2$ error bound:}} \quad 
 % \begin{equation}
    \|H_{\theta_{true}}-H_{\theta_{\star}}\|_{H_2}^2 \le 
     \frac{L(\ln(N))^{r+1}}{m_{\rvw} \sqrt{N}} \mathcal{C}_{or}(\delta,\sigma)
   \label{col1:eq1}
\end{align}
Then the following holds: \textbf{(A)} \eqref{bound:emp} holds
with probability at least $1-(2+r_{\text{Lip},1})\delta$,  \\
\textbf{(B)} \eqref{bound:oracle} holds with probability at least $1-2\delta$ for single draw and MEP, and with probability at least $1-(2+r_{\text{Lip},1})\delta$ for MAP, \\
\textbf{(C)} \eqref{cor:param:est1:eq1} holds with probability at least $1-2\delta$ for single draw and MEP, and with probability at least $1-(2+r_{\text{Lip},1})\delta$ for MAP, 
if \textbf{(SConv)} of Assumption \ref{as:oracle} holds,  \\
\textbf{(D)} \eqref{col1:eq1} holds with probability at least $1-2\delta$ for single draw and MEP,
and with probability at least $1-(2+r_{\text{Lip},1})\delta$ for MAP,
if \textbf{(Real)} and \textbf{(PE)} of Assumption  \ref{as:oracle} holds, $\ell$ is the quadratic loss, and 
 $H_{\theta}$ the transfer function of the predictor  corresponding to $\theta$, and denote by $\|H\|_{H_2}$ the $H_2$ norm of the transfer function $H$. \\
%with probability at least $1-(2+r_{\text{Lip},1})\delta$, 
%\begin{align}
%\label{bound:oracle} \\
% 
%\end{align}
%with probability at least $1-(2+r_{\text{Lip},1})\delta$, 
%\begin{align}
%
%\end{align}
%i
 Here, the probability is taken over data for MEP and MAP,  and over the data and over all samples $\theta_{\star}$ drawn from $\rho_{\mathrm{Gibbs}}$ for the single draw.
 \footnote{The precise meaning of probability over data and samples drawn from $\rho_{\mathrm{Gibbs}}$ is explained in Remark \ref{rem:jointProb}}
%Then \eqref{bound:emp} holds with probability at least at least $1-(2+r_{\text{Lip},1})\delta$,
%\eqref{bound:oracle} holds with p \eqref{cor:param:est1:eq1}.
%Moreover, if \textbf{(SConv)} of Assumption \ref{as:oracle} holds, then \eqref{cor:param:est1:eq1} holds 
%with probability at least  $1-(2+r_{\text{Lip},1})\delta$,

%Then 
%for any $\delta \in (0,0.5)$ and $\sigma > 0$ and $N \ge K_{\delta}$ 
%each of the  
 %for single draw:
%a transfer function
%$\|\cdot\|_{H_2}$ denotes the $H_2$ norm of a transfer function. 
%For MEP and single draw learning algorithms, assume that  Assumption \ref{as:oracle} either \textbf{(TL1)} or Assumption \ref{as:lipschitz:param} holds, and for MAP assume that 
%Assumption \ref{as:lipschitz:param} holds.
%Define
%$K=1$ for MEP and single draw, and $K=2$ for MAP.
%$Cfor MAP assume that Assumption \ref{as:oracle} \textbf{(CL2)} holds.
%L_{\Theta}(\delta)$ be defined as in   Corollary \ref{lem:pac}, and
% if Assumption \ref{as:lipschitz:param} holds, or let $L_{\Theta}(\delta)=L_{\Theta}$ if Assumption \ref{as:oracle} \textbf{(TL1)} holds.
%let $C(\theta)$ as in Corollary \ref{lem:pac}, and 
  %\textbf{Oracle inequality for Gibbs posterior} for any prior $\pi$;:
  %\begin{equation}
  %\label{lem:bound:oracle}
  %  \frac{2L(\ln(N))^r}{\sqrt{N}} \left(\KL(\rho^{\train} ||\pi)+\ln\frac{2C_{\delta}}{\delta} + 2\mathcal{C}(\delta,\epsilon))\right)
  %\end{equation}
  %with probability at least $1-2\delta$ over data. \\
 % \textbf{Oracle bound for MEP}
  % Corollary 2.6 \cite{alquier2021userfriendly})}
    %if $\mathcal{L}(\theta)$ is a convex function of $\theta$ and 
%$\theta_{\star}$ is ME or a single draw, then %for any prior $\pi$,
%Then each of the following inequalities holds
 % \begin{align}
   
 %& \textbf{Oracle bound :} \quad 
 %\begin{equation}
       
 %\label{bound:oracle}  \\
 % %\end{equation}
% &\text{\textbf{Eestimation error:}  if \textbf{(SConv)} of Assumption \ref{as:oracle} holds: } \quad
%     \\
  %\end{equation}
%& \text{\textbf{$H_2$ error bound:} ,} \nonumber \\ 
 % \begin{equation}
%  &  \|H_{\theta_{true}}-H_{\theta_{\star}}\|_{H_2}^2 \le 
%     \frac{L(\ln(N))^{r+1}}{m_{\rvw} \sqrt{N}} \mathcal{C}_{or}(\delta,\sigma)
%   \label{col1:eq1}
%\end{equation} 
%\end{align}
 %Mwhere and
\end{corollary}
\begin{table}[ht]
\centering
\caption{Constants for learning bounds under different algorithm choices (Corollary \ref{lem:bound:mean})}\label{tab:constants:sysid}
\begin{tabular}{|l|c|c|}
\hline
\textbf{Constant} & \textbf{MEP / Single Draw} & \textbf{MAP} \\
\hline
$L_{\Theta}(\delta)$ & \multicolumn{2}{c|}{$\begin{array}{ll}
  L_{\alpha}L_{\ell} \|\alpha_g\|_{\ell_1}\sqrt{n_\rve}c_{\rve} & \text{bounded noise} \\
  \frac{1}{16} L_{\alpha}\sqrt{n_\rve} \frac{1}{K} & \text{Lipschitz/quadratic}
  \end{array}$} \\
\hline
$C(\theta)$ & \multicolumn{2}{c|}{$\!\!\!\begin{array}{rl} 
  \frac{1}{2} \big (\sigma^2 \mathrm{trace} (P^{-1}) -n_{\theta}+ (\theta_m-\theta)^TP^{-1}(\theta_m-\theta)+
   \ln \frac{\det(P)}{\sigma^{2n_\theta}} + n_\theta \big) & \!\! \textbf{(Gauss)} \\
     \ln\frac{vol(\Theta)}{(\sqrt{\pi \sigma^2}^{n_{\theta}} \Gamma(n_{\theta}/2+1)}+ \frac{n_{\theta}}{2}  & \!\! \textbf{(Uni)} 
   \end{array}
   $\!\! } \\
\hline 
 $r_{\text{Lip},1}$, 
 $r_{\text{Lip},2}$
 & \multicolumn{2}{c|}{\(r_{\text{Lip},1}=\left\{\begin{array}{rl} 0 & \e_g\mbox{ is bounded }  \\
   1 & \mbox{ otherwise }
   \end{array}\right. \quad 
   r_{\text{Lip},2}=\left\{\begin{array}{rl} 2 & \ell \mbox{ is quadratic and $\e_g$ is unbounded }  \\
   0 & \mbox{ otherwise }
   \end{array}\right.
   \)
}\\
\hline
$\mathcal{C}_{L,emp}$ & \multicolumn{2}{c|}{$L_{\Theta}(\delta) \sigma+2r_{\text{Lip},1}(C_{u,1} + C_{u,2})$} \\
\hline
$\mathcal{C}_{L,true}$ & \multicolumn{2}{c|}{$\mathcal{C}_{L,emp}+2r_{\text{Lip},2}\frac{L}{512}$} \\
$\mathcal{C}_{L}(\delta,\epsilon)$ & \multicolumn{2}{c|}{$\ln\left(\frac{C_{\delta}}{\delta}\right)+ \mathcal{C}_{L,emp}+C(\theta_{\star})+ C(\delta,\epsilon)$} \\
\hline
$C_{em}(\theta_{\star},\delta)$ & $\mathcal{C}_L(\delta,1)+2r_{\text{Lip},1}\ln\left(\frac{1}{\delta}\right)$ & $\mathcal{C}_L(\delta,1)+\mathcal{C}_{L,true} +2r_{\text{Lip},1}\ln\left(\frac{1}{\delta}\right) $ \\
%$C_{em}(\theta_{\star},\delta)$ & 
%$\ln\left(\frac{C_{\delta}}{\delta^{2r_{\text{Lip},1}+1}}\right)+ \mathcal{C}_{L,emp}+C(\theta_{\star})+ C(\delta,1)$ &
%$\mathcal{C}_{L,emp}+\mathcal{C}_{L,true}+C(\theta_{\star})+\ln\left(\frac{C_{\delta}}{\delta^{2r_{\text{Lip},1}+1}}\right) + \mathcal{C}(\delta,1)$ \\
\hline
$C(\theta_{true},\theta_{\star})$ & $\mathcal{C}_{L,true}-2\mathcal{C}_{L,emp}$ & $2\mathcal{C}_{L,true}-C(\theta_{true})+C(\theta_{\star})+\ln\left(\frac{\pi(\theta_{true})}{\pi(\theta_{\star})}\right)$ \\
%$\mathcal{C}_{L,true}+ 2C(\theta_{true})$ &
%$2 (\mathcal{C}_{L,true}+\mathcal{C}_{L,emp})+C(\theta_{true})+C(\theta_{\star})+\ln\left(\frac{\pi(\theta_{true})}{\pi(\theta_{\star})}\right)$ \\
\hline
$\mathcal{C}_{or}(\delta,\sigma)$ & \multicolumn{2}{c|}{$C(\theta_{true},\theta_{\star})+\mathcal{C}_L(\frac{\delta}{2},1)+\mathcal{C}_L(\frac{\delta}{2},-1)$} \\
  %$C(\theta_{true},\theta_{\star},\delta)+2\ln\frac{C_{\delta}2^{2r_{\text{Lip},1}+1}}{\delta^{2r_{\text{Lip},1}+1}} + \mathcal{C}(\frac{\delta}{2},1)+\mathcal{C}(\frac{\delta}{2},-1)$} \\
\hline
\end{tabular}
\end{table}
%\textbf{Discussion on the bounds}
In summary, for any outcome $\theta_{\star}$ obtained via MAP, MEP, or a single draw from the Gibbs posterior, Corollary~\ref{lem:bound:mean} yields $O\left(\frac{(\ln N)^{r+1}}{\sqrt{N}}\right)$ bounds on: \textbf{(1)} the generalization gap \eqref{bound:emp}, \textbf{(2)} the excess true error over the optimal model \eqref{bound:oracle}, and \textbf{(3)} the parameter estimation error \eqref{cor:param:est1:eq1}--\eqref{col1:eq1}. The constants for the learning bounds, as given by Corollary \ref{lem:bound:mean}, are in Table \ref{tab:constants:sysid}. The influence of system properties and convergence rates is discussed in Section~\ref{sec:conclusions}.

\textbf{Intuition behind the proof:}
The proof of Corollary \ref{lem:bound:mean} is in Appendix \ref{sec:proof_single_draw}.  For MAP, \eqref{bound:emp}--\eqref{bound:oracle} are derived from Corollary~\ref{thm:bounded:alt:col} by constructing suitable data-dependent posteriors and applying oracle inequality techniques from \cite[Theorem 4.1]{alquier2021userfriendly}. For single draw and MEP, similar arguments apply using \cite[Theorems 2.7 and Corollary 2.6]{alquier2021userfriendly} combined with the moment generating function bounds from  the proof of Theorem~\ref{thm:main}. Finally, the parameter estimation bounds \eqref{cor:param:est1:eq1}--\eqref{col1:eq1} follow from the oracle inequality \eqref{bound:oracle} by relating the parameter error to the loss difference through strong convexity (for \eqref{cor:param:est1:eq1}) or persistence of excitation (for \eqref{col1:eq1}).


  
\section{Discussion on the bound  and conclusions}
\label{sec:conclusions}
In this paper, we derived PAC-Bayesian error bounds for stochastic LTI systems with inputs and sub-Gaussian noise.
\par
The convergence rates are fastest for bounded noise, slower for Lipschitz losses, and slowest for quadratic losses with general sub-Gaussian noise (see Table~\ref{rates}).
All bounds in Theorem~\ref{thm:main}, Corollary~\ref{thm:bounded:alt:col}, and Corollary~\ref{lem:bound:mean} improve with increased system stability, lower sub-Gaussian noise norm, and depend logarithmically on the confidence level $\delta$ via $\ln(1/\delta)$.
\par
Specifically, the bounds increase monotonically with the sub-Gaussian norm $\|\e_g\|_{\psi_2}$ of the noise, and with the $\ell_1$ norms and mixing coefficients of the Markov parameters of both the predictors and the data generator ($\|\alpha_\theta\|_{\ell_1}$, $\|\alpha_g\|_{\ell_1}$, $\theta_{\infty}(\alpha_\theta)$, $\theta_{\infty}(\alpha_g)$). Smaller $\ell_1$ norms and mixing coefficients correspond to more stable systems and weaker dependence on past inputs, making learning easier. Similarly, a smaller sub-Gaussian norm (i.e., lower noise variance) yields tighter bounds.
\par 
In summary, the bounds are tighter for less noisy and more stable systems, aligning with intuition: learning is easier when the system is less noisy and more stable.

Additionally, if the goal is to bound not only the generalization gap but also the parameter estimation error (in terms of the $H_2$ distance between the true and learned systems \eqref{col1:eq1}, or the Euclidean distance between parameters \eqref{cor:param:est1:eq1}), the bounds also depend on the \emph{persistence of excitation} of the input $\w$ ($m_{\rvw}$) and the \emph{identifiability} constant ($m_{\Theta}$). Larger values of $m_{\rvw}$ (richer input) and $m_{\Theta}$ (greater separation between models) yield tighter bounds. The bound \eqref{col1:eq1} can also be used to bound the system matrices (up to similarity), see Remark~\ref{matrix:trans}, Appendix~\ref{app:indent}.

%The bound \eqref{bound:emp} is a 
%bounds the  generalization gap, i.e. the difference between 
%the true error and the empirical error achieved
%by the learning algorithm.  The bound \eqref{bound:oracle} is  
%$O((\ln(N))^{r+1}/\sqrt{N})$  oracle bound on the difference between the true error of the learned model and the best achievable true error.
%Both bounds converge to zero as $N \rightarrow \infty$,  showing  that the empirical error converges to the true error and the learned system is asymptotically optimal.

%If we assume strong convexity of the true error, then \eqref{cor:param:est1:eq1} provides  a 
% $O((\ln(N))^{r+1}/\sqrt{N})$ bound on the parameter estimation error. In turn, strong convexity of the true error is often done when studying asymptotic behavior of the parameter estimation error for classical algorithms \cite[Theorem 9.1]{LjungBook}, 
% and it can be viewed as an identifiability assumption, see Remark \ref{app:strong convexity}, Appendix \ref{app:indent}.

 %which, in addition, ensures that the problem of optimizing the true error, and by
 %extension, the empirical error, is well-posed. 
% In the absence of strong convexity, but assuming persistence of exicitation of
% the signal $\w$,  we can provide a $O((\ln(N))^{r+1}/\sqrt{N})$ bound on the square  $H_2$ distance between the true and the learned model. 

  
\begin{table}[ht]
\centering
\caption{Summary of convergence rates for different noise and loss function types \label{rates}}
\begin{tabular}{|c|c|c|}
\hline
\textbf{Scenario} & \textbf{Empirical PAC-Bayesian Bound} & \textbf{PAC/Oracle/parameter est./$H_2$ bounds} \\
\hline
\textbf{Bounded noise} & $O\left(\frac{1}{\sqrt{N}}\right)$ & $O\left(\frac{\ln(N)}{\sqrt{N}}\right)$ \\
\hline
\textbf{Lipschitz loss} & $O\left(\frac{\ln(N)}{\sqrt{N}}\right)$ & $O\left(\frac{(\ln(N))^2}{\sqrt{N}}\right)$
 \\
\hline
\textbf{Quadratic loss} & $O\left(\frac{(\ln(N))^{3/2}}{\sqrt{N}}\right)$ & $O\left(\frac{(\ln(N))^{5/2}}{\sqrt{N}}\right)$ 
\\
\hline
\end{tabular}
\end{table}

For parameter estimation error, our obtained rate $O((\ln N)^{r+1}/\sqrt{N})$ is more conservative than classical asymptotic theory \cite{LjungBook,CainesBook} (which yields  $O(1/\sqrt{N})$ asymptotic consistency) and existing finite-sample bounds \cite{oymak2021revisiting,lale2020logarithmic,Simchowitz_Foster_2020,Ziemann1}(which achieve $O(\ln N/\sqrt{N})$) . However, this is the expected price for generality: our bound applies to a broad class of learning algorithms, provides non-asymptotic guarantees with high probability, and simultaneously bounds the generalization gap (properties not offered by prior results). In addition, most existing finite-sample bounds are algorithm-specific (e.g., for least-squares only) and do not address the generalization gap.

%Compared to classical system identification theory (e.g., \cite{LjungBook,CainesBook}), our $O((\ln N)^{r+1}/\sqrt{N})$ bound on the squared parameter estimation error is slightly more conservative than the typical asymptotic $O(1/N)$ rates.
%
%This is expected: classical results are asymptotic and concern mostly parameter estimation error,  while our PAC-Bayesian bounds are non-asymptotic, hold with high probability, and also bound the generalization gap. 

%Compared to existing finite-sample bounds for LTI system identification (e.g., \cite{oymak2021revisiting,lale2020logarithmic,Simchowitz_Foster_2020}), our results differ in several key aspects: (1) we allow for a wide range of learning algorithms including PEM, maximum likelihood, Bayesian, and regularized FIR methods), not just least-squares or specific estimators; (2) we assume the data generator is in a stationary regime, rather than starting from a fixed initial state, which can be problematic,
%see Remark \ref{rem:fin_samp:zero_init}. (3) we provide generalization gap bounds in addition to parameter estimation error.  The trade-off is that our $O(((\ln N)^{r+1}/\sqrt{N})$ bound for the squared parameter estimation error is slightly slower than the $O(\ln(N)/N)$ rate in the existing finite-sample bounds, but the results are more broadly applicable. See Remark \ref{rem:fin_samp:zero_init}.

% and provide stronger (high-probability, non-asymptotic) guarantees.
%that the rate  but the results are more broadly applicable and provide stronger (high-probability, non-asymptotic) guarantees.




%\textbf{Comparison with system identification}
%The results of Corollary \ref{lem:bound:mean} are consistent with classical asymptotic results for maximum likelihood and PEM methods \cite{LjungBook,CainesBook}, which typically yield $O(1/\sqrt{N})$ convergence rates. Our finite-sample PAC-Bayesian bounds are slightly more conservative, with rates of $O((\ln N)^{r+1}/\sqrt{N})$ for the generalization gap and $O((\ln N)^{(r+1)}/\sqrt{N})$ for the squared parameter or $H_2$ error. 
%In contrast, classical results focus on asymptotic convergence in probability under more restrictive assumptions. Our bounds also explicitly relate the parameter estimation error to the empirical error, which is not typically addressed in classical system identification theory. 
%\end{remark}
%(\ref{eq:T0sys2})}, 
%$f_{\theta_{\st$, 
%see \eqref{param:eq4}-\eqref{param:eq6} of the supplementary material for more details.
%\par
%These bounds involve either the average
%the empirical error $\E_{f \sim \rho^{\train}} \hat{\mathcal{L}}_N(f)$ or the empirical
%error $\hat{\mathcal{L}}_N(f_{\theta_{\star}})$.
%Intuitively, \eqref{lqg:rem:eq1}-\eqref{lqg:rem:eq2} establishes
%an upper bound on the distance between \emph{all}
%the true and estimated Markov parameters ($H_2$ distance).
%in terms of 
%the difference between the empirical error of the estimated and the true parameter. 
%\par
%\textbf{Comparison with other finite-sample bounds}
%They apply to  any learning algorithm which can be presented as sampling/take average 
%from a  posterior, see \cite{pillonetto2022regularized} for examples.  
%Recent works such as \cite{oymak2021revisiting,lale2020logarithmic,Simchowitz_Foster_2020} provide finite-sample parameter estimation bounds for specific algorithms, typically assuming the data generator starts from a zero initial state and often neglecting the stochastic component. In contrast, our results: \textbf{(1)} apply to a broad class of learning algorithms (including PEM, maximum likelihood, Bayesian, and regularized FIR methods), \textbf{(2)} assume the data generator is in a stationary regime, \textbf{(3)} address both deterministic and stochastic components, and \textbf{(4)} provide generalization gap bounds. The trade-off is that our parameter estimation bound is $O\left(\frac{(\ln N)^{r+1}}{\sqrt{N}}\right)$, which is more conservative than the $O\left(\frac{\ln N}{N}\right)$ rates in the cited works.
%However, in contrast to the cited results, they apply to a wide variety of algorithms. 
%That is, the conservativeness of the PAC-Bayesian bound is the price to pay for its generality. 
%Moreover, PAC-Bayesian bounds allow us not only to bound the parameter estimation error, but
%also to bound the generalization loss.
% with a term of order $O(\frac{1}{\sqrt{N}})$.
%In turn, $O(\frac{1}{N})$ bounds on the $H_2$ norm of the difference between the true and estimated system do not translate automatically to 
%$O(\frac{1}{N})$ bounds on the generalization error in the presence of 
%noise. 
%That is, the presented bounds can be used to bound quantities not addressed by \cite{oymak2021revisiting,lale2020logarithmic,Simchowitz_Foster_2020}, and the presented bounds apply to a wider range of learning algorithms.
%The price to pay for this generality is that the bounds are more conservative when applied to the specific 
%setting of
%\cite{oymak2021revisiting,lale2020logarithmic,Simchowitz_Foster_2020}.
%From a practical point of view, the difference between %$O(\frac{1}{N})$ and 
%$O(\frac{1}{\sqrt{N}})$ need not always be critical, 
%as for relatively low $N$, higher constants can mask the advantages of the faster rates. A detailed numerical comparison remains topic of future research.
%\end{remark}



%For data generated by an LTI system with gaussian noise, the error bound is asymptotically bounded from below, which indicates that the bound is not
%For data generated by an LTI system with bounded innovation noise, the error bound converges to zero at the rate $O(\frac{1}{\sqrt{N}})$,
%which is comparable to most of PAC-Bayesian bounds.

%\textbf{Future research}
%Future research will be directed towards extending these results to more general state-space representations and using the results of the paper for deriving oracle inequalities \cite{alquier2021userfriendly}.

\bibliographystyle{plain}
\bibliography{bib}
\appendix
\newpage
\section{Proofs of Section \ref{sect:Problem_formulation}}
\label{app:sect:Problem_formulation}
\begin{remark}[LTI in innovation form]
  \label{rem:innovation}
   Let us assume that $\e_g$ has zero mean. It follows
   that $\e_g(t)$ is the innovation process of $[\rvy^T(t)\;\rvu^T(t)]^T$
   and \eqref{eq:generator} is in the so called 
   (forward) innovation form, see
   \cite[Section 8.5.2]{Katayama:05}.
   Intuitively, it means that $\e_g$
   is the difference between $[\rvy^T(t)\;\rvu^T(t)]^T$ and its best linear prediction based on its own past values.
  \end{remark} 
\begin{remark}[Equivalence of \eqref{eq:generator} and \eqref{eq:T0sys1}]
  \label{rem:equiv_gen:proof}
Let us assume that $\e_g$ and $\e_s$ are both zero mean.
More precisely, if there is no feedback from $\rvy$ to $\rvu$ according to Definition 17.1.1. of \cite{LindquistBook},
then \eqref{eq:T0sys1} holds and by \cite{eringis2021optimal} the matrices $A_0,B_0,C_0,K_0,C_0,D_0$ can be computed from the matrices $A_g,K_g,C_g$. 

%\textbf{Existence of \eqref{eq:T0sys1} $\implies$ Assumption 4.1}
Conversely, if \eqref{eq:T0sys1} holds and $\rvu(t)$ is generated by an ARMA process, and there is no feedback from $\rvy(t)$ to $\rvu(t)$, then 
Assumption \ref{as:generator} holds and $A_g,K_g,C_g$  can be computed from those of \eqref{eq:T0sys1} and the ARMA representation of $\rvu(t)$, again by reversing the equations from
\cite{eringis2021optimal}. As a passing remark, if there is a feedback from $\rvy(t)$ to $\rvu(t)$ and  $\rvu(t)$ is generated by an LTI systems driven by $\y(t)$, i.e.,  
\begin{equation}\label{eq:u_gen}
 %   \begin{align}
    \rvs_\rvu(t+1)=A_u\rvs_\rvu(t)+K_u\y(t), \quad 
    %\label{eq:u_gen_x} 
    \rvu(t)=C_u\rvs_\rvu(t) 
    %\label{eq:u_gen}
%\end{align}
\end{equation}
then again the generator described in Assumption \ref{as:generator} can be obtained by adding \eqref{eq:u_gen} to \eqref{eq:T0sys1}. 
\end{remark} 
\begin{remark}[Methods for stable parametrisations]
\label{rem:parameter:methods}
 The construction of \cite[Section 3.2.3]{RalfPeeters} relies on canonical forms based on structure indices, to which the diffeomorphism from \cite{HanzonStable1}
 between general and stable parametersations are applied. Intuitively, the latter transformation maps $(A,B,C,D)$ to $(\frac{A}{\mu},B,C,D)$, where $\mu$ is the maximum of moduli of the eigenvalues of $A$. 
 The construction of \cite[Section  4.6]{Ribarits1} uses Ober's balanced canonical forms, and it relies on fixing the observability and controllability Grammians to being equal and diagonal. 
 That is, stability is assured by using diagonal quadratic Lyapunov functions. 
 In this case, by fixing upper bounds on the diagnonals of the Grammians and on the norms of the system matrices occurring in the parameterisation, we ensure existence of constants $M,C,\gamma$ mentioned in Assumption \ref{as:parameterisation}. 
\end{remark}
\begin{proof}[Proof of Lemma \ref{l:ihp}]
We have
%\begin{align}\label{ytt0}
$
\y_\theta(t\,|\,t_0)=\sum_{k=t_0}^{t-1} H_\theta(t-1-k)\,\w(k)+D_\theta \w(t)=\sum_{k=0}^{t-1-t_0} H_\theta(k)\,\w(t-1-k)+D_\theta \w(t)
$
%\end{align}
with the impulse response $H_\theta(k)\triangleq C_\theta A_\theta^{k}B_\theta$ for $k\ge 0$. Because $A_\theta$ is Schur, there exist $M>0$ and $\rho\in(0,1)$ such that
$\|A_\theta^{k}\|_2\le M\rho^k$, hence $\sum_{k=s}^\infty \|H_\theta(k)\|_2\leq\|C_\theta\|_2\|B_\theta\|_2M\frac{\rho^{s}}{1-\rho} \rightarrow 0$ as $s \rightarrow \infty$. 
For $t_1<t_2\leq t$ we get, using stationarity of $\w$ ($\sigma_\w^2=\bE\|\w(s)\|_2^2$) and setting $\beta=\rho^2$, that 
\begin{subequations}\label{ytt2-ytt1}
\begin{align}
\bE[\left\|\y_\theta(t\,|\,t_2)-\y_\theta(t\,|\,t_1)\right\|_2^2]
&=\bE\left[
\left\|\sum_{m=t-t_2}^{t-1-t_1}H_\theta(m)\w(t-1-m)\right\|_2^2
\right]\\
&\leq(t_2-t_1)\bE\left[\sum_{m=t-t_2}^{t-1-t_1}
\|H_\theta(m)\|_2^2\|\w(t-1-m)\|_2^2
\right]\\
&\leq\sigma_\w^2(\|C_\theta\|_2\|B_\theta\|_2M)^2(t_2-t_1)\sum_{m=t-t_2}^{t-1-t_1}\beta^m\\
&=\sigma_\w^2(\|C_\theta\|_2\|B_\theta\|_2M)^2\frac{\beta^t}{1-\beta}
(t_2-t_1)(\beta^{-t_2}-\beta^{-t_1})
\end{align}
\end{subequations}
% \begin{align}\label{ytt2-ytt1}
% \bE[\left\|\y_\theta(t\,|\,t_2)-\y_\theta(t\,|\,t_1)\right\|_2^2]
% %\le
% %\sigma_w\Big(\sum_{k=t_1}^{t_2-1}\|H_\theta(t-1-k)\|\Big)^2
% \le
% \sigma_w^2\Big(\sum_{m=t-t_2}^{t-1-t_1}\|H_\theta(m)\|_2\Big)^2
% \leq
% \sigma_w^2\Big(\sum_{m=t-t_2}^{\infty}\|H_\theta(m)\|_2\Big)^2
% %\leq
% %\sigma_w\Big(K\frac{\rho^{t-t_2}}{1-\rho}\Big)^2
% \end{align}
%with the right-hand side tending to $0$ as $t_2\to-\infty$. 
Therefore $\{\y_\theta(t\,|\,t_0)\}_{t_0\le t}$ is Cauchy and hence converges in mean square to
%\begin{align}\label{yt}
$
\y_\theta(t)= \sum_{k=-\infty}^{t-1} H_\theta(t-1-k)\,\w(k) \;+\; D_\theta \w(t)=\sum_{k=0}^{\infty} H_\theta(k)\,\w(t-1-k)+D_\theta \w(t).
$
%\end{align}
This representation shows that $\y_\theta$ is the output of a causal, time-invariant, stable LTI filter
(with impulse response $\{H_\theta(k)\}$ and feed-through $D_\theta$) driven by the stationary input $\w$. Hence $\y_\theta$ is stationary, since stable LTI filtering preserves stationarity. Moreover, it follows that
\begin{align}\label{ytytt0}
\y_\theta(t) - \y_\theta(t\mid t_0) 
=  \sum_{k = t - t_0}^{\infty} H_\theta(k)\, \w(t-1-k).
\end{align}
so $\y_\theta(t)-\y_\theta(t\,|\,0)\to 0$ in mean square as $t\to\infty$ by arguments as in \eqref{ytt2-ytt1}. 


Finally, consider
\begin{align}\label{LE}
\left|\bE\big[\ell\big(\y(t),\y_\theta(t)\big)\big]-\bE\big[\ell\big(\y(t),\y_\theta(t\,|\,t_0)\big)\big]\right|
\le \bE\big[\left|\ell\big(\y(t),\y_\theta(t)\big)-\ell\big(\y(t),\y_\theta(t\,|\,t_0)\big)\right|\big]
\end{align}
If $\ell$ is $L_\ell$-Lipschitz, then
the right hand side of \eqref{LE} is bounded by $L_\ell \bE[\|\y_\theta(t) - \y_\theta(t\mid t_0) \|_2]$, 
and hence, as $\y_\theta(t)-\y_\theta(t\,|\,0)\to 0$ in mean square as $t\to\infty$ or 
as $t_0 \to -\infty$, the right hand side of \eqref{LE} also tends to $0$, i.e., 
\[
   \bE\big[\left|\ell\big(\y(t),\y_\theta(t)\big)-\ell\big(\y(t),\y_\theta(t\,|\,t_0)\big)\right|\big] \to 0.
\]
as $t \to \infty$ or $t_0 \to -\infty$.

If $\ell$ is quadratic, then 
% then notice that from the properties of mean-squared convergence and the fact
% that $\y_{\theta}(t\,|\,0)\to \y_\theta(t)$ in mean square as $t_0 \to -\infty$
% it follows that $\y(t) - \y_{\theta}(t,|,t_0) \to \y(t) - \y_{\theta}(t)$ in mean square sense as $t_0 \to -\infty$.
% The latter means convergence in the Hilber-space of squar integrable random variables with
%Mthe norm $\|z\|=\sqrt{\bE[\|z\|_2^2]}$. But if a sequence of elements has a limit in a Hilber-space, then 
%the sequence of the square norms of those elements converges to the square norm of the limit, as the norm is a continuous function.
%Note that $\bE[\ell(\y(t),\y_\theta(t))] = \|\y(t) - \y_\theta(t)\|^2$ and $\bE[\ell(\y(t),\y_\theta(t\,|\,t_0))] = \|\y(t) - \y_\theta(t\,|\,t_0)\|^2$, and hence $\bE[\ell(\y(t),\y_\theta(t\,|\,t_0))] \to \bE[\ell(\y(t),\y_\theta(t))]$ as $t_0 \to -\infty$.
%Finally, 
 %it follows that %$\y(t)-\y_\theta(t\,|\,0)\to 0$ in mean square as $t\to\infty$.
\begin{align*}
  & |\ell\big(\y(t),\y_\theta(t)\big)-\ell\big(\y(t),\y_\theta(t\,|\,t_0)\big)|=
    \left|\|\y(t)-\y_\theta(t)\|_2^2 - \|\y(t)-\y_\theta(t\,|\,t_0)\|_2^2\right|= \\
  & |2\y^T(t)(-\y_\theta(t)+\y_\theta(t\,|\,t_0))+\left(\|\y_\theta(t)\|_2^2-\|\y_\theta(t\,|\,t_0)\|_2^2\right)| \le \\
  & 2\|\y^T(t)\|_2 \|\y_\theta(t)-\y_\theta(t\,|\,t_0)\|_2 + \left|\|\y_\theta(t)\|_2-\|\y_\theta(t\,|\,t_0)\|_2\right| \left(\|\y_\theta(t)\|_2+\|\y_\theta(t\,|\,t_0)\|_2\right)\\
\end{align*}
Hence, using the Cauchy-Schwartz  inequality
\begin{align*} 
%& \left|\bE\left[\|\y(t)-\y_\theta(t)\|_2\right] - \bE\left[|\y(t)-\y_\theta(t\,|\,t_0)\|_2^2\right]\right|  \le 
%& \bE\left[|\|\y(t)-\y_\theta(t)\|_2^2 - \|\y(t)-\y_\theta(t\,|\,t_0)\|_2^2|\right]  \le  \\
& \bE\left[|\ell\big(\y(t),\y_\theta(t)\big)-\ell\big(\y(t),\y_\theta(t\,|\,t_0)\big)|\right]  \le  \\
%& 2\bE\left[\|\y^T(t)\|_2 \|\y_\theta(t)-\y_\theta(t\,|\,t_0)\|_2\right] 
%+ \bE\left[\left|\|\y_\theta(t)\|_2-\|\y_\theta(t\,|\,t_0)\|_2\right|\right] = \\
& 2\bE\left[\|\y^T(t)\|_2 \|\y_\theta(t)-\y_\theta(t\,|\,t_0)\|_2\right] 
+ \bE\left[\left|\|\y_\theta(t)\|_2-\|\y_\theta(t\,|\,t_0)\|_2\right|\big(\|\y_\theta(t)\|_2+\|\y_\theta(t\,|\,t_0)\|_2\big)\right] \le \\
%  & \bE\left[\|\y^T(t)\|_2 \|\y_\theta(t)-\y_\theta(t\,|\,t_0)\|_2\right] 
%+ \bE\left[\|\y_\theta(t)-\y_\theta(t\,|\,t_0)\|_2 \big(\|\y_\theta(t)\|_2+\|\y_\theta(t\,|\,t_0)\|_2\big)\right] \le \\
& 
2\sqrt{\bE\left[\|\y_\theta(t) - \y_\theta(t\mid t_0) \|_2^2\right]}
\sqrt{\bE\left[\|\y(t)\|_2^2\right]}+ \\
&
\sqrt{\bE\left[\|\y_\theta(t)-\y_\theta(t\,|\,t_0)\|_2^2)\right]}
\sqrt{(\bE[\|\y_\theta(t)\|_2^2] +2\bE[\|\y_\theta(t)\|_2\|\y_\theta(t\,|\,t_0)\|_2]+\bE[\|\y_\theta(t\,|\,t_0)\|_2^2]\big)}
\end{align*}
Note that $\bE\left[\|\y_\theta(t) - \y_\theta(t\mid t_0) \|_2^2\right] \rightarrow 0$ as $t_0 \rightarrow -\infty$
or  $t \rightarrow \infty$, since $\y_{\theta}(t\mid t_0)- \y_{\theta}(t \mid t_0) \rightarrow 0$ in the mean-square
sense as $t_0 \to -\infty$ or $t \to \infty$. Moreover, 
\[ 2\bE[\|\y_\theta(t)\|_2\|\y_\theta(t\,|\,t_0)\|_2] \le  \bE[\|\y_\theta(t)\|_2^2] + \bE[\|\y_\theta(t\,|\,t_0)\|_2^2], \]
and  $\bE[\|\y_\theta(t)\|_2^2]$ is constant due to stationarity of $\y_{\theta}(t)$. 
Now for sufficiently large $t > T_{\epsilon}$ or $-t_0 > T_{\epsilon}$, $\bE\left[\|\y_\theta(t) - \y_\theta(t\mid t_0) \|_2^2\right] < \epsilon^2$, and it follows that
\begin{align*} 
\bE[\|\y_{\theta}(t \mid t_0)\|_2^2] &\le \bE[(\|\y_\theta(t)\|_2+\|\y_{\theta}(t\mid t_0)-\y_{\theta}(t)\|_2)^2] \\
&\le 2(\bE[\|\y_\theta(t)\|_2^2] + \bE[\|\y_\theta(t) - \y_\theta(t\mid t_0) \|_2^2]) < 2\bE[\|\y_\theta(t)\|_2^2] + 2\epsilon^2 
\end{align*}
Hence, it follows that for $t > T_{\epsilon}$ or $-t_0 > T_{\epsilon}$,
\[ \sqrt{(\bE[\|\y_\theta(t)\|_2^2] +2\bE[\|\y_\theta(t)\|_2]\|\y_\theta(t\,|\,t_0)\|_2]+\bE[\|\y_\theta(t\,|\,t_0)\|_2^2]\big)}
\le \sqrt{6\bE[\|\y_\theta(t)\|_2^2] + 4\epsilon^2} 
\]
and 
$\sqrt{6\bE[\|\y_\theta(t)\|_2^2] + 4\epsilon^2}=C$ does not depend on $t$.
Hence, if $t > T_{\epsilon}$ or $-t_0 > T_{\epsilon}$, then
\begin{align*} 
\bE\big[\left|\ell\big(\y(t),\y_\theta(t)\big)-\ell\big(\y(t),\y_\theta(t\,|\,t_0)\big)\right|\big]
\leq \epsilon (2\bE[\|\y(t)\|_2^2]+C)
% \\
% =
% \bE\left[\left|\|\y(t)-\y_\theta(t)\|_2^2 - \|\y(t)-\y_\theta(t\,|\,t_0)\|_2^2\right|\right]   \le 
%    \epsilon (\bE[\|\y(t)\|_2^2]+C)
\end{align*}
from which  it follows that 
$\bE\big[\left|\ell\big(\y(t),\y_\theta(t)\big)-\ell\big(\y(t),\y_\theta(t\,|\,t_0)\big)\right|\big] \to 0$
as either $t_0 \to -\infty$ or $t \to \infty$.
%and by $E\|\y_\theta(t) - \y_\theta(t\mid t_0) \|_2$ if $\ell$ is quadratic. 
%In any case, arguments as in \eqref{ytytt0} show that 


That is, in all cases, 
\begin{align}\label{Lcon}\displaystyle 
\lim_{t\to\infty} E\big[\ell\big(\y(t),\y_\theta(t\,|\,0)\big)\big]
= E\big[\ell\big(\y(t),\y_\theta(t)\big)\big]%=\mathcal L(\theta)
=\lim_{t_0\to -\infty} E\big[\ell\big(\y(t),\y_\theta(t\,|\,t_0)\big)\big]
\end{align}
proving the lemma since $\mathcal L(\theta)$ is the right hand side of \eqref{Lcon} by definition.
\end{proof}

\section{Proofs of Section \ref{sec:asm}}
\label{app:constants}

\begin{remark}[$\alpha$ as Markov-parameters]
    \label{norms:rem}
   If $\alpha$ is the sequence of Markov parameters of a stable deterministic LTI system  determined
  by matrices $(A,B,C,D)$, i.e.,
  $\alpha(0)=D$, $\alpha(k)=CA^{k-1}B$, $k \ge 1$, then the $\ell_1$ norm $\|\alpha\|_{\ell_1}$ and the WDMC $\theta_{\infty}(\alpha)$ are both finite. In this case,
  $\|\alpha\|_{\ell_1}$ is the classical $\ell_1$ norm of the underlying LTI system.  
  Moreover, if $M >0$, $\gamma \in [0,1)$ are such that
  $\|A^k\|_2 < M \gamma^k$, then 
  $\|\alpha\|_{\ell_1} \le \|D\|_2+M\|B\|_2\|C\|_2/(1-\gamma)$ and 
  $\theta_{\infty}(\alpha) \le \|D\|+ \gamma M\|B\|_2\|C\|_2/(1-\gamma)^2$.
  In turn, such $M$
  always exist if $A$ is a Schur matrix, and $\gamma$ can be chosen to be 
  the spectral radius of $A$. In particular, $\|\alpha\|_{\ell_1}$, $\theta_{\infty}(\alpha)$
  decrease as the spectral radius of $A$ decreases.
  \end{remark}
  Finally, as promised before Lemma \ref{alpha:lemma2}, we can state that all the
  constants from Definition \ref{def:constants} are well-defined. 
\begin{lemma}
\label{alpha:lemma1.5}
  If Assumption \ref{as:generator} and \ref{as:parameterisation} hold, then
  the quantities $\|\alpha_{\theta}\|_{\ell_1}$ and $\theta_{\infty}(\alpha_{\theta})$, $\|\bar{\alpha}\|_{\ell_1}$,
  $\theta_{\infty}(\bar{\alpha})$, $G_e(\theta)$, $G_{e,1}(\theta)$, $G_e$, $G_{e,1}$ are well defined 
  real numbers. 
\end{lemma}
\begin{proof}
  From Remark \ref{norms:rem} and the assumption that $A_g$ is Schur it follows that 
  $\|\alpha_g\|_{\ell_1}$ and $\theta_{\infty}(\alpha_g)$ are well-defined. 
   By Assumption \ref{as:parameterisation} there exists constants $M,\gamma,C$ such that 
   for all $\theta \in \Theta$, $\|A_{\theta}^k\|_2 \le M \gamma^k$, $\|B_{\theta}\|_2 \le C$, and $\|C_{\theta}\|_2 \le C$. 
   Therefore, $\bar{\alpha}(k)=\sup \|\alpha_{\theta}(k)\|_2 \le M \gamma^k C^2$ and hence,
   by using the argument of Remark \ref{norms:rem},  
   $\|\alpha_{\theta}\|_{\ell_1} \le \sum_{k=0}^{\infty} \|\alpha\|_{\ell_1} \le  \sum_{k=0}^{\infty} \bar{\alpha}(k)=\|\bar{\alpha}\|_{\ell_1} \le \frac{MC^2}{1-\gamma} < \infty$ and
   $\theta_{\infty}(\alpha_{\theta})=\sum_{k=0}^{\infty} k \|\alpha_{\theta}(k)\|_2 \le \sum_{k=0}^{\infty} k \bar{\alpha}(k)=\theta_{\infty}(\bar{\alpha}) \le \frac{MC^2}{(1-\gamma)^2} < \infty$.
   Therefore, the quantities $\|\alpha_{\theta}\|_{\ell_1}$ and $\theta_{\infty}(\alpha_{\theta})$, $\|\bar{\alpha}\|_{\ell_1}$,
   $\theta_{\infty}(\bar{\alpha})$ defining $G_e(\theta)$, $G_{e,1}(\theta)$, $G_e$, $G_{e,1}$
   are well-defined, and hence so are $G_e(\theta)$, $G_{e,1}(\theta)$, $G_e$, $G_{e,1}$
\end{proof}
\begin{lemma}
\label{alpha:lemma2}
%\[
%\begin{aligned*}
%&  \alpha_{w}(k)= \left\{ \begin{array}{rl} \alpha_{g}(k) & \rvw = \begin{bmatrix} \rvy \rvu \end{bmatrix} \\
%            \begin{bmatrix} 0 & I_{n_{\rvu}} \end{bmatrix} \alpha_{g}(k) & \rvw=\rvu 
 %   \end{array}\right. \\
 %   & \alpha_{y}(k)=\begin{bmatrix} I_{n_{\rvy}} & 0 \end{bmatrix} \alpha_{g}(k)  \\
 % \end{aligned*}
%\]\left(\sum_{l=0}^{\min\{k,t-t_0\}} \alpha_{\theta}(k)\alpha_w(k-l)\right)
Let $\alpha_{g,y}(k)$ and $\alpha_{g,w}(k)$ be the matrix formed by the first $n_\rvy$ and the last $n_\rvw$ 
rows of $\alpha_g(k)$ respectively, where $n_{\rvw}$ is the number of entries of $\w(t)$. Let $c_{\theta,s}(k)=\left(\sum_{l=0}^{\min\{k,s\}} \alpha_{\theta}(k)\alpha_{g,w}(k-l)\right)$
and $c_{\theta}(k)=\sum_{l=0}^{k} \alpha_{\theta}(k)\alpha_{g,w}(k-l)$. 
\begin{equation}
\label{alpha:lemma:eq1}
\begin{aligned}
%& \hat{\rvy}_{\theta}(t)=(c_{\theta} \star \e_g)(t), ~~
%\hat{\rvy}_{\theta}(t\mid t_0)=(c_{\theta, t-t_0} \star \e_g)(t), ~~  \rvy(t)=(\alpha_{g,y} \star \e_g)(t), ~~ 
& \begin{bmatrix} \y(t)
    \\ \hat{\rvy}_{\theta}(t)
\end{bmatrix}=\left(\begin{bmatrix} \alpha_{g,y} \\ c_{\theta} \end{bmatrix} \star \e_g\right)(t), \quad  
\begin{bmatrix} \y(t) \\ \hat{\rvy}_{\theta}(t \mid t_0) \end{bmatrix}=\left(\begin{bmatrix} \alpha_{g,y} \\ c_{\theta,t-t_0} \end{bmatrix} \star \e_g\right)(t),   \\
& \y(t)-\hat{\rvy}_{\theta}(t) = ((\alpha_{g,y} - c_{\theta}) \star \e_g)(t), \quad
\y(t)-\hat{\rvy}_{\theta}(t \mid t_0) = ((\alpha_{g,y} - c_{\theta,t-t_0}) \star \e_g)(t)
 % \begin{bmatrix}
 %    \hat{\y}_\theta(t) \\                %\sum_{k=0}^{\infty} c_{\theta,k}(k)  \e_g(t-k), ~ 
 %   \y(t)-\hat{\y}_\theta(t \mid t_0)
 % \end{bmatrix} = \sum_{k=0}^{\infty}  \begin{bmatrix} c_{\theta,k}(k) \\ c_{\theta,t-t_0}  \end{bmatrix} \e_g(t-k) \\
%\widetilde{\alpha}_{k,e,t-t_0} 
%\left(\sum_{l=0}^{\min\{k,t-t_0\}} \alpha_{\theta}(k)\alpha_w(k-l)\right)
% \e_g(t-k) \\ 
 %& \begin{bmatrix} \rvy(t) \\ \rvw(t) \end{bmatrix}=\sum_{k=0}^{\infty} \begin{bmatrix} \alpha_{g,y}(k) \\ \alpha_{g,w}(k) \end{bmatrix} \e_g(t-k).
%&\y(t)-\hat{\y}_\theta(t) =\sum_{k=0}^{\infty} \alpha_{k,e} \e_g(t-k), ~
%&\y(t)-\hat{\y}_\theta(tt \mid t_0)=\sum_{k=0}^{\infty} \alpha_{k,e,t-t_0} \e_g(t-k) \\ 
%& \y(t) - \hat{\y}_\theta(t) =\sum_{k=0}^{\infty} \begin{bmatrix} I_{n_\rvy} & - I_{n_{\rvy}} \end{bmatrix} \alpha_{k,e} \e_g(t-k), ~
%&\y(t) - \hat{\y}_\theta(t) =\sum_{k=0}^{\infty} \begin{bmatrix} I_{n_\rvy} & - I_{n_{\rvy}} \end{bmatrix} \alpha_{k,e,t-t_0} \e_g(t-k), ~
\end{aligned}
\end{equation}
Moreover,  the constants from Definition \ref{def:constants} satisfy:
%\begin{equation}
%\label{alpha:lemma:eq2}
\begin{align}
& \max\{\|c_{\theta}\|_{\ell_1}, \|c_{\theta,s}\|_{\ell_1}\} \le \|\alpha_{\theta}\|_{\ell_1}\|\alpha_g\|_{\ell_1}, \label{alpha:lemma:eq2a} \\ 
 & \max\{\left\|\begin{bmatrix} \alpha_{g,y} \\ c_{\theta} \end{bmatrix}\right \|_{\ell_1}, \left \|\begin{bmatrix} \alpha_{g,y} \\ c_{\theta,t-t_0} \end{bmatrix}\right\|_{\ell_1}, \|\alpha_{g,y}-c_{\theta}\|_{\ell_1}, \|\alpha_{g,y}-c_{\theta,s}\|_{\ell_1}\} \le G_e(\theta) \le G_e \label{alpha:lemma:eq2b} \\
& \max\{
  \theta_{\infty}\left(\begin{bmatrix} \alpha_{g,y} \\ c_{\theta} \end{bmatrix}\right),
  \theta_{\infty}\left(\begin{bmatrix} \alpha_{g,y} \\ c_{\theta,t-t_0} \end{bmatrix}\right),
\theta_{\infty}(\alpha_{g,y}-c_{\theta}), \theta_{\infty}(\alpha_{g,y}-c_{\theta,s})\} \le G_{e,1}(\theta) \label{alpha:lemma:eq2c} \\
& \sum_{t=t_0}^{N+t_0}\|c_{\theta,t-t_0}-c_{\theta}\|_{\ell_1} \le \theta_{\infty}(\alpha_{\theta})\|\alpha_g\|_{\ell_1} \label{alpha:lemma:eq2d}
  \end{align}
%\end{equation}
%$\max\{\theta_{\infty}(c_{\theta}), \theta_{\infty}(\alpha_{g,y}-c_{\theta,s})\} \le G_{e,1}(\theta)$
%$\sum_{k=0}^{\infty} \|c_{\theta,k}(k)\|_2 \le G_f(\theta)G_g$
%$\sum_{k=0}^{\infty} \|c_{\theta,t-t_0}(k)\|_2 \le G_f(\theta)G_g$,
%$\sum_{k=0}^{\infty} \|\alpha_{g,y}-c_{\theta,k}(k)\|_2 \le G_e(\theta)$,
%$\sum_{k=0}^{\infty} (k+1) \|\alpha_{g,y}-c_{\theta,k}(k)\|_2 \le G_{e,1}(\theta)$,
%\[ 
%  \begin{aligned}
  % & \max\{\bE[\|\hat{\y}_{\theta}(t)\|_2],\bE[\|\hat{\y}_{\theta}(t\mid t_0)\|_2]\} \le \bar{G}_f(\theta) G_e \bE[\|\e_g(t)\|], \\
   %&  \le \bar{G}_f(\theta) G_g \bE[\|\e_g(t)\|], ~ \\
   %\bE[\|\y(t)-\hat{\y}_{\theta}(t)\|_2] \le G_e(\theta) \bE[\|\e_g(t)\|], \\
   %& \bE[\|\y(t)-\hat{\y}_{\theta}(t \mid t_0)\|_2] \le G_e(\theta)G_g \bE[\|\e_g(t)\|], 
%  \end{aligned}
%\]
%\[
%  \begin{aligned}
%   & \sum_{k=0}^{\infty} \|\widetilde{\alpha}_{k,e}\|_2 \le  
%   \sum_{k=0}^{\infty} \|\alpha_{k}-\theta)\|_2\|\alpha_{k,w}\|_2 \\
%    & \bar{G}_{f,1}(\theta)\|\Sigma_{gen}\|_{\ell_1}, ~
%   \sum_{k=0}^{\infty} \|\widetilde{\alpha}_{k,e,t-t_0}\| \le
%    \bar{G}_{f,1}(\theta)\|\Sigma_{gen}\|_{\ell_1}
%  \end{aligned}
%\]
In particular, the two LTI filters in \eqref{alpha:lemma:eq1} are stationary. %\cite[Section 5.6]{gikh96}.
\end{lemma}
\begin{proof}[Proof of Lemma  \ref{alpha:lemma2}]
  Let $C_w$ denote the matrix formed by the last $n_w$ rows of $C_g$. From Assumption \ref{as:generator} and the proof of Lemma~\ref{l:ihp} we get that
  \begin{align*}
  \w(t)&=\sum_{k=1}^{\infty} C_w A_g^{k-1} K_g \e_g(t-k)+ \e_g(t)=\sum_{k=0}^{\infty} \alpha_{g,w}(k)\e_g(t-k)
  =(\alpha_{g,w}\star\e_g)(t)\\
  \y(t)&=\sum_{k=0}^{\infty} \alpha_{g,y}(k)\e_g(t-k)\\
  \hat{\y}_{\theta}(t \mid t_0)&=\sum_{k=1}^{t-t_0} \hat{C}_{\theta}\hat{A}_{\theta}^{k-1}\hat{B}_{\theta}\w(t-k)+\hat{D}_{\theta}\w(t)=\sum_{k=0}^{t-t_0} \alpha_{\theta}(k)\w(t-k)\\ 
  \hat{\y}_{\theta}(t)&=\lim_{t_0 \rightarrow -\infty} \hat{\y}_{\theta}(t \mid t_0)=\sum_{k=0}^{\infty} \alpha_{\theta}(k)\w(t-k) 
  \end{align*}
  It then follows using the new index $n=k+j\geq 0$ and the standard argument involving Cauchy products of series, 
  % From Assumption \ref{as:generator} it follows that
  % $\w(t)=\sum_{k=1}^{\infty} C_w A_g^{k-1} K_g \e_g(t-k)+ \e_g(t)=\sum_{k=0}^{\infty} \alpha_w(k)\e_g(t-k)$,
  % and $\y(t)=\sum_{k=0}^{\infty} \alpha_y(k)\e_g(t-k)$.
  % From the definition of a predictor \eqref{eq:predictor} it follows that
  % $\hat{\y}_{\theta}(t \mid t_0)=\sum_{k=1}^{t-t_0} \hat{C}_{\theta}\hat{A}_{\theta}^{k-1}\hat{B}_{\theta}\w(t-k)+\hat{D}_{\theta}\w(t)=\sum_{k=0}^{t-t_0} \alpha_{\theta}(k)\w(t-k)$. Since by Lemma \ref{l:ihp},
  % $\hat{\y}_{\theta}(t)=\lim_{t_0 \rightarrow -\infty} \hat{\y}_{\theta}(t \mid t_0)$, it follows that
  % \(
  %    \hat{\y}_{\theta}(t)=\sum_{k=0}^{\infty} \alpha_{\theta}(k)\w(t-k)
  % \) and the convergence of the right-hand side is in the mean square sense. 
  % From $\w(t)=\sum_{k=0}^{\infty} \alpha_w(k)\e_g(t-k)$ and $\hat{\y}_{\theta}(t \mid t_0)=\sum_{k=0}^{t-t_0} \alpha_{\theta}(k)\w(t-k)$ it then follows using the standard argurment involving Cauchy products of series, 
  \begin{align*}
     \hat{\y}_{\theta}(t \mid t_0)
     &=\sum_{k=0}^{t-t_0} \alpha_{\theta}(k)\sum_{j=0}^{\infty} \alpha_{g,w}(j)\e_g(t-k-j)\\
     % &=
     %   \sum_{k=0}^{\infty} \left(\sum_{i,l=0, l+i=k, l+i=k, i \le t-t_0}^{\infty} \alpha_{\theta}(i)\alpha_{g,w}(l)\right)\e_g(t-k)\\
      &= \sum_{n=0}^{\infty } \left(\sum_{k=0}^{\min(n,t-t_0)} \alpha_{\theta}(k)\alpha_{g,w}(n-k)\right)\e_g(t-n)\\
      &=\sum_{n=0}^{\infty} c_{\theta, t-t_0}(n)\e_g(t-n)=(c_{\theta, t-t_0} \star \e_g)(t)
  \end{align*}
  Notice that since $A_g$ and $A_{\theta}$ are Schur, the series 
  $\sum_{k=0}^{\infty} \alpha_{\theta}(k)$ and $\sum_{k=0}^{\infty} \alpha_{g,w}(k)$ are absolutely convergent. 
  With a similar argument, using that absolutely convergent series can be rearranged, it follows that
  \begin{align*}
     \hat{\y}_{\theta}(t)
     &=\sum_{k=0}^{\infty} \alpha_{\theta}(k)\sum_{j=0}^{\infty} \alpha_{g,w}(j)\e_g(t-k-j)\\
     % =
     %   \sum_{k=0}^{\infty} \left(\sum_{i,l=0, l+i=k}^{\infty} \alpha_{\theta}(i)\alpha_{g,w}(l)\right)\e_g(t-k)= \\
      &= \sum_{n=0}^{\infty } \left(\sum_{k=0}^{n} \alpha_{\theta}(k)\alpha_{g,w}(n-k)\right)\e_g(t-n)\\
      &=
      \sum_{n=0}^{\infty} c_{\theta}(n)\e_g(t-n)=(c_{\theta} \star \e_g)(t)
  \end{align*}
  From this the first  claim of the lemma, i.e., \eqref{alpha:lemma:eq1}, follows.
  The first inequality of the second claim, i.e., 
  \eqref{alpha:lemma:eq2a}, follows
  by first noticing that 
  \begin{align*}
    & \|c_{\theta}(k) \|_2 \le \sum_{i=0}^{k} \|\alpha_{\theta}(i)\|_2 \|\alpha_{g,w}(k-i)\|_2, \\
     & \|c_{\theta,t-t_0}(k) \|_2 \le \sum_{i=0}^{\min(k,t-t_0)} \|\alpha_{\theta}(i)\|_2 \|\alpha_{g,w}(k-i)\|_2 
     \le \sum_{i=0}^{k} \|\alpha_{\theta}(i)\|_2 \|\alpha_{g,w}(k-i)\|_2
  \end{align*}
  %and by using the well-known property of the Cauchy-product of two series and noticing that 
  Since $\sum_{i=0}^{k} \|\alpha_{\theta}(i)\|_2 \|\alpha_{g,w}(k-i)\|_2$ is the  $k$'th coefficient of the Cauchy-product of 
  $\sum_{k=0}^{\infty} \|\alpha_{\theta}(k)\|_2$ and $\sum_{k=0}^{\infty} \|\alpha_{g,w}(k)\|_2$, we obtain  \eqref{alpha:lemma:eq2a} by  
    \begin{align}\label{alpha:lemma2:pf2}
      \|c\|_{\ell_1}=\sum_{k=0}^{\infty}\|c(k)\|_2
  \le \sum_{i=0}^{\infty} \|\alpha_{\theta}(i)\|_2 \sum_{j=0}^{\infty} \|\alpha_{g,w}(j)\|_2    
      =\|\alpha_{\theta}\|_{\ell_1} \|\alpha_{g,w}\|_{\ell_1}
       \leq\|\alpha_{\theta}\|_{\ell_1} \|\alpha_{g}\|_{\ell_1}
  \end{align}
  % \begin{equation}
  %   \label{alpha:lemma2:pf2}
  %   \begin{aligned}
  %     & \sum_{k=0}^{\infty} \|c_{\theta}(k)\|_2 \le \sum_{i=0}^{\infty} \|\alpha_{\theta}(i)\|_2 \sum_{j=0}^{\infty} \|\alpha_{g,w}(j)\|_2
  %     = \\
  %     & (\sum_{i=0}^{\infty} \|\alpha_{\theta}(i)\|_2)(\sum_{j=0}^{\infty} \|\alpha_{g,w}(j)\|_2)=\|\alpha_{\theta}\|_{\ell_1} \|\alpha_{g,w}\|_{\ell_1}
  % \end{aligned}
  % \end{equation}
 with $c$ denoting either $c_{\theta}$ or $c_{\theta,t-t_0}$. Concerning the second inequality \eqref{alpha:lemma:eq2b}, notice that
  \begin{subequations}\label{alpha:lemma2:pf0}    
  \begin{align}
    & \!\!\!\!\!\left\|\begin{bmatrix} \alpha_{g,y}(k) \\ c_{\theta}(k) \end{bmatrix}\right\|_2 \le \|\alpha_{g,y}(k)\|_2 + \|c_{\theta}(k)\|_2, \quad
    \|\alpha_{g,y}(k) - c_{\theta}(k)\|_2 \le \|\alpha_{g,y}(k)\|_2 + \|c_{\theta}(k)\|_2 \\
   & \left\|\begin{bmatrix} \alpha_{g,y}(k) \\ c_{\theta,t-t_0}(k) \end{bmatrix}\right\|_2 \le \|\alpha_{g,y}(k)\|_2 + \|c_{\theta,t-t_0}(k)\|_2,\\
    &\|\alpha_{g,y}(k) - c_{\theta,t-t_0}(k)\|_2 \le \|\alpha_{g,y}(k)\|_2 + \|c_{\theta,t-t_0}(k)\|_2 
  \end{align}
  \end{subequations}
  and hence,  
\begin{subequations}\label{alpha:lemma2:pf1}    
  \begin{align}
  &  \max\left\{\left\|\begin{bmatrix} \alpha_{g,y} \\ c_{\theta} \end{bmatrix}\right\|_{\ell_1}, \|\alpha_{g,y}-c_{\theta}\|_{\ell_1}\right\}
    \le \|\alpha_{g,y}\|_{\ell_1} + \|c_{\theta}\|_{\ell_1} \\
   & \max\left\{\left\|\begin{bmatrix} \alpha_{g,y} \\ c_{\theta,t-t_0} \end{bmatrix}\right\|_{\ell_1}, \|\alpha_{g,y}-c_{\theta,t-t_0}\|_{\ell_1}\right\}
    \le \|\alpha_{g,y}\|_{\ell_1} + \|c_{\theta,t-t_0}\|_{\ell_1}
  \end{align}
\end{subequations}
  From \eqref{alpha:lemma2:pf1} and \eqref{alpha:lemma2:pf2}, the inequality of \eqref{alpha:lemma:eq2b} follows
  by using the definition of $G_{e}(\theta)$ and the fact that $\max\{\|\alpha_{g,y}\|_{\ell_1}, \|\alpha_{g,w}\|_{\ell_1}\} \le \|\alpha_{g}\|_{\ell_1}$.

 Concerning \eqref{alpha:lemma:eq2c}, notice that by definining $d_1(i)=i \|\alpha_{\theta}(i)\|_2$ and $d_2(i)=i \|\alpha_{g,w}(i)\|_2$,
  \begin{equation*}
    %\label{alpha:lemma2:pf3}
  \begin{aligned}
     & \theta_{\infty}(c_{\theta,t-t_0})=\sum_{k=0}^{\infty} k \|c_{\theta,t-t_0}(k)\|_2 \le \sum_{k=0}^{\infty} k \sum_{i=0}^{\min(k,t-t_0)} \|\alpha_{\theta}(i)\|_2 \|\alpha_{g,w}(k-i)\|_2 = \\
     & \sum_{k=0}^{\infty} \sum_{i=0}^{\min(k,t-t_0)} i \|\alpha_{\theta}(i)\|_2 \|\alpha_{g,w}(k-i)\|_2  +  \sum_{k=0}^{\infty} \sum_{i=0}^{\min(k,t-t_0)} (k-i) \|\alpha_{\theta}(i)\|_2 \|\alpha_{g,w}(k-i)\|_2 =  \\
     & \sum_{k=0}^{\infty} \sum_{i=0}^{\min(k,t-t_0)} d_1(i) \|\alpha_{g,w}(k-i)\|_2 
     + \sum_{k=0}^{\infty} \sum_{i=0}^{\min(k,t-t_0)} d_2(k-i) \|\alpha_{\theta}(i)\|_2 \le \\
     & \sum_{k=0}^{\infty} \sum_{i=0}^{k} d_1(i) \|\alpha_{g,w}(k-i)\|_2 
     + \sum_{k=0}^{\infty} \sum_{i=0}^{k} d_2(k-i) \|\alpha_{\theta}(i)\|_2 
  \end{aligned}
\end{equation*}
Similarly,
\begin{equation*}
\begin{aligned}
  & \theta_{\infty}(c_{\theta})=\sum_{k=0}^{\infty} k \|c_{\theta}(k)\|_2 \le 
   \sum_{k=0}^{\infty} k \sum_{i=0}^{k} \|\alpha_{\theta}(i)\|_2 \|\alpha_{g,w}(k-i)\|_2 = \\
     & \sum_{k=0}^{\infty} \sum_{i=0}^{k} i \|\alpha_{\theta}(i)\|_2 \|\alpha_{g,w}(k-i)\|_2  +  \sum_{k=0}^{\infty} \sum_{i=0}^{k} (k-i) \|\alpha_{\theta}(i)\|_2 \|\alpha_{g,w}(k-i)\|_2 = \\
     & \sum_{k=0}^{\infty} \sum_{i=0}^{k} d_1(i) \|\alpha_{g,w}(k-i)\|_2 
     + \sum_{k=0}^{\infty} \sum_{i=0}^{k} d_2(k-i) \|\alpha_{\theta}(i)\|_2 \le \\
     & \left(\sum_{k=0}^{\infty} d_1(k)\right)
        \left(\sum_{k=0}^{\infty} \|\alpha_{g,w}(k)\|_2 \right)+
      \left(\sum_{k=0}^{\infty} d_2(k)\right)
        \left(\sum_{k=0}^{\infty} \|\alpha_{\theta}(k)\|_2 \right) = \\
     &   \theta_{\infty}(\alpha_{\theta})\|\alpha_{g,w}\|_{\ell_1}+
        \theta_{\infty}(\alpha_{g,w})\|\alpha_{\theta}\|_{\ell_1} \le  \theta_{\infty}(\alpha_{\theta})\|\alpha_{g}\|_{\ell_1}+
        \theta_{\infty}(\alpha_{g})\|\alpha_{\theta}\|_{\ell_1}
  \end{aligned}
\end{equation*}
In particular, 
\begin{equation}
    \label{alpha:lemma2:pf4}
\max\{\theta_{\infty}(c_{\theta}), \theta_{\infty}(c_{\theta,t-t_0})\} \le \theta_{\infty}(\alpha_{\theta})\|\alpha_g\|_{\ell_1} +\theta_{\infty}(\alpha_{g})\|\alpha_{\theta}\|_{\ell_1}
\end{equation}
Moreover, using \eqref{alpha:lemma2:pf0} it follows that 
\begin{align*}
     & \max\left\{\theta_{\infty}\left(\begin{bmatrix} \alpha_{g,y} \\ c_{\theta} \end{bmatrix}\right),
  \theta_{\infty}\left(\begin{bmatrix} \alpha_{g,y} \\ c_{\theta,t-t_0} \end{bmatrix}\right),
\theta_{\infty}(\alpha_{g,y}-c_{\theta}), \theta_{\infty}(\alpha_{g,y}-c_{\theta,t-t_0
})\right\}  \le \\
& \theta_{\infty}(\alpha_{g,y})+\max\{\theta_{\infty}(c_{\theta}), \theta_{\infty}(c_{\theta,t-t_0})\}
\end{align*}
from which inequality \eqref{alpha:lemma:eq2c} follows by using \eqref{alpha:lemma2:pf4}. Finally, notice that 
\[
   c_{\theta}(k)-c_{\theta,t-t_0}(k)=\left\{
    \begin{array}{ll}
       0 & k \le t-t_0 \\
       \sum_{i=t-t_0+1}^{k} \alpha_{\theta}(i)\alpha_{g,w}(k-i) & k > t-t_0
   \end{array}\right.
\] 
Hence, by setting $d_3(k)=\|\alpha_{\theta}(k+t-t_0+1)\|_2$
  \begin{align*}
        \|c_{\theta}-c_{\theta,t-t_0}\|_{\ell_1}
       &=
      \sum_{k=t-t_0+1}^{\infty} \|\sum_{i=t-t_0+1}^{k} \alpha_{\theta}(i)\alpha_{g,w}(k-i)\|_2\\
       &\le  \sum_{k=t-t_0+1}^{\infty} \sum_{i=t-t_0}^{k} \|\alpha_{\theta}(i)\|_2 \|\alpha_{g,w}(k-i)\|_2\\
       &= \sum_{k=t-t_0+1}^{\infty} \sum_{i=0}^{k-(t-t_0+1)} d_3(i)
       \|\alpha_{g,w}(k-i-(t-t_0+1))\|_2 \\
       &=\sum_{k=0}^{\infty}\sum_{i=0}^{k} d_3(i)
       \|\alpha_{g,w}(k-i)\|_2=
       \left(\sum_{k=t-t_0+1}^{\infty} \|\alpha_{\theta}(k)\|_2\right) \cdot
       \|\alpha_{g,w}\|_{\ell_1}
  \end{align*}
     By noticing that 
     \[
       \begin{aligned}
        & \sum_{t=t_0}^{N+t_0} \sum_{k=t-t_0+1}^{\infty} \|\alpha_{\theta}(k)\|_2 \le
        \sum_{k=1}^{\infty} \min\{k,N\} \|\alpha_{\theta}(k)\|_2
        \le 
        \sum_{k=1}^{\infty} k \|\alpha_{\theta}(k)\|_2=\theta_{\infty}(\alpha_{\theta})
       \end{aligned}
     \]
the last inequality \eqref{alpha:lemma:eq2d} follows. Finally, the stationarity of the two LTI filters is a concequence of the standard result that stable LTI filters maps stationary process to stationary process, see e.g., \cite[Section 5.6]{gikh96}.


%%    Let
%%	$A_g,K_g,C_g$ be the matrices of the data generator  from Assumption \ref{as:generator}. 
%%  \textcolor{cyan}{For every $\theta \in \theta$}, let us define
%%	the matrices $A_{\theta,e},K_{\theta,e},C_{\theta,e},D_{\theta,e},\textcolor{black}{\widetilde{C}_{\theta,e},\widetilde{D}_{\theta,e}}$ as
%%    \begin{color}{red}
%%    \begin{equation}
%%    \label{error-sys-mat}    
%%	\begin{aligned}
%%	   & A_{\theta,e}=
%%	   \begin{bmatrix} A_g & 0 \\ \hat{B}_\theta C_w & \hat{A}_\theta  \end{bmatrix}, 
%%	   &K_{\theta,e}=
%%	    \begin{bmatrix} K_g \\
%%	    \hat{B}_w
%%	    \end{bmatrix},\\
%%         &C_{\theta,e}=
%%	         \begin{bmatrix} C_1-\hat{D}_\theta C_w &  -\hat{C}_\theta  \end{bmatrix}, & D_{\theta,e}= I-\hat{D}_w, \\
%%        & \bar{C}_{\theta,e}=\begin{bmatrix} \hat{D}_\theta C_w & \hat{C}_\theta  \end{bmatrix}, &
%%          \bar{D}_{\theta,e}=\hat{D}_w
%%        \\
%%        & \bar{C}_e = \begin{bmatrix} C_1 & 0 \end{bmatrix}, & D_e=\begin{bmatrix} I_{n_{\y}} &  0 \end{bmatrix} \\
%%        & \widetilde{C}_{\theta,e}=\begin{bmatrix} \bar{C}_e \\ \bar{C}_{\theta,e} \end{bmatrix}, &  
%%        \widetilde{D}_{\theta,e}=\begin{bmatrix} D_e \\ \hat{D}_w \end{bmatrix}
%%	    \end{aligned}
%%    \end{equation}
%%    \end{color}
%%	where $C_g=\begin{bmatrix} C_1^T & C_2^T \end{bmatrix}^T$ and $C_1$ has $n_\y$ rows and $C_2$ has $n_\rvu$ rows; and
%%	\[ 
%%       (C_w,\hat{B}_w,\hat{D}_w)\hspace{-1pt}=\hspace{-1pt}\left \{ \hspace{-5pt}\begin{array}{ll} 
%%         (C_2,\begin{bmatrix} 0 & \hat{B}_\theta  \end{bmatrix}, \begin{bmatrix} 0 & \hat{D}_\theta  \end{bmatrix}) & 
%%	\mbox{if $\w=\mathbf{u}$ }, \\ 
%%	%(C_w,\hat{B}_w,\hat{D}_w)=
%% (C_g,\hat{B}_\theta ,\hat{D}_\theta )  & \hspace{-4pt}\mbox{ if $\w=\begin{bmatrix} \y \\ \mathbf{u}  \end{bmatrix}$} \end{array}\right.
%% \]
%% The matrices $(A_{\theta,e},K_{\theta,e},C_{\theta,e},D_{\theta,e})$ 
%% \textcolor{red}{respectively $(A_{\theta,e},K_{\theta,e},\widetilde{C}_{\theta,e},\widetilde{D}_{\theta,e})$}
%% represent the LTI system driven
%%  by the innovation process $\e_g$ of $(\y,\w)$, output
%%  of which is $\y(t)-\hat{\y}_\theta(t)$ respectively $\hat{\y}_{\theta}(t)$, i.e.,
%%\begin{equation}
%%\label{error-sys1}
%%\begin{split}
%%   \tilde{\x}(t+1)&=A_{\theta,e}\tilde{\x}(t)+K_{\theta,e}\e_g(t), \\
%%   \y(t)-\hat{\y}_\theta(t) &=C_{\theta,e}\tilde{\x}(t)+D_{\theta,e}\e_g(t) \\
%%   \y(t) &=\bar{C}_{e}\tilde{\x}(t)+\bar{D}_{e}\e_g(t) \\
%%   \hat{\y}_\theta(t) &=\bar{C}_{\theta,e}\tilde{\x}(t)+\bar{D}_{\theta,e}\e_g(t) \\
%%   \begin{bmatrix} \hat{\y}_\theta(t) \\ \hat{\y}_\theta(t) \end{bmatrix} &=\widetilde{C}_{\theta,e}\tilde{\x}(t)+\widetilde{D}_{\theta,e}\e_g(t) \\
%%\end{split}  
%%\end{equation}
%%with $\tilde{\x}=[\x^T~\hat{\x}^T]^T$. 
%%\begin{color}{red}
%%Moreover, for $t \ge t_0$, by setting
%%$\tilde{\x}(t \mid t_0)=[\x^T(t_0)~\hat{\x}^T(t\mid t_0)]^T$ and 
%%it holds that
%%\begin{equation}
%%\label{error-sys1t0}
%%\begin{split}
%%   \tilde{\x}(t+1 \mid t_0)&=A_{\theta,e}\tilde{\x}(t \mid t_0)+K_{\theta,e}\e_g(t), \\
%%   \y(t)-\hat{\y}_\theta(t\mid t_0) &=C_{\theta,e}\tilde{\x}(t\mid t_0)+D_{\theta,e}\e_g(t) \\
%%   \y(t) &=\bar{C}_{e}\tilde{\x}(t\mid t_0)+\bar{D}_{e}\e_g(t) \\
%%   \hat{\y}_\theta(t\mid t_0) &=\bar{C}_{\theta,e}\tilde{\x}(t\mid t_0)+\bar{D}_{\theta,e}\e_g(t) \\
%%   \begin{bmatrix} \y(t) \\ \hat{\y}_\theta(t \mid t_0) \end{bmatrix} &=\widetilde{C}_{\theta,e}\tilde{\x}(t \mid t_0)+\widetilde{D}_{\theta,e}\e_g(t) \\
%%\end{split}  
%%\end{equation}
%%\end{color}
%%%Next, choose $\hat{M}(\theta)>1$, such that $\|\hat{A}_\theta^k\|_2\leq \hat{M}(\theta)\hat{\gamma}_\theta^k$ (hence $0<\hat{\gamma}_\theta<1$).
%%%\begin{lemma}
%%%\label{alpha:lemma}
%%From the consruction of $A_{\theta,e},K_{\theta,e},C_{\theta,e},D_{\theta,e},\widetilde{C}_{\theta,e},\widetilde{D}_{\theta,e}$, it follows that by induction on $k$, it follows that
%%\begin{align*}
%%  & c_{\theta}(k)=\left\{\begin{array}{ll}
%%            \bar{C}_{\theta,e}A_{\theta,e}^{k-1}K_{\theta,e} & k \ge 1 \\
%%            \bar{D}_e & k=0 
%%    \end{array}\right. \\
%%  & c_{\theta,t-t_0}(k)=\left\{\begin{array}{ll}
%%            \bar{C}_{\theta,e}A_{\theta,e}^{k-1}K_{\theta,e} & k > t-t_0 \\
%%           \bar{D}_{\theta,e}   & k = t-t_0 \\
%%           D_{\theta}C_w A_g^{k-t-t_0-1}K_g & k < t-t_0 
%%    \end{array}\right., \quad k \ge 1
%%\end{align*}
%%
%%
%%
%%Then 
%%\begin{equation}
%%\label{alpha:lemma:eq1}
%%\begin{aligned}
%%&\hat{\y}_\theta(t)=\sum_{k=0}^{\infty}  \left(\alpha_{\theta} \star \alpha_w\right)(k)\e_g(t-k), ~
%%&\hat{\y}_\theta(t \mid t_0)=\sum_{k=0}^{\infty} \widetilde{\alpha}_{k,e,t-t_0} \e_g(t-k) \\ 
%%&\y(t)-\hat{\y}_\theta(t) =\sum_{k=0}^{\infty} \alpha_{k,e} \e_g(t-k), ~
%%&\y(t)-\hat{\y}_\theta(t \mid t_0)=\sum_{k=0}^{\infty} \alpha_{k,e,t-t_0} \e_g(t-k) \\ 
%%%& \y(t) - \hat{\y}_\theta(t) =\sum_{k=0}^{\infty} \begin{bmatrix} I_{n_\rvy} & - I_{n_{\rvy}} \end{bmatrix} \alpha_{k,e} \e_g(t-k), ~
%%%&\y(t) - \hat{\y}_\theta(t) =\sum_{k=0}^{\infty} \begin{bmatrix} I_{n_\rvy} & - I_{n_{\rvy}} \end{bmatrix} \alpha_{k,e,t-t_0} \e_g(t-k), ~
%%\end{aligned}
%%\end{equation}
%%Moreover, 
%%\[
%%   \begin{aligned}
%%  & \widetilde{\alpha}_{k,e}=\left\{\begin{array}{rl}
%%            \widetilde{C}_{\theta,e}(A_{\theta,e}^{k-1}{B}_{\theta,e} & k > 0 \\
%%           \widehat{D}_{\theta,e}   & k = 0
%%    \end{array}\right. \\
%% & \alpha_{k,e}(\theta)=\left\{\begin{array}{rl}
%%            C_{\theta,e}(A_{\theta,e}^{k-1}{B}_{\theta,e} & k > 0 \\
%%           D_{\theta,e}   & k = 0
%%    \end{array}\right. \\
%%    \end{aligned}
%%\]
%\end{lemma}
\end{proof}





\section{Proof of Theorem \ref{thm:main}}
\label{sect:pf:thm_main}
  The proof of Theorem \ref{thm:main} is organized as follows. 
  First, we define an auxiliary quantity called \emph{the infinite past empirical error}
\[ V_N(\theta)=%\lim_{t_0\rightarrow-\infty}
\frac{1}{N}\sum_{t=0}^{N-1}\ell(\rvy(t),\hat{\rvy}_\theta(t)) \]
It follows that $V_N(\theta)$ is an unbiased estimate of $\mathcal{L}(\theta)$, i.e., $\bE[V_N(\theta)]=\mathcal{L}(\theta)$.
%and as it will be shown in \textcolor{red}{Lemma \ref{lemma:moments_|V-Lhat|}}, that under some mild assumptions,
%$V_N(\theta)$ converges to the empirical error in moments, as $N \rightarrow \infty$.

As the first step, we show the following general theorem.
\begin{theorem}%[\cite{nips-16,alquier-15}]
\label{thm:initial}
    Given a parameter set $\Theta$, a prior distribution $\pi$ over $\Theta$, $\delta\in[0,0.5)$, $\lambda>0$ and $\epsilon \in \{-1,1\}$. Then with probability at least $1-2\delta$, the following holds
    \begin{align}
        \forall \rho\in \mathcal{M}_\pi:\epsilon E_{\theta\sim\rho} \mathcal{L}(\theta)
        \leq \epsilon E_{\theta\sim\rho}\hat{\mathcal{L}}_N(\theta) +        
        \tilde{\mathbb{B}}_N(\rho,\pi,\lambda,\epsilon) \label{eq:initial:tildern} \\
    %\end{equation}
    %with
    %\begin{align}       
    \tilde{\mathbb{B}}_N(\rho,\pi,\lambda,\epsilon)
        \triangleq \frac{1}{\lambda}\left ( \KL(\rho\|\pi) + \ln\dfrac{1}{\delta} + \Psi(\lambda,\pi,N,\epsilon)\right ) \label{eq:initial:def:tildern} \\
    %\end{align}
    %and
    %\begin{multline}
        \Psi(\lambda,\pi,N,\epsilon)\triangleq \frac{1}{2} \Big (\ln E_{\theta\sim\pi}\bE[e^{2\epsilon\lambda(V_N(\theta)-\hat{\mathcal{L}}_N(\theta))}] \nonumber \\
        +\ln E_{\theta\sim\pi}\bE[e^{2\epsilon\lambda (\mathcal{L}(\theta)-V_N(\theta))}] \Big ),\label{eq:Psi:initial}
    \end{align}
    with $\KL(\rho\|\pi)=E_{\theta\sim\rho}\ln\frac{\rho(\theta)}{\pi(\theta)}$ the Kullback–Leibler divergence. 
\end{theorem}
In the theorem above, the moment generating functions 
$\bE[e^{2\epsilon\lambda (\mathcal{L}(\theta)-V_N(\theta))}]$ and  $\bE[e^{2\epsilon\lambda(V_N(\theta)-\hat{\mathcal{L}}_N(\theta))}]$
in the right-hand side are interpreted as $\infty$ whenever $2\epsilon \lambda$ falls outside of their domain of definition. In this case, $\Psi(\lambda,\pi,N,\epsilon)$ is taken to mean $\infty$ and  the bounds are trivially true. 
\begin{proof}[Proof of Theorem \ref{thm:initial}]
    By \cite[Theorem 3]{nips-16}, for any measurable functions $X(\theta,\omega),Y(\theta,\omega)$,  $\tilde{\lambda}>0$, and $\epsilon \in \{-1,1\}$, with probability at least $1-\delta$ we have:
    \begin{multline}
        \forall \rho\in \mathcal{M}_{\pi}:\ E_{\theta\sim\rho}X(\theta,\omega)
        \leq E_{\theta\sim\rho}Y(\theta,\omega)+ \frac{1}{\tilde{\lambda}}\left (\KL(\rho\|\pi)+\ln\dfrac{1}{\delta}+ \Psi_{X,Y}(\tilde{\lambda}) \right ),
        \label{pf:thm:initial:eq1}
    \end{multline}
    with $\Psi_{X,Y}(\tilde{\lambda})=\ln E_{\theta\sim\pi}\bE[e^{\tilde{\lambda}(X(\theta,\omega)-Y(\theta,\omega))}]$.
    Now apply the above theorem for $X=\epsilon \mathcal{L}(\theta),Y=\epsilon V_N(\theta)$, and $X=\epsilon V_N(\theta),Y=\epsilon \hat{\mathcal{L}}_N(\theta)$. Then by applying a union bound on the two obtained probabilistic inequalities, we obtain with probability at least $1-2\delta$
    \begin{multline*}
        \forall \rho \in \mathcal{M}_{\pi}:\ \epsilon E_{\theta\sim\rho}\mathcal{L}(\theta)\leq \epsilon  E_{\theta\sim\rho}\hat{\mathcal{L}}_N(\theta)\\
        + \frac{1}{\tilde{\lambda}}\left (2\KL(\rho\|\pi)+2\ln\dfrac{1}{\delta}+ \Psi_{\epsilon V_N, \epsilon \hat{\mathcal{L}}_N}(\tilde{\lambda}) +\Psi_{\epsilon \mathcal{L},\epsilon V_N}(\tilde{\lambda})
        \right ) 
        %\\+ \frac{1}{\tilde{\lambda}}\left (\KL(\rho\|\pi)+\ln\dfrac{1}{\delta}+ \Psi_{\mathcal{L},V_N}(\tilde{\lambda}) \right )
    \end{multline*}
     With some algebraic manipulation and $\tilde{\lambda}=2\lambda$, we obtain the statement of the theorem. 
\end{proof}
For the case of bounded innovation noise or Lipschitz loss function, we
then show that $\bE[e^{2\lambda(V_N(\theta)-\hat{\mathcal{L}}_N(\theta))}]$
and $\bE[e^{2\lambda (\mathcal{L}(\theta)-V_N(\theta))}]$ can be both upper bounded by
$\mathcal{C}_i(\theta,\lambda,N,\delta,\epsilon)$, from which the statement of the theorem follows.
For the case of quadaratic loss function, we will show that it can be reduced to the case of Lipschitz loss function
with probability $\delta/3$ and then we apply the result for Lipschitz loss with $\delta/3$ instead of $\delta$.



\subsection{Proof of Theorem \ref{thm:main}: bounded case}
For the sake of readability, we repeat below the basic assumptions for the case of bounded data.
\begin{assumption}[Bounded noise and Lipschitz loss]\label{as:bounded}
    The innovation noise $\rve_g(t)$ from \eqref{eq:generator} is essentially bounded , i.e.,
    \begin{align}
        \|\rve_g(t)\|_\infty \leq c_\rve,\; \forall t\in\sZ
    \end{align}
and the loss function $\ell$ is either $L_{\ell}$-Lipschitz or quadratic.
\end{assumption}
The following Lemma is obvious in the $L_\ell$-Lipschits case, and follow by applying Cauchy–Schwarz in the quadratic case.
\begin{lemma}
  \label{lemma:bounded:quadratic_loss}
   Let $B=\{(z,z') \in \mathbb{R}^{n_\rvy} \times \mathbb{R}^{n_\rvy} \mid \|z-z'\|_2 \le G_e(\theta)\sqrt{n_{\rve}}c_{\rve}\}$. Then 
   $(\y(t), \y_{\theta}(t)), (\y(t),\y_{\theta}(t \mid 0)) \in B$ 
   for all $t$, $t_0 \le t$, and for all
   $(z_1,z_2),(z_3,z_4) \in B$, 
   \[ \ell(z_1,z_2)-\ell(z_3,z_4)\le L(\theta)\left(\|z_1-z_3\|_2 +\|z_2-z_4\|_2 \right) \]
   where $L(\theta)$ is as in Theorem \ref{thm:main}, i.e., $L(\theta)=2G_e(\theta)\sqrt{n_{\rve}}c_{\rve}$ if $\ell$-quadratic, and $L(\theta)=L_{\ell}$ if $\ell$ is $L_{\ell}$-Lipschitz.
\end{lemma}
\begin{remark}[Norm loss $L_{\ell}$-Lipschitz]
\label{rem:lipschitz2}
     Note that if the loss $\ell$ is the $p$ norm of the difference, i.e., $\ell(y,\hat{y})=\|y-\hat{y}\|_p$, then 
    $\ell(y,\hat{y})\leq C_p\|y-\hat{y}\|_2$ with $C_p=\sqrt{n_{\rvy}}$ if $p=1$, and $C_p=1$
     if $p=\infty$ or $p=2$.   
    %it satisfies Assumption \ref{as:loss} with $\widetilde{\ell}(z)=\|z\|_p$.
    %and $L(\theta)=C_p$ is such that $\|z\|_p \le C_p \|z_2\|$. In particular, $L(\theta)=1$ if  $p=2$, $L(\theta)=\sqrt{2}$ if $p=1$, and
    %$L(\theta)=\sqrt{n_{\rvy}}$ if $p=\infty$. 
     %If $\ell$ is the quadratic loss, i.e., $\ell(y,\hat{y})=\|y-\hat{y}\|_2^2$,
     %then with $\widetilde{\ell}(z)=z^2$ 
     %it satisfies \eqref{as:Lipschitz} with $L(\theta)=2C_e(\theta)=2G_e(\theta)\sqrt{n_{\rve}}c_{\rve}$. 
     %2(1+\bar{G}_{f}(\theta))\|\Sigma_{gen}\|_{\ell_1}\sqrt{n_{\rve}}c_{\rve} \ge 2C_{\ell}(\theta)$.
\end{remark}
%\begin{assumption}[Locally lipshcitz loss function]\label{as:Lipschitz}
%  %Assume that $C_{\ell}(\theta)=C(\theta)+C_g$  and 
%  Assume that 
%the loss function $\ell$ is of the form $\ell(y,\hat{y})=\widetilde{\ell}(y-\hat{y})$  where $\widetilde{\ell}: \mathbb{R}^{n_\rvy} \rightarrow \mathbb{R}$ is such that for any $C > 0$
%$\widetilde{\ell}$ $L(C)$-Lipschitz if restricted to $B(C)=\{ z \in \mathbb{R}^{n_\rvy} \mid \|z\|_2  \le C\}$, i.e.,
%for all $z,z^{'} \in B: |\widetilde{\ell}(z)-\widetilde{\ell}(z^{'})\| \le L(C) \|z-z^{'}\|_2$.
%\end{assumption}
%Moreover, in the sequel we will use moment generating functions (MGF), to which we assign the value $\infty$  whenever they  are not finite. Note that in these cases the derived bounds will be true but vacuous. Later however we provide bounds (on certain intervals) for the MGFs thus ensuring finiteness \textcolor{red}{(see Lemma~\ref{lemma:mgf(Vn-hatL)} and \ref{lem:mgf})}.
\begin{lemma}\label{lemma:bounded_mgf(V-hatL)}Let assumptions \ref{as:generator}, \ref{as:parameterisation} and \ref{as:bounded} hold. Then
    \begin{align*}
        \bE[e^{\left(\epsilon \lambda|V_N(\theta)-\hat{\mathcal{L}}_N(\theta)|\right)}] 
        \leq e^{\left(\frac{L(\theta)\lambda}{N} \theta_{\infty}(\alpha_{\theta})\|\alpha_g\|_{\ell_1} (c_\rve\sqrt{n_\rve})\right)}
        =e^{C_{b,1}(\lambda,\theta,L(\theta),c_\rve,N)}
        % \left \{
        %\frac{2\lambda K^2}{N} \bar{G}_{f,1}(\theta)\bar{G}_{f,2}(\theta)\right \} \right )}
    \end{align*}
    %with $K\triangleq\|\Sigma_{gen}\|_{\ell_1} (c_\rve\sqrt{n_\rve})$, and
%    \begin{align}
%        \bar{G}_{f,1}(\theta)&\triangleq \hat{M}_\theta\|\hat{B}_\theta\|\|\hat{C}_\theta\| (1-\hat{\gamma}_\theta)^{-1}\\%\dfrac{\hat{M}_\theta\|\hat{B}_\theta\|\|\hat{C}_\theta\|}{1-\hat{\gamma}_\theta}\\
%        \bar{G}_{f,2}(\theta)&\triangleq (1+\|\hat{D}_\theta\|+\bar{G}_{f,1}(\theta))(1-\hat{\gamma}_\theta)^{-1}
%    \end{align}
\end{lemma}
\begin{proof}[Proof \ref{lemma:bounded_mgf(V-hatL)} ]
        Using the Lipschitz property of the loss function it follows that 
    \begin{equation}
        \label{lemma:bounded_mgf(V-hatL):eq1}
       \begin{aligned}
        & \epsilon (V_N(\theta)-\hat{\mathcal{L}}_N(\theta)) \le 
      |V_N(\theta) - \hat{\mathcal{L}}_N(\theta)| \le 
    \frac{1}{N} \sum_{t=0}^{N-1}|\ell(\y(t),\hat{\rvy}_{\theta}(t)) 
    - \ell(\rvy(t),\hat{\rvy}_{\theta}(t \mid 0))|
  \le \\
  & %\frac{L(\theta)}{N} \sum_{t=0}^{N-1}
    \frac{L(\theta)}{N} \sum_{t=0}^{N-1} \|\hat{\rvy}_{\theta}(t)-\hat{\rvy}_{\theta}(t \mid 0)\|_2
       \end{aligned}
    \end{equation}
    From Lemma \ref{alpha:lemma2} it follows that 
    \[ 
       \hat{\rvy}_{\theta}(t)-\hat{\rvy}_{\theta}(t \mid 0) = 
       ((c_{\theta}-c_{\theta,t})\star\e_g)(t)
       =\sum_{k=0}^{\infty} 
       (c_{\theta}-c_{\theta,t})(k)\e_g(t-k)
    \]
    and hence 
    \[ 
       \sum_{t=0}^{N-1}\|\hat{\rvy}_{\theta}(t)-\hat{\rvy}_{\theta}(t \mid 0)\|_2 \le \sum_{t=0}^{N-1} \sum_{k=0}^{\infty} \|(c_{\theta}-c_{\theta,t})(k)\|_{2} \|\e_g(t-k)\|_2 \le 
       \sum_{t=0}^{N-1} \|c_{\theta}-c_{\theta,t}\|_{\ell_1} \sqrt{n_\rve} c_\rve
    \]
    From \eqref{alpha:lemma:eq2d} it follows that 
    \[
       \sum_{t=0}^{N-1} \|c_{\theta}-c_{\theta,t}\|_{\ell_1} \le \theta_{\infty}(\alpha_{\theta}) \|\alpha_g\|_{\ell_1}
    \]
    and hence, 
    \[
   \sum_{t=0}^{N-1} \|\hat{\rvy}_{\theta}(t)-\hat{\rvy}_{\theta}(t \mid 0)\|_2 \le \theta_{\infty}(\alpha_{\theta}) \|\alpha_g\|_{\ell_1}\sqrt{n_\rve} c_\rve
    \]
    from which the statement of the lemma follows using the definition of $\mathcal{C}_{b,1}(\lambda,\theta,L(\theta),c_\rve,N)$.
%%    %Note that 
%%    \[
%%      \begin{aligned}
%%      & (\widetilde{\alpha}_{k,e}-\widetilde{\alpha}_{k,e,t})= \\
%%      & \sum_{l=0}^{k} \alpha_{l}(\theta)\alpha_{k-l,w} - \sum_{l=0}^{\min\{k,t\}} \alpha_{l}(\theta)\alpha_{k-l,w} \\
%%      & = \left\{\begin{array}{rl}
%%           0   & k \le t \\
%%           \sum_{l=t+1}^{k} \alpha_{l}(\theta)\alpha_{k-l,w} & k > t
%%            \end{array}
%%           \right.
%%      \end{aligned} 
%%    \]
%%    and hence
%%    \[
%%      \begin{aligned}
%%    &  \|\hat{\rvy}_{\theta}(t)-\hat{\rvy}_{\theta}(t \mid 0)\|_2  \le  \\
%%    &   \sum_{k=t+1}^{\infty} \|\widetilde{\alpha}_{k,e}-\widetilde{\alpha}_{k,e,t}\|_2 
%%       \|\e_g(t-k)\|_2 \le \\
%%    &    \sum_{k=t+1}^{\infty} \|\sum_{l=t+1}^{k} \alpha_{l}(\theta)\alpha_{k-l,w}\|_2 c_{\rve} \sqrt{n_{\rve}} \\
%%    &    \left(\sum_{k=t+1}^{\infty} \sum_{l=t+1}^{k} \|\alpha_{l}(\theta)\|_2\|\alpha_{k-l,w}\|_2 \right)  c_{\rve} \sqrt{n_{\rve}} \\
%%      \end{aligned}
%%    \]
%%    If we define $\beta_k=\|\alpha_{k+t+1}\|_2$, then 
%%    \[ 
%%       \begin{aligned}
%%       & \sum_{l=t+1}^{k} \|\alpha_{l}(\theta)\|_2\|\alpha_{k-l,w}\|_2 = \\
%%       & \sum_{l=0}^{k-t-1} \beta_k \|\alpha_{k-l-t-1,w}|_2 =:\widetilde{c}_{k-t-1},
%%       \end{aligned}
%%    \]
%%     i.e., the series $\sum_{k=t+1}^{\infty} \sum_{l=t+1}^{k} \|\alpha_{l}(\theta)\|_2\|\alpha_{k-l,w}\|_2=\sum_{k=0}^{\infty} \widetilde{c}_k$ is the Cauchy-product of
%%     the series
%%     $\sum_{k=0}^{\infty} \beta_k=\sum_{k=t+1} \|\alpha_{k}(\theta)\|_2$ and $\sum_{k=0}^{\infty} \|\alpha_{g,w}\|_2 = \|\Sigma_{gen}\|_{\ell_1}$, from which 
%%     by the well-known property of the Cauchy-product and the fact that
%%     \[ 
%%       \sum_{k=0}^{\infty} \beta_k=\sum_{k=t+1} \|\alpha_{k}(\theta)\|_2 \le \bar{G}_f(\theta) < +\infty
%%     \]
%%     it follows that 
%%     \[
%%       \sum_{k=t+1}^{\infty} \sum_{l=t+1}^{k} \|\alpha_{l}(\theta)\|_2\|\alpha_{k-l,w}\|_2=
%%       \left(\sum_{k=t+1}^{\infty} \|\alpha_{k}(\theta)\|_2\right) \cdot \|\Sigma_{gen}\|_{\ell_1}
%%     \]
%%     Hence, 
%%     \[
%%       \begin{aligned}
%%        & \|\hat{\rvy}_{\theta}(t)-\hat{\rvy}_{\theta}(t \mid 0)\|_2 \le \sum_{k=t+1}^{\infty} \|\sum_{l=t+1}^{k} \alpha_{l}(\theta)\alpha_{k-l,w}\|_2 c_{\rve} \sqrt{n_{\rve}} \\
%%        & \left(\sum_{k=t+1}^{\infty} \|\alpha_{k}(\theta)\|_2\right) \cdot \|\Sigma_{gen}\|_{\ell_1} c_{\rve} \sqrt{n_{\rve}}
%%       \end{aligned}
%%     \]
%%     Therefore, 
%%     \begin{equation}
%%        \label{lemma:bounded_mgf(V-hatL):eq2}
%%        \begin{aligned}
%%        & \sum_{t=0}^{N-1} \|\hat{\rvy}_{\theta}(t)-\hat{\rvy}_{\theta}(t \mid 0)\|_2) \le \\
%%        & \left( \sum_{t=0}^{N-1} \sum_{k=t+1}^{\infty} \|\alpha_{k}(\theta)\|_2\right) \cdot \|\Sigma_{gen}\|_{\ell_1} c_{\rve} \sqrt{n_{\rve}} \\
%%        \end{aligned}
%%    \end{equation}
%%     By noticing that 
%%     \[
%%       \begin{aligned}
%%        & \sum_{t=0}^{N} \sum_{k=t+1}^{\infty} \|\alpha_{k}(\theta)\|_2 \le \\
%%        & \sum_{k=1}^{\infty} \min\{k,N\} \|\alpha_{k}(\theta)\|_2
%%        \le \\
%%        & \sum_{k=1}^{\infty} k \|\alpha_{k}(\theta)\|_2=\bar{G}_{f,2}(\theta)
%%       \end{aligned}
%%     \]
%%     By combining the latter inequality with  \eqref{lemma:bounded_mgf(V-hatL):eq1} and  \eqref{lemma:bounded_mgf(V-hatL):eq2},
%%     it follows that 
%%     \[  
%%        \left|V_N(\theta) - \hat{\mathcal{L}}_N(\theta)\right| \le \frac{L(\theta)}{N} \bar{G}_{f,2}(\theta) \|\Sigma_{gen}\|_{\ell_1} c_{\rve} \sqrt{n_{\rve}} 
%%    \]
%%    from which the statement follows by multiplying the inequality by $\lambda$ and using the monotinicity of the exponential function and
%%    taking expectations.
%%        %\begin{multline}
%%        %    \bE[e^{\lambda(V_N(\theta)-\hat{\mathcal{L}_N}(\theta))}]\leq \bE[e^{\lambda|V_N(\theta)-\hat{\mathcal{L}}_N(\theta)|}] \\
%%        %    = 1 + \sum_{r=1}^\infty \frac{\lambda^r}{r!} \bE[|V_N(\theta)-\hat{\mathcal{L}}_N(\theta)|^r] \label{eq:PowerSeriesplus}
%%        %\end{multline}
%%        %With moments from Lemma \ref{lemma:newE|V-Lhat|^r}, and $K=\|\Sigma_{gen}\|_{\ell_1} (c_\rve\sqrt{n_\rve})$ it follows that
%%        %\begin{align} \label{eq:newE|V-Lhat|^r_in_mgf}
%%        %    \bE[\|V_N(\theta)-\hat{\mathcal{L}}_N(\theta)\|^r] \leq  \textcolor{red}{(\bar{G}_{f,1}(\theta)K)^{r}} \left (\frac{2K}{N} \bar{G}_{f,2}(\theta) \right )^r
%%        %\end{align}
%%        %Using \eqref{eq:newE|V-Lhat|^r_in_mgf} in \eqref{eq:PowerSeriesplus} then yields the statement of the lemma.
%%    % Thus
%%    % \begin{multline*}
%%    %     \bE[e^{\lambda|V_N(\theta)-\hat{\mathcal{L}}(\theta)|}] \\
%%    %     \leq 1 + \bar{G}_{f,1}(\theta) K\left (\exp \left \{\frac{2\lambda K}{N} \bar{G}_{f,2}(\theta)\right \}-1 \right )
%%    % \end{multline*}
%%    % and therefore the statement of the lemma holds.
\end{proof}
\begin{lemma}\label{lemma:bounded:mgf(L-V)}%[see Lemma \ref{cor:KL(L-V):alt2}] 
Let assumptions \ref{as:generator}, \ref{as:parameterisation}  and \ref{as:bounded} hold. Then
%for any \textcolor{cyan}{$\epsilon \in \{1,-1\}$} following holds
\begin{equation}
 \bE[e^{\epsilon \lambda(\mathcal{L}(\theta)-V_N(\theta))}] \leq 
  e^{\frac{(L(\theta))^2\lambda^2}{2N}(G_e(\theta)+2G_{e,1}(\theta))(c_{\rve}\sqrt{n_{\rve}})^2}
=e^{\mathcal{C}_{b,2}(\lambda,\theta,L(\theta),c_\rve,N)}
        %e^{\frac{2\lambda^2 }{N} c_\rve^4 n_\rve^2( G_e(\theta)  + 2G_{e,1}(\theta))^2G_e(\theta)^2}%  e^{\frac{\lambda^2}{2N} (G_e(\theta)+G_{e,1}(\theta))^2C^2 (4G_e(\theta)C+1)^2}  
    \end{equation}
   % where  $C=c_\rve\sqrt{n_\rvu+n_\rvy}$
   % \[ G_{e,1}(\theta)=\|D_e\|_2+\sum_{k=1}^{\infty} (k+1) \|C_eA_e^{k-1}K_e\|_2 \]
 \end{lemma}
\begin{proof}[Proof of Lemma \ref{lemma:bounded:mgf(L-V)} ]
 For each $\Sigma(\theta) \in \mathcal{F}$, consider
 $\rvz_t=\left[\y^T(t)~~\hat{\y}_\theta^T(t)\right]^T$.
 %\begin{color}{red
 Then by Lemma \ref{alpha:lemma2},  \( \rvz_t=\left(\begin{bmatrix} \alpha_{g,y} \\ c_{\theta} \end{bmatrix} \star \e_g\right)(t) \)
 and $\rvz_t$ is stationary. Moreover, as $\left\|\begin{bmatrix}\alpha_{g,y} \\ c_{\theta} \end{bmatrix}\right\|_{\ell_1} \leq G_e(\theta) <\infty$ and $\theta_{\infty}\left(\begin{bmatrix} \alpha_{g,y} \\ c_{\theta} \end{bmatrix}\right) \leq G_{e,1}(\theta) <\infty$, it follows that $\rvz_t$ is a Bernoulli-shift. In particular,  $\|\rvz_t\|_\infty \leq G_e(\theta) c_\rve \sqrt{n_\rve}$, and by \cite[Proposition 4.2]{alquier2013prediction} 
 $\rvz_t$ is a weakly dependent process in the terminology of
 \cite{alquier2013prediction}, with $\|\rvz_0\|_\infty \leq G_e(\theta) c_\rve \sqrt{n_\rve}$, and the mixing coefficient $\theta_{\infty,N}(1)$ of $\rvz_t$ satisfies
 $\theta_{\infty,N}(1) < 2G_{e,1}(\theta)c_\rve \sqrt{n_\rve}$. 

If $\ell$ is $L_\ell$-Lipschitz, then let $B=\mathbb{R}^{n_{\rvy}} 
\times \mathbb{R}^{n_{\rvy}}$. If $\ell$ is quadratic, then consider the set $B$ from Lemma \ref{lemma:bounded:quadratic_loss}.
% and consider the restruction of $h$ in $B^{2N}$. 
Define  the function $h(z_0,z_0',\ldots,z_{N-1},z_{N-1}')=\epsilon \frac{1}{L(\theta)} \sum_{i=0}^{N-1} \ell(z_i,z_i')$, $z_0,\ldots, z_{N-1}, z_0',\ldots, z_{N-1}' \in \mathbb{R}^{n_{\rvy}}$ on $B^N$.
%If $\ell$ is $L_\ell$ Lipschitz, then $L(\theta)=L_{\ell}$
We claim that 
 $h$ has Lipschitz constant $1$ in the following sense:
\[ |h(z_0,z_0',\ldots,z_{N-1}, z_{N-1}')-h(\tilde{z}_0, \tilde{z}_0',\ldots,\tilde{z}_{N-1}, \tilde{z}_{N-1}')| \le 
\sum_{i=0}^{N-1}  \|z_i-\tilde{z}_i\|_2 + \|z_i'-\tilde{z}_i'\|_2 \]
Indeed, using Lemma \ref{lemma:bounded:quadratic_loss}, %$k=2G_e(\theta) c_\rve \sqrt{n_\rve}$. 
%That is,
\begin{align*}
    & |h(z_0,z_0',\ldots,z_{N-1}, z_{N-1}')-h(\tilde{z}_0, \tilde{z}_0',\ldots,\tilde{z}_{N-1}, \tilde{z}_{N-1}')|
    =\left |\frac{\epsilon}{L(\theta)}\sum_{i=0}^{N-1} \ell(z_i,z_i')-\ell(\tilde{z}_i,\tilde{z}_i') \right |\\
    &\leq \frac{1}{L(\theta)} \sum_{i=0}^{N-1} L(\theta)  (\|z_i-\tilde{z}_i\|_2 + \|z_i'-\tilde{z}_i'\|_2) =
     \sum_{i=0}^{N-1}  \|z_i-\tilde{z}_i\|_2 + \|z_i'-\tilde{z}_i'\|_2
\end{align*}
%If $\ell$ is quadratic, then 

If $\ell$ is $L_{\ell}$-Lipschitz, then $h$ is already defined on  $(\mathbb{R}^{n_{\rvy}} \times \mathbb{R}^{n_{\rvy}})^N$. If $\ell$ is quadratic, 
we can use Kirszbraun's theorem for scalar valued functions
\cite[2.10.44,page 202]{Federer1996},
from which it then follows that $h$ can be extended to a $1$-Lipschitz map on $(\mathbb{R}^{n_{\rvy}} \times \mathbb{R}^{n_{\rvy}})^N$
without changing its values on $B^N$.  In fact, such an 
extension can be defined as 
$h_{ext}(z)=\inf_{x \in B^N} \{h(x) + \|z-x\|_2\}$.
By abuse of notation this extension will be denote by $h$ too. 
Notice that by Lemma \ref{lemma:bounded:quadratic_loss}, $(\y(t),\y_{\theta}(t)) \in B$ for all $t$, hence
$h(\y(0),\y_{\theta}(0),\ldots,\y(N-1),\y_{\theta}(N-1)) = \frac{\epsilon}{L(\theta)} \sum_{i=0}^{N-1} \ell(\y(i),\y_{\theta}(i))=\frac{\epsilon}{L(\theta)} N V_N(\theta)$.

%Now notice that $|\|\rvz_i\|_2+\|\rvz_i'\|_2\big |\leq 2\|\rvz_t\|_2\leq 2G_e(\theta) c_\rve \sqrt{n_\rve}$, and by reverse triangle inequality we have
%\begin{multline}
%    |h(\rvz_0,\ldots,\rvz_{N-1})-h(\rvz_0',\ldots,\rvz_{N-1}')|\\
%    \leq \frac{2G_e(\theta) c_\rve \sqrt{n_\rve}}{k} \sum_{i=0}^{N-1} \|\rvz_i-\rvz_i'\|_2
%\end{multline}
%thus with $k=2G_e(\theta) c_\rve \sqrt{n_\rve}$, $h$ is $1$-Lipschitz

Notice that, 
\[\epsilon \lambda(\bE[V_N(\theta)]- V_N(\theta))=\frac{\lambda L(\theta)}{N} \left ( \bE[h]-h \right ).\]

 Now we can apply \cite[Theorem 6.6]{alquier2013prediction}, from which it follows that for any $q\geq 0$
 \begin{align}
     \bE[e^{q(\bE[h]-h)}]\leq e^{\frac{q^2N}{2} (\|\rvz_t\|_\infty+\theta_{\infty,N}(1))^2}
 \end{align}
By  
choosing $q=\frac{\lambda L(\theta)}{N}$ and using the definition of $\mathcal{C}_{b,2}(\lambda,\theta,L(\theta),c_\rve,N)$, we obtain the statement of the lemma .
\end{proof}
\begin{proof}[Proof of Theorem \ref{thm:main}: bounded processes]
  The statement of the theorem follows from Theorem \ref{thm:initial}
  by replacing $\bE[e^{2\lambda \epsilon(V_N(\theta)-\hat{\mathcal{L}}_N(\theta))}]$ and $\bE[e^{2\lambda \epsilon (\mathcal{L}(\theta)-V_N(\theta))}]$ with their upper bound from
  Lemma \ref{lemma:bounded_mgf(V-hatL)}  and Lemma \ref{lemma:bounded:mgf(L-V)} respectively. 
\end{proof}

\subsection{Proof of Theorem \ref{thm:main}: unbounded processes and Lipschitz loss functions}
\label{lipschitz:unbounded}

%\begin{assumption}[Cramer's condition on the noise]
%\label{as:unbounded_noise}
 %The noise process $\e_g$ is sub-Gaussian and 
% The loss $\ell$ function is $L$-Lipschitz, and 
%The moment generating function
 
%If $\e_g$ is sub-Gaussian and the pseudo-maximal eigenvalue of its covariance matrix is $\mu_{max}(Q_e)$.
%\end{assumption}
Recall that according to our assumptions, the 2-norm of the noise process $\e_g$ is subgaussian. 
Bounded processes have this property, however, in this section we will not
assume that the process $\e_g$ is bounded. 
%We start by recalling the following property of sub-gaussian processes. 

Recall that we also consider Lipschitz losses.
\begin{assumption}[Lipschitz loss]
  \label{as:unbounded_lipschitz}
We assume that the loss function $\ell$ is $L_{\ell}$-Lipschitz, i.e., 
for all $y,\hat{y},y',\hat{y}' \in \mathbb{R}^{n_\rvy}$,
$|\ell(y,\hat{y})-\ell(y',\hat{y}')| \le L_{\ell} (\|y-y'\|_2+\|\hat{y}-\hat{y}'\|_2)$.
\end{assumption}


The main idea is to use Theorem \ref{thm:initial} and to derive bounds for 
moment generating functions of 
$\bE[ 2\lambda \epsilon (V_N(\theta)-\hat{\mathcal{L}}_N(\theta))]$ and $\bE[2\lambda \epsilon (\mathcal{L}(\theta)-V_N(\theta))]$. 
In order to achieve this, we will reduce the problem of bounding 
$\bE[ 2\lambda \epsilon (V_N(\theta)-\hat{\mathcal{L}}_N(\theta))]$ and $\bE[2\lambda \epsilon (\mathcal{L}(\theta)-V_N(\theta))]$
to the same problem for bounded noises, %obtained unbounded signals to that of bounded ones, arising 
by truncating the unbounded noise
$\e_g$  which drives the data generator. %In turn, the latter 
This problem was solved in Lemma \ref{lemma:bounded_mgf(V-hatL)} and Lemma \ref{lemma:bounded:mgf(L-V)}, which thus can be re-used. 

To this end, let 
$C > 0$ be any constant and let 
$\bar{\e}_g(t)$ be the truncated version of $\e_g(t)$, i.e., 
\begin{align}\label{e-}
\bar{\e}_g(t)=\e_g(t) \chi_{\{\|\e_g(t)\| < C\}}(t)
:=\e_g(t) \chi(\|\e_g(t)\| < C). 
\end{align}
Let $\bar{\rvw}(t)=\begin{bmatrix} \bar{\rvy}^T(t) \\ \bar{\rvu}^T(t) \end{bmatrix}$ be the output of the data generator, if $\e_g(t)$ is replaced by $\bar{\e}_g(t)$, and let 
$\hat{\bar{\rvy}}_{\theta}(t \mid t_0)$ be the output of    the predictor 
\eqref{eq:predictor}
parametrized by $\theta$
with $\rvw$ replaced by $\bar{\rvw}$, and let 
$\hat{\bar{\rvy}}_{\theta}(t)$ 
be the limit 
$\lim_{t_0 \rightarrow -\infty} \hat{\bar{\rvy}}_{\theta}(t \mid t_0)$.
Note that Lemma \ref{l:ihp} applies, as $\bar{\e}_g(t)$ satisfies all the assumptions of
Assumption \ref{as:generator}.
%i.e.,
%$\hat{\bar{\rvy}}(t)$ is the output of \eqref{eq:formal_predictor} with $\rvw$ being replaced by $\bar{\rvw}$.
Let $\bar{V}_N(\theta)=\frac{1}{N} \ell(\bar{\rvy}(t),\hat{\bar{\rvy}}_{\theta}(t))$ be the corresponding infinite past empirical error, 
$\hat{\bar{\mathcal{L}}}_N(\theta)=\frac{1}{N}\sum_{t=0}^{N-1} \ell(\bar{\rvy}(t),\hat{\bar{\rvy}}_{\theta}(t \mid 0))$
the finite past empirical error, and 
$\bar{\mathcal{L}}(\theta)=\bE[\ell(\bar{\rvy}(t),\hat{\bar{\rvy}}(t))]$ the true error for the model corresponding to $\theta \in \Theta$, if it is fed with the truncated signal $\bar{\rvw}$
instead of $\rvw$. 

It turns out that we can bound the difference between these quantities and their non-truncated counterparts.
The basic steps are as follows:
First, using Lipschitz properties of the loss function, we bound the moment generating functions of
the differences $\sup_{\theta \in \Theta} |\hat{\mathcal{L}}_N(\theta) - \hat{\bar{\mathcal{L}}}_N(\theta)|$,
$\sup_{\theta \in \Theta}|V_N(\theta) - \bar{V}_N(\theta)|$ by the corresponding quantity for the $\ell_2$ loss, i.e.; by
the moment generating functions of 
$\sup_{\theta \in \Theta} \sum_{t=0}^{N-1} \|\hat{\rvy}_{\theta}(t \mid 0)-\hat{\bar{\rvy}}_{\theta}(t \mid t_0)\|_2$
and $\sup_{\theta \in \Theta} \sum_{t=0}^{N-1} \|\rvy(t)-\bar{\rvy}(t)\|_2$.
\begin{lemma}
\label{vl-bl:lemma0}
\begin{align}
    \bE[ & e^{\lambda \sup_{\theta \in \Theta}  |\hat{\mathcal{L}}_N(\theta) - \hat{\bar{\mathcal{L}}}_N(\theta)|}]
   \le  \nonumber \\
   %\bE[e^{\frac{\lambda L_{\ell}}{N} \sum_{t=1}^{N}
   % \|\y(t)-\bar{\y}(t)\|_2} e^{\frac{\lambda L_{\ell}}{N} \sup_{\theta \in \Theta}\sum_{t=0}^{N-1} \|\hat{\rvy}_{\theta}(t \mid 0)-\hat{\bar{\rvy}}_{\theta}(t \mid t_0)\|_2}]
    %\le \\
    %&
     & \sqrt{\bE[e^{\frac{2 \lambda L_{\ell}}{N} \sum_{t=0}^{N-1}
    (\|\y(t)-\bar{\y}(t)\|_2}]}
    \sqrt{\bE[e^{\frac{2\lambda L_{\ell}}{N} \sup_{\theta \in \Theta} 
    \sum_{t=0}^{N-1} \|\hat{\rvy}_{\theta}(t \mid 0)-\hat{\bar{\rvy}}_{\theta}(t \mid t_0)\|_2}]}
\label{vl-vbl:lemma1:pf:eq0} \\
%\begin{equation}
%\begin{aligned*}
    \bE[ & e^{\lambda \sup_{\theta \in \Theta}  |V_N(\theta) - \bar{V}_N(\theta)|}] 
   \le \nonumber \\
   %\bE[e^{\frac{\lambda L_{\ell}}{N} \sum_{t=0}^{N-1}
   % \|\y(t)-\bar{\y}(t)\|_2}e^{\frac{\lambda L_{\ell}}{N} \sup_{\theta \in \Theta}\sum_{t=0}^{N-1} \|\hat{\rvy}_{\theta}(t)-\hat{\bar{\rvy}}_{\theta}(t)\|_2}]
    %\le  \\
    %& 
    & \sqrt{\bE[e^{\frac{2 \lambda L_{\ell}}{N} \sum_{t=0}^{N-1}
    (\|\y(t)-\bar{\y}(t)\|_2}] }\sqrt{\bE[e^{\frac{2 \lambda L_{\ell}}{N} \sup_{\theta \in \Theta} \sum_{t=0}^{N-1} \|\hat{\rvy}_{\theta}(t)-\hat{\bar{\rvy}}_{\theta}(t)\|_2}]}
\label{vl-vbl:lemma1:pf:eq3}
%\end{aligned*}
\end{align}
\end{lemma}
Then, using the technique of \cite{alquier2012pred} we bound the moment generating functions of 
the truncation errors $\sup_{\theta \in \Theta} \sum_{t=0}^{N-1} \|\hat{\rvy}_{\theta}(t \mid 0)-\hat{\bar{\rvy}}_{\theta}(t \mid t_0)\|_2$
and $\sup_{\theta \in \Theta} \sum_{t=0}^{N-1} \|\rvy(t)-\bar{\rvy}(t)\|_2$. To this end, we use the following bound on the 
moment generating function of $\|\e_g(t)\|_2$.
\begin{lemma}[Subgaussian noise]
\label{lemma:subgauss:Cramer}
If
$\|\e_g(t)\|_2$ is sub-Gaussian  then
%satisfies Cramer's condition, i.e.,
$\Psi_{\e_a}(\lambda)=\bE[e^{\lambda \|\e_g(0)\|_2}]$ 
%of the noise $\e_g(t)$
 is well-defined for all $\lambda > 0$.
%the moment generating function$\Psi_{\e}(\lambda)=\bE[e^{\lambda \|\e_g(0)\|_2}]$,
%is wel-defined for all $\lambda > 0$.
In fact, %there exists a constant $K_e \le C \sigma_e$, where $C$ is a universal
%constant such that
%with $C_{1,e}=\mu_{max}(Q_e)\sqrt{n_e!}$,
%and $C_{2,e}=\sqrt{n_e+2} \mu_{max}(Q_e)$
%$K_e=\max\{\mu_{max}(Q_e)^2n_e, C_{2,e}+C_{1,e}\}$
\begin{align*}
& \Psi_{\e_g}(\lambda) \le  \bar{\Psi}_{\rve_g}(\lambda)=e^{(\lambda^2 14^2 \|\rve_g\|_{\psi_2}^2+3\lambda \|\rve_g\|_{\psi_2})}
%1+\lambda C_{1,e} (1+\lambda C_{2,e})e^{\lambda^2 \mu_{max}(Q_e)^2 n_\e} \le  
\end{align*}
\end{lemma}
Using this bound, the moment generating function of the  truncation error can be bounded as follows:
\begin{lemma}
 \label{basic_lemma:lipschitz2}
 For any $\lambda < 0.5N$ 
\begin{align*} 
  \bE[e^{\frac{\lambda}{N} \sum_{t=0}^{N-1}\|\y(t)-\bar{\y}(t)\|_2}] \le e^{\mathcal{C}(\|\alpha_{g,y}\|_{\ell_1},\lambda,N,C)} \\
  \bE[e^{\frac{\lambda}{N} \sup_{\theta \in \Theta}\sum_{t=0}^{N-1} \|\hat{\y}(t)-\hat{\bar{\y}}(t)\|_2}] \le e^{\mathcal{C}(\|\alpha_{g,w}\|_{\ell_1}\|\bar{\alpha}\|_{\ell_1},\lambda,N,C)} \\
\bE[e^{\frac{\lambda}{N} \sup_{\theta \in \Theta}\sum_{t=0}^{N-1} \|\hat{\y}(t \mid 0  )-\hat{\bar{\y}}(t \mid 0 )\|_2}] \le e^{\mathcal{C}(\|\alpha_{g,w}\|_{\ell_1}\|\bar{\alpha}\|_{\ell_1},\lambda,N,C)}
\end{align*}
where
\begin{equation}
\label{basic_lemma:lipschitz2:theo:eq1}
   \mathcal{C}(s,\lambda,N,C)=\lambda \frac{\bar{\Psi}_{\rve_g}(s)s^2C}{(e^{sC}-1)}+
     \frac{\lambda^2}{N} 2\bar{\Psi}_{\rve_g}( s)
\end{equation}
\end{lemma}
From this we derive
\begin{lemma}
\label{vl-vbl:lemma1}
For any $\lambda < N/4L_{\ell}$ %the following holds
\begin{align*} 
%& \max\{\bE[e^{\lambda |\hat{\mathcal{L}}_N(\theta) - \hat{\bar{\mathcal{L}}}_N(\theta)|}],
%\bE[e^{\lambda\|V_N(\theta) - \bar{V}_N(\theta)|}], \\
%& e^{\lambda |\mathcal{L}(\theta) - \mathcal{\bar{L}}(\theta)|}\} \le  \\
%& \prod_{k=0}^{\infty} \sqrt{\Psi_{\rve_g}(\frac{2\lambda L c_k(\theta)}{N},\bar{G}_{f}(\theta)\|\Sigma_{gen}\|_{\ell_1},C)}  \times \\
%& \prod_{k=0}^{\infty} \sqrt{\Psi_{\rve_g}(\frac{c_{k,g} 2L\lambda}{N},\|\Sigma_{gen}\|_{\ell_1},C)} \le \\
%& e^{\frac{1}{2} \left( C_2(2 L\lambda,N,\theta)+C_1(2\lambda,N,\theta)\right)}=:\Psi_{\e_g,1}(\lambda,C,\theta) \\
%& C_1(\lambda,N)=\|\Sigma_{gen}\|_{\ell_1} \lambda \frac{\Psi_{\e_g}(\|\Sigma_{gen}\|_{\ell_1}) \|\Sigma_{gen}\|_{\ell_1} C}{(e^{\|\Sigma_{gen}\|_{\ell_1}C}-1)}+ \\
%&
%        \frac{\lambda^2}{N} (\|\Sigma_{gen}\|_{\ell_1})^2 \frac{2\Phi(\|\Sigma_{gen}\|_{\ell_1})}{ \|\Sigma_{gen}\|_{\ell_1}^2}
& \max\{\bE[e^{\lambda \sup_{\theta \in \Theta} |\hat{\mathcal{L}}_N(\theta) - \hat{\bar{\mathcal{L}}}_N(\theta)|}],
\bE[e^{\lambda \sup_{\theta \in \Theta} |V_N(\theta) - \bar{V}_N(\theta)|}], 
e^{\lambda \sup_{\theta \in \Theta} |\mathcal{L}(\theta) - \mathcal{\bar{L}}(\theta)|}\} \le  \\
%& \prod_{k=0}^{\infty} \sqrt{\Psi_{\rve_g}(\frac{2\lambda L \sup_{\theta \in \Theta} c_k(\theta) }{N},G_e,C)}  \times \\
%& \prod_{k=0}^{\infty} \sqrt{\Psi_{\rve_g}(\frac{c_{k,g} 2L\lambda}{N},\|\Sigma_{gen}\|_{\ell_1},C)} \le \\
& e^{\frac{1}{2} \left( \mathcal{C}(\|\bar{\alpha}\|_{\ell_1} \|\alpha_{g,w}\|_{\ell_1}, 2L_{\ell}\lambda,N)+
\mathcal{C}(\|\alpha_{g,y}\|_{\ell_1}, 2L_{\ell}\lambda,N)\right)} %\Psi_{\e_g,1}(\lambda,C) \\
%& e^{\frac{1}{2} \left( C_2(2 L\lambda,N)+C_1(2\lambda,N)\right)}=:\Psi_{\e_g,1}(\lambda,C) \\
%& C_1(\lambda,N)=\|\Sigma_{gen}\|_{\ell_1} \lambda \frac{\Psi_{\e_g}(\|\Sigma_{gen}\|_{\ell_1}) \|\Sigma_{gen}\|_{\ell_1} C}{(e^{\|\Sigma_{gen}\|_{\ell_1}C}-1)}+ \\
%&
%        \frac{\lambda^2}{N} (\|\Sigma_{gen}\|_{\ell_1})^2 \frac{2\Phi(\|\Sigma_{gen}\|_{\ell_1})}{ \|\Sigma_{gen}\|_{\ell_1}^2}
%\bE[e^{\lambda |V_N(\theta) - \bar{V}_N(\theta)|}]
%\le \prod_{k=0}^{\infty} \Psi_{\rve_g}(\frac{\lambda L c_k(\theta)}{N},G_{f,3}(\theta),C) \Psi_{\rve_g}(\frac{c_{k,g} \lambda}{N},\|\Sigma_{gen}\|_{\ell_1},C)\\
%e^{\lambda |\mathcal{L}(\theta) - \mathcal{L}(\theta)|}
%\le \prod_{k=0}^{\infty} \Psi_{\rve_g}(\frac{\lambda L c_k(\theta)}{N},G_{f,3}(\theta),C) \Psi_{\rve_g}(\frac{c_{k,g} \lambda}{N},\|\Sigma_{gen}\|_{\ell_1},C)
\end{align*}
%where 
%\[ c_{k,q}=\sum_{i=\max\{0,k-N+1\}}^{k} \|\alpha_{k,g}\|_2.  \]
\end{lemma}
In order to go back to bounding the moment generating functions of 
$\epsilon (V_N(\theta)-\hat{\mathcal{L}}_N(\theta))$ and $\epsilon (\mathcal{L}(\theta)-V_N(\theta))$, 
we use the previous lemma, and the following inequality which we derive using the Cauchy-Scwhartz inequality: %we show that 
\begin{lemma}
  \label{basic_lemma:lipschitz3.1}
  For $\epsilon=\{-1,1\}$
   \begin{align}
 %\le  \\
      %\bE[e^{\lambda \sup_{\theta \in \Theta} |\bar{V}_N(\theta)-V_N(\theta)|} e^{\lambda \sup_{\theta \in \Theta} |\hat{\bar{\mathcal{L}}}_N(\theta)-\hat{\mathcal{L}}_N(\theta)|} e^{\lambda (\bar{V}_N(\theta)-\hat{\bar{\mathcal{L}}}_N(\theta))}] 
      %\nonumber \\
 %   & \le \sqrt{\bE[e^{2\lambda (\bar{V}_N(\theta)-\hat{\bar{\mathcal{L}}}_N(\theta))}]}
     %\sqrt{\bE[e^{2\lambda \sup_{\theta \in \Theta} |\bar{V}_N(\theta)-V_N(\theta)|} e^{2\lambda \sup_{\theta \in \Theta}  |\hat{\bar{\mathcal{L}}}_N(\theta)-\hat{\mathcal{L}}_N(\theta)|}]} \nonumber \\
     %& \le 
      & \bE[e^{\lambda \epsilon (\mathcal{L}(\theta)-V_N(\theta))}]   \le \nonumber  \\
      & \sqrt{\bE[e^{2\lambda \epsilon (\bar{\mathcal{L}}(\theta)-\bar{V}_N(\theta))}]}
     \sqrt{\sqrt{\bE[e^{4\lambda \sup_{\theta \in \Theta} |\bar{V}_N(\theta)-V_N(\theta)|}]} 
     \sqrt{e^{4\lambda \sup_{\theta \in \Theta} |\bar{\mathcal{L}}(\theta)-\mathcal{L}(\theta)|}]}}
     \label{eq:basic_lemma_lipschitz3:pf:eq0.1}  \\
        & \bE[e^{\lambda \epsilon (V_N(\theta)-\hat{\mathcal{L}}_N(\theta))}]   
      \le  \nonumber \\
      & \sqrt{\bE[e^{2\lambda \epsilon (\bar{V}_N(\theta)-\hat{\bar{\mathcal{L}}}_N(\theta))}]}
     \sqrt{\sqrt{\bE[e^{4\lambda \sup_{\theta \in \Theta} |\bar{V}_N(\theta)-V_N(\theta)|}]} \sqrt{e^{4\lambda \sup_{\theta \in \Theta} |\hat{\bar{\mathcal{L}}}_N(\theta)-\hat{\mathcal{L}}_N(\theta)|}]}}
     \label{eq:basic_lemma_lipschitz3:pf:eq0} 
    \end{align}
\end{lemma}
If we apply Lemma~\ref{vl-vbl:lemma1} %\ref{basic_lemma:lipschitz2} 
to the right-hand sides of the latter inequalities, then we obtain.
\begin{lemma}
\label{basic_lemma:lipschitz3}
For $\epsilon \in \{1,-1\}$, for $\lambda< \frac{N}{16L_{\ell}}$,
\begin{align*}
 & \bE[e^{\lambda \epsilon (V_N(\theta)-\hat{\mathcal{L}}_N(\theta))}]   
 \le \\
 %& \sqrt{\Psi_{\e_g,1}(4\lambda,C)}
 %\sqrt{\bE[e^{2\lambda (\bar{V}_N(\theta)-\hat{\bar{\mathcal{L}}}_N(\theta))}]} = \\
 &e^{\frac{1}{4}(\mathcal{C}(\|\bar{\alpha}\|_{\ell_1}\|\alpha_{g,w}\|_{\ell_1}, 8L_{\ell}\lambda,N,C)+\mathcal{C}(\|\alpha_{g,y}\|_{\ell_1}, 8L_{\ell}\lambda,N,C))}
 \sqrt{\bE[e^{2\lambda \epsilon (\bar{V}_N(\theta)-\hat{\bar{\mathcal{L}}}_N(\theta))}]} 
 %\le \sqrt{\Psi_{\e_g,1}(2\lambda,C)} 
 \\
 & \bE[e^{\lambda \epsilon (\mathcal{L}(\theta)-V_N(\theta))}]   
 \le e^{\frac{1}{4}(\mathcal{C}(\|\bar{\alpha}\|_{\ell_1}\|\alpha_{g,w}\|_{\ell_1}, 8L_{\ell}\lambda,N,C)+\mathcal{C}(\|\alpha_{g,y}\|_{\ell_1}, 8L_{\ell}\lambda,N,C))}
 %& \sqrt{\Psi_{\e_g,1}(4\lambda,C)}
 %\sqrt{\Psi_{\e_g,1}(2\lambda,C)}
 %\sqrt{\bE[e^{2\lambda (\hat{\bar{\mathcal{L}}}(\theta)-\bar{V}_N(\theta)})]}=\\
% &  e^{0.25(C_2(8L\lambda,N)+C_2(4L\lambda,C)+C_1(8L\lambda,N)+C_1(4L\lambda,N))}
 \sqrt{\bE[e^{2\lambda \epsilon (\hat{\bar{\mathcal{L}}}(\theta)-\bar{V}_N(\theta)})]}
\end{align*}
\end{lemma}
Finally, Lemma \ref{lemma:bounded_mgf(V-hatL)} and Lemma \ref{lemma:bounded:mgf(L-V)} 
derived for the bounded case can be applied to the moment generating functions of
$\epsilon (\hat{\bar{\mathcal{L}}}(\theta)-\bar{V}_N(\theta))$ and $\epsilon (\bar{V}_N(\theta)-\hat{\bar{\mathcal{L}}}_N(\theta))$, for a suitable choice of $C$. %the result follows. 
The detailed proof is presented below. 
\begin{proof}[Proof of Theorem \ref{thm:main}: unbounded noise and Lipschitz loss]
 Let $C=\frac{2\ln(N)}{K\|\alpha_g\|_{\ell_1}}$ where $K=\min\{1, \|\bar{\alpha}\|_{\ell_1}\}$. Then
 \[ C\|\alpha_g\|_{\ell_1} \ge 2\ln(N), \quad C\|\bar{\alpha}\|_{\ell_1}\|\alpha_g\|_{\ell_1} \ge 2\ln(N) \]
and hence for $N \ge 2$,
 \begin{align*}
    e^{C\|\alpha_g\|_{\ell_1}}-1 \ge  e^{2\ln(N)} - 1 = N^2-1 \ge N, \quad
    e^{C\|\alpha_g\|_{\ell_1}\|\bar{\alpha}\|_{\ell_1}}-1 \ge  e^{2\ln(N)} - 1 = N^2 -1 \ge N
 \end{align*}
 and therefore
 \[
     \frac{1}{e^{C\|\alpha_g\|_{\ell_1}}-1} \le \frac{1}{N}, \quad 
     \frac{1}{e^{C\|\alpha_g\|_{\ell_1}\|\bar{\alpha}\|_{\ell_1}}-1} \le \frac{1}{N}
 \]
 Therefore, as $C\|\bar{\alpha}\|_{\ell_1}\|\alpha_g\|_{\ell_1} = 2\ln(N) \max\{1,\|\bar{\alpha}\|_{\ell_1}\}$,
 \begin{align*}
    & \mathcal{C}(\|\alpha_g\|_{\ell_1}, 8L_{\ell}\lambda,N,C) \le
      \frac{8L_{\ell}\lambda}{N}\Psi_{\e_g}(\|\alpha_g\|_{\ell_1})\|\alpha_g\|_{\ell_1}^2\frac{2\ln(N)}{K\|\alpha_g\|_{\ell_1}} + \frac{ (8L_{\ell}\lambda)^2}{N}2\Psi_{\e_g}(\|\alpha_g\|_{\ell_1}) \le \\
      & \frac{8L_{\ell}\lambda}{N}\Psi_{\e_g}(\|\alpha_g\|_{\ell_1})\frac{\|\alpha_g\|_{\ell_1}}{\min\{1,\|\bar{\alpha}\|_{\ell_1}\}}2\ln(N) + \frac{ 128L_{\ell}^2\lambda^2}{N}\Psi_{\e_g}(\|\alpha_g\|_{\ell_1})
      \\
    & \mathcal{C}(\|\bar{\alpha}\|_{\ell_1}\|\alpha_g\|_{\ell_1}, 8L_{\ell}\lambda,N,C)
     \le \\
     &\frac{8L_{\ell}\lambda}{N}\Psi_{\e_g}(\|\bar{\alpha}\|_{\ell_1}\|\alpha_g\|_{\ell_1})(\|\bar{\alpha}\|_{\ell_1}\|\alpha_g\|_{\ell_1})^2C + \frac{ (8L_{\ell}\lambda)^2}{N}2\Psi_{\e_g}(\|\bar{\alpha}\|_{\ell_1}\|\alpha_g\|_{\ell_1}) = \\
    & \frac{8L_{\ell}\lambda}{N} \Psi_{\e_g}(\|\bar{\alpha}\|_{\ell_1}\|\alpha_g\|_{\ell_1})\|\bar{\alpha}\|_{\ell_1}\|\alpha_g\|_{\ell_1}2\ln(N)\max\{1,\|\bar{\alpha}\|_{\ell_1}\}+ \frac{128 L_{\ell}^2\lambda^2}{N}\Psi_{\e_g}(\|\bar{\alpha}\|_{\ell_1}\|\alpha_g\|_{\ell_1})
 \end{align*}
 Hence, 
 \begin{equation}
  \label{thm:main:pf:unbounded:eq-1}
 \begin{aligned}
    & \frac{1}{4}(\mathcal{C}(\|\bar{\alpha}\|_{\ell_1}\|\alpha_g\|_{\ell_1}, 8L_{\ell}\lambda,N,C)+\mathcal{C}(\|\alpha_g\|_{\ell_1}, 8L_{\ell}\lambda,N,C))
    \le  \\
     & \frac{4\lambda \ln(N) L_{\ell}}{N} \left( 
      \underbrace{\bar{\Psi}_{\e_g}(\|\bar{\alpha}\|_{\ell_1}\|\alpha_g\|_{\ell_1})\|\bar{\alpha}\|_{\ell_1}\max\{1,\|\bar{\alpha}\|_{\ell_1}\}\|\alpha_g\|_{\ell_1} + \bar{\Psi}_{\e_g}(\|\alpha_g\|_{\ell_1})\frac{\|\alpha_g\|_{\ell_1}}{\min\{1,\|\bar{\alpha}\|_{\ell_1}\}}}_{C_{u,2}}\right)  \\
    & +\frac{32\lambda^2L_{\ell}^2}{N}\underbrace{\left(\bar{\Psi}_{\e_g}(\|\bar{\alpha}\|_{\ell_1}\|\alpha_g\|_{\ell_1})+ \bar{\Psi}_{\e_g}(\|\alpha_g\|_{\ell_1}) \right)}_{C_{u,1}}=\bar{\mathcal{C}}_{ub}(\lambda, L_{\ell},N)
 \end{aligned}
\end{equation}
  From Lemma \ref{lemma:bounded_mgf(V-hatL)} and Lemma \ref{lemma:bounded:mgf(L-V)} applied to 
  $\bar{V}_N(\theta)$, $\bar{\mathcal{L}}(\theta)$, $\hat{\bar{\mathcal{L}}}_N(\theta)$, with  $c_{\rve}=C=\frac{2\ln(N)}{K\|\alpha_g\|_{\ell_1}}$
  and $L(\theta)=L_{\ell}$,
  it follows that 
  \begin{align*}
    & \sqrt{\bE[e^{2\lambda \epsilon (\bar{V}_N(\theta)-\hat{\bar{\mathcal{L}}}_N(\theta))}]}
    \le \sqrt{e^{\mathcal{C}_{b,1}(2\lambda,\theta,L_{\ell},\frac{2\ln(N)}{K\|\alpha_g\|_{\ell_1}},N)}}=e^{\frac{1}{2}\mathcal{C}_{b,1}(2\lambda,\theta,L_{\ell},\frac{2\ln(N)}{K\|\alpha_g\|_{\ell_1}},N)}
      %\left(\frac{L_{\ell}2\lambda}{N} \theta_{\infty}(\alpha_{\theta})\|\alpha_g\|_{\ell_1} (\frac{2\ln(N)}{K}\sqrt{n_\rve})\right)}} =
    %e^{\left(\frac{L_{\ell}\lambda}{N} \theta_{\infty}(\alpha_{\theta})\|\alpha_g\|_{\ell_1} (\frac{2\ln(N)}{K}\sqrt{n_\rve})\right)}
    \\
    & \sqrt{\bE[e^{\epsilon 2\lambda(\bar{\mathcal{L}}(\theta)-\bar{V}_N(\theta))}]} \leq 
    \sqrt{e^{\mathcal{C}_{b,2}(2\lambda,\theta,L_{\ell},\frac{2\ln(N)}{K\|\alpha_g\|_{\ell_1}},N)}}=e^{\frac{1}{2}\mathcal{C}_{b,2}(2\lambda,\theta,L_{\ell},\frac{2\ln(N)}{K\|\alpha_g\|_{\ell_1}},N)}
 % \sqrt{e^{\frac{(L_{\ell})^24\lambda^2}{2N}(G_e(\theta)+2G_{e,1}(\theta))(\frac{2\ln(N)}{K}\sqrt{n_{\rve}})^2}} =
 % e^{}
  %\frac{(L_{\ell})^2\lambda^2}{N}(G_e(\theta)+2G_{e,1}(\theta))(\frac{2\ln(N)}{K}\sqrt{n_{\rve}})^2}
  \end{align*} 
 Using Lemma \ref{basic_lemma:lipschitz3} with $\|\alpha_{g,y}\|_{\ell_1}$ and $\|\alpha_{g,w}\|_{\ell_1}$ replaced by their upper bound $\|\alpha_{g}\|_{\ell_1}$ it follows that 
 for any $\lambda  < L_{\ell}/16N$, 
 \begin{equation}
  \label{thm:main:pf:unbounded:eq1}
 \begin{aligned}
    & \bE[e^{2\lambda \epsilon (V_N(\theta)-\hat{\mathcal{L}}_N(\theta))}] \le 
    e^{\bar{C}_{ub}(2\lambda,L_{\ell},N)+ \frac{1}{2}\mathcal{C}_{b,1}(4\lambda,\theta,L_{\ell},\frac{2\ln(N)}{K\|\alpha_g\|_{\ell_1}},N)}=
    e^{\mathcal{C}_1(2\lambda,\theta,N,\delta,\epsilon)}
    \\
    %\left(\frac{2L_{\ell}\lambda}{N} \theta_{\infty}(\alpha_{\theta})\|\alpha_g\|_{\ell_1} (\frac{2\ln(N)}{K}\sqrt{n_\rve})\right)} \\
    &  \bE[e^{2\lambda \epsilon (\mathcal{L}(\theta)-V_N(\theta))}]  \le e^{\bar{C}_{ub}(2\lambda,L_{\ell},N)+\frac{1}{2}\mathcal{C}_{b,2}(4\lambda,\theta,L_{\ell},\frac{2\ln(N)}{K\|\alpha_g\|_{\ell_1}},N)}=e^{\mathcal{C}_2(2\lambda,\theta,N,\delta,\epsilon)}
 \end{aligned}
\end{equation}
 The statement of the theorem follows from Theorem \ref{thm:initial}
  by replacing $\bE[e^{2\lambda \epsilon(V_N(\theta)-\hat{\mathcal{L}}_N(\theta))}]$ and $\bE[e^{2\lambda \epsilon (\mathcal{L}(\theta)-V_N(\theta))}]$ with their upper bound from \eqref{thm:main:pf:unbounded:eq1}. 

 %The the statement of Theorem \ref{thm:main} follows from Theorem \ref{thm:init}
%\begin{align*}
%   & \Big (\ln E_{\theta\sim\pi}\bE[e^{2\epsilon\lambda(V_N(\theta)-\hat{\mathcal{L}}_N(\theta))}] 
%        +\ln E_{\theta\sim\pi}\bE[e^{2\epsilon\lambda (\mathcal{L}(\theta)-V_N(\theta))}] \Big ) \le
%        %& \Big( \ln E_{\theta \sim \pi} e^{\bar{C}_{ub}(\lambda,L_{\ell},N)+\mathcal{C}_{b,1}(\lambda,\theta,L_{\ell},C,N)}  +
%               \ln E_{\theta \sim \pi} e^{\bar{C}_{ub}(2\lambda,L_{\ell},N)+\mathcal{C}_{b,2}(\lambda,\theta,L_{\ell},C,N)} \Big) = \\
%      %& \ln (e^{\bar{C}_{ub}(\lambda,L_{\ell},N)}  E_{\theta \sim \pi}  e^{\mathcal{C}_{b,1}(\lambda,\theta,L_{\ell},C,N)}) +
%      M  \ln (e^{\bar{C}_{ub}(2\lambda,L_{\ell},N)}  E_{\theta \sim \pi}  e^{\mathcal{C}_{b,2}(\lambda,\theta,L_{\ell},C,N)})
%\end{align*}      
% & \bE[e^{\lambda (V_N(\theta)-\hat{\mathcal{L}}_N(\theta))}]   
% \le \sqrt{\bE[e^{2\lambda (\bar{V}_N(\theta)-\hat{\bar{\mathcal{L}}}_N(\theta))}]} 
% e^{(\lambda L  C_{11} \frac{\ln(N+1)}{N} + (L\lambda)^2C_{12}\frac{1}{N})} \\
% %psi(8L\lambda,G_fG_g,K,N)+\psi(8L\lambda,G_g,K,N))} \\
% & \bE[e^{\lambda (\mathcal{L}(\theta)-V_N(\theta))}]   
 %\le \sqrt{\bE[e^{2\lambda (\hat{\bar{\mathcal{L}}}(\theta)-\bar{V}_N(\theta)})]} e^{(\lambda L C_{21}\frac{\ln(N+1)}{N} + (\lambda L)^2 C_{22}\frac{1}{N})} \\
 %0.25(\psi(8L\lambda,G_fG_g,K,N)+\psi(4L\lambda,G_fG_g,K,N)+\psi(8L\lambda,G_g,K,N)+\psi(4L\lambda,G_g,K,N))} \\
 %& C_{11}=\frac{2(\Psi_{\e_g}(G_fG_g)G_fG_g+\Psi_{\e_g}(G_g)G_g}{K} \\
 %& C_{12}=16(\Psi_{\e_g}(G_fG_g)+\Psi_{\e_g}(G_g)) \\
 %& C_{21} =\frac{3(\Psi_{\e_g}(G_fG_g)G_fG_g+\Psi_{\e_g}(G_g)G_g)}{K} \\
 %& C_{22}= 24(\Psi_{\e_g}(G_fG_g)+\Psi_{\e_g}(G_g))
 %& \psi_b(s,K_1,K)=s \frac{\Psi_{\e_g}(K_1)K_1\ln(N+1)}{NK}+\frac{s^2}{N}\Psi_{\e_g}(K_1)
%\end{align*}
\end{proof}

%%\begin{itemize}
%%\item 
%%      \begin{equation}
%%      \label{unbounded:trunc:eq1}
%%      \begin{aligned} 
%%      \left|\hat{\bar{L}}_N(\theta) - \hat{L}_N(\theta)\right| \le \frac{L_{\ell}}{N} \sum_{t=0}^{N-1} \|\hat{\bar{\rvy}}_{\theta}(t \mid 0) - \hat{\rvy}_{\theta}(t \mid 0)\|_2 \\
%%      \left|\bar{V}_N(\theta) - V_N(\theta)\right| \le \frac{L_{\ell}}{N} \sum_{t=0}^{N-1} \|\bar{\rvy}(t) - \rvy(t)\|_2
%%      \end{aligned}
%%    \end{equation}
%%\item Bound the difference 
%%      \begin{equation}
%%      \label{unbounded:trunc:eq2}
%%      \begin{aligned}
%%      \|\hat{\bar{\rvy}}_{\theta}(t) - \hat{\rvy}_{\theta}(t)\|_2  \le \sum_{k=0}^{\infty} \|c_{\theta}(k)\|_2 \|e_g(t)\|_2 \chi(\|\e_g(t)\|_2 > C) \\
%%      \|\hat{\bar{\rvy}}_{\theta}(t \mid 0) - \hat{\rvy}_{\theta}(t \mid 0)\|_2  \le \sum_{k=0}^{\infty}  \|c_{\theta,t}(k)\|_2 \|e_g(t)\|_2 \chi(\|\e_g(t)\|_2 > C) \\
%%      \end{aligned}
%%      \end{equation}
%%\item By defining $c(k)=\sup_{\theta \in \Theta} \|c_{\theta}(k)\|_2$ and $c_t(k)=\sup_{\theta \in \Theta} \|c_{\theta,t}(k)\|_2$,
%%and use the fact that $\|e_g(t)\|_2 \chi(\|\e_g(t)\|_2 > C) \le \|e_g(t)\|_2$ are independent to
%%we can further bound moment generating function of the differences in \eqref{unbounded:trunc:eq2} as follows:
%%\begin{equation}
%%      \label{unbounded:trunc:eq2}
%%      \begin{aligned}
%%      \bE[e^{\frac{\lambda}{N}\sup_{\theta \in \Theta}\|\hat{\bar{\rvy}}_{\theta}(t) - \hat{\rvy}_{\theta}(t)\|_2}] \le \prod_{k=0}^{\infty} \bE[e^{\frac{\lambda}{N}\|c_{\theta}(k)\|_2 \|e_g(t)\|_2 \chi(\|\e_g(t)\|_2 > C)}] \\
%%      \bE[e^{\frac{\lambda}{N}\sup_{\theta \in \Theta}\|\hat{\bar{\rvy}}_{\theta}(t \mid 0) - \hat{\rvy}_{\theta}(t \mid 0)\|_2}] \le \prod_{k=0}^{\infty} \bE[e^{\frac{\lambda}{N}\|c_{\theta,t}(k)\|_2 \|e_g(t)\|_2 \chi(\|\e_g(t)\|_2 > C)}] \\
%%      \end{aligned}
%%      \end{equation}
%%\item We bound $\bE[e^{s \|e_g(t)\|_2 \chi(\|\e_g(t)\|_2 > C)}]$ by a function $\Psi_{\rve_g}(s,\bar{\lambda},C)$ 
%%for $s \in [0,0.5\bar{\lambda})$.
%%\item We show that for all $\lambda < N \|c\|_{\ell_1}$
%%      \begin{equation}
%%      \label{unbounded:trunc:eq3}
%%  \begin{aligned}
%%      \bE[e^{s\sup_{\theta \in \Theta}\|\hat{\bar{\rvy}}_{\theta}(t) - \hat{\rvy}_{\theta}(t)\|_2}] \le 
%%       \prod_{k=0}^{\infty} \Psi_{\rve_g}(\frac{\lambda}{N}\|c(k)\|_2, \|c\|_{\ell_1},C)
%%       \le e^{} \\
%%      \bE[e^{s\sup_{\theta \in \Theta}\|\hat{\bar{\rvy}}_{\theta}(t \mid 0) - \hat{\rvy}_{\theta}(t \mid 0)\|_2}] \le \prod_{k=0}^{\infty} \Psi_{\rve_g}(\frac{\lambda}{N}\|c_t(k)\|_2, \|c\|_{\ell_1},C) \\
%%      \end{aligned}
%%      \end{equation}
%% \item Finally, we show that with $C=\frac{2\ln(N)}{K}$, $K=\min\{G_e,G_g\}$
%%\end{itemize}

The rest of the section is devoted to the proof of the lemmas used above. 

First, we prove Lemma \ref{lemma:subgauss:Cramer}
\begin{proof}[Proof of Lemma \ref{lemma:subgauss:Cramer}]
  %From \cite{jin2019short} it follows that $\|\e_g(t)\|_2$ is sub-Gaussian.
  
  
  First note that $\z(t):=\|\e_g(t)\|_2 - \bE[\|\e_g(t)\|_2]$ is zero mean sub-Gaussian,
  and using sub-additivity of the sub-Gaussian norm we get $\|\z(t)\|_{\psi_2} \le \|\|\e_g(t)\|_2\|_{\psi_2}+ \|\bE[\|\e_g(t)\|_2]\|_{\psi_2}=\|\rve_g\|_{\psi_2}+ \|\bE[\|\e_g(t)\|_2]\|_{\psi_2}$. Now from the proof of \cite[Proposition 2.6.1]{vershynin2026highdimensional}
  it follows that $\bE[\|\e_g(t)\|_2] \le 3\|\e_g\|_{\psi_2}$, and by \cite[Exercise 2.24]{vershynin2026highdimensional},
  $\|\bE[\|\e_g(t)\|_2]\|_{\psi_2} \le \bE[\|\e_g(t)\|_2]/\sqrt{\ln 2} \le 2 \bE[\|\e_g(t)\|_2]$. Hence $\|\z(t)\|_{\psi_2} \le 7\|\rve_g\|_{\psi_2}$.  
  
  % satisfies 
  % $\|\rve_g\|_{\psi_2}+ 3 \cdot 2\|\rve_g\|_{\psi_2} \le 7\|\rve_g\|_{\psi_2}$. 
  % Here we used that by \cite[Exercise 2.24]{vershynin2026highdimensional},
  % $\|\bE[\|\e_g(t)\|_2]\|_{\psi_2} \le \bE[\|\e_g(t)\|_2]/\sqrt{\ln 2} \le 2 \bE[\|\e_g(t)\|_2]$ and hence $\|\bE[\|\e_g(t)\|_2]\|_{\psi_2} \le 6\|\rve_g\|_{\psi_2}$.
  Then from the proof of
  \cite[Proposition 2.6.1]{vershynin2026highdimensional} it follows that 
  $\bE[e^{\lambda (\|\e_g(t)\|_2 - \bE[\|\e_g(t)\|_2])}] \le e^{\lambda^2 K^2}$ 
  and  $K=\sqrt{\frac{3}{2}} \|\z(t)\|_{\psi_2} \le 14 \|\rve_g\|_{\psi_2}$. It then follows that 
  $\Psi_{\e}(\lambda) \le e^{(\lambda^2 K^2+ \lambda \bE[\|\e_g(t)\|_2])}$ and
  % by using $K \le 14\|\rve_g\|_{\psi_2}$ and $\bE[\|\e_g(t)\|_2] \le 3\|\rve_g\|_{\psi_2}$, 
  the statement of the lemma follows. 
  %^{(\lambda^2 14^2\sigma_e^2+3\lambda\sigma_e)}
  %and noticing that by \cite[Proposition 2.6.1]{vershynin2026highdimensional}
  %$\bE[\|\e_g(t)\|_2] \le C_2 \sigma_e$ for some universal constant $C_2$,
  %we obtain the desired result.
\end{proof}
In order to present the proof of Lemma \ref{vl-bl:lemma0}-\ref{basic_lemma:lipschitz3}, we need 
the following sequence of results.
%follows by using 
%where we use that
%First, we e
\begin{lemma}
\label{noise:boundC-1}
  Let $\z$ be positive valued random variable such that for all $s \le \bar{\lambda}$
  its moment generating function is well-defined. Then for any $C > 0$,
  \[ \bE[\z \chi(\z > C)] \le \frac{C\bE[e^{\bar{\lambda} \z}]}{e^{\bar{\lambda} C}-1} \]
\end{lemma}
Note that, as in \eqref{e-}, we use $\chi(\z > C)$ as a shorthand for the indicator function $\chi_{\{\z > C\}}(\z)$. This practices will be used in the sequel.
\begin{proof}[Proof of Lemma \ref{noise:boundC-1}]
  Let us consider $\phi(x)=\frac{e^x-1}{x}$. It then follows
  that $\phi$ is a monotonically increasing function of $x$ for $x > 0$.
  Hence, 
  \[ 
    \frac{\phi(\bar{\lambda} \z)}{\phi(\bar{\lambda}C)} \ge  \chi(\z > C).
  \]
   Indeed, if $\z > C$, then 
   $\frac{\phi(\bar{\lambda} \z)}{\phi(\bar{\lambda}C)} \ge 1 = \chi(\z > C)$. If $\z \le  C$, then 
   $1 > \frac{\phi(\bar{\lambda} \z)}{\phi(\bar{\lambda}C)} \ge 0  = \chi(\z > C)$. 
   It then follows that 
   \[ 
      \z  \frac{\phi(\bar{\lambda} \z)}{\phi(\bar{\lambda}C)} \ge \z \chi(\z > C)
   \]
  and by taking expectations it follows that
  \[
     \begin{aligned}
       \bE[\z \chi(\z > C))] \le
       \frac{1}{\phi(\bar{\lambda}C)} \bE[\z \phi(\bar{\lambda} \z)]
     \end{aligned}  
  \]
  Note that
  \[
     \z \phi(\bar{\lambda} \z)=\frac{e^{\bar{\lambda} \z}-1}{\bar{\lambda}}
  \]
  and hence by using the inequality $e^x-1 \le e^{x}$ for $x \ge 0$ it follows that
  \[ 
     \bE[\z \phi(\bar{\lambda} \z)]=\frac{\bE[e^{\bar{\lambda} \z}]-1}{\bar{\lambda}} \le \frac{\bE[e^{\bar{\lambda} \z}]}{\bar{\lambda}}
  \]
  By noticing that
  \[
     \bar{\lambda}\phi(\bar{\lambda}C)=\frac{e^{\bar{\lambda}C}-1}{ C}
  \]
  the statement of the lemma follows. 
  \end{proof}
\begin{lemma}
\label{noise:boundC}
%With notation as in Lemma~\ref{lemma:subgauss:Cramer}
\begin{align*} 
     \bE[e^{s \|e_g(t)\|_2 \chi(\|\e_g(t)\|_2 > C)}]
     \le 
     \left(1+ s 
         \frac{\bar{\Psi}_{\rve_g}(\bar{\lambda})C}{e^{\bar{\lambda}C}-1} + s^2
         \frac{2\bar{\Psi}_{\rve_g}(\bar{\lambda})}{\bar{\lambda}^2}\right):=\Psi_{\rve_g}(s,\bar{\lambda},C)
 \end{align*}
 for all $s \in [0,0.5\bar{\lambda})$. 
\end{lemma}
\begin{proof}[Proof Lemma \ref{noise:boundC}]
   First note that 
   \begin{equation}
     \label{noise:boundC:eq1}
     \begin{aligned}
     & \bE[e^{s \|\e_g(t) \|_2 \chi(\|\e_g(t)\|_2 > C))}] =  1 + s  \bE[\|\e_g(t)\|
     _2\chi(\|\e_g(t)\|_2 > C))] + \\
      & \sum_{k=2}^{\infty} \frac{\bE[\|\e_g(t)\|_2^k \chi(\|\e_g(t)\|_2 > C)] s^k]}{k!}.
     \end{aligned}
    \end{equation}
and $\bE[\|\e_g(t)\|_2^k \chi(\|\e_g(t)\|_2 > C) \le \bE[\|\e_g(t)\|_2^k]$.

As in Lemma~\ref{lemma:subgauss:Cramer} let $\Psi_{\rve_g}(\bar{\lambda})=
\bE[e^{\bar{\lambda}\|\e_g(0)\|_2}]=
\bE[e^{\bar{\lambda}\|\e_g(t)\|_2}]$ be the moment generating function for $\|\e_g(t)\|_2$. Then
\[ 
%\bar{\Psi}_{\e_g}(\bar{\lambda}) \ge 
\Psi_{\rve_g}(\bar{\lambda})
%=\bE[e^{\bar{\lambda}\|\e_g(t)\|_2}]
=\sum_{k=0}^{\infty} \frac{\bE[\|\e_g(t)\|_2^k]}{k!} 
\bar{\lambda}^k
\]
and we get 
\[
    \frac{\bE[\|\e_g(t)\|_2^k]}{k!} \le \frac{\Psi_{\rve_g}(\bar{\lambda})}{\bar{\lambda}^k}.
\]
It follows that for all $k \ge 2$,
\[
   \frac{\bE[\|\e_g(t)\|_2^k \chi(\|\e_g(t)\|_2 > C)]s^k}{k!} \le 
   \frac{\bE[\|\e_g(t)\|_2^k]s^k}{k!} \le  
   \Psi_{\rve_g}(\bar{\lambda})\frac{s^k}{\bar{\lambda}^k}.
\]
Therefore, for all $s \in [0,0.5\bar{\lambda})$. 
\begin{equation}
\label{noise:boundC:eq2}
  \begin{aligned}
     & \sum_{k=2}^{\infty} \frac{\bE[\|e_g(t)\|_2^k \chi(\|\e_g(t)\|_2 > C)] s^k]}{k!} \le 
      \sum_{k=2}^{\infty}  \Psi_{\rve_g}(\bar{\lambda})\frac{s^k}{\bar{\lambda}^k}=
     \Psi_{\rve_g}(\bar{\lambda}) \frac{s^2}{\bar{\lambda}^2} \frac{1}{1-\frac{s}{\bar{\lambda}}} 
     = \\
     & \Psi_{\rve_g}(\bar{\lambda}) \frac{s^2}{\bar{\lambda} (\bar{\lambda}-s)} \le 
     \Psi_{\rve_g}(\bar{\lambda}) \frac{s^2}{0.5 \bar{\lambda}^2}=
     2\Psi_{\rve_g}(\bar{\lambda}) \frac{s^2}{\bar{\lambda}^2} 
  \end{aligned}
\end{equation}
where in the last inequality we used that as $s < 0.5\bar{\lambda}$,
$\bar{\lambda}-s \ge 0.5\bar{\lambda}$.
Finally, use Lemma \ref{noise:boundC-1} with $\z=\|\e_g(t)\|_2$ to obtain
%and noticing that $\bE[e^{\bar{\lambda} \|\e_g(t)\|_2}] \le \bar{\Psi}_{\rve_g}(\bar{\lambda})$,
%it then follows that 
 \begin{equation}
 \label{noise:boundC:eq3}
 \bE[\|\e_g(t)\|_2\chi(\|\e_g(t)\|_2 > C))]
 \le \frac{\Psi_{\rve_g}(\bar{\lambda}) C}{e^{\bar{\lambda}C}-1} 
    % \frac{1}{\phi(\bar{\lambda}C)} \bE[\|e_g(t)\|_2 \phi(\bar{\lambda} \|e_g(t)\|_2)]
    % \le \frac{\bar{\Psi}_{\rve_g}(\bar{\lambda}) C}{e^{\bar{\lambda}C}-1}
  \end{equation}
The result now follows by combining \eqref{noise:boundC:eq1}, \eqref{noise:boundC:eq2} and 
\eqref{noise:boundC:eq3}, and then use Lemma~\ref{lemma:subgauss:Cramer}. 
\end{proof}
We will also need the following technical lemma
\begin{lemma}
\label{basic_lemma:lipschitz:lemma_aux}
  Let $a_k \ge 0$ be a sequence of positive numbers such that the series $\sum_{k=0}^{\infty} a_k$ is convergent. For any $N$, define
  $c_{k,N}=\sum_{l=\max\{0,k-N+1\}}^k a_l$. Then  $\sum_{k=0}^{L} c_{k,N} \le N \sum_{k=0}^{L} a_k$, 
  the series $\sum_{k=0}^{\infty} c_{k,N}$ is convergent, and $\sum_{k=0}^{\infty} c_{k,N} \le N \sum_{k=0}^{\infty} a_k$. 
\end{lemma}
\begin{proof}[Proof of Lemma \ref{basic_lemma:lipschitz:lemma_aux}]
     We will prove by induction on $L \ge N-1$ that
     \begin{equation} 
     \label{basic_lemma:lipschitz:lemma_aux:eq1}
       \sum_{k=0}^{L} c_{k,N} = \sum_{k=0}^{L-N+1} N a_k + \sum_{k=L-N+2}^{L} (L-k+1) a_k 
     \end{equation}
     For $L=N-1$, by induction on $N$, it can be shown that 
     \[ \sum_{k=0}^{N-1} c_{k,N}=\sum_{k=0}^{N-1} \sum_{l=0}^{k} a_l = \sum_{k=0}^{N-1} (N-k) a_k.  \]
     
     %Note that $\sum_{k=0}^{N} (N-k) a_k=\sum_{k=0}^{0=L-N+1} N a_k + \sum_{k=1}^{N-1} \underbrace{(N-k)}_{L-k+1} a_k$, i.e., 
     which yield \eqref{basic_lemma:lipschitz:lemma_aux:eq1} for $L=N-1$. 
     Assume \eqref{basic_lemma:lipschitz:lemma_aux:eq1} holds for $N-1\leq L'\leq L$. Notice that $c_{L+1,N}=\sum_{l=L-N+2}^{L} a_l$ and
     hence
     \[
       \begin{aligned}
         & \sum_{k=0}^{L+1} c_{k,N}=\sum_{k=0}^{L} c_{k,N} + c_{L+1,N}=  \sum_{k=0}^{L-N+1} N a_k + \sum_{k=L-N+2}^{L} (L-k+1) a_k + \sum_{l=L-N+2}^{L} a_l= \\
         & \sum_{k=0}^{L-N+1} N a_k + \sum_{k=L-N+2}^{L} (L-k+2) a_k=\sum_{k=0}^{L-N+1} N a_k + (L-(L-N+2)+2)a_{L-N+2} +  \\
         & \sum_{k=L-N+3}^{L} ((L+1)-k+1) a_k=  
          \sum_{k=0}^{L-N+2} N a_k +\sum_{k=L-N+3}^{L} ((L+1)-k+1) a_k 
       \end{aligned} 
     \]
     The last expression is the right-hand side of 
     \eqref{basic_lemma:lipschitz:lemma_aux:eq1} for $L$ being replaced by $L+1$, thus proving \eqref{basic_lemma:lipschitz:lemma_aux:eq1}.
    Note that for $k=L-N+2,\ldots, L$ we have $L-k+1 \le N$, and hence 
     \eqref{basic_lemma:lipschitz:lemma_aux:eq1} implies that $\sum_{k=0}^{L} c_{k,N} \le N \sum_{k=0}^{L} a_k$. Since 
     $a_k$ and $c_k$ are nonnegative, and $\sum_{k=0}^{\infty} a_k$ is convergent, the rest of the lemma follows.
\end{proof}

\begin{lemma}
\label{basic_lemma:lipschitz}
% Let $\mathbf{I}$ be an index set, and let $\{\rvz_i\}_{i \in \mathbf{I}}$ be a collection of stochastic process
% such that for all $i \in \mathbb{I}$, $\rvz_{i}(t)=(\beta_{i,t} \star \e_g)(t)$ and
% $\bar{\rvz}(t)=(\beta_{i,t} \star \bar{\e}_g)(t)$
% %\sum_{k=0}^{\infty} \beta_{i,t}(k) \bar{\e}_g(t-k)$, for 
% suitable matrices $\{\beta_{i,t}(k)\}_{k=0}^{\infty}$, $i \in \mathbf{I}$, 
% such that %$\sup_{i \in \mathbf{I}} \|\beta_i\|_{\ell_1} <  +\infty$, where
% %for 
% $\|\bar{\beta}\|_{\ell_1}< +\infty$, where $\bar{\beta}(k)=\sup_{i \in \mathbf{I},t \in \mathbb{Z}} \|\beta_{i,t}(k)\|_2$.
% %$\bar{H}(\underline{\beta}):=\sum_{k=0}^{\infty}   +\infty$. 
Let $\mathbf{I}$ be an index set, and $\{\beta_{i,t}(k)\}_{k=0}^{\infty}$, $i \in \mathbf{I}$, be a sequence of matrices such that $\|\bar{\beta}\|_{\ell_1}< \infty$, where $\bar{\beta}(k)=\sup_{i \in \mathbf{I},t \in \mathbb{Z}} \|\beta_{i,t}(k)\|_2$. Let $\rvz_{i}(t)=(\beta_{i,t} \star \e_g)(t)$ and
$\bar{\rvz}_i(t)=(\beta_{i,t} \star \bar{\e}_g)(t)$.
Then for $\lambda < 0.5N$
 \begin{align*}
 & \bE[e^{\frac{\lambda}{N} \sup_{i \in \mathbf{I}}\sum_{t=0}^{N-1} \|\rvz(t)-\bar{\rvz}_i(t)\|_2}]
  \le \prod_{k=0}^{\infty} 
  \bar{\Psi}_{\rve_g}\left(\frac{c(k) \lambda}{N},\|\bar{\beta}\|_{\ell_1},C\right) \le 
   e^{\mathcal{C}(\|\bar{\beta}\|_{\ell_1},\lambda,N,C)}  \\
   & \mathcal{C}(s,\lambda,N,C)=\lambda \frac{\bar{\Psi}_{\rve_g}(s)s^2C}{(e^{sC}-1)}+
     \frac{\lambda^2}{N} 2\bar{\Psi}_{\rve_g}( s)
  %(1+\frac{\lambda c_{k,g}}{N} 
  %    \frac{\Phi(\|\Sigma_{gen}\|C) \|\Sigma_{gen}\|}{e^{\|\Sigma_{gen}\|C}-1} + \frac{\lambda^2 c_{k,g}^2}{N^2}
  %    \frac{2\Phi(\|\Sigma_{gen}\|C)}{\|\Sigma_{gen}\|^2})
 \end{align*} 
 where 
 %$c_{k}=\sup_{i \in \mathbf{I}} c_k(i)$ and  
 $c(k)=\sum_{l=\max\{0,k-N+1\}}^{k} \|\bar{\beta}(l)\|_2$
 and the infinite product on the right-hand side is well-defined. 
 %$\alpha_{k,g}=C_gA_g^{k-1}K_g$, $k > 1$
%$\alpha_{0,g}=I_{n_y}$.
\end{lemma}
\begin{proof}[Proof of Lemma \ref{basic_lemma:lipschitz}]
  Note that 
  \begin{align}
    %& c_{i,t}(k) \le \sum_{l=0}^{\infty} \|\beta_{i,t}(l)\|_2 \le \|\bar{\beta}\|_{\ell_1} \label{basic_lemma:lipschitz:eq0.1} \\
    &  c(k) 
    %= \sum_{l=\max\{0,k-N+1\}}^{k} \bar{\beta}(l) 
    \le \|\bar{\beta}\|_{\ell_1}
      \label{basic_lemma:lipschitz:eq0.2} 
   \end{align}
    and by Lemma \ref{basic_lemma:lipschitz:lemma_aux}
    \begin{equation}
      \label{basic_lemma:lipschitz:eq1}
      \begin{aligned}
     & \sum_{k=0}^{\infty} c(k) \le N \sum_{k=0}^{\infty} \|\bar{\beta}(k)\|_2 = N \|\bar{\beta}\|_{\ell_1} \\
     %& \sum_{k=0}^{\infty} \sup_{i \in \mathbf{I}} c_{i}(k) \le N \sum_{k=0}^{\infty} \sup_{i \in \mathbf{I}}  \|\beta_i(k)\|_2 = N \|\bar{\beta}\|_{\ell_1} \\
      \end{aligned}
    \end{equation}
    %Indeed, the first inequality is a direct consequence of Lemma \ref{basic_lemma:lipschitz:lemma_aux}, while the second one follows
    %from the statement of Lemma \ref{basic_lemma:lipschitz:lemma_aux} 
    %and \eqref{basic_lemma:lipschitz:eq0.2}.
    Note now that 
      \begin{align*}
      & \sum_{t=0}^{N-1}  \|\rvz_{i,t}(t)-\bar{\rvz}_{i,t}(t)\|_2  \\
      & \le
      \sum_{t=0}^{N-1}  \sum_{k=0}^{\infty} \|\beta_{i,t}(k)\|_2 \|\e_g(t-k)-\bar{\e}_g(t-k)\|_2 \\
      & \le \sum_{t=0}^{N-1}  \sum_{k=0}^{\infty} \|\bar{\beta}(k)\|_2 \|\e_g(t-k)\chi(\|\e_g(t-k)\|_2 > C)\|_2 \\
      & =\sum_{k=0}^{\infty} c(k) \|\e_g(N-k-1) \chi(\|\e_g(N-k-1)\|_2 > C)\|_2 
        \end{align*}    
        Hence,
        \begin{align*}
      &\bE[ e^{\frac{\lambda}{N} \sup_{i \in \mathbf{I}, t \in \mathbb{Z}} \sum_{t=1}^{N} \| \|\rvz_{i,t}(t)-\bar{\rvz}_{i,t}(t)\|_2}]
      \le \\
      %& \bE[e^{\frac{\lambda}{N} \sup_{i \in \mathbf{I}, t \in \mathbb{Z}}\sum_{k=0}^{\infty} c_{i,t}(k) \|e_g(N-k-1) \chi(\|\e_g(N-k-1)\|_2 > C)\|_2}] \\
      & \le \bE[e^{\frac{\lambda}{N} \sum_{k=0}^{\infty} 
      %\sup_{i \in \mathbf{I}, t \in \mathbb{Z}} c_{i,t}(k) 
      c(k) \|\e_g(N-k-1) \chi(\|\e_g(N-k-1)\|_2 > C)\|_2}] \\
      & = \lim_{L \rightarrow \infty} 
      \prod_{k=0}^{L} \bar{\Psi}_{\rve_g}\left(\frac{c(k) \lambda}{N},\|\bar{\beta}\|_{\ell_1},C\right)
        \end{align*}
        where we applied  Lemma \ref{noise:boundC}
        to $\bE[e^{\frac{\lambda}{N} c(k)\|e_g(N-k-1) \chi(\|\e_g(N-k-1)\|_2 > C)\|_2}]$
        with $\bar{\lambda}=\|\bar{\beta}\|_{\ell_1}$ and $s=\frac{c(k)\lambda}{N} < 0.5\|\bar{\beta}\|_{\ell_1}$ since $\lambda<0.5N$ by assumption.

        Note that 
        \[
        \begin{aligned}
        &  d_L=\ln\left(\prod_{k=0}^{L} \bar{\Psi}_{\rve_g}\left(\frac{c(k) \lambda}{N},\|\bar{\beta}\|_{\ell_1},C\right) \right)=\sum_{k=0}^{L} r_k \\
        & r_k:=\ln \left(1+\frac{c(k) \lambda}{N} \frac{\bar{\Psi}_{\rve_g}(\|\bar{\beta}\|_{\ell_1})\|\bar{\beta}\|_{\ell_1}C}{e^{\|\bar{\beta}\|_{\ell_1}C}-1} + \left(\frac{c(k) \lambda}{N}\right)^2 \frac{2\bar{\Psi}_{\rve_g}(\|\bar{\beta}\|_{\ell_1})}{\|\bar{\beta}\|_{\ell_1}^2} \right) \geq 0
        \end{aligned}
        \]
        %As $c(k) \ge 0$, $\|\bar{\beta}\|_{\ell_1} > 0$, $r_k \ge 0$. 
        We show that the series $\sum_{k=0}^{\infty} r_k$ 
        is bounded, more precisely, 
        \[
        \sum_{k=0}^{L} r_k \le 
      \|\bar{\beta}\|_{\ell_1} \lambda \frac{\bar{\Psi}_{\rve_g}(\|\bar{\beta}\|_{\ell_1})\|\bar{\beta}\|_{\ell_1}C}{e^{\|\bar{\beta}\|_{\ell_1}C}-1} + \frac{\lambda^2}{N} (\|\bar{\beta}\|_{\ell_1})^2 \frac{2\bar{\Psi}_{\rve_g}(\|\bar{\beta}\|_{\ell_1})}{\|\bar{\beta}\|_{\ell_1}^2} = \mathcal{C} (\|\bar{\beta}\|_{\ell_1},\lambda,N,C)
        \]
        From this it follows that $d_L$ is convergent, and from
        \[
      \prod_{k=0}^{ L} \bar{\Psi}_{\rve_g}\left(\frac{c(k) \lambda}{N},\|\bar{\beta}\|_{\ell_1},C\right) = e^{d_L}
        \]
        it follows that  
        \begin{equation}
       \label{basic_lemma:lipschitz:eq1.3}
       \begin{aligned}
        & \lim_{L \rightarrow \infty} 
        \prod_{k=0}^{ L} \bar{\Psi}_{\rve_g}\left(\frac{c(k) \lambda}{N},\|\bar{\beta}\|_{\ell_1},C\right) \le  e^{\mathcal{C}(\|\bar{\beta}\|_{\ell_1},\lambda,N,C)}
       \end{aligned}
        \end{equation}
        i.e., the infinite product is well-defined, and the statement of the lemma holds. 

        In order to show \eqref{basic_lemma:lipschitz:eq1.3}, note that using $\ln(1+x) \le x$, it follows that 
        \[
        \begin{aligned}
      & r_k \le \frac{c(k) \lambda}{N} \frac{\bar{\Psi}_{\rve_g}(\|\bar{\beta}\|_{\ell_1})\|\bar{\beta}\|_{\ell_1}C}{e^{\|\bar{\beta}\|_{\ell_1}C}-1} + \left(\frac{c(k) \lambda}{N}\right)^2 \frac{2\bar{\Psi}_{\rve_g}(\|\bar{\beta}\|_{\ell_1})}{\|\bar{\beta}\|_{\ell_1}^2}
        \end{aligned}
        \]
        Hence,
        \[
        \begin{aligned}
      & \sum_{k=0}^{L} r_k  \le \frac{\lambda \sum_{k=0}^{L}c(k)}{N} \frac{\bar{\Psi}_{\rve_g}(\|\bar{\beta}\|_{\ell_1})\|\bar{\beta}\|_{\ell_1}C}{e^{\|\bar{\beta}\|_{\ell_1}C}-1} +  \left(\frac{\sum_{k=0}^{L} (c(k))^2}{N^2}\right) \frac{2\bar{\Psi}_{\rve_g}(\|\bar{\beta}\|_{\ell_1})}{\|\bar{\beta}\|_{\ell_1}^2}
        \end{aligned}
        \]
        
        From \eqref{basic_lemma:lipschitz:eq0.2}
        it follows that
        $c(k) \le \|\bar{\beta}\|_{\ell_1}$.
        %Moreover, 
        %\[ (\frac{c_{i,t}(k)}{N})^2=\left(\frac{1}{N} \sum_{l=\max\{0,k-N+1\}}^k \beta_{i,t}(l) \right)^2 \le \frac{1}{N} \sum_{l=\max\{0,k-N+1\}}^k \beta_{i,t}(l)^2 \le \frac{1}{N} \sum_{l=\max\{0,k-N+1\}}^k (\bar{\beta}(l))^2
        %\]
        Hence
        \[
         \left(\frac{1}{N} c(k)\right)^2 \le \frac{1}{N^2} c(k)  \|\bar{\beta}\|_{\ell_1}
        \]
        and as from \eqref{basic_lemma:lipschitz:eq1} it follows that $\sum_{k=0}^{L} c(k) \le N \|\bar{\beta}\|_{\ell_1}$, we get that 
        % as by Lemma \ref{basic_lemma:lipschitz:lemma_aux}
        %
        \begin{align*}
        & \sum_{k=0}^{L} \left(\frac{1}{N} c(k)\right)^2 \le 
        \frac{1}{N^2} \sum_{k=0}^{L}  c(k)
        \|\bar{\beta}\|_{\ell_1} \le 
        \frac{1}{N} \|\bar{\beta}\|_{\ell_1}^2
      \end{align*}
      This completes the proof.
    \end{proof}
Lemma \ref{basic_lemma:lipschitz} provides the basis for several
technical results. 

\begin{proof}[Proof of Lemma \ref{basic_lemma:lipschitz2}]
  From Lemma \ref{alpha:lemma2} it follows that 
  $\y(t)=(\alpha_{g,y} \star \e_g)(t)$ and 
  $\bar{\y}(t)=(\alpha_{g,y} \star \bar{\e}_g)(t)$.
  By setting $I=\{g\}$, $\beta_{i,t}=\alpha_{g,y}$, $i=g$
  and using Lemma \ref{basic_lemma:lipschitz}, the first inequality follows.

  Similar, from Lemma \ref{alpha:lemma2} it follows that 
  $\hat{\y}(t)=(c_{\theta} \star \e_g)(t)$ and 
  $\hat{\bar{\y}}(t)=(c_{\theta} \star \bar{\e}_g)(t)$.
  By setting $I=\Theta$, $\beta_{i,t}=c_{\theta}$, $i=\theta \in \Theta$, 
  it follows that 
  \[ \bE[e^{\frac{\lambda}{N} \sup_{\theta \in \Theta}\sum_{t=0}^{N} \|\hat{\y}(t)-\hat{\bar{\y}}(t)\|_2}] \le e^{\mathcal{C}(\|\bar{c}\|_{\ell_1},\lambda,N,C)} 
  \]
  where $\bar{c}(k)=\sup_{\theta \in \Theta} \|c_{\theta}(k)\|_2$.
  Note that by definition 
  $c_{\theta}(k)=\sum_{l=0}^{k} \alpha_{\theta}(l) \alpha_{g,w}(k-l)$ and hence 
  \[
     \|c_{\theta}(k)\|_2 \le \sum_{l=0}^{k} \|\alpha_\theta(l)\|_2 \|\alpha_{g,w}(k-l)\|_2
    \le \sum_{l=0}^{k} \|\bar{\alpha}(l)\|_2 \|\alpha_{g,w}(k-l)\|_2
  \]
  In particular, 
  \[
      \bar{c}(k) \le \sum_{l=0}^{k} \|\bar{\alpha}(l)\|_2 \|\alpha_{g,w}(k-l)\|_2
  \]
  from which by the well-known property of Cauchy products of series (the infinite sum of the Cauchy product of two absolutely convergent series is equal to the product of their infinite sums) it follows
  that
  \[ \|\bar{c}\|_{\ell_1} \le \|\bar{\alpha}\|_{\ell_1} \|\alpha_{g,w}\|_{\ell_1}.
  \]
  By noticing that $\mathcal{C}(\beta,\lambda,N,C)$ is increasing in $\beta$, it 
  follows that $\mathcal{C}(\|\bar{c}\|_{\ell_1},\lambda,N,C) \le \mathcal{C}(\|\bar{\alpha}\|_{\ell_1} \|\alpha_{g,w}\|_{\ell_1},\lambda,N,C)$, from which the second inequality follows. 

%  and using Lemma \ref{basic_lemma:lipschitz}, the second inequality follows.

 Finally, the proof of the third inequality is similar to the second one. 
 By Lemma \ref{alpha:lemma2} it follows that 
 $\hat{\y}(t\mid 0)=(c_{\theta,t} \star \e_g)(t)$ and 
 $\hat{\bar{\y}}(t\mid 0)=(c_{\theta,t} \star \bar{\e}_g)(t)$.
 By applying lemma \ref{basic_lemma:lipschitz} to 
 $I=\Theta$, $\beta_{i,t}=c_{\theta,t}$  it follows that 
 \[ \bE[e^{\frac{\lambda}{N} \sup_{\theta \in \Theta}\sum_{t=0}^{N} \|\hat{\y}(t \mid 0)-\hat{\bar{\y}}(t \mid 0)\|_2}] \le e^{\mathcal{C}(\|\bar{\bar{c}}\|_{\ell_1},\lambda,N,C)} 
  \]
  where $\bar{\bar{c}}(k)=\sup_{\theta \in \Theta, t \ge 0} \|c_{\theta,t}(k)\|_2$.
  By noticing that 
  \begin{align*}
     & \|c_{\theta,t}(k)\|_2=\|\sum_{l=0}^{\min\{k,t\}} \alpha_{\theta}(k)\alpha_{g,w}(k-l)\|_2 \le \sum_{l=0}^{\min\{k,t\}} \|\alpha_{\theta}(k)\|_2 \|\alpha_{g,w}(k-l)\|_2  \\
     & \le \sum_{l=0}^{k} \|\bar{\alpha}(k)\|_2 \|\alpha_{g,w}(k-l)\|_2
     \le \sum_{l=0}^{k} \|\bar{\alpha}(k)\|_2 \|\alpha_{g,w}(k-l)\|_2,
  \end{align*}
from which it follows that 
\[
\|\bar{\bar{c}}\|_{\ell_1} \le \|\bar{\alpha}\|_{\ell_1} \|\alpha_{g,w}\|_{\ell_1}
\]
From the monotonicity of $\mathcal{C}(s,\lambda,N,C)$ in $s$, it follows that
\[
\mathcal{C}(\|\bar{\bar{c}}\|_{\ell_1},\lambda,N,C) \le \mathcal{C}(\|\bar{\alpha}\|_{\ell_1} \|\alpha_{g,w}\|_{\ell_1},\lambda,N,C),
\]
which completes the proof.
%% Consider the LTI tem \eqref{error-sys1} and \eqref{error-sys1t0} whose output is 
%% $\hat{\y}_{\theta}(t)$ and $\hat{\rvy}(t \mid t_0)$ respectively. By replacing $\e_g$ by $\bar{\e}_g$ it follows that
%%\begin{equation}
%%\label{error-sys1:bar}
%%\begin{split}
%%   \tilde{\bar{\x}}(t+1)&=A_{\theta,e}\tilde{\bar{\x}}(t)+K_{\theta,e}\bar{e}_g(t), \\
%%   \hat{\bar{\y}}_\theta(t)&=\widetilde{C}_{\theta,e}\tilde{\bar{\x}}(t)+\widetilde{D}_{\theta,e}\bar{\e}_g(t) \\
%%\end{split}
%%\end{equation}
%%and for $t \ge t_0$,
%%\begin{equation}
%%\label{error-sys1t0:bar}
%%\begin{split}
%%   \tilde{\bar{\x}}(t+1 \mid t_0)&=A_{\theta,e}\tilde{\bar{\x}}(t \mid t_0)+K_{\theta,e}\bar{e}_g(t), \\
%%   \hat{\bar{\y}}_\theta(t)&=\widetilde{C}_{\theta,e}\tilde{\bar{\x}}(t \mid t_0)+\widetilde{D}_{\theta,e}\bar{\e}_g(t) \\
%%\end{split}
%%\end{equation}
%%where $\tilde{\bar{\x}}(t)=[\bar{x}(t) ~ \hat{\bar{\x}}(t)]^T$ and
%% $\tilde{\bar{\x}}(t \mid t_0)=[\bar{x}(t \mid t_0) ~ \hat{\bar{\x}}(t \mid t_0)]^T$  and $\bar{\x}(t)$ is the state of the data generator \eqref{eq:generator} if $\e_g$ is replaced by $\bar{\e}_g(t)$. 
%%
%%It then follows that 
%%\[
%% \begin{split}
%%& \hat{\rvy}_{\theta}(t)=\sum_{k=0}^{\infty} \widetilde{\alpha}_{k,e}(\theta) \e_g(t-k) \\
%%& \hat{\bar{\rvy}}_{\theta}(t)=\sum_{k=0}^{\infty} \widetilde{\alpha}_{k,e}(\theta) \bar{\e}_g(t-k) \\
%%& \hat{\rvy}_{\theta}(t \mid t_0)=\sum_{k=0}^{\infty} \widetilde{\alpha}_{k,e,t-t_0}(\theta) \e_g(t-k) \\
%%& \hat{\bar{\rvy}}_{\theta}(t \mid t_0)=\sum_{k=0}^{\infty} \widetilde{\alpha}_{k,e,t-t_0}(\theta) \bar{\e}_g(t-k)
%% \end{split}
%%\]
%%Notice moreover that 
%%\begin{equation}
%%\label{basic_lemma:lipschitz2:pf:eq1}
%%  \begin{split}
%%     &   \|\widetilde{\alpha}_{k,e}(\theta)\|_2 \le 
%%         \|\sum_{l=0}^{k} \alpha_{k}(\theta) \alpha_{k,w}\|_2 \le \\
%%         & \sum_{l=0}^{k} \|\alpha_k(\theta)\|_2 \|\alpha_{k,w}\|_2 \le 
%%          \sum_{l=0}^{k} \|\alpha_k(\theta)\|_2 \|\alpha_{k,g}\|_2 \\
%%  & \|\widetilde{\alpha}_{k,e,t-t_0}(\theta)\|_2 \le 
%%         \|\sum_{l=0}^{\min\{k,t-t_0\}} \alpha_{k}(\theta) \alpha_{k,w}\|_2\le \\
%%         & \sum_{l=0}^{k} \|\alpha_k(\theta)\|_2 \|\alpha_{k,w}\|_2 \le 
%%          \sum_{l=0}^{k} \|\alpha_k(\theta)\|_2 \|\alpha_{k,g}\|_2
%%  \end{split}
%%\end{equation}
%%It then follows that 
%%\begin{equation}
%%\label{basic_lemma:lipschitz2:pf:eq2}
%%\begin{split}
%%& \sum_{i=\max\{0,k-N+1\}}^{k} \|\widetilde{\alpha}_{i,e}(\theta)\|_2 \le c_k(\theta) \\
%%& \sum_{i=\max\{0,k-N+1\}}^{k} \|\widetilde{\alpha}_{i,e,t-t_0}(\theta)\|_2 \le c_k(\theta)  \\
%%& \sum_{i=0}^{\infty} \|\widetilde{\alpha}_{i,e}(\theta)\|_2 \le 
%%   \sum_{k=0}^{\infty} \sum_{l=0}^{k} \|\alpha_k(\theta)\|_2 \|\alpha_{k,w}\|_2  \\
%%   & \le \left( \sum_{k=0}^{\infty} \|\alpha_k(\theta)\|_2  \right)
%%   \left( \sum_{k=0}^{\infty} \|\alpha_{k,e}\|_2  \right) \\
%%   & =\bar{G}_{f,1}(\theta) \|\Sigma_{gen}\|_{\ell_1} \\
%%& \sum_{i=0}^{\infty} \|\widetilde{\alpha}_{i,e,t-t_0}(\theta)\|_2 \le 
%%   \sum_{k=0}^{\infty} \sum_{l=0}^{k} \|\alpha_k(\theta)\|_2 \|\alpha_{k,w}\|_2 
%%   \\
%%   & \le \bar{G}_{f,1}(\theta) \|\Sigma_{gen}\|_{\ell_1} \\
%%\end{split}
%%\end{equation}
%%Then the statements of the lemma follows from
%%Lemma \ref{basic_lemma:lipschitz}, by 
%%taking $\rvz(t)=\hat{\bar{\rvy}}_{\theta}(t)$
%%and $\alpha_k=\widetilde{\alpha}_{k,e}(\theta)$
%%and then $\rvz(t)=\hat{\bar{\rvy}}_{\theta}(t \mid t_0)$
%%and $\alpha_k=\widetilde{\alpha}_{k,e,t-t_0}(\theta)$.
%%and then using the inequalities \eqref{basic_lemma:lipschitz2:pf:eq2}.
%%% by noticing that $\wideltilde{G}_e(\theta)=\sum_{k=0}^{\infty} \|\widetilde{\alpha}_{k,e}\|_2$.
\end{proof}



\begin{proof}[Proof of Lemma \ref{vl-bl:lemma0}]
   Notice that  
\begin{align*}
   &  |\hat{\mathcal{L}}_N(\theta) - \hat{\bar{\mathcal{L}}}_N(\theta)| \le 
    \frac{1}{N} \sum_{t=0}^{N-1}|\ell(\y(t),\hat{\rvy}_{\theta}(t \mid 0) 
    - \ell(\bar{\rvy}(t),\hat{\bar{\rvy}}_{\theta}(t \mid 0)|
  \le \\
  & \frac{L_{\ell}}{N} \sum_{t=0}^{N-1}
    \|\y(t)-\bar{\y}(t)\|_2+\frac{L_{\ell}}{N} \sum_{t=0}^{N-1} \|\hat{\rvy}_{\theta}(t \mid 0)-\hat{\bar{\rvy}}_{\theta}(t \mid t_0)\|_2
\end{align*}
and hence 
\begin{align*}
\sup_{\theta \in \Theta} |\hat{\mathcal{L}}_N(\theta) - \hat{\bar{\mathcal{L}}}_N(\theta)| \le 
 \frac{L_{\ell}}{N} \sum_{t=0}^{N-1}
    (\|\y(t)-\bar{\y}(t)\|_2+\frac{L_{\ell}}{N}  \sup_{\theta \in \Theta} \sum_{t=0}^{N-1} \|\hat{\rvy}_{\theta}(t \mid 0)-\hat{\bar{\rvy}}_{\theta}(t \mid t_0)\|_2)
\end{align*}
Then \begin{align*}
   & \bE[e^{\lambda \sup_{\theta \in \Theta} |\hat{\mathcal{L}}_N(\theta) - \hat{\bar{\mathcal{L}}}_N(\theta)|}]
   \le \bE[e^{\frac{\lambda L_{\ell}}{N} \sum_{t=0}^{N-1}
    \|\y(t)-\bar{\y}(t)\|_2} e^{\frac{\lambda L_{\ell}}{N} \sup_{\theta \in \Theta}\sum_{t=0}^{N-1} \|\hat{\rvy}_{\theta}(t \mid 0)-\hat{\bar{\rvy}}_{\theta}(t \mid t_0)\|_2}]
    \le \\
    & \sqrt{\bE[e^{\frac{2 \lambda L_{\ell}}{N} \sum_{t=0}^{N-1}
    (\|\y(t)-\bar{\y}(t)\|_2}]}
    \sqrt{\bE[e^{\frac{2\lambda L_{\ell}}{N} \sup_{\theta \in \Theta} 
    \sum_{t=0}^{N-1} \|\hat{\rvy}_{\theta}(t \mid 0)-\hat{\bar{\rvy}}_{\theta}(t \mid t_0)\|_2}]}
\end{align*}
from which \eqref{vl-vbl:lemma1:pf:eq0} follows.
Similarly, 
\begin{align*}
   &  |V_N(\theta) - \bar{V}_N(\theta)| \le 
    \frac{1}{N} \sum_{t=0}^{N-1}|\ell(\y(t),\hat{\rvy}_{\theta}(t)) 
    - \ell(\bar{\rvy}(t),\hat{\bar{\rvy}}_{\theta}(t))|
  \le \\
  & \frac{L_{\ell}}{N} \sum_{t=0}^{N-1}
    \|\y(t)-\bar{\y}(t)\|_2+\frac{L_{\ell}}{N} \sum_{t=0}^{N-1} \|\hat{\rvy}_{\theta}(t)-\hat{\bar{\rvy}}_{\theta}(t)\|_2
\end{align*}
and hence
\begin{align*}
  & \sup_{\theta \in \Theta} |V_N(\theta) - \bar{V}_N(\theta)| \le 
    \frac{L_{\ell}}{N} \sum_{t=0}^{N-1}
    (\|\y(t)-\bar{\y}(t)\|_2+\frac{L_{\ell}}{N} \sup_{\theta \in \Theta} \sum_{t=0}^{N-1} \|\hat{\rvy}_{\theta}(t)-\hat{\bar{\rvy}}_{\theta}(t)\|_2)
\end{align*}
so that
%\begin{equation}
%\label{vl-vbl:lemma1:pf:eq3}
\begin{align*}
   & \bE[e^{\lambda \sup_{\theta \in \Theta}|V_N(\theta) - \bar{V}_N(\theta)|}]
   \le \bE[e^{\frac{\lambda L_{\ell}}{N} \sum_{t=0}^{N-1}
    \|\y(t)-\bar{\y}(t)\|_2}e^{\frac{\lambda L_{\ell}}{N} \sup_{\theta \in \Theta}\sum_{t=0}^{N-1} \|\hat{\rvy}_{\theta}(t)-\hat{\bar{\rvy}}_{\theta}(t)\|_2}]
    \le  \\
    & \sqrt{\bE[e^{\frac{2 \lambda L_{\ell}}{N} \sum_{t=0}^{N-1}
    (\|\y(t)-\bar{\y}(t)\|_2}]}\sqrt{\bE[e^{\frac{2 \lambda L_{\ell}}{N} \sup_{\theta \in \Theta} \sum_{t=0}^{N-1} \|\hat{\rvy}_{\theta}(t)-\hat{\bar{\rvy}}_{\theta}(t)\|_2}]}
\end{align*}
from which \eqref{vl-vbl:lemma1:pf:eq3} follows.
%\end{equation}
\end{proof}

\begin{proof}[Proof of Lemma \ref{vl-vbl:lemma1}]
%\begin{align*}
%   &  |\hat{\mathcal{L}}_N(\theta) - \hat{\bar{\mathcal{L}}}_N(\theta)| \le 
%    \frac{1}{N} |\ell(\y(t),\hat{\rvy}_{\theta}(t \mid 0) 
%    - \ell(\bar{\rvy}(t),\hat{\bar{\rvy}}_{\theta}(t \mid 0)|
%  \le \\
%  & \frac{L_{\ell}}{N} \sum_{t=0}^{N-1}
%    (\|\y(t)-\bar{\y}(t)\|_2+\frac{L_{\ell}}{N} \sum_{t=0}^{N-1} \|\hat{\rvy}_{\theta}(t \mid 0)-\hat{\bar{\rvy}}_{\theta}(t \mid t_0)\|_2)
%\end{align*}
From Lemma \ref{basic_lemma:lipschitz2} it follows that
%$\rvz(t)=\y(t)$ and $\alpha_k=\begin{bmatrix} I_{n_{\rvy}} & 0 \end{bmatrix}\alpha_{k,e}$ and noticing that $\|\alpha_k\|_2 \le \|\alpha_{k,g}\|$ it follows that 
\begin{equation} 
\label{vl-vbl:lemma1:pf:eq1}
\begin{aligned}
    & \sqrt{\bE[e^{\frac{2 \lambda L_{\ell}}{N} \sum_{t=0}^{N-1}
    \|\y(t)-\bar{\y}(t)\|_2}]} \le e^{\frac{1}{2} \mathcal{C}(\|\alpha_{g,y}\|_{\ell_1}, 2L_{\ell}\lambda,N,C)}\\
    %& \prod_{k=0}^{\infty} \Psi_{\rve_g}(\frac{c_{k,g} \lambda}{N},\|\Sigma_{gen}\|_{\ell_1},C)
\end{aligned}    
\end{equation}
 and %rom Lemma \ref{basic_lemma:lipschitz2} it follows that 
 \begin{equation}
\label{vl-vbl:lemma1:pf:eq2}
\begin{aligned}
& \sqrt{\bE[e^{\frac{2\lambda L_{\ell}}{N} \sup_{\theta \in \Theta}\sum_{t=0}^{N-1} \|\hat{\rvy}_{\theta}(t \mid 0)-\hat{\bar{\rvy}}_{\theta}(t \mid t_0)}]} \le  e^{\frac{1}{2} \mathcal{C}(\|\bar{\alpha}\|_{\ell_1} \|\alpha_{g,w}\|_{\ell_1}, 2L_{\ell}\lambda,N,C)}
%& \prod_{k=0}^{\infty} \Psi_{\rve_g}(\frac{\lambda L c_k(\theta)}{N},\bar{G}_{f}(\theta)\|\Sigma_{gen}\|_{\ell_1})
\end{aligned}
 \end{equation}
 from which by using \eqref{vl-vbl:lemma1:pf:eq0} it follows that %the statement of the lemma follows. 
 \begin{align*}
    & \bE[e^{\lambda \sup_{\theta \in \Theta} |\hat{\mathcal{L}}_N(\theta) - \hat{\bar{\mathcal{L}}}_N(\theta)|}]
    \le   e^{\frac{1}{2} (\mathcal{C}(\|\alpha_{g,y}\|_{\ell_1}, 2L_{\ell}\lambda,N,C) + \mathcal{C}(\|\bar{\alpha}\|_{\ell_1} \|\alpha_{g,w}\|_{\ell_1}, 2L_{\ell}\lambda,N,C))}  
    %& \prod_{k=0}^{\infty} \sqrt{\Psi_{\rve_g}(\frac{\lambda L c_k(\theta)}{N},\bar{G}_{f}(\theta)\|\Sigma_{gen}\|_{\ell_1},C)} \times \\
    %& \prod_{k=0}^{\infty} \sqrt{\Psi_{\rve_g}(\frac{c_{k,g} \lambda}{N},\|\Sigma_{gen}\|_{\ell_1},C)}
 \end{align*}   
Likewise,from Lemma \ref{basic_lemma:lipschitz2} it follows that 
\begin{align*}
    & \sqrt{\bE[e^{\frac{2 \lambda L_{\ell}}{N} \sup_{\theta \in \Theta}\sum_{t=0}^{N-1} \|\hat{\rvy}_{\theta}(t)-\hat{\bar{\rvy}}_{\theta}(t)\|_2}]} \le  e^{\frac{1}{2} \mathcal{C}(\|\bar{\alpha}\|_{\ell_1} \|\alpha_{g,y}\|_{\ell_1}, 2L_{\ell}\lambda,N,C)} 
    %& \prod_{k=0}^{\infty} \Psi_{\rve_g}(\frac{\lambda L c_k(\theta)}{N},\bar{G}_{f}(\theta)\|\Sigma_{gen}\|_{\ell_1})
\end{align*}
By combining this with \eqref{vl-vbl:lemma1:pf:eq1}, it follows by \eqref{vl-vbl:lemma1:pf:eq3} that 
\begin{equation}
\label{vl-vbl:lemma1:pf:eq4}
\begin{aligned}
& \bE[e^{\lambda \sup_{\theta \in \Theta}|V_N(\theta) - \bar{V}_N(\theta)|}] \le e^{\frac{1}{2} (\mathcal{C}(\|\alpha_{g,y}\|_{\ell_1}, 2L_{\ell}\lambda,N,C) + \mathcal{C}(\|\bar{\alpha}\|_{\ell_1} \|\alpha_{g,w}\|_{\ell_1}, 2L_{\ell}\lambda,N,C))}  
%& \prod_{k=0}^{\infty} \sqrt{\Psi_{\rve_g}(\frac{\lambda L c_k(\theta)}{N},\bar{G}_{f}(\theta)\|\Sigma_{gen}\|_{\ell_1},C)} \times \\
%& \prod_{k=0}^{\infty} \sqrt{\Psi_{\rve_g}(\frac{c_{k,g} \lambda}{N},\|\Sigma_{gen}\|_{\ell_1},C)}
\end{aligned}
\end{equation}
Finally, as $\mathcal{L}(\theta)=\bE[V_N(\theta)]$ 
and $\bar{\mathcal{L}}(\theta)=\bE[\bar{V}_N(\theta)]$, it follows that
\begin{align*}
  \sup_{\theta \in \Theta} |\mathcal{L}(\theta) - \bar{\mathcal{L}}(\theta)| &= 
   \sup_{\theta \in \Theta} |\bE[V_N(\theta) - \bar{V}_N(\theta)]| \\
   &\le \sup_{\theta \in \Theta} \bE[|V_N(\theta) - \bar{V}_N(\theta)|] \le    \bE[\sup_{\theta \in \Theta} |V_N(\theta) - \bar{V}_N(\theta)|]
\end{align*}
and hence
\begin{equation}
  \label{vl-bl:lemma1:pf:eq5}
   e^{\lambda \sup_{\theta \in \Theta} |\mathcal{L}(\theta) - \bar{\mathcal{L}}(\theta)|} \le e^{\lambda \bE[\sup_{\theta \in \Theta} |V_N(\theta) - \bar{V}_N(\theta)|]}
\end{equation}
As the function 
$\phi:x \mapsto e^{\lambda x}$ is convex, by Jensen's inequality it follows that 
\begin{align*}
& e^{\lambda \bE[\sup_{\theta \in \Theta} |V_N(\theta) - \bar{V}_N(\theta)|]}=
\phi(\bE[\sup_{\theta \in \Theta} |V_N(\theta) - \bar{V}_N(\theta)|]) \le \\
&  \bE[\phi(\sup_{\theta \in \Theta} |V_N(\theta) - \bar{V}_N(\theta)|)]
   = \bE[e^{\lambda \sup_{\theta \in \Theta} |V_N(\theta) - \bar{V}_N(\theta)|}]
\end{align*}
and then the statement of the lemma follows from 
%\eqref{vl-vbl:lemma1:pf:eq3}-
\eqref{vl-vbl:lemma1:pf:eq4} and \eqref{vl-bl:lemma1:pf:eq5}.
%\begin{align*}
%|\hat{\bar{\mathcal{L}}}_N(\theta)-\bar{V}_N(\theta)| \le
%\frac{L}{N} \sum_{t=0}^{N-1} \|\hat{\bar{\y}}_{\theta}(t \mid 0)-\hat{\bar{\y}}_{\theta}(t \mid 0)\|_2
%\end{align*}
%Note that
%$\hat{\bar{\y}}(t \mid 0)$
%\begin{align*}
%    \|\hat{\bar{\y}}_{\theta}(t \mid 0)-\hat{\bar{\y}}_{\theta}(t \mid 0)\|_2=\sum_{k=0}^{tt}
%\end{align*}
\end{proof}

\begin{proof}[Proof of Lemma \ref{basic_lemma:lipschitz3.1}]
   Note that
 \[
    \begin{aligned}
     & \epsilon (V_N(\theta)-\hat{\mathcal{L}}_N(\theta)) =
     \epsilon(\hat{\bar{\mathcal{L}}}_N(\theta)-\hat{\mathcal{L}}_N(\theta)) +\epsilon(\bar{V}_N(\theta)-\hat{\bar{\mathcal{L}}}_N(\theta))
     +\epsilon    (V_N(\theta)-\bar{V}_N(\theta)) \\
     & \le \sup_{\theta \in \Theta} |\hat{\mathcal{L}}_N(\theta)-\hat{\bar{\mathcal{L}}}_N(\theta)| 
     +\epsilon (\bar{V}_N(\theta)-\hat{\bar{\mathcal{L}}}_N(\theta))
     +\sup_{\theta \in \Theta} |\bar{V}_N(\theta)-V_N(\theta)| 
    \end{aligned}  
 \]
Hence, by using monotonicity of $x \mapsto e^{x}$ and 
Cauchy-Schwartz inequality for expectations,
    \begin{align*}
        & \bE[e^{\lambda\epsilon (V_N(\theta)-\hat{\mathcal{L}}_N(\theta))}]   
 \le 
      \bE[e^{\lambda \sup_{\theta \in \Theta} |\bar{V}_N(\theta)-V_N(\theta)|} e^{\lambda \sup_{\theta \in \Theta} |\hat{\bar{\mathcal{L}}}_N(\theta)-\hat{\mathcal{L}}_N(\theta)|} e^{\lambda \epsilon(\bar{V}_N(\theta)-\hat{\bar{\mathcal{L}}}_N(\theta))}] 
      \nonumber \\
     & \le \sqrt{\bE[e^{2\lambda\epsilon (\bar{V}_N(\theta)-\hat{\bar{\mathcal{L}}}_N(\theta))}]}
     \sqrt{\bE[e^{2\lambda \sup_{\theta \in \Theta} |\bar{V}_N(\theta)-V_N(\theta)|} e^{2\lambda \sup_{\theta \in \Theta}  |\hat{\bar{\mathcal{L}}}_N(\theta)-\hat{\mathcal{L}}_N(\theta)|}]} \nonumber \\
     & \le \sqrt{\bE[e^{2\lambda\epsilon (\bar{V}_N(\theta)-\hat{\bar{\mathcal{L}}}_N(\theta))}]}
     \sqrt{\sqrt{\bE[e^{4\lambda \sup_{\theta \in \Theta} |\bar{V}_N(\theta)-V_N(\theta)|}]} \sqrt{e^{4\lambda \sup_{\theta \in \Theta} |\hat{\bar{\mathcal{L}}}_N(\theta)-\hat{\mathcal{L}}_N(\theta)|}]}}
     %\label{eq:basic_lemma_lipschitz3:pf:eq0}
    \end{align*}
    proving \eqref{eq:basic_lemma_lipschitz3:pf:eq0}. A similar argument proves \eqref{eq:basic_lemma_lipschitz3:pf:eq0.1}.  
\end{proof}
\begin{proof}[Proof of Lemma \ref{basic_lemma:lipschitz3}]
 
 From Lemma~\ref{vl-vbl:lemma1} %\ref{basic_lemma:lipschitz2} 
 it follows that 
 \[
   \begin{aligned}
      & \sqrt{\bE[e^{4\lambda \sup_{\theta \in \Theta} |\bar{V}_N(\theta)-V_N(\theta)|}]} \le e^{\frac{1}{4}(\mathcal{C}(\|\bar{\alpha}\|_{\ell_1}\|\alpha_{g,w}\|_{\ell_1}, 8L_{\ell}\lambda,N,C)+\mathcal{C}(\|\alpha_{g,y}\|_{\ell_1}, 8L_{\ell}\lambda,N,C))} \\
      &  \sqrt{e^{4\lambda \sup_{\theta \in \Theta} |\hat{\bar{\mathcal{L}}}_N(\theta)-\hat{\mathcal{L}}_N(\theta)|}]} \le e^{\frac{1}{4}(\mathcal{C}(\|\bar{\alpha}\|_{\ell_1}\|\alpha_{g,w}\|_{\ell_1}, 8L_{\ell}\lambda,N,C)+\mathcal{C}(\|\alpha_{g,y}\|_{\ell_1}, 8L_{\ell}\lambda,N,C))}
    \end{aligned}    
 \]
 and hence 
\begin{multline*} 
\sqrt{\sqrt{\bE[e^{4\lambda |\bar{V}_N(\theta)-V_N(\theta)|}]} \sqrt{e^{4\lambda |\hat{\bar{\mathcal{L}}}_N(\theta)-\hat{\mathcal{L}}_N(\theta)|}]}}\\ 
\le e^{\frac{1}{4}(\mathcal{C}(\|\bar{\alpha}\|_{\ell_1}\|\alpha_{g,w}\|_{\ell_1}, 8L_{\ell}\lambda,N,C)+\mathcal{C}(\|\alpha_{g,y}\|_{\ell_1}, 8L_{\ell}\lambda,N,C))}
\end{multline*}
By using \eqref{eq:basic_lemma_lipschitz3:pf:eq0} the first statement of the lemma follows.
The proof of the second claim is similar. 
% That is,
% \[
%     \begin{aligned}
%      & (\epsilon)(\mathcal{L}(\theta)-V_N(\theta)) = (\epsilon)(\mathcal{L}(\theta)-\bar{\mathcal{L}}(\theta)) +\epsilon(\bar{\mathcal{L}}(\theta)-\bar{V}_N(\theta))
%      +\epsilon(\bar{V}_N(\theta)-V_N(\theta)) \\
%      & \le \sup_{\theta \in \Theta} |\mathcal{L}(\theta)-\bar{\mathcal{L}}(\theta)| +\epsilon (\bar{\mathcal{L}}(\theta)-\bar{V}_N(\theta))
%      +\sup_{\theta \in \Theta} |\bar{V}_N(\theta)-V_N(\theta)| \
%     \end{aligned}  
%  \]
% Hence, by using monotonicity of $x \mapsto e^{x}$ and  Cauchy-Schwartz inequality for expectations,
%     \begin{align}
%         & \bE[e^{\lambda \hat{\mathcal{L}}(\theta)-V_N(\theta)}]   
%  \le 
%         \bE[e^{\lambda \sup_{\theta \in \Theta} |\bar{V}_N(\theta)-V_N(\theta)|} e^{\lambda \sup_{\theta \in \Theta} |\bar{\mathcal{L}}(\theta)-\mathcal{L}(\theta)|} e^{\lambda (\bar{\mathcal{L}}(\theta)-\bar{V}_N(\theta))}] \nonumber \\
%      & \le \sqrt{\bE[e^{2\lambda (\bar{\mathcal{L}}(\theta)-\bar{V}_N(\theta))}]} 
%     \sqrt{\bE[e^{2\lambda \sup_{\theta \in \Theta} |\bar{V}_N(\theta)-V_N(\theta)|} e^{2\lambda \sup_{\theta \in \Theta} |\bar{\mathcal{L}}(\theta)-\mathcal{L}(\theta)|}] } \nonumber \\
%      & = \sqrt{\bE[e^{2\lambda (\bar{\mathcal{L}}(\theta)-\bar{V}_N(\theta))}]} 
%       \sqrt{\bE[e^{2\lambda \sup_{\theta \in \Theta} |\bar{V}_N(\theta)-V_N(\theta)|}]} \sqrt{e^{2\lambda \sup_{\theta \in \Theta} |\hat{\bar{\mathcal{L}}}_N(\theta)-\hat{\mathcal{L}_N(\theta)|}]} }
%      \label{eq:basic_lemma_lipschitz3:pf:eq0.11}
%     \end{align}    
%  From Lemma \ref{basic_lemma:lipschitz2} it follows that 
%  \[
%     \begin{aligned}
%         & \sqrt{\bE[e^{2\lambda \sup_{\theta \in \Theta} |\bar{V}_N(\theta)-V_N(\theta)|}]} \le  e^{\frac{1}{4} (\mathcal{C}(\|\bar{\alpha}\|_{\ell_1}\|\alpha_g\|_{\ell_1}, 8L_{\ell}\lambda,N,C)+\mathcal{C}(\|\alpha_g\|_{\ell_1}, 8L_{\ell}\lambda,N,C))}\\
%        & \sqrt{e^{2\lambda \sup_{\theta \in \Theta} |\bar{\mathcal{L}}(\theta)-\mathcal{L}(\theta)|}]} \le e^{\frac{1}{4}(\mathcal{C}(\|\bar{\alpha}\|_{\ell_1}\|\alpha_g\|_{\ell_1}, 8L_{\ell}\lambda,N,C)+\mathcal{C}(\|\alpha_g\|_{\ell_1}, 8L_{\ell}\lambda,N,C))}
%     \end{aligned}    
%  \]
%  By using \eqref{eq:basic_lemma_lipschitz3:pf:eq0.11} the second statement of the lemma follows.
\end{proof}

%It turns out that we can bound the MGF of 
%$V_N(\theta)-\bar{V}_{\theta}$ and
%$\hat{\mathcal{L}}_N(\theta)-\hat{\bar{\mathcal{L}}}_N(\theta)$. 
%%\begin{theorem}\label{thm:unbounded_alt} Let assumptions \ref{as:generator}, \ref{as:parameterisation}, \ref{as:unbounded_lipschitz}, \ref{as:unbounded_noise} hold.  For any density $\pi$ on hypothesis class $\mathcal{F}$, %any $\delta\in(0,0.5]$, and
%%any $\lambda>0$, 
%%%the following inequality holds with probability at least $1-2\delta$,
%%%    \begin{multline}
%%%        \forall \hat{\rho}\in\mathcal{M}_\pi: E_{f\sim\hat{\rho}} \mathcal{L}(\theta) \leq E_{\theta\sim \hat{\rho} }\hat{\mathcal{L}}_N(\theta)+\bar{\mathbb{B}}_N(\lambda,\hat{\rho},\delta)
%%%    \end{multline}
%%%     where
%%    \begin{align}
%% \Psi(\lambda,\pi,N) & \le  \bar{\Psi}(\pi,\lambda,N) \mbox{ where } \nonumber \\
%%     \bar{\Psi}(\pi\lambda,N) & \triangleq  \frac{1}{2} \left ( \ln \underset{f\sim\pi}{E} \bar{\Psi}_{1}(\theta,\lambda)+ \ln \underset{f\sim\pi}{E} \bar{\Psi}_{2}(\theta,\lambda)\right ) \nonumber \\
%%        \bar{\Psi}_{1}(\theta,\lambda) & \triangleq
%%        %\exp\left(K_{f,1}(\theta,\lambda,N)+K_{f,2}(\lambda,N)\right) 
%%        \exp\left(\lambda L  \left(C_{11}+\frac{2 C_{b,1}(\theta)}{K}\right) \frac{\ln(N+1)}{N} + \right. \nonumber \\
%%        & \left.  (L\lambda)^2C_{12}\frac{1}{N}\right)
%%        \nonumber \\
%%         %\bar{\Psi}_{2}(\theta) & \triangleq \exp\left(K_{\infty,1}(\theta,\lambda,N)+K_{\infty,2}(\lambda,N)\right) \nonumber \\
%%         \bar{\Psi}_{2}(\theta,\lambda) & \triangleq \exp\left((L\lambda) C_{21}\frac{\ln(N+1)}{N} + 
%%         (L\lambda)^2C_{22}\frac{1}{N} + \right.  \nonumber  \\
%%         &  \left. (L\lambda)^2 \frac{2 C_{b,2}(\theta)}{K^2} \frac{(\ln(N+1))^2}{N} \right)
%%         %(K_{\infty,1}(\theta,\lambda,N)+K_{\infty,2}(\lambda,N)\right) 
%%         \nonumber 
%%         %K_{f,1}(\theta,\lambda,N)& \triangleq \frac{2L \lambda\bar{G}_{f,2}(\theta)\|\Sigma_{gen}\|_{\ell_1} \ln(N+1)\sqrt{n_\rve}}{N K} \nonumber \\
%%%
%%%         K_{f,2}(\lambda,N) & \triangleq 0.25 \left(\phi_{b}(8L\bar{G}_e\lambda,G_{e},K,N)+\right. \nonumber \\
%%%         & \left. \phi_{b}(8L\lambda,\|\Sigma_{gen}\|_{\ell_1},K,N)\right) \nonumber \\
%%       % & \bar{\mathbb{B}}_N(\lambda,\hat{\rho},\delta) \triangleq \frac{1}{\lambda}\left [\KL(\hat{\rho}||\pi)+\ln\frac{1}{\delta} + \bar{\Psi}_{\pi}(\lambda,N)\right ] \nonumber \\
%%        %\left (\exp \left \{\textcolor{red}{\frac{4\lambda K^2}{N}\bar{G}_{f,1}(\theta)} \bar{G}_{f,2}(\theta)\right \}\right ) \nonumber \\ %\left ( 1-C_{1,2}(\theta)+ C_{1,2}(\theta) e^{\frac{\lambda}{N}2 C_{1,1}(\theta)}  \right ) \nonumber \\
%%%         K_{\infty,1}(\theta,\lambda,N) & \triangleq \left(\frac{(4 L)^2\lambda^2 n_{\rve} (\ln(N+1))^2 (G_e(\theta)+2G_{e,1}(\theta)}{2N K^2}\right) \nonumber \\
%%%         K_{\infty,2}(\lambda,N) & \triangleq 0.25 \left(\phi_{b}(8L\lambda\bar{G}_e,G_e,K,N)+ \right. \nonumber \\
%%%         &\left. \phi_{b}(4L\lambda\bar{G}_e,G_{e},K,N)+ \right. \nonumber \\
%%%        &  \left. \phi_{b}(8L\lambda,\|\Sigma_{gen}\|_{\ell_1},K,N)+\right. \nonumber \\
%%%        &  \left. \phi_{b}(4L\lambda,\|\Sigma_{gen}\|_{\ell_1},K,N) \right) \nonumber \\
%%%        & K=\min\{\|\Sigma_{gen}\|_{\ell_1}, G_{e}\} \nonumber \\
%%%        \phi_{b}(s,K_1,K_2,N) & \triangleq \frac{s\Psi_{\e_g}(K_1)\ln(N+1)K_1}{NK_2}+\frac{s^2}{N} 2\Psi_{\e_g}(K_1)
%%%        \nonumber 
%%    \end{align}
%%    where $C_{i,j}$ is as in Lemma \ref{basic_lemma:lipschitz3v0} and $C_{b,i}(\theta)$ is as in Theorem \ref{thm:bounded_alt}, $i,j=1,2$. 
%%  In particular, \eqref{eq:initial:tildern} holds with probability $1-2\delta$ over data, if in the definition of
%%    $\tilde{\mathbb{B}}(\rho,\pi,\lambda)$
%%    we replace $\Psi(\pi,\lambda,N)$ by
%%    $\bar{\Psi}(\pi,\lambda,N)$.
%%%    \begin{align}
%%%        & C_{1,1}(\theta)\triangleq 2\|\Sigma_{gen}\|_{\ell_1} C \bar{G}_{f,2}(\theta), \nonumber \\
%%%        &C_{1,2}(\theta)\triangleq \|\Sigma_{gen}\|_{\ell_1}C \bar{G}_{f,1}(\theta)    \nonumber 
%%%%        & C_2(\theta)\triangleq 8(G_e(\theta)+G_{e,1}(\theta))^2C^2 (4G_e(\theta)C+1)^2 \nonumber 
%%%    \end{align}
%%\end{theorem}
%Intuitively,  terms $C_{i,j}$,$i,j=1,2$
%represent the bounds for the worst-case difference 
%for the choice of the truncation 
%In contrast, parameter dependent terms $C_{b,i}(\theta)$, $i=1,2$ are taken from
%Theorem \ref{thm:bounded_alt} for the bounded noise, and they represent bounds on the MGF $\sqrt{\bE[e^{2\lambda (\bar{V}_N(\theta)-\hat{\bar{\mathcal{L}}}_N(\theta))}]}$ and $\sqrt{\bE[e^{2\lambda  (\bar{\mathcal{L}}(\theta)-\hat{V}_N(\theta))}]}$
%when the noise is assumed to be bounded by the truncation constant
%$C=\frac{\ln(N+1)}{K}$.
% and they do not depend on the parameter $\theta$. 

%The terms $C_{i,j}$,$i,j=2$, which do not depend on the parameter $\theta$, bound the moment generating function of the worst-case difference $\bE[e^{\lambda \sup_{\theta \in \Theta} |\hat{\mathcal{L}}_N(\theta) - \hat{\bar{\mathcal{L}}}_N(\theta)|}]$, and 
%$\bE[e^{\lambda \sup_{\theta \in \Theta} \|V_N(\theta) - \bar{V}_N(\theta)|}]$.
%for the truncation constant $C=\frac{\ln(N+1)}{K}$.
%& e^{\lambda \sup_{\theta \in \Theta} |\mathcal{L}(\theta) - \mathcal{\bar{L}}(\theta)|}\}






%%Similarly to Corollary \ref{thm:bounded:alt:col} we can simplify the general bound
%%\begin{corollary}[Convergence $O(\frac{\ln(N)}{\sqrt{N}})$]
%%\label{thm:ubounded:alt:col}
%%Define the
%%constants 
%%\[
%%\begin{split}
%% %\bar{G}_{f,2}&=\sup_{\theta \in \Theta} \bar{G}_{f,2}(\theta),  \\
%% %& G_{e}=\sup_{\theta \in \Theta} G_e(\theta), \quad G_{e,1}=\sup_{\theta \in \Theta} G_{e,1}(\theta), ~ \\C_{11}+\frac{2 C_{b,1}(\theta)}{K}
%% \mathcal{C}_1 \triangleq (C_{12}+C_{22}), \quad 
%% \mathcal{C}_2 \triangleq  (C_{11}+C_{21}+\frac{2 C_{b,1}}{K}), \quad  
%% \mathcal{C}_3 \triangleq  \frac{2C_{b,2}}{K^2}  
%% %0.5\mathcal{C}_{ub,11}+0.5\mathcal{C}_{ub,12} \\
%% %\mathcal{C}_2 &=0.5\mathcal{C}_{b,2}+0.5\mathcal{C}_{ub,21}+0.5\mathcal{C}_{ub,22} \\
%%  \\ %+0.5\mathcal{C}_{ub,21}+0.5\mathcal{C}_{ub,22} \\
%% %  \mathcal{C}_{b,3}&=\frac{\frac{4L^2}{2}(G_e+2G_{e,1}) n_{\rve}}{K^{2}} \\ 
%% % \mathcal{C}_{b,2}= & \left(L \bar{G}_{f,2}\|\Sigma_{gen}\|_{\ell_1}(\frac{\ln(N+1)}{K}\sqrt{n_\rve})\right) \\
%% %  \mathcal{C}_{ub,i1}&=0.25 \left(\bar{\phi}_{b,1}(8L \bar{G}_e,G_{e},K)+\right. \\
%%  %       & \left. \phi_{b,1}(8L \|\Sigma_{gen}\|_{\ell_1},K)\right) i=1,2\\
%%  %\mathcal{C}_{ub,i2}&=0.25 \left(\bar{\phi}_{b,i}(8L\bar{G}_e,G_e,K)+ \right.  \\
%%  %       &\left. \bar{\phi}_{b,i}(4L\bar{G}_e,G_{e},K)+ \right.  \\
%%  %      &  \left. \bar{\phi}_{b,i}(8L\|\Sigma_{gen}\|_{\ell_1},K)+\right. \\
%%  %      &  \left. \bar{\phi}_{b,i}(4L\|\Sigma_{gen}\|_{\ell_1},K) \right), ~ i=1,2  \\
%%%\bar{\phi}_{b,1}(s,K_1,K_2) &=\Psi_{\e_g}(K_1)s^2, \\
%%%\bar{\phi}_{b,2}(s,K_1,K_2) &=\frac{s\Psi_{\e_g}(K_1)K_1}{K_2}
%%%& \mathcal{C} =4K\bar{G}_{f,2}+ \ln(\max{\bar{G}_{f,i}K,1})+8 c_\rve^4 n_\rve^2( G_e  + G_{e,1}^2G_e^2) \\
%%\end{split}
%%\]
%%where $C_{b,1},C_{b,2}$ are as in Corollary \ref{thm:bounded:alt:col}
%%%then 
%%%with  and 
%%%$\mathcal{C}_2=4K\bar{G}_{f,2}$, $\mathcal{C}=\mathcal{C}_1+\mathcal{C}_2$ 
%%%it follows that
%%For any prior $\pi$, and %for $\lambda=\sqrt{N}$, 
%%for  $\lambda =\frac{\sqrt{N}}{L \ln(N+1)}$,  and $N \ge 2$
%%% in \eqref{thm:unbounded:alt:col:eq1},
%%\begin{equation}
%%\label{thm:unbounded:alt:col:eq1}
%%\begin{split}
%%     & \frac{1}{\lambda} \Phi(\lambda,\pi,N) \le \frac{1}{\lambda}
%%      \bar{\Phi}(\pi,\lambda,N) \le 2LC_{ub}\frac{\ln(N)}{\sqrt{N}}
%%       \\
%%  & C_{ub}=\mathcal{C}_1+\mathcal{C}_2+\mathcal{C}_3
%%      %& \frac{(L\lambda)^2\mathcal{C}_1}{N}+\frac{(L\lambda) \mathcal{C}_2\ln(N+1)}{N}+ 
%%      %\frac{(\lambda L)^2 \mathcal{C}_3(\ln(N+1))^2}{N}
%%\end{split}      
%%\end{equation}
%%In particular, for any prior $\pi$
%%and for any $\delta \in (0,0.5)$, 
%%%$\Phi(\lambda,\pi,N) \le 2L C_{ub} \frac{1}{N}
%%\begin{equation}
%%\label{catoni:unbounded1}
%%\begin{split}
%%   \forall \rho \in\mathcal{M}_\pi: & E_{\theta \sim \rho } \mathcal{L}(\theta)  \leq E_{\theta\sim \theta } \hat{\mathcal{L}}_N(\theta)+ \\
%%  %& \frac{L \ln(N+1)}{\sqrt{N}}(\KL(\rho ||\pi)+\ln\frac{1}{\delta}) + \\ 
%%  %& \frac{L\mathcal{C}_1}{\sqrt{N}}+\frac{L\mathcal{C}_2 \ln(N+1)}{\sqrt{N}}+ \frac{L\mathcal{C}_3(\ln(N+1)}{\sqrt{N}} \\
%%  %& \le   E_{\theta\sim \theta } \hat{\mathcal{L}}_N(\theta)
%%  & + \frac{\ln(N)}{\sqrt{N}} 2L(\KL(\rho ||\pi)+\ln\frac{1}{\delta}+C_{ub}) \\
%%\end{split}
%%\end{equation}
%%holds with probability at least $1-2\delta$.
%%%\begin{color}{red}
%%%    Moreover, for quadratic loss the constant $\mathcal{C}$ is of the form
%%%    \[
%%%      \begin{aligned}
%%%        & \mathcal{C}=\left(2G_e(\theta)\bar{G}_{f,2}\|\Sigma_{gen}\|_{\ell_1}(c_\rve\sqrt{n_\rve})^2+ \right. \\
%%%        & \left. 2(2G_e(\theta))^2(G_e(\theta)+2G_{e,1}(\theta))(c_{\rve}\sqrt{n_{\rve}})^4 \right). 
%% %     \end{aligned} 
%% %   \]
%%\end{corollary} 
%%%That is, 
%%
%%
%%%$V_N(\theta)-\hat{\mathcal{L}}_N(\theta)$ and 
%%%$\mathcal{L}(\theta)-V_N(\theta)$ by
%%%by $\mathcal{O}(c^r)$ for a suitable constant $c$.
%%%}
%%%$\bE[\|\rve_g(t)\|_2^r]\leq (c_\rve\sqrt{n_\rve})^r$ by Lemma \ref{lemma:bounded_e_moments}. This gives us the desired property that the moments are $\mathcal{O}(c^r)$. One could also obtain $\mathcal{O}(c^r)$ bounds for $\sigma(r)$. 
%%\textcolor{cyan}{The assumptions above allow us to bound the MGF involved in \eqref{pf:thm:initial:eq1} 
%%for all $\lambda>0$, under Assumption \ref{as:bounded}.} Moreover, we can get tighter results, since we can now apply more general tools, i.e., generalisation of the Hoeffding's inequality \cite[Theorem 6.6]{alquier2013prediction}, and recursive Cauchy–Schwarz inequality, see Lemma \ref{lem:repeatedCa




\subsection{Proof of Theorem \ref{thm:main}: quadratic loss}
\label{thm:main:quadratic}
  The core of the proof is to to show that the prediction errors
  $\|\y(t)-\y_{\theta}(t)\|_2$ and $\|\y(t) - \y_{\theta}(t\mid 0)\|_2$ 
  are sub-Gaussian and hence
  they are bounded with high probablity as follows.
  \begin{lemma}
  \label{lem:quadratic1}
  For $N \ge 1$ and $\delta \in (0,0.5)$
  \begin{equation}
\label{lem:quadratic1:eq1}
\begin{aligned}
     & \bP\left(\forall \theta \in \Theta,~  \forall t=0,\ldots, N-1: \quad 
     \right. \\
     & \left. 
          \max\{\|\y(t)-\y_{\theta}(t)\|_2, \|\y(t)-\y_{\theta}(t\mid 0)\|_2\} \le 6G_e (\|\e_g\|_{\psi_2})\sqrt{\ln\left(\frac{3N^2}{2\delta}\right)} \right) \ge 1-\frac{2\delta}{3}
\end{aligned}
  \end{equation}
  and 
  \begin{equation}
\label{lem:quadratic1:eq2}
   \bE[\|\y(t)-\y_{\theta}(t)\|_2^2 \chi(\|\y(t)-\y_{\theta}(t)\|_2^2 > (6G_e\|\e_g\|_{\psi_2})^2\ln\left(\frac{3N^2}{2\delta}\right))] \le 
   \frac{2 \cdot 6^2 G_e^2 \|\e_g\|_{\psi_2}^2\ln(\frac{3N^2}{2\delta})}{N}
\end{equation}
\end{lemma}
We then replace the quadratic loss with a loss which is Lipschitz for 
the case when $\|y(t)-\y_{\theta}(t\mid 0)\|_2$ are bounded. More precisely, let us define the loss function $\ell_C$ as follows:
\[ \ell_C(y, \hat{y}) = \min\{\|y-\hat{y}\|_2^2, C^2\} \]
\begin{lemma}
\label{lem:quadratic2}
For any $C > 0$, $\ell_C$ is $2C$-Lipschitz.
\end{lemma}
The proof then reduces the quadratic case to the case of Lipschitz loss $\ell_C$
for a suitable $C$ and it goes as follows.
\begin{proof}[Proof of Theorem \ref{thm:main}:quadratic case]
      Let us set $C=6G_e\|\e_g\|_{\psi_2}\sqrt{\ln(\frac{3N^2}{2\delta})}$ and let us consider the 
      true and empirical errores for the loss function $\ell_C$:
      \begin{equation}
        \label{eq:initial:tildern1:quadratic0} 
         {}^C\hat{\mathcal{L}}_{N}(\theta) = \frac{1}{N} \sum_{t=0}^{N-1} \ell_C(y(t), \hat{y}_{\theta}(t \mid 0)), \quad 
         {}^C \mathcal{L}(\theta) = \bE[\ell_C(y(t), \hat{y}_{\theta}(t))]
      \end{equation}
      By applying Theorem \ref{thm:main} for the case of Lipschitz loss function $\ell_C$  and
      with $\frac{2\delta}{3}$ instead of $\delta$, it follows
      that for all $\epsilon \in \{-1,1\}$,
      for all $0 < \lambda  < \frac{N}{16 \cdot 2C}$
        \begin{equation}
        \label{eq:initial:tildern1:quadratic} 
    \begin{aligned}
         \forall \rho & \in \mathcal{M}_\pi: \epsilon E_{\theta\sim\rho} {}^C\mathcal{L}(\theta)
         \leq   \epsilon E_{\theta\sim\rho}{}^C\hat{\mathcal{L}}_N(\theta) +  
          \frac{1}{\lambda}\left ( \KL(\rho\|\pi) + \ln\dfrac{3}{2\delta} + \right. \\
          &  \left. \frac{1}{2} 
          \left(\ln E_{\theta\sim\pi} \exp(\mathcal{\bar{C}}_1(2\lambda,\theta,N,\frac{2\delta}{3},\epsilon))+
          \ln E_{\theta \sim \pi} \exp(\mathcal{\bar{C}}_2(2\lambda,\theta,N,\frac{2\delta}{3},\epsilon))\right)
           %\Psi(\lambda,\pi,N,\delta,\epsilon)
           \right )
    \end{aligned}
    \end{equation}
    holds with probability at least $1-\frac{4\delta}{3}$, where  for $i=1,2$.
    \begin{equation}
    \label{thm:main:quadratic:eq2}
       \bar{C}_i(\lambda,\theta,N,\delta,\epsilon)= 
       \mathcal{C}_{ub,i}(\lambda,\theta,12G_e\|\e_g\|_{\psi_2}\sqrt{\ln\left(\frac{3N^2}{2\delta}\right)},N)
    \end{equation}
    By using Lemma \ref{lem:quadratic1}, \eqref{eq:initial:tildern1:quadratic}
    and the union bound, it follows that 
   \begin{equation}
        \label{eq:initial:tildern1:quadratic22} 
    \begin{aligned}
         \forall \rho & \in \mathcal{M}_\pi: \epsilon E_{\theta\sim\rho} {}^C\mathcal{L}(\theta)
         \leq   \epsilon E_{\theta\sim\rho}{}^C\hat{\mathcal{L}}_N(\theta) +  
          \frac{1}{\lambda}\left ( \KL(\rho\|\pi) + \ln\dfrac{3}{2\delta} + \right. \\
          &  \left. \frac{1}{2} 
          \left(\ln E_{\theta\sim\pi} \exp(\mathcal{\bar{C}}_1(2\lambda,\theta,N,\frac{2\delta}{3},\epsilon))+
          \ln E_{\theta \sim \pi} \exp(\mathcal{\bar{C}}_2(2\lambda,\theta,N,\frac{2\delta}{3},\epsilon))\right) \right.\\
          \text{and}&\\
          %& \left. \mbox{ and } 
          &\forall \theta \in \Theta,~\forall t=0,\ldots,N-1: \quad
          \|\y(t)-\hat{\y}_{\theta}(t)\|_2 \le C, \|\y(t)-\hat{\y}_{\theta}(t \mid 0)\|_2 \le C
           %\Psi(\lambda,\pi,N,\delta,\epsilon)
           %\right )
    \end{aligned}
    \end{equation}
    holds with probability at least  $1-(\frac{4\delta}{3}+\frac{2\delta}{3}) = 1-2\delta$ over data.

    Notice that if $\omega \in \Omega$ is such that
     \[ \forall \theta \in \Theta, t=0,\ldots,N-1: \quad
          \|\y(t)(\omega)-\hat{\y}_{\theta}(t)(\omega)\|_2 \le C, \|\y(t)(\omega  )-\hat{\y}_{\theta}(t \mid 0)(\omega)\|_2 \le C
    \]
    then $\ell_C(\y(t)(\omega),\hat{\y}_{\theta}(t \mid 0)(\omega)) = \|\y(t)(\omega)-\hat{\y}_{\theta}(t \mid 0)(\omega)\|_2^2$ and hence
    \[ \hat{\mathcal{L}}_N(\theta)(\omega) = {}^C\hat{\mathcal{L}}_N(\theta)(\omega) \]
    Therefore, the event defined by  \eqref{eq:initial:tildern1:quadratic22} equals the event defined by 
     \begin{equation}
        \label{eq:initial:tildern1:quadratic2} 
    \begin{aligned}
         \forall \rho & \in \mathcal{M}_\pi: \epsilon E_{\theta\sim\rho} {}^{C}\mathcal{L}(\theta)
         \leq   \epsilon E_{\theta\sim\rho}{} \hat{\mathcal{L}}_N(\theta) +  
          \frac{1}{\lambda}\left ( \KL(\rho\|\pi) + \ln\dfrac{3}{2\delta} + \right. \\
          &  \left. \frac{1}{2} 
          \left(\ln E_{\theta\sim\pi} \exp(\mathcal{\bar{C}}_1(2\lambda,\theta,N,\dfrac{2\delta}{3},\epsilon))+
          \ln E_{\theta \sim \pi} \exp(\mathcal{\bar{C}}_2(2\lambda,\theta,N,\dfrac{2\delta}{3},\epsilon)) \right)  \right. \\
          \text{and}&\\%& \left. \mbox{ and } 
          &\forall \theta \in \Theta,~\forall t=0,\ldots,N-1: \quad
          \|\y(t)-\hat{\y}_{\theta}(t)\|_2 \le C, \|\y(t)-\hat{\y}_{\theta}(t \mid 0)\|_2 \le C
           %\Psi(\lambda,\pi,N,\delta,\epsilon)
    \end{aligned}
    \end{equation}

  For $\epsilon=-1$
  we will show \eqref{eq:initial:tildern1:quadratic2} implies
\begin{equation}
        \label{eq:initial:tildern1:quadratic3} 
    \begin{aligned}
         \forall \rho & \in \mathcal{M}_\pi: \epsilon E_{\theta\sim\rho} \mathcal{L}(\theta)
         \leq   \epsilon E_{\theta\sim\rho} \hat{\mathcal{L}}_N(\theta) +  
          \frac{1}{\lambda}\left ( \KL(\rho\|\pi) + \ln\dfrac{3}{2\delta} + \right. \\
          &  \left. \frac{1}{2} 
          \left(\ln E_{\theta\sim\pi} \exp(\mathcal{\bar{C}}_1(2\lambda,\theta,N,\frac{2\delta}{3},\epsilon))+
          \ln E_{\theta \sim \pi} \exp(\mathcal{\bar{C}}_2(2\lambda,\theta,N,\frac{2\delta}{3},\epsilon))\right) \right)
          %& \mbox{ and } \forall \theta \in \Theta, t=0,\ldots,N-1: \quad
          %\|\y(t)-\hat{\y}_{\theta}(t)\|_2 \le C, \|\y(t)-\hat{\y}_{\theta}(t \mid 0)\|_2 \le C
           %\Psi(\lambda,\pi,N,\delta,\epsilon)
          % \right )
    \end{aligned}
    \end{equation}
   and as \eqref{eq:initial:tildern1:quadratic2} holds with probability at least $1-2\delta$, it follows that \eqref{eq:initial:tildern1:quadratic3} also holds with probability at least $1-2\delta$. 
   To this end,   notice that 
    \begin{multline}
    \label{thm:main:quadratic:eq24}
     %\begin{aligned}
        \mathcal{L}(\theta) = \bE{[\|\y(t)-\hat{\y}_{\theta}(t)\|_2^2]}= 
       \bE[\|\y(t)-\hat{\y}_{\theta}(t)\|_2^2\chi(\|\y(t)-\hat{\y}_{\theta}(t)\|_2 \le C)]+ \\
        \bE[\|\y(t)-\hat{\y}_{\theta}(t)\|_2^2\chi(\|\y(t)-\hat{\y}_{\theta}(t)\|_2 > C)]
     %\end{aligned}
    \end{multline}
    and that $\bE[\|\y(t)-\hat{\y}_{\theta}(t)\|_2^2\chi(\|\y(t)-\hat{\y}_{\theta}(t)\|_2 > C)] \ge C^2\bP(\|\y(t)-\hat{\y}_{\theta}(t)\|_2 > C)$. Hence
    \begin{equation}
    \label{thm:main:quadratic:eq210}
    \begin{aligned}
     \mathcal{L}(\theta)& 
    % = \bE{[\|\y(t)-\hat{\y}_{\theta}(t)\|_2^2]}=  \bE[\|\y(t)-\hat{\y}_{\theta}(t)\|_2^2\chi(\|\y(t)-\hat{\y}_{\theta}(t)\|_2 \le C)]+ \\
    % & \bE[\|\y(t)-\hat{\y}_{\theta}(t)\|_2^2\chi(\|\y(t)-\hat{\y}_{\theta}(t)\|_2 > C)] 
    \ge  
         \bE[\|\y(t)-\hat{\y}_{\theta}(t)\|_2^2\chi(\|\y(t)-\hat{\y}_{\theta}(t)\|_2 \le C)]+ 
          C^2\bP(\|\y(t)-\hat{\y}_{\theta}(t)\|_2 > C)\\
          &= 
    \bE[\ell_C(\y(t),\hat{\y}_{\theta}(t))]={}^C \mathcal{L}(\theta)
    \end{aligned}
    \end{equation}
  i.e., $\epsilon \mathcal{L}(\theta) \le \epsilon {}^C \mathcal{L}(\theta)$
  for $\epsilon=-1$, from which it follows that \eqref{eq:initial:tildern1:quadratic2}
  implies \eqref{eq:initial:tildern1:quadratic3}.

   
  For $\epsilon=1$, we will show that \eqref{eq:initial:tildern1:quadratic2} implies
   \begin{equation}
        \label{eq:initial:tildern1:quadratic4} 
    \begin{aligned}
         \forall \rho & \in \mathcal{M}_\pi: \epsilon E_{\theta\sim\rho} \mathcal{L}(\theta)
         \leq   \epsilon E_{\theta\sim\rho} \hat{\mathcal{L}}_N(\theta) +   \frac{72G_e^2\|\e_g\|_{\psi_2}^2\ln\left(\frac{3N^2}{2\delta}\right)}{N} + %\\
         %& 
          \frac{1}{\lambda}\left ( \KL(\rho\|\pi) + \ln\dfrac{3}{2\delta} + 
          \right. \\
          &\left. 
          \frac{1}{2} 
          \left(\ln E_{\theta\sim\pi} \exp(\mathcal{\bar{C}}_1(2\lambda,\theta,N,\dfrac{2\delta}{3},\epsilon))+
          \ln E_{\theta \sim \pi} \exp(\mathcal{\bar{C}}_2(2\lambda,\theta,N,\dfrac{2\delta}{3},\epsilon))\right) 
          %& \mbox{ and } \forall \theta \in \Theta, t=0,\ldots,N-1: \quad
          %\|\y(t)-\hat{\y}_{\theta}(t)\|_2 \le C, \|\y(t)-\hat{\y}_{\theta}(t \mid 0)\|_2 \le C
           %\Psi(\lambda,\pi,N,\delta,\epsilon)
           \right )
    \end{aligned}
    \end{equation}
   First note that 
   \[ \frac{72G_e^2\|\e_g\|_{\psi_2}^2\ln\left(\frac{3N^2}{2\delta}\right)}{N}=\frac{1}{\lambda}\ln\left(\exp\left(\lambda \frac{72G_e^2\|\e_g\|_{\psi_2}^2\ln\left(\frac{3N^2}{2\delta}\right)}{N}\right)\right) 
   \]
   and hence after algebraic manipulations %\eqref{eq:initial:tildern1:quadratic4}
   %becomes \eqref{eq:initial:tildern1}, i.e., 
   \begin{align*}
    & \frac{72G_e^2\|\e_g\|_{\psi_2}^2\ln\left(\frac{3N^2}{2\delta}\right)}{N} +
    \frac{1}{\lambda}\left ( \KL(\rho\|\pi) + \ln\dfrac{3}{2\delta} + 
            \right. \\
          &\left. 
          \frac{1}{2} 
          \left(\ln E_{\theta\sim\pi} \exp(\mathcal{\bar{C}}_1(2\lambda,\theta,N,\dfrac{2\delta}{3},\epsilon))+
          \ln E_{\theta \sim \pi} \exp(\mathcal{\bar{C}}_2(2\lambda,\theta,N,\dfrac{2\delta}{3},\epsilon))\right)\right) = \\
    & \frac{1}{\lambda}\left ( \KL(\rho\|\pi) + \ln\dfrac{3}{2\delta} + 
          %  \right. \\
          %&\left. 
          \frac{1}{2} 
          \left(\ln E_{\theta\sim\pi} \exp(\mathcal{C}_1(2\lambda,\theta,N,\delta,\epsilon))+
          \ln E_{\theta \sim \pi} \exp(\mathcal{C}_2(2\lambda,\theta,N,\delta,\epsilon))\right) \right)
   \end{align*}
    where for $i=1,2$, 
   \[ \mathcal{C}_i(\lambda,\theta,N,\delta,\epsilon)=
      \lambda \frac{72G_e^2\|\e_g\|_{\psi_2}^2\ln\left(\frac{3N^2}{2\delta}\right)}{N}
      + \mathcal{C}_{ub,i}(\lambda,\theta,12G_e\|\e_g\|_{\psi_2}\sqrt{\ln\left(\frac{3N^2}{2\delta}\right)},N)
   \]   
   % Since \eqref{eq:initial:tildern1:quadratic2}
   % holds with probability $1-2\delta$, and it implies  \eqref{eq:initial:tildern1:quadratic4}, it follows that \eqref{eq:initial:tildern1} also holds with probability at least $1-2\delta$.


   To show that \eqref{eq:initial:tildern1:quadratic2} implies   \eqref{eq:initial:tildern1:quadratic4}, notice that 
   \begin{equation}
    \label{thm:main:quadratic:eq22}
     \begin{aligned}
      & \mathcal{L}(\theta) = \bE{[\|\y(t)-\hat{\y}_{\theta}(t)\|_2^2]}= \\
      & \bE[\ell_C(\y(t),\hat{\y}_{\theta}(t))\chi(\|\y(t)-\hat{\y}_{\theta}(t)\|_2 \le C)]+\bE[\|\y(t)-\hat{\y}_{\theta}(t)\|_2^2\chi(\|\y(t)-\hat{\y}_{\theta}(t)\|_2 > C)] \le \\
      & \bE[\ell_C(\y(t),\hat{\y}_{\theta}(t))]+\bE[\|\y(t)-\hat{\y}_{\theta}(t)\|_2^2\chi(\|\y(t)-\hat{\y}_{\theta}(t)\|_2 > C)] \\
      &= {}^C \mathcal{L}(\theta)+ \bE[\|\y(t)-\hat{\y}_{\theta}(t)\|_2^2\chi(\|\y(t)-\hat{\y}_{\theta}(t)\|_2 > C)]
    \end{aligned}
    \end{equation}
    Using %Lemma \ref{lem:quadratic2}, 
    \eqref{lem:quadratic1:eq2} it
    follows that 
    \begin{equation}
    \label{thm:main:quadratic:eq21}
     \begin{aligned}
       & \mathcal{L}(\theta) \le 
       {}^C \mathcal{L}(\theta)+ 72G_e^2\|\e_g\|_{\psi_2}^2\frac{\ln\left(\frac{3N^2}{2\delta}\right)}{N}
    \end{aligned}
    \end{equation}
    Hence, by adding  $72G_e^2\|\e_g\|_{\psi_2}^2\frac{\ln\left(\frac{3N^2}{2\delta}\right)}{N}$ to both sides of the inequality \eqref{eq:initial:tildern1:quadratic2}, it follows that \eqref{eq:initial:tildern1:quadratic4} is implied by \eqref{eq:initial:tildern1:quadratic2}. Combining the two case it follows that for quadratic loss and $\epsilon\in\{-1,1\}$, the inequality \eqref{eq:initial:tildern1} also holds with probability at least $1-2\delta$ over data.
  \end{proof}




We will conclude by presenting the proof of Lemma \ref{lem:quadratic1}-\ref{lem:quadratic2}. To this end, we need the following technical result. 
  \begin{lemma}
    \label{lem:subGaussian1}
%     Let $\mathbf{I}$ be an index set, and let $\{\rvz_i\}_{i \in \mathbf{I}}$ be a collection of stochastic process
% such that for all $i \in \mathbb{I}$, $\rvz_{i}(t)=(\beta_{i,t} \star \e_g)(t)$ and
% $\bar{\rvz}(t)=(\beta_{i,t} \star \bar{\e}_g)(t)$
% %\sum_{k=0}^{\infty} \beta_{i,t}(k) \bar{\e}_g(t-k)$, for 
% suitable matrices $\{\beta_{i,t}(k)\}_{k=0}^{\infty}$, $i \in \mathbf{I}$, 
% such that %$\sup_{i \in \mathbf{I}} \|\beta_i\|_{\ell_1} <  +\infty$, where
% %for 
% $\|\bar{\beta}\|_{\ell_1}< +\infty$, where $\bar{\beta}(k)=\sup_{i \in \mathbf{I},t \in \mathbb{Z}} \|\beta_{i,t}(k)\|_2$.
% %$\bar{H}(\underline{\beta}):=\sum_{k=0}^{\infty}   +\infty$. 
% Then $\|\z_i(t)\|_{2}$, $i \in \mathbb{I}$ and 
% $\z(t)=\sup_{i \in \mathbf{I}} \|\rvz_i(t)\|_2$ are sub-Gaussian with 
% sub-Gaussian norm 
    Let $\mathbf{I}$ be an index set, $\{\beta_{i,t}(k)\}_{k=0}^{\infty}$, $i \in \mathbf{I}$, be a sequence of matrices such that $\|\bar{\beta}\|_{\ell_1}< \infty$, where $\bar{\beta}(k)=\sup_{i \in \mathbf{I},t \in \mathbb{Z}} \|\beta_{i,t}(k)\|_2$, and let $\rvz_{i}(t)=(\beta_{i,t} \star \e_g)(t)$. 
    %Let $\rvz_{i}(t)=(\beta_{i,t} \star \e_g)(t)$ and
%$\bar{\rvz}(t)=(\beta_{i,t} \star \bar{\e}_g)(t)$. 
Then $\|\z_i(t)\|_{2}$, $i \in \mathbb{I}$ and 
$\z(t)=\sup_{i \in \mathbf{I}} \|\rvz_i(t)\|_2$ are sub-Gaussian with 
sub-Gaussian norm 
\begin{equation}
\label{lem:subGaussian1:eq0}
    \|\|\z_i(t)\|_2\|_{\psi_2} \le \|\z(t)\|_{\psi_2} \le 6\|\bar{\beta}\|_{\ell_1}\|\e_g\|_{\psi_2}
\end{equation}
and the moments and the moment generating function of $\z^2(t)=\sup_{i \in \mathbf{I}} \|\rvz_i(t)\|_2^2$ satisfy
\begin{align}
   & \bE[\z^{2r}(t)] \le \|\bar{\beta}\|_{\ell_1}^{2r} (3\|\e_g\|_{\psi_2})^{2r} 2r!
      \label{lem:subGaussian1:eq1} \\
   & \bE[e^{\lambda \z^2(t)}] \le \frac{1+\lambda \|\bar{\beta}\|_{\ell_1}^2 (3\|\e_g\|_{\psi_2})^2}{1-\lambda \|\bar{\beta}\|_{\ell_1}^2 (3\|\e_g\|_{\psi_2})^2},  \quad |\lambda| < \frac{1}{\|\bar{\beta}\|_{\ell_1}^2 (3\|\e_g\|_{\psi_2})^2}
   \label{lem:subGaussian1:eq2}
\end{align}
Moreover, for any $C > 0$,
\begin{equation}
\label{lem:subGaussian1:eq4}
   \bE[\|\z_i(t)\|_2^2 \chi(\|\z_i(t)\|_2^2 > C^2)] \le 
   \frac{2C^2}{e^{\frac{C^2}{(6\|\bar{\beta}\|_{\ell_1} \|\e_g\|_{\psi_2})^2}}-1}
\end{equation}
\end{lemma}
\begin{proof}[Proof Lemma \ref{lem:subGaussian1}]
  Notice that 
  \begin{align*}
      \z(t) &= \sup_{i \in \mathbf{I}} \|\sum_{k=0}^{\infty} \beta_{i,t}(k) \e_g(t-k)\|_2\\ 
      &\le \sup_{i \in \mathbf{I}}\sum_{k=0}^{\infty} \|\beta_{i,t}(k)\|_2 \|\e_g(t-k)\|_2 \le \sum_{k=0}^{\infty} \|\bar{\beta}(k)\|_2 \|\e_g(t-k)\|_2 
     % \\
     % &  \|\bar{\beta}\|_{\ell_1} \|\e_g\|_{\psi_2} \le 
     % \sum_{k=0}^{\infty} \|\bar{\beta}(k)\|_2 \|\e_g(t-k)\|_2
  \end{align*}
  Hence, for any $r \ge 1$, by using the arithmetic-geometric mean inequality we get
  $a_1 \cdots a_k \le \frac{1}{k} \sum_{i=1}^k a_i^k$, $a_1,\ldots, a_k \ge 0$, 
  \begin{align*}
    \z^{2r}(t) \le \left( \sum_{k=0}^{\infty} \|\bar{\beta}(k)\|_2 \|\e_g(t-k)\|_2 \right)^{2r}=\sum_{k_1,\ldots,k_{2r}=0}^{\infty} \prod_{j=1}^{2r} \|\bar{\beta}(k_j)\|_2 \|\e_g(t-k_j)\|_2 \le \\
    \sum_{k_1,\ldots,k_{2r}=0}^{\infty} \frac{1}{2r} \left(\sum_{j=1}^{2r} \|\bar{\beta}(k_j)\|_2^{2r} \right)  \frac{1}{2r} \sum_{j=1}^{2r}\|\e_g(t-k_j)\|_2^{2r}
  \end{align*}
  By taking expectations, using the fact that as $\|\e_g(t)\|_2$ is sub-Gaussian,
  and the proof of \cite[Proposition 2.6.1]{vershynin2026highdimensional} we obtain
  \begin{equation} 
   \label{lem:subGaussian1:pf:eq1} 
      \bE[\|\e_g(t-k_j)\|_2^{2r}]=\bE[\|\e_g(t)\|_2^{2r}] \le 2r(r-1)! (3\|\e_g\|_{\psi_2})^{2r}=2 r! (3\|\e_g\|_{\psi_2})^{2r}
  \end{equation} 
  it follows that
  \begin{align*}
    & \bE[\z^{2r}(t)] \le \sum_{k_1,\ldots,k_{2r}=0}^{\infty} \left(\sum_{j=1}^{2r} \|\bar{\beta}(k_j)\|_2^{2r} \right)  \frac{1}{2r} \sum_{j=1}^{2r}\bE[\|\e_g(t-k_j)\|_2^{2r}] = \\
    & \sum_{k_1,\ldots,k_{2r}=0}^{\infty} \left(\sum_{j=1}^{2r} \|\bar{\beta}(k_j)\|_2^{2r} \right) \bE[\|\e_g(t)\|_2^{2r}] \le 
    \sum_{k_1,\ldots,k_{2r}=0}^{\infty} \left(\sum_{j=1}^{2r} \|\bar{\beta}(k_j)\|_2^{2r} \right) 2 r! (3\|\e_g\|_{\psi_2})^{2r}= \\
    & (\sum_{k=0}^{\infty} \|\bar{\beta}(k)\|_2)^{2r} 2 r! (3\|\e_g(t)\|_{\psi_2})^{2r}=
    \|\bar{\beta}\|_{\ell_1}^{2r} 2 r! (3\|\e_g\|_{\psi_2})^{2r}
  \end{align*}
  i.e., \eqref{lem:subGaussian1:eq1} holds.
  By noticing that
  \begin{align*}
     & \bE[e^{\lambda \z^2(t)}]=1+\sum_{r=1}^{\infty} \frac{\lambda^r \bE[\z^{2r}(t)]}{r!} \le 1+\sum_{r=1}^{\infty} \frac{\lambda^r \|\bar{\beta}\|_{\ell_1}^{2r} (3\|\e_g\|_{\psi_2})^{2r} 2r!}{r!} = \\
     &  1+ 2\lambda \|\bar{\beta}\|_{\ell_1}^2(3\|\e_g(t)\|_{\psi_2})^2 \sum_{r=0}^{\infty} (\lambda \|\bar{\beta}\|_{\ell_1}^2 (3\|\e_g\|_{\psi_2})^2)^r =
      1 + \frac{2 \lambda \|\bar{\beta}\|_{\ell_1}^2(3\|\e_g(t)\|_{\psi_2})^2}{1- \lambda \|\bar{\beta}\|_{\ell_1}^2 (3\|\e_g\|_{\psi_2})^2}= \\
      & \frac{1+\lambda \|\bar{\beta}\|_{\ell_1}(3\|\e_g(t)\|_{\psi_2})^2}{1- \lambda \|\bar{\beta}\|_{\ell_1}^2 (3\|\e_g\|_{\psi_2})^2}
\end{align*}
it follows that \eqref{lem:subGaussian1:eq2} holds. 
Finally, for $\lambda = \frac{1}{4\|\bar{\beta}\|_{\ell_1}^2 (3\|\e_g\|_{\psi_2})^2}$, the bound becomes
\[
\bE[e^{\lambda \z^2(t)}] \le \frac{1+\frac{1}{4}}{1-\frac{1}{4}} = \frac{5/4}{3/4} = \frac{5}{3} \le 2
\]
from which by the definition of sub-Gaussian norm of $\z(t)$ it follows
\[
\|\z(t)\|_{\psi_2} \le \sqrt{\frac{1}{\lambda}} = 6\|\bar{\beta}\|_{\ell_1} \|\e_g\|_{\psi_2}
\]
i.e., \eqref{lem:subGaussian1:eq0} holds.
In particular, the sub-Gaussian norm of $\z(t)$ is bounded
and hence by \cite[Proposition 2.6.1]{vershynin2026highdimensional} 
$\z(t)$ is sub-Gaussian. Moreover, as 
$\z^2(t) \ge \|\z_i(t)\|_2^2$, it follows that the sub-Gaussian norm of
$\|\z_i(t)\|_2$ is also bounded and  it is also sub-Gaussian. 

Finally, from Lemma \ref{noise:boundC-1} applied to $\|\z_i(t)\|_2$
it follows that 
\begin{equation}
   \label{lem:subGaussian1:pf:eq01} 
\bE[\|\z_i(t)\|_2^2 \chi(\|\z_i(t)\|_2^2 > C^2)] \le 
\frac{\bE[e^{\bar{\lambda}\|\z_i(t)\|_2^2}]C^2}{e^{C^2\bar{\lambda}}-1}
\end{equation}
if $\bar{\lambda}$ is such that the moment generating function 
$\bE[e^{s\|\z_i(t)\|_2^2}]$ is well-defined on $[0,\bar{\lambda}]$.
Note that $\|\z_i(t)\|^2_2 \le \z^2(t)$ and hence 
$\bE[e^{s\|\z_i(t)\|^2}] \le \bE[e^{s\z^2(t)}]$. Now  $\bar{\lambda}$ can be chosen as $\bar{\lambda}=\frac{1}{4\|\bar{\beta}\|_{\ell_1}^2 (3\|\e_g\|_{\psi_2})^2}$
and hence $\bE[e^{\bar{\lambda}\|\z_i(t)\|^2}] \le \bE[e^{\bar{\lambda}\z^2(t)}] \le 2$.  Combining \eqref{lem:subGaussian1:pf:eq01} and the bound on $\bE[e^{\bar{\lambda}\|\z_i(t)\|^2}]$, we obtain \eqref{lem:subGaussian1:eq4}.
\end{proof}

\begin{proof}[Proof of Lemma \ref{lem:quadratic1}]
    Let $\mathbb{I}=\Theta \times \{1,2\}$,
    $\z_{(\theta,1)}(t)=\y(t)-\y_{\theta}(t)$, 
    $\z_{(\theta,2)}(t)=\y(t)-\y_{\theta}(t \mid 0)$,
    $\beta_{(\theta,1),t}=\alpha_{g,y}-c_{\theta}$, and $\beta_{(\theta,2),t}=\alpha_{g,y}-c_{\theta,t}$.
    Note that from the proof of Lemma \ref{alpha:lemma2} it follows
    that 
    \begin{align*}
     \|(\alpha_{g,y}+c_{\theta})(k)\|_2 
    &\le \|\alpha_{g,y}(k)\|_2+\sum_{l=0}^{k} \|\alpha_{\theta}(l)\|_2\|\alpha_{g,w}(k-l)\|_2 \\
    &\le\|\alpha_{g,y}(k)\|_2+ \sum_{l=0}^{k} \bar{\alpha}(l)\|\alpha_{g,w}(k-l)\|_2 \\
     \|(\alpha_{g,y}+c_{\theta,t})(k)\|_2 
     &\le \|\alpha_{g,y}(k)\|_2+\sum_{l=0}^{\min\{k,t\}} \|\alpha_{\theta}(l)\|_2\|\alpha_{g,w}(k-l)\|_2\\ 
     &\le \|\alpha_{g,y}(k)\|_2+\sum_{l=0}^{k} \bar{\alpha}(l)\|\alpha_{g,w}(k-l)\|_2.
    \end{align*}
    Hence $\bar{\beta}(k)=\sup_{i \in \mathbb{I}} \|\beta_{i,t}(k)\|_2 \le \|\alpha_{g,y}(k)\|_2+\sum_{l=0}^{k} \bar{\alpha}(l)\|\alpha_{g,w}(k-l)\|_2$ and 
    \begin{align*}
    & \|\bar{\beta}\|_{\ell_1} \le \sum_{k=0}^{\infty} \|\alpha_{g,y}(k)\|_2 + \sum_{k=0}^{\infty} \sum_{l=0}^{k} \bar{\alpha}(l)\|\alpha_{g,w}(k-l)\|_2 \le 
\|\alpha_g\|_{\ell_1}+\|\bar{\alpha}\|_{\ell_1}\|\alpha_{g}\|_{\ell_1}=G_e
    \end{align*}
    From this and Lemma \ref{lem:subGaussian1} and Lemma \ref{alpha:lemma2}
    it follows that for any $t=0,\ldots,N-1$
    $\sup_{\theta \in \Theta} \max\{\|\y(t)-\y_{\theta}(t)\|_2, \|\y(t)-\y_{\theta}(t \mid 0)\|_2\}$ are sub-Gaussian random variable
    whose sub-Gaussian norms are upper bounded by 
    $6G_e \|\e_g\|_{\psi_2}$.
    Hence, using the sub-Gaussian tail bound (proof of (iii)$\implies$(i) in \cite[Proposition 2.6.1]{vershynin2026highdimensional}), we obtain the desired result
    \begin{multline*}
        \bP\left( 
          \sup_{\theta \in \Theta} \max\{\|\y(t)-\y_{\theta}(t)\|_2, \|\y(t)-\y_{\theta}(t \mid 0)\|_2\} \ge 6G_e \|\e_g\|_{\psi_2}\sqrt{\ln(\frac{3N^2}{2\delta})} \right) \le \\
           %e^{-\frac{\ln(\frac{3N^2}{2\delta})}{6^2G_e^2 (\|\e_g\|_{\psi_2})^2}} = 
           =2e^{-\ln(\frac{3N^2}{2\delta})} = \frac{2}{N}\frac{2\delta}{3N} \le \frac{2\delta}{3N} 
    \end{multline*}
    Therefore, by applying a union bound over $t=0,\ldots, N-1$, the first inequality
    \eqref{lem:quadratic1:eq1} of 
    the lemma follows. 

    In order to obtain the second inequality, we use 
    \eqref{lem:subGaussian1:eq4},
    %Lemma~\ref{noise:boundC-1},
    %\eqref{lem:quadratic1:eq2}, 
    and notice that 
    \begin{align*}
        \frac{2C^2}{e^{\frac{C^2}{(6G_e \|\e_g\|_{\psi_2})^2}}-1} &=
          \frac{2(6G_e \|\e_g\|_{\psi_2})^2\ln(\frac{3N^2}{2\delta})}
          {e^{\ln(\frac{3N^2}{2\delta})}-1} = 
         \frac{2 (6G_e\|\e_g\|_{\psi_2})^2\ln(\frac{3N^2}{2\delta}) 2\delta}{3N^2-2\delta}\\
         &\le 72 G_e^2\|\e_g\|_{\psi_2}^2\frac{\ln(\frac{3N^2}{2\delta})}{N}
    \end{align*}
  where we used that $2\delta \le 1$, $3N^2-2\delta \ge N^2 \ge N$, i.e.,
  \eqref{lem:quadratic1:eq2}.
\end{proof}


\begin{proof}[Proof of Lemma \ref{lem:quadratic2}]
  Notice that
\begin{equation}
  \label{lem:quadratic2:eq1}
  \begin{aligned}
   |\ell_C(y_1, \hat{y}_1) - & \ell_C(y_2, \hat{y}_2)|=  \\
   & \left \{  
    \begin{array}{ll}
    |\|y_1-\hat{y}_1\|_2^2 - \|y_2-\hat{y}_2\|_2^2| & \|y_1-\hat{y}_1\|_2^2\leq C^2, \quad \|y_2-\hat{y}_2\|_2^2 \le C^2 \\
    |\|y_2 - \hat{y}_2\|_2^2 - C^2| & \|y_1-\hat{y}_1\|_2^2 > C^2, \quad  \|y_2-\hat{y}_2\|_2^2 \le C^2 \\
    |\|y_1 - \hat{y}_1\|_2^2 - C^2| & \|y_2-\hat{y}_2\|_2^2 > C^2, \quad \|y_1-\hat{y}_1\|_2^2 \le C^2 \\
    0  & \|y_1-\hat{y}_1\|_2^2 > C^2, \quad  \|y_2-\hat{y}_2\|_2^2 > C^2 \\ 
    \end{array}
   \right.
  \end{aligned}
\end{equation}
 For each of the cases above we will show that the righ-hand side of 
 \eqref{lem:quadratic2:eq1} is bounded by $2C (\|y_1-y_2\|_2+\|\hat{y}_1-\hat{y}_2\|_2)$.

 First, consider the case when $\|y_1-\hat{y}_1\|_2^2\leq C^2$ and $\|y_2-\hat{y}_2\|_2^2 \le C^2$.
 Then by repeated use of the triangle inequality,
 \begin{align*}
 & |\|y_1-\hat{y}_1\|_2^2 - \|y_2-\hat{y}_2\|_2^2|= 
 |\|y_1-\hat{y}_1\|_2 - \|y_2-\hat{y}_2\|_2|(\|y_1-\hat{y}_1\|_2 + \\
 & \|y_2-\hat{y}_2\|_2) \le \|(y_1-y_2)-(\hat{y}_1-\hat{y}_2)\|_2 (2C) \le 
 2C (\|y_1-y_2\|_2+\|\hat{y}_1-\hat{y}_2\|_2)
 \end{align*}

  Consider the case when $\|y_1-\hat{y}_1\|_2^2 > C^2, \|y_2-\hat{y}_2\|_2^2 \le C^2$.
  Then by the definition of $\ell_C$,
 %  \begin{align*}
 % &  |\ell_C(y_1, \hat{y}_1) - \ell_C(y_2, \hat{y}_2)| = \|\|y_1-\hat{y}_1\|_2^2 - C^2| C^2 \|y_1-\hat{y}_1\|_2^2  = \\
 %  & (C - \|y_1-\hat{y}_1\|_2 )(\|y_1-\hat{y}_1\|_2 + C) \le 2C (\|y_1-\hat{y}_1\|_2 - C) \le \\
 %  &  2C (\|y_2-\hat{y}_2\|_2-\|y_1-\hat{y}_1\|_2)=
 %  2C\|(y_2-y_1)-(\hat{y}_2-\hat{y}_1)\|_2 \le
 %  2C (\|y_1-y_2\|_2+\|\hat{y}_1-\hat{y}_2\|_2)
 %  \end{align*}
   \begin{align*}
 &  |\ell_C(y_1, \hat{y}_1) - \ell_C(y_2, \hat{y}_2)| = |C^2 - \|y_2-\hat{y}_2\|_2^2| = \\
  & (C - \|y_2-\hat{y}_2\|_2 )(\|y_2-\hat{y}_2\|_2 + C) \le (C-\|y_2-\hat{y}_2\|_2) 2C  \le \\
  &  2C (\|y_1-\hat{y}_1\|_2-\|y_2-\hat{y}_2\|_2)=
  2C\|(y_2-y_1)-(\hat{y}_2-\hat{y}_1)\|_2 \le
  2C (\|y_1-y_2\|_2+\|\hat{y}_1-\hat{y}_2\|_2)
  \end{align*}
  The case when $\|y_2-\hat{y}_2\|_2^2 > C^2, \|y_1-\hat{y}_1\|_2^2 \le C^2$ is symmetric.

  Finally, consider the case when $\|y_1-\hat{y}_1\|_2^2 > C^2, \|y_2-\hat{y}_2\|_2^2 > C^2$.
  Then by the definition of $\ell_C$,
  \begin{align*}
  |\ell_C(y_1, \hat{y}_1) - \ell_C(y_2, \hat{y}_2)| &= |C^2 - C^2| = 0 \le 2C (\|y_1-y_2\|_2+\|\hat{y}_1-\hat{y}_2\|_2)
  \end{align*}
  This completes the proof.
  \end{proof}
  


  \subsection{Proof of Corollary \ref{thm:bounded:alt:col}}
  \label{sec:proof_corollary}
    The corollary follows from Theorem \ref{thm:main} by choosing $\lambda$ as follows:
    \begin{equation}
      \label{thm:bounded:alt:col:lambda}
      \lambda=\frac{\sqrt{N}}{L(\ln(N))^r} = \left\{\begin{array}{ll}
           \sqrt{N} & \e_g \mbox{ is bounded } \\
           \frac{\sqrt{N}}{16L_{\ell}\ln(N)} & \ell \mbox{ is } L_{\ell} \mbox{-Lipschitz} \\
           \frac{\sqrt{N}}{192G_e\|\e_g\|_{\psi_2}\sqrt{1+\ln\left(\frac{3}{2\delta}\right)}(\ln(N))^{3/2}} & \ell \mbox{is quadratic}
      \end{array}\right. 
    \end{equation}
    and upper bounding the right-hand side of \eqref{eq:initial:tildern1}, i.e., by showing that
    \begin{equation}
    \label{thm:bounded:alt:col:pf:eq11}
     \begin{aligned}
          & \frac{1}{2} \left(   \ln E_{\theta \sim \pi} \exp(\mathcal{C}_1(\lambda,\theta,N,\delta,\epsilon)) +
           \ln E_{\theta \sim \pi} \exp(\mathcal{C}_2(\lambda,\theta,N,\delta,\epsilon)) 
          \right)  \le 
          \mathcal{C}(\delta,\epsilon)
          %& \le \frac{1}{2} \left( \mathcal{C}_{b,m,2}(1,\bar{\alpha},G_e+2G_{e,1},L_b,c_{\e}) + \mathcal{C}_{b,m,2}(1,\bar{\alpha},G_e+2G_{e,1},L_b,c_{\e}) \right) = 
      \end{aligned}
    \end{equation}
    
    
    
    
    To this end, note that 
    from Lemma \ref{alpha:lemma2}
    \begin{equation}
      \label{thm:bounded:alt:col:pf:eq01}
      \begin{aligned}
       & \|\alpha_{\theta}\|_{\ell_1} \le \|\bar{\alpha}\|_{\ell_1}, \quad \theta_{\infty}(\alpha_{\theta}) \le \theta_{\infty}(\bar{\alpha}), \quad
        G_e(\theta) \le G_e, \quad  G_{e,1}(\theta) \le G_{e,1}
      \end{aligned}
    \end{equation} 
    In addition, notice that %for any $\lambda \le \bar{\lambda}\sqrt{N}$, $L_1 \le L_2$, and for any $i=1,2$, 
     the functions
    $\mathcal{C}_{m,b,i}(\lambda,\alpha,X,L,c)$ 
    %are monotonically increasing in $\lambda,X,L,c$ and in $\alpha$ in the sense that 
    has the following property;
    if $\lambda \le \bar{\lambda}\sqrt{N}$, $L_1 \le L_2$, $X_1 \le X_2$, $c_1 \le c_2$, and $\|\alpha(k)\|_2 \le \|\beta(k)\|_2$ for all $k \in \mathbb{N}$, then 
    \begin{equation}
      \label{thm:bounded:alt:col:pf:eq2}
      \begin{aligned}
       \mathcal{C}_{m,b,i}(\lambda,\alpha,X_1,L_1,c_1) \le N\mathcal{C}_{m,b,i}(\bar{\lambda},\beta,X_2,L_2,c_2)
       %\mathcal{C}_{m,b,i}(\lambda,\alpha_{\theta},G_e(\theta)+2G_{e,1}(\theta),L_1,c) 
       %\le  \mathcal{C}_{m,b,i}(\bar{\lambda},\alpha_{\theta},G_e+2G_{e,1},L_2,c_{\e})
      \end{aligned}
    \end{equation}



    \paragraph{Bounded case}
      Note that $L(\theta) \le L_b$,
      any by using 
      by using \eqref{thm:bounded:alt:col:pf:eq01}, it follows that 
       $\|\alpha_{\theta}(k)\|_2 \le \bar{\alpha}(k)=\|\bar{\alpha}(k)\|_2$, $G_e(\theta) \le G_e$, $G_{e,1}(\theta) \le G_{e,1}$, 
      and hence,  by using \eqref{thm:bounded:alt:col:pf:eq2}, it follows that for all $i=1,2$
      \begin{align*}
      \mathcal{C}_i(2\lambda,\theta,N,\delta,\epsilon)=& \frac{\mathcal{C}_{m,b,i}(N,\alpha_{\theta},G_e(\theta)+2G_{e,1}(\theta),L(\theta),c_{\e})}{N} \le \mathcal{C}_{m,b,i}(1,\bar{\alpha},G_e+2G_{e,1},L_b,c_{\e})
       %\frac{\sqrt{N}}{N} L(\theta)c_e\sqrt{n_e}\theta_{\infty}(\alpha_{\theta})\|\alpha_g\|_{\ell_1}  \\
      %& \le 
      %  \frac{1}{\sqrt{N}} L_bc_e\sqrt{n_{\e}}\theta_{\bar{\alpha}}\|\alpha_g\| = 
        %\sqrt{N} L_b c_e \sqrt{n_{\e}} \theta_{\bar{\alpha}} \|\alpha_g\|_{\ell_1}\\
       % \frac{1}{\sqrt{N}} \mathcal{C}_{b,m,1}(1,\bar{\alpha},G_e+2G_{e,1},L_b,c_{\e}) \\
       % & \le \mathcal{C}_{b,m,1}(1,\bar{\alpha},G_e+2G_{e,1},L_b,c_{\e}) \\
      %\mathcal{C}_2(\lambda,\theta,N,\delta,\epsilon)=& \frac{\mathcal{C}_{m,b,2}(\sqrt{N},\alpha_{\theta},G_{e}(\theta)+2G_{e,1}(\theta),L(\theta),c_{\e})}{N}=\frac{N}{N} \frac{L(\theta)^2}{2} c_{\e}^2 n_e (G_e(\theta)+2G_{e,1}(\theta)) \\
       %&\le L_b^2n_e(G_e+2G_{e,1}) = \mathcal{C}_{b,m,2}(1,\bar{\alpha},L_b,c_{\e})
      \end{align*}
      Hence for $i=1,2$,
      \begin{align*}
        \ln E_{\theta \sim \pi} \exp(\mathcal{C}_i(2\lambda,\theta,N,\delta,\epsilon)) &\le 
        \ln E_{\theta \sim \pi} \exp(\mathcal{C}_{b,m,i}(2,\bar{\alpha},G_e+2G_{e,1},L_b,c_{\e})) 
        \\
        & = \mathcal{C}_{b,m,i}(2,\bar{\alpha},G_e+2G_{e,1},L_b,c_{\e})
      \end{align*}
      %proving the bounded case.
      and hence \eqref{thm:bounded:alt:col:pf:eq11}.


      \paragraph{Lipschitz loss case:}
         First, notice that for $N \ge 2$, $\ln(N) \ge 1$ and hence  
         \begin{equation}
          \label{thm:bounded:alt:col:pf:eq3}
         \begin{aligned}
            \mathcal{\bar{C}}_{ub}(2\lambda,L_{\ell},N)= &\frac{1}{N}\left(\frac{2\sqrt{N}}{16 L_{\ell}\ln(N)} \right)^2 \left(32 L_{\ell}^2C_{u,1}\right)
            + 4\frac{1}{N} \frac{2\sqrt{N}}{16L_{\ell}\ln(N)} L_{\ell}C_{u,2} \ln(N) \\
            & = \frac{1}{2\ln(N)^2}C_{u,1}+\frac{1}{2} C_{u,2} \frac{1}{\sqrt{N}} \le C_{u,1}+C_{u,2}
         \end{aligned}
        \end{equation}
        Moreover, using \eqref{thm:bounded:alt:col:pf:eq2}, 
        \begin{equation}
          \label{thm:bounded:alt:col:pf:eq4}
         \begin{aligned}
            %\mathcal{\bar{C}}_{ub}(\lambda,L_{\ell},N) \\
           \mathcal{C}_{b,i}(2\lambda, & \theta,L_{\ell},\frac{2\ln(N)}{\|\alpha_g\|_{\ell_1}K},N)  =
           \mathcal{C}_{b,i}(\frac{2\sqrt{N}}{16L_{\ell}\ln(N)}, \theta,L_{\ell},\frac{2\ln(N)}{\|\alpha_g\|_{\ell_1}K},N)  \\
           &= \frac{\mathcal{C}_{b,m,i}(\frac{2}{16L_{\ell}}\frac{\sqrt{N}}{\ln(N)},\alpha_{\theta},G_{\theta}+2G_{e,1}(\theta),L_{\ell},\frac{2\ln(N)}{\|\alpha_g\|_{\ell_1}K})}{N} \\
           & \le \mathcal{C}_{b,m,i}(\frac{1}{8L_{\ell}\ln(N)},\bar{\alpha},G_e+2G_{e,1},L_{\ell}, \frac{2\ln(N)}{\|\alpha_g\|_{\ell_1}K}) \\
           & = \left\{ \begin{array}{ll}
             \frac{1}{8L_{\ell}\ln(N)} L_{\ell} \frac{2\ln(N)}{\|\alpha_g\|_{\ell_1}K} \sqrt{n_{\e}}\theta_\infty(\bar\alpha)\|\alpha_g\|_{\ell_1}  & i=1 \\
             \frac{1}{(8L_{\ell}\ln(N))^2} \frac{L_{\ell}^2}{2} (G_e+2G_{e,1}) \left(\frac{2\ln(N)}{\|\alpha_g\|_{\ell_1}K}\right)^2 n_{\e}  & i=2 \end{array} \right. \\
           & = \left\{ \begin{array}{ll}
             \frac{1}{4K} \sqrt{n_{\e}}\theta_\infty(\bar\alpha)& i=1 \\
             \frac{1}{16} \frac{(G_e+2G_{e,1})}{2} \left(\frac{1}{\|\alpha_g\|_2K}\right)^2 n_{\e} & i=2 \end{array} \right. \\
             & \le  \mathcal{C}_{b,m,i}(\frac{1}{4},\bar{\alpha}, G_e+2G_{e,1},1, \frac{1}{\|\alpha_g\|_{\ell_1}K})
         \end{aligned}
        \end{equation}
     From this it then follows that 
     \begin{equation}
      \label{thm:bounded:alt:col:pf:eq5}
         \mathcal{C}_{ub,i}(\frac{2}{16L_{\ell}}\frac{\sqrt{N}}{\ln(N)},\theta,L_{\ell},N) \le 
           C_{u,1}+C_{u,2}+\frac{1}{2}\mathcal{C}_{b,m,i}(\frac{1}{4}, \bar{\alpha}, G_e+2G_{e,1},1, \frac{1}{\|\alpha_g\|_{\ell_1}K})
     \end{equation}
     and hence
     \begin{align*}
          & \frac{1}{2} \left(   \ln E_{\theta \sim \pi} \exp(\mathcal{C}_1(2\lambda,\theta,N,\delta,\epsilon)) +
           \ln E_{\theta \sim \pi} \exp(\mathcal{C}_2(2\lambda,\theta,N,\delta,\epsilon)) 
          \right)  \\
          & \le \frac{1}{2} \sum_{i=1}^{2} \left(C_{u,1}+C_{u,2}+ \frac{1}{2}\mathcal{C}_{b,m,i}(1/4, \bar{\alpha}, G_e+2G_{e,1}, 1, \frac{1}{\|\alpha_g\|_{\ell_1}K})\right) = \mathcal{C}(\delta,\epsilon)
      \end{align*}
      so \eqref{thm:bounded:alt:col:pf:eq11} holds.
      %Finally, notice that for $L=16L_{\ell}$ and $r=1$, 
      %\[ 
      %\frac{1}{\lambda} = \frac{\ln(N)4L_{\ell}}{\sqrt{N}}=2L \frac{(\ln(N))^{r}}{\sqrt{N}}
      %\]
     
      \paragraph{Quadratic loss case:}
           %œLet us set  $L_{\ell}=12G_e\|\e_g\|_{\psi_2}\sqrt{\ln\left(\frac{3N^2}{2\delta}\right)}$.
           %that $\lambda \le \frac{\sqrt{N}}{16 L_{\ell}\ln(N)}$.
            Notice that as $N^2 \ge 4$, it follows that $\ln(N^2) \ge 1$, and  hence
        \begin{equation} 
        \label{thm:bounded:alt:col:pf:eq6.0}
           \ln\left(\frac{3N^2}{2\delta}\right)=\left(\ln(N^2)+\ln\left(\frac{3}{2\delta}\right)\right)
            \le \ln(N^2)\left(1+\ln\left(\frac{3}{2\delta}\right)\right)
        \end{equation}
        %and hence 
        %\[ L_{\ell} \le 12G_e\|\e_g\|_{\psi_2} \sqrt{2\ln(N)} \sqrt{1+\ln\left(\frac{3}{2\delta}\right)} \le \frac{L}{16}(\ln(N))^{1/2} \]
        %therefore
        %\begin{equation}
        %  \label{thm:bounded:alt:col:pf:eq6}
        %   \lambda=\frac{\sqrt{N}}{16 L(\ln N)^{3/2}} 
           %\bar{\lambda}=\frac{\sqrt{N}}{16L_{\ell}\ln(N)}
            %\ge 
            %\frac{\sqrt{N}}{192G_e\|\e_g\|_{\psi_2} \sqrt{2} (\ln(N))^{\frac{3}{2}} 
            %\sqrt{1+\ln\left(\frac{3}{2\delta}\right)}} = \sqrt{N}{L (\ln(N))^{\frac{3}{2}}}
        %\end{equation}
       % and finally for $r=3/2$,
        %\[ \frac{1}{\bar{\lambda}}  \le \frac{1}{\lambda}=L \frac{(\ln(N))^{r}}{\sqrt{N}}
        %\frac{192G_e\|\e_g\|_{\psi_2} \sqrt{2} (\ln(N))^{\frac{3}{2}} 
         %   \sqrt{1+\ln\left(\frac{3}{2\delta}\right)}}{\sqrt{N}}=\frac{L\sqrt{2}(\ln(N))^{\frac{3}{2}}}{\sqrt{N}}
        %M\]
       %Finally, %from the definition of $L$ and 
        By using \eqref{thm:bounded:alt:col:pf:eq6.0}
        and  applying \eqref{thm:bounded:alt:col:pf:eq5} with
      $L_{\ell}=\frac{L(\ln(N))^{1/2}}{16}$ and noticing that $\lambda=\frac{\sqrt{N}}{16 L_{\ell} \ln(N)}$
       %12G_e\|\e_g\|_{\psi_2}\sqrt{\ln\left(\frac{3N^2}{2\delta}\right)}$ %\le \frac{L}{16}$ 
       %and noticing that $\mathcak{C}_{b,m,i}()
       it follows that
      % for $i=1,2$,
       \begin{equation}
        \label{thm:bounded:alt:col:pf:eq7}
        \begin{aligned}
           \mathcal{C}_{ub,i}(2\lambda, & \theta,12G_e\|\e_g\|_{\psi_2}\sqrt{\ln\left(\frac{3N^2}{2\delta}\right)},N)  \
           \le 
           \mathcal{C}_{ub,i}(2\lambda, \theta, \frac{L(\ln(N))^{1/2}}{16},N)  \\
            & \le C_{u,1}+C_{u,2}+\frac{1}{2}\mathcal{C}_{b,m,i}\left(1/4, \bar{\alpha}, G_e+2G_{e,1}, 1, \frac{1}{\|\alpha_g\|_{\ell_1}K}\right)
       \end{aligned}
       \end{equation}
       %Moreover, \eqref{thm:bounded:alt:col:pf:eq6.0} and the fact that as $N \ge K_{\delta}=\max\{2, \frac{9}{2\delta}\}$, 
       Moreover, from \eqref{thm:bounded:alt:col:pf:eq6.0} and the fact that 
       %as $N \ge K_{\delta}=\max\{2, \frac{9}{2\delta}\}$, 
       $\ln(N) \ge 1$ and $\ln\left(\frac{3N^2}{2\delta}\right) \ge 1$, it follows that
       %it follows that
       %it follows that
       %it follows that
       \begin{equation}
        \label{thm:bounded:alt:col:pf:eq8}
       \begin{aligned}
         2\lambda & \frac{\ln\left(\frac{3N^2}{2\delta}\right)}{N} 72G^2_e\|\e_g\|_{\psi_2}^2\max\{0,\epsilon\}
          \\
         &= \frac{2\sqrt{N}}{192 G_e\|\e_g\|_{\psi_2} \sqrt{2}\ln(N) \sqrt{1+\ln\left(\frac{3}{2\delta}\right)}\sqrt{\ln(N)}}
            (72G^2_e\|\e_g\|_{\psi_2}^2\max\{0,\epsilon\})\frac{\ln\left(\frac{3N^2}{2\delta}\right)}{N}   \\
        & \leq  \frac{2\sqrt{N}}{192 G_e\|\e_g\|_{\psi_2}\ln(N) \sqrt{1+\ln\left(\frac{3}{2\delta}\right)}\sqrt{\ln(N^2)}}
            (72G^2_e\|\e_g\|_{\psi_2}^2\max\{0,\epsilon\}) \frac{(1+\ln\left(\frac{3}{2\delta}\right))\ln(N^2)}{N} \\
         & = \frac{2 \cdot 72G^2_e\|\e_g\|_{\psi_2}^2\max\{0,\epsilon\}\sqrt{1+\ln\left(\frac{3}{2\delta}\right)}\sqrt{\ln(N^2)}}{192 G_e\|\e_g\|_{\psi_2}\ln(N)\sqrt{N}} \\
        %  & \frac{2\sqrt{N}}{192 G_e\|\e_g\|_{\psi_2} \sqrt{2}\ln(N) \sqrt{1+\ln\left(\frac{3}{2\delta}\right)}\sqrt{\ln(N)}}
        %     (72G^2_e\|\e_g\|_{\psi_2}^2\max\{0,\epsilon\})\frac{\ln\left(\frac{3N^2}{2\delta}\right)}{N} \le  \\
        % &   \frac{2\sqrt{N}}{192 G_e\|\e_g\|_{\psi_2}\ln(N) \sqrt{1+\ln\left(\frac{3}{2\delta}\right)}\sqrt{\ln(N^2)}}
        %     (72G^2_e\|\e_g\|_{\psi_2}^2\max\{0,\epsilon\}) \frac{(1+\ln\left(\frac{3}{2\delta}\right))\ln(N^2)}{N} \\
        %  & \le \frac{2 \cdot 72G^2_e\|\e_g\|_{\psi_2}^2\max\{0,\epsilon\}\sqrt{1+\ln\left(\frac{3}{2\delta}\right)}\sqrt{\ln(N^2)}}{192 G_e\|\e_g\|_{\psi_2}\ln(N)\sqrt{N}}\\   
         &= 
         \frac{6}{16} \frac{2G_e\|\e_g\|_{\psi_2}\sqrt{2}(\ln(N))^{1/2} \sqrt{1+\ln\left(\frac{3}{2\delta}\right)}\max\{0,\epsilon\}}{\ln(N) \sqrt{N}} \\
         & = \frac{6}{8 \cdot 192} \frac{L (\ln(N))^{1/2}}{\sqrt{N}\ln(N)} \max\{0,\epsilon\}\\ 
         % =
         %  \frac{1}{256} \frac{\sqrt{2} L}{\sqrt{N}(\ln(N))^{1/2}} \max\{0,\epsilon\} \\
          & = \frac{1}{128} \frac{L}{\sqrt{N}(\ln(N))^{1/2}} \max\{0,\epsilon\}
       \end{aligned}
      \end{equation}
       Hence,  for $i=1,2$,
       \begin{align*}
            \mathcal{C}_i(2\lambda, & \theta,N,\delta,\epsilon) \le \frac{L}{128} \max\{0,\epsilon\} + \\
               & C_{u,1}+C_{u,2}+\frac{1}{2}\mathcal{C}_{b,m,i}\left(1/4, \bar{\alpha}, G_e+2G_{e,1}, 1, \frac{1}{\|\alpha_g\|_2K}\right)
       \end{align*}
       and thus 
      \begin{align*}
          & \frac{1}{2} \left(   \ln E_{\theta \sim \pi} \exp(\mathcal{C}_1(\lambda,\theta,N,\delta,\epsilon)) +
           \ln E_{\theta \sim \pi} \exp(\mathcal{C}_2(\lambda,\theta,N,\delta,\epsilon)) 
          \right)  \\
          & \le \frac{L}{128} \max\{0,\epsilon\} + \frac{1}{2} \sum_{i=1}^{2} \left(C_{u,1}+C_{u,2}+ \frac{1}{2}\mathcal{C}_{b,m,i}\left(1/4, \bar{\alpha}, G_e+2G_{e,1}, 1, \frac{1}{\|\alpha_g\|_2K}\right)\right)  \\
          & = \mathcal{C}(\delta,\epsilon)
      \end{align*}
   which completes the proof. 



 \subsection{Uniform bounds derived from Corollary \ref{thm:bounded:alt:col}}
 \label{sec:uniform_bounds}
 In order to present the uniform bounds, we need the following assumptions from Assumption \ref{as:oracle}:
 \textbf{(LipModel)} and either 
\textbf{(Gauss)} or \textbf{(Uni)}.
%% \begin{assumption}[\textcolor{cyan}{Lipschitz parameterisation}]
%%  \label{as:lipschitz:param}
%%  The map $\theta \mapsto \alpha_{\theta}$ is $L_{\alpha}$-Lipschitz w.r.t.
%%  to the $\|\cdot\|_{\ell_1}$ norm, i.e.,
%%  \(
%%    \|\alpha_{\theta_1}-\alpha_{\theta_2}\|_2 \le L_{\alpha} \|\theta_1-\theta_2\|_{\ell_1}, \quad \forall \theta_1,\theta_2 \in \Theta.
%%  \)
%%  \end{assumption}
%%  \begin{assumption}[Prior]
%%    \label{as:prior}
%%      Assume that either 
%%    \textbf{(Gauss)} $\Theta=\mathbb{R}^{n_{\theta}}$ and $\pi$ is
%%   Gaussian $\mathcal{N}(\theta_m,P)$, or
%%   \textbf{(Uni)} 
%%   $\Theta$ is compact and $\pi$ is the uniform density on $\Theta$
%%  \end{assumption}  
%  In order to present the uniform bounds, we need the following assumptions from Assumption \ref{as:oracle}:
%  \textbf{(LipModel)} and either 
% \textbf{(Gauss)} or \textbf{(Uni)}.
%% \begin{assumption}[\textcolor{cyan}{Lipschitz parameterisation}]
%%  \label{as:lipschitz:param}
%%  The map $\theta \mapsto \alpha_{\theta}$ is $L_{\alpha}$-Lipschitz w.r.t.
%%  to the $\|\cdot\|_{\ell_1}$ norm, i.e.,
%%  \(
%%    \|\alpha_{\theta_1}-\alpha_{\theta_2}\|_2 \le L_{\alpha} \|\theta_1-\theta_2\|_{\ell_1}, \quad \forall \theta_1,\theta_2 \in \Theta.
%%  \)
%%  \end{assumption}
%%  \begin{assumption}[Prior]
%%    \label{as:prior}
%%      Assume that either 
%%    \textbf{(Gauss)} $\Theta=\mathbb{R}^{n_{\theta}}$ and $\pi$ is
%%   Gaussian $\mathcal{N}(\theta_m,P)$, or
%%   \textbf{(Uni)} 
%%   $\Theta$ is compact and $\pi$ is the uniform density on $\Theta$
%%  \end{assumption}  
%  In order to present the uniform bounds, we need the following assumptions from Assumption \ref{as:oracle}:
%  \textbf{(LipModel)} and either 
% \textbf{(Gauss)} or \textbf{(Uni)}.
%% \begin{assumption}[\textcolor{cyan}{Lipschitz parameterisation}]
%%  \label{as:lipschitz:param}
%%  The map $\theta \mapsto \alpha_{\theta}$ is $L_{\alpha}$-Lipschitz w.r.t.
%%  to the $\|\cdot\|_{\ell_1}$ norm, i.e.,
%%  \(
%%    \|\alpha_{\theta_1}-\alpha_{\theta_2}\|_2 \le L_{\alpha} \|\theta_1-\theta_2\|_{\ell_1}, \quad \forall \theta_1,\theta_2 \in \Theta.
%%  \)
%%  \end{assumption}
%%  \begin{assumption}[Prior]
%%    \label{as:prior}
%%      Assume that either 
%%    \textbf{(Gauss)} $\Theta=\mathbb{R}^{n_{\theta}}$ and $\pi$ is
%%   Gaussian $\mathcal{N}(\theta_m,P)$, or
%%   \textbf{(Uni)} 
%%   $\Theta$ is compact and $\pi$ is the uniform density on $\Theta$
%%  \end{assumption}  
 \begin{corollary}[Uniform [PAC bounds, Example 2.1,2.2 \cite{alquier2021userfriendly}]
 \label{lem:pac}
    %Assume that
    % $\theta \mapsto \alpha_{\theta}$ is $L_{\alpha}$-Lipschitz w.r.t.
%%    %to the $\|\cdot\|_2$ norm, and that
%%  % $\mathcal{L}(\theta)$ and $\hat{\mathcal{L}}_N(\theta)$ are  $L$-Lipschitz functions of $\theta$,
  With the Assumption  \textbf{(LipModel)} and either 
\textbf{(Gauss)} or \textbf{(Uni)} from Assumption \ref{as:oracle} the following holds.
  Let $L_{\Theta}(\delta)$, $\mathcal{C}_{L,emp}$, $\mathcal{C}_{L,true}$, $C(\theta)$,  $r_{\text{Lip},1}$, $r_{\text{Lip},2}$
   be defined as in 
   Table~\ref{tab:constants:sysid},
   %Corollary \ref{lem:bound:mean}, 
   and let 
   $L$, $C_{\delta}$, %$K_{\delta}$, 
   $\mathcal{C}(\delta,\epsilon)$, $r$  be defined as in 
   Corollary \ref{thm:bounded:alt:col}
   For any $\delta \in (0,0.5)$, $\lambda,\sigma  > 0
   $, $N\geq 3$, %$$N \ge \max\{K_{\sigma},4\}$, 
   $\epsilon \in \{-1,1\}$, 
   \begin{equation}
   \label{lem:pac:eq1:uniform}
   \begin{aligned}
     \forall \theta \in \Theta: &
    \epsilon \mathcal{L}(\theta)  \le 
    \epsilon \hat{\mathcal{L}}_{N}(\theta) + \\
     %\frac{C(\theta)}{\lambda}  \\
     & \frac{L(\ln(N))^{r+1}}{\sqrt{N}} \left(\mathcal{C}_{L,emp}+\mathcal{C}_{L,true}+C(\theta)+\ln \left(\frac{C_{\delta}}{\delta^{2r_{\text{Lip},1}+1}}\right) + \mathcal{C}(\delta,\epsilon)\right)
       %2\sigma L_{\Theta}(\delta) +  4r_{\text{Lip},1} \left(C_{ub,1}+C_{ub,2}\right) \right. \\
       %& \left. +2r_{\text{Lip},2}\frac{L}{256} 
    % +) 
     %\Psi(\lambda,\pi,N)}{\lambda}
    \end{aligned}
   \end{equation}
    holds with probability at least $1-(2+r_{\text{Lip},1})\delta$ over data.
%%   $r_{\text{Lip},1}=0$ when $\e_g$ is bounded and $r_{\text{Lip},1}=1$   in all other cases, $r_{\text{Lip},2}=1$ when $\ell$ is quadratic, and $r_{\text{Lip},2}=0$ otherwise, and
%%   \begin{equation}
%%    \label{lem:pac:eq2}
%%   %\begin{split}
%%    C(\theta)=\left\{ \begin{array}{rl} 
%%    \frac{1}{2} \big (\sigma^2 \mathrm{trace} (P^{-1}) %}{\sqrt{N}}
%%   -n_{\theta}+ (\theta-\theta_m)^TP^{-1}(\theta-\theta_m)+
%%   \ln \frac{\det(P)}{\sigma^2} + \frac{1}{2} \big) & 
%%   \mbox{ if \textbf{(Gauss)} holds} \\
%%       \ln\frac{vol(\Theta)}{(\sqrt{\pi \sigma^2}^{n_{\theta}} \Gamma(n_{\theta}/2+1)}
%%      + \frac{n_{\theta}}{2}  & 
%%   \mbox{ if \textbf{(Uni)} holds} 
%%   \end{array}\right.
%%  \end{equation}
%%    %$L$ is defined as in Corollary \ref{thm:bounded:alt:col}, 
%%    and $L_{\Theta}(\delta)$ is defined as follows:
%%    \begin{equation}
%%    \label{lem:pac:eq3}
%%    L_{\Theta}(\delta)=\left\{\begin{array}{ll}
%%    L_{\alpha}L_{\ell} \|\alpha_g\|_{\ell_1}\sqrt{n_\rve}c_{\rve} & \text{if } \|e_g(t)\|_{\infty} \le c_{\rve} \text{ (bounded noise)} \\[0.5em]
%%    \frac{1}{16} L_{\alpha}\sqrt{n_\rve} \frac{1}{K} & \text{if } \ell \text{ is } L_{\ell}\text{-Lipschitz} \\[0.5em]
%%    L_{\alpha} \frac{1}{16} \sqrt{n_\rve} \frac{1}{K} & \text{if } \ell \text{ is quadratic}
%%    \end{array}\right.
%%  \end{equation}
%     \|\alpha_g\|_{\ell_1} \sqrt{n_\rve} \frac{2(\ln(N))^3}{\|\alpha_g\|_2^2}
 \end{corollary}
The proof of Corollary \ref{lem:pac} is based on a number of technical lemmas, which will also be used for
the proof of Corollary \ref{lem:bound:mean}.
\begin{lemma}[$\KL$ divergence]
  \label{pac:kl}
   Assume that either \textbf{(Gauss)} or \textbf{(Uni)} holds.
   For every $\theta \in \Theta$, define the density $\rho_{\theta}$ as follows
  \[ \rho_{\theta}=\left\{ \begin{array}{rl}
                           \mathcal{N}(\theta,\frac{\sigma^2}{N}I) & \mbox{ if \textbf{(Gauss)} holds} \\
                            \mathcal{U}(\theta,\frac{\sigma}{\sqrt{N}}I)  & \mbox{ if \textbf{(Uni)} holds}
                         \end{array} \right. 
  \]
   where $\mathcal{N}(m,\sigma^2I)$ denotes the Gaussian distribution with mean $m$ and variance $\sigma^2I$, and
   $\mathcal{U}(m,r)$ denotes the uniform  distribution on the ball with center in $m$ and radius $r$.
  Then for $N \ge 3$,
  \[ \KL(\rho_{\theta} \mid \pi) \le C(\theta)\ln(N)\]
\end{lemma}
\begin{proof}[Proof of Lemma \ref{pac:kl}]
  If \textbf{(Gauss)} holds, then 
  using KL divergence for Gaussian densities (\cite[Lemma 11.3.4]{LindquistBook}) it follows
  that 
  \begin{align*}
     \KL&(\rho_{\theta} \mid \pi)  = \\
     & \frac{1}{2} \left(\mathrm{trace}\left(\frac{\sigma^2}{N}P^{-1} \right) +
          (\theta_m-\theta)^TP^{-1}(\theta_m-\theta) - n_{\theta}  + \ln \frac{\det(P)N^{n_{\theta}}}{\sigma^{2n_{\theta}}} \right)  \\
      & =\frac{1}{2} \left(\mathrm{trace}\left(\frac{\sigma^2}{N} P^{-1} \right) +
          (\theta_m-\theta)^TP^{-1}(\theta_m-\theta) - n_{\theta}  + \ln \frac{\det(P)}{\sigma^{2n_{\theta}}}+ n_{\theta}\ln(N)\right)\\
        & \le \ln(N) C(\theta)  
  \end{align*}
  where we used the fact that $\ln(N) \ge 1$   and $\frac{1}{N} \le 1$ if $N \ge 3$.
  Likewise, if \textbf{(Uni)} holds,  and we denote by
  $C_2$ the volume of the ball of radius $\frac{\sigma}{\sqrt{N}}$ centered at $\tilde{\theta}$, then
\begin{align*}
    & KL(\rho_{\theta} \mid \pi) = \int_{\|\tilde{\theta}-\theta\|_2 \le \frac{\sigma^2}{\sqrt{N}}} 
     \ln\left(\frac{\rho_{\theta}(\tilde{\theta})}{\pi(\tilde{\theta})}\right) \frac{1}{C_2} dm(\tilde{\theta})  \\
     & =\int_{\|\tilde{\theta}-\theta\|_2 \le \frac{\sigma^2}{\sqrt{N}}}  \frac{1}{C_2} \ln\left(\frac{\mathrm{vol}(\Theta)}{C_2}\right) dm(\tilde{\theta})
    = \ln\left(\frac{\mathrm{vol}(\Theta)}{C_2}\right)
\end{align*}
From the formula for the volume of a high-dimensional ball \cite{BallVolume} it follows that
\[
  C_2= \frac{\pi^{n_\theta/2}}{\Gamma(n_\theta/2+1)} \left(\frac{\sigma}{\sqrt{N}}\right)^{n_\theta}
\]
and hence 
\begin{align*}
  & \KL(\rho_{\theta} \mid \pi) = \ln\left(\frac{\mathrm{vol}(\Theta)}{C_2}\right)=
  \ln\left(\frac{\mathrm{vol}(\Theta)}{\frac{\pi^{n_\theta/2}}{\Gamma(n_\theta/2+1)} \left(\frac{\sigma}{\sqrt{N}}\right)^{n_\theta}}\right) \\
  & =\ln\left(\frac{\mathrm{vol}(\Theta) \Gamma(n_\theta/2+1)}{\pi^{n_\theta/2} (\sigma)^{n_\theta}}\right) + 
     \ln \left(\left(\sqrt{N}\right)^{n_\theta}\right) \\
     & = \ln\left(\frac{\mathrm{vol}(\Theta) \Gamma(n_\theta/2+1)}{\sqrt{\pi\sigma^2}^{n_\theta}}\right) + 
       \frac{n_{\theta}}{2} \ln(N)
       \le C(\theta)\ln(N)
\end{align*}
%
%  it follows that
%$\KL(\rho_{\theta} \mid \pi)=C(\theta_{\star})$.
%Assume that $\pi$ is the uniform density of $\Theta$, then
%let $\rho$ be the uniform density centered at the 
%intersection of $\Theta$ with the ball %$B_{\frac{\sigma^2}{\sqrt{N}}}(\theta_{\star})$
%centered in $\theta_{\star}$
%and having radius $\frac{\sigma^2}{\sqrt{N}}$.
%$E_{\theta \sim \rho} \|\theta - \theta_{\star}\|^2 < \frac{\sigma^4}{N}$. Moreover, from
%the formula for KL-divergence of uniform densities and the 
%$\KL(\rho|\pi)=C(\theta_{\star})$. 
\end{proof}
Another ingredient is the following:
\begin{lemma}[Lipschitz continuity of the true and empirical errores]
  \label{pac:lipsch:loss}
   %If $\ell_g$ is bounded and $\ell$ is $L_{\ell}$-Lipschitz, then the true and empirical errores are Lipschitz continuous
   %in $\theta$, with the Lipschitz constant $L_{\Theta}(\delta)$, 
   %Consider the inequalities
   %\begin{equation}
   % \label{pac:lipschitz:loss:eq1}
   % \begin{aligned}
   %|\mathcal{L}(\theta_1) - \mathcal{L}(\theta_2)| \le L_{\Theta}(\delta)2\ln(N) \|\theta_1 - \theta_2\| \quad \forall \theta_1, \theta_2 \in \Theta \\
   %|\hat{\mathcal{L}}_N(\theta_1) - \hat{\mathcal{L}}_N(\theta_2)| \le L_{\Theta}(\delta)2\ln(N) \|\theta_1 - \theta_2\| \quad \forall \theta_1, \theta_2 \in \Theta
   %\end{aligned}
   % \label{pac:lipschitz:loss:eq2}
   %\end{equation}
   Using the constants defined in Corollary \ref{thm:bounded:alt:col}-\ref{lem:pac}, 
   the following inequality holds
   \begin{equation}
   \label{pac:lipschitz:loss:eq31}
   \begin{aligned}
     & \forall \theta_1,\theta_2 \in \Theta:\\
          &|\hat{\mathcal{L}}_N(\theta_1) - \hat{\mathcal{L}}_N(\theta_2)| \le L_{\Theta}(\delta) L(\ln(N))^r \|\theta_1 - \theta_2\|_2
           +  \\
           & \qquad\qquad\qquad 2r_{\text{Lip,1}}\frac{L\ln(N)^{r}}{\sqrt{N}}\left(C_{u,1} + C_{u,2}+\ln\left(\frac{1}{\delta}\right)\right), 
   \end{aligned}
   \end{equation}  
   with probability $1-r_{\text{Lip},1}\delta$ over data, 
   where 
   $$ r_{\text{Lip},1}=\left\{\begin{array}{ll} 0 & \e_g \mbox{ is bounded } \\ 1 & \mbox{otherwise} \end{array}\right.$$ 
   and 
   \begin{equation}
   \label{pac:lipschitz:loss:eq32}
   \begin{aligned}
       & \forall ~ \theta_1,\theta_2 \in \Theta: \quad \\
          & |\mathcal{L}(\theta_1) - \mathcal{L}(\theta_2)| \le L_{\Theta}(\delta) L(\ln(N))^r \|\theta_1 - \theta_2\|_2
          + 2r_{\text{Lip},1}\frac{L\ln(N)^{r}}{\sqrt{N}}\left(C_{u,1} + C_{u,2}\right) + \\
           & \qquad\qquad\qquad 2r_{\text{Lip},2}  \frac{L}{512} \frac{L(\ln(N))^r}{\sqrt{N}}
   \end{aligned}
   \end{equation}
    where $$ r_{\text{Lip},2}=\left\{\begin{array}{ll} 2 & \ell \mbox{ is quadratic and $\e_g$ is unbounded} \\ 0 & \mbox{otherwise} \end{array}\right. .$$
   Moreover, the inequalities below 
   \begin{equation}
   \label{pac:lipschitz:loss:eq3}
   \begin{aligned}
     & \mbox{\eqref{catoni1} from Corollary \ref{thm:bounded:alt:col} holds, and }  \\
     & \forall \theta_1,\theta_2 \in \Theta: \\
         & |\hat{\mathcal{L}}_N(\theta_1) - \hat{\mathcal{L}}_N(\theta_2)| \le L_{\Theta}(\delta) L(\ln(N))^r \|\theta_1 - \theta_2\|_2
           +  \\
            & \qquad\qquad\qquad 2r_{\text{Lip},1}\frac{L\ln(N)^{r}}{\sqrt{N}}\left(C_{u,1} + C_{u,2}+\ln\left(\frac{1}{\delta}\right)\right), \\
           \text{ and } &\\
       & \forall ~ \theta_1,\theta_2 \in \Theta: \quad \\
          & |\mathcal{L}(\theta_1) - \mathcal{L}(\theta_2)| \le L_{\Theta}(\delta) L(\ln(N))^r \|\theta_1 - \theta_2\|_2
          + 2r_{\text{Lip},1}\frac{L\ln(N)^{r}}{\sqrt{N}}\left(C_{u,1} + C_{u,2}\right) + \\
           & \qquad\qquad\qquad 2r_{\text{Lip},2}  \frac{L}{512} \frac{L(\ln(N))^r}{\sqrt{N}}
           % 144 G_e^2\|e_g\|_{\psi_2}^2 \frac{\ln(N)}{N}\left (1+\ln(\frac{3}{2\delta})\right)
   \end{aligned}
   \end{equation}
   hold simultaneously  with probability at least $1-(2+r_{\text{Lip},1})\delta$.
  
\end{lemma}
\begin{proof}[Proof of Lemma \ref{pac:lipsch:loss}]
  \textbf{Bounded noise case:}
  If $\ell$ is $L_{\ell}$-Lipschitz and $\e_g$ is bounded, then using (\textbf{LipModel}) and the argument of
  the proof from Lemma \ref{alpha:lemma2}, it follows that 
  \begin{align*}
    & |\ell(\y(t),\y_{\theta_1}(t)) - \ell(\y(t),\y_{\theta_2}(t))| \le L_{\ell} \|\y_{\theta_1}(t) - \y_{\theta_2}(t)\|_2 \\
    & \le L_{\ell}\sum_{k=0}^{\infty} \|c_{\theta_1}(k)-c_{\theta_2}(k)\|_{2} \|\e_g(t-k)\|_{2} 
    \le L_{\ell} \sum_{k=0}^{\infty} \sum_{l=0}^{k} \|\alpha_{\theta_1}(l)-\alpha_{\theta_2}(l)\|_2\|\alpha_{g,w}(k-l)\|_2 c_{\e}\sqrt{n_\e} \\
    & \le L_{\ell}\|\alpha_{\theta_1}-\alpha_{\theta_2}\|_{\ell_1}\|\alpha_{g}\|_{\ell_1}c_{\e}\sqrt{n_\e}=
          \underbrace{L_{\alpha}L_{\ell}\|\alpha_{g}\|_{\ell_1}c_{\e}\sqrt{n_\e}}_{L_{\Theta}(\delta)} \|\theta_1-\theta_2\|_2
  \end{align*}
  From this it then follows that 
    \begin{align}
   & |\mathcal{L}(\theta_1) - \mathcal{L}(\theta_2)| \le L_{\Theta}(\delta) \|\theta_1 - \theta_2\| \quad \forall \theta_1, \theta_2 \in \Theta 
    \label{pac:lipschitz:loss:eq1} \\
   & |\hat{\mathcal{L}}_N(\theta_1) - \hat{\mathcal{L}}_N(\theta_2)| \le L_{\Theta}(\delta) \|\theta_1 - \theta_2\| \quad \forall \theta_1, \theta_2 \in \Theta
    \label{pac:lipschitz:loss:eq2}
   \end{align}
   By noticing that $r=0$ and $L=1$ for the bounded case, 
   %$2\ln(N) = \ln(N^2) \ge 1$ for $N \ge 2$, it follows that 
   %$L_{\Theta}(\delta) 2\ln(N) \ge L_{\Theta}(\delta)$, from which 
   the claim of the lemma follows. 

  \textbf{Unbounded noise case:}
  Assume that the noise is not necessarily bounded and $\ell$  is $L_{\ell}$-Lipschitz. From Subsection \ref{lipschitz:unbounded},
  recall the definition of truncated empirical and true errors 
  $\bar{\mathcal{L}}_N(\theta)$ and $\bar{\mathcal{L}}(\theta)$, corresponding to the 
  output  $\bar{\y}(t)$ and $\bar{\y}_{\theta}(t)$ and $\bar{\y}_{\theta}(t \mid 0)$ generated for the truncated noise
  $\bar{\e}_g(t)=\e_g(t)\chi(\e_g(t) < C)$ for $C=\frac{\ln(N)}{K\|\alpha_g\|_{\ell_1}}$.
  Using Chernoff's inequality and Lemma \ref{vl-vbl:lemma1}, it follows that 
  \begin{equation}
    \label{pac:lipschitz:loss:eq05}
  \begin{aligned}
     \bP(\forall \theta \in \Theta: \left|\hat{\mathcal{L}}_N(\theta) - \bar{\mathcal{L}}_N(\theta)\right| > c) \le  
     \frac{e^{\lambda\frac{1}{2} \left(  \mathcal{C}(\|\bar{\alpha}\|_{\ell_1} \|\alpha_g\|_{\ell_1}, 2L_{\ell}\lambda,N)+
\mathcal{C}(\|\alpha_g\|_{\ell_1}, 2L_{\ell}\lambda,N)\right)}}{e^{\lambda c}}
  \end{aligned}
  \end{equation}
  and hence 
   \begin{equation}
    \label{pac:lipschitz:loss:eq5}
  \begin{aligned}
     & \bP\left(\forall \theta \in \Theta:  \left|\hat{\mathcal{L}}_N(\theta) - \bar{\mathcal{L}}_N(\theta)\right| \le 
      \frac{1}{2\lambda} \left(\mathcal{C}(\|\bar{\alpha}\|_{\ell_1} \|\alpha_g\|_{\ell_1}, 2L_{\ell}\lambda,N)+ \right. \right. \\
& \left. \left. \mathcal{C}(\|\alpha_g\|_{\ell_1}, 2L_{\ell}\lambda,N) + \ln\left(\frac{1}{\delta }\right)\right)\right)
      \ge 1-\delta
  \end{aligned}
  \end{equation}
   Similarly, from Lemma \ref{vl-vbl:lemma1} it follows that
   \begin{equation}
    \label{pac:lipschitz:loss:eq5.1}
    \forall \theta \in \Theta: \quad
    |\mathcal{L}(\theta) - \mathcal{\bar{L}}(\theta)| \le  
   \frac{1}{2\lambda} \left( \mathcal{C}(\|\bar{\alpha}\|_{\ell_1} \|\alpha_g\|_{\ell_1}, 2L_{\ell}\lambda,N)+
    \mathcal{C}(\|\alpha_g\|_{\ell_1}, 2L_{\ell}\lambda,N)\right) 
   %\mathcal{C}(\|\bar{\alpha}\|_{\ell_1} \|\alpha_g\|_{\ell_1}, 2L_{\ell}\lambda,N)
   \end{equation}
   From \eqref{thm:main:pf:unbounded:eq-1} in Subsection \ref{lipschitz:unbounded} for $C=\frac{\ln(N)}{K\|\alpha_g\|_{\ell_1}}$ it follows that
   \begin{align*}
   \frac{1}{2} \left( \mathcal{C}(\|\bar{\alpha}\|_{\ell_1} \|\alpha_g\|_{\ell_1}, 2L_{\ell}\lambda,N)+
    \mathcal{C}(\|\alpha_g\|_{\ell_1}, 2L_{\ell}\lambda,N)\right) = 2\bar{\mathcal{C}}_{ub}(\frac{\lambda}{8}, L_{\ell},N)
   \end{align*}
   and hence by  setting $\lambda=\frac{\sqrt{N}}{16 L_{\ell}\ln(N)}$  it follows that
   \begin{align}
     & \frac{1}{\lambda} (\frac{1}{2} \left(\mathcal{C}(\|\bar{\alpha}\|_{\ell_1} \|\alpha_g\|_{\ell_1}, 2L_{\ell}\lambda,N)+
\mathcal{C}(\|\alpha_g\|_{\ell_1}, 2L_{\ell}\lambda,N)\right) \nonumber \\
     & \le (\frac{1}{2} C_{u,1}+\frac{1}{8} C_{u,2}) \frac{16 L_{\ell}\ln(N)}{\sqrt{N}}
     \le (C_{u,1}+C_{u,2}) \frac{16L_{\ell}\ln(N)}{\sqrt{N}}
    \label{pac:lipschitz:loss:eq6}
   \end{align}
    and hence, appplying it to \eqref{pac:lipschitz:loss:eq5}-\eqref{pac:lipschitz:loss:eq5.1} and recalling that
    $L=16L_{\ell}$ results in
    \begin{equation}
    \label{pac:lipschitz:loss:eq07}
      \begin{aligned}
      & \bP\left(\forall \theta \in \Theta: \left|\hat{\mathcal{L}}_N(\theta) - \bar{\mathcal{L}}_N(\theta)\right| \le (C_{u,1}+C_{u,2}+\ln(\frac{1}{\delta})) \frac{16 L_{\ell}\ln(N)}{\sqrt{N}}\right) \ge 1-\delta \\
      &  \forall \theta_1,\theta_2 \in \Theta: |\mathcal{L}(\theta_1) - \mathcal{L}(\theta_2)| \le  (C_{u,1}+C_{u,2}) \frac{L\ln(N)}{\sqrt{N}}\\
   %\frac{1}{2\lambda} \left( \mathcal{C}(\|\bar{\alpha}\|_{\ell_1} \|\alpha_g\|_{\ell_1}, 2L_{\ell}\lambda,N\right)
      \end{aligned}
    \end{equation}
     Note that  \eqref{pac:lipschitz:loss:eq1}-\eqref{pac:lipschitz:loss:eq2} holds for $\bar{\mathcal{L}}_N(\theta)$ and 
     $\bar{\hat{\mathcal{L}}}_N(\theta)$ with $L_{\Theta}(\delta)=L_{\alpha}L_{\ell}\sqrt{n_\e}\frac{\ln(N)}{K\|\alpha_g\|_{\ell_1}}$ and hence 
     by using \eqref{pac:lipschitz:loss:eq07} and the observation
     \begin{align*}
       & |\hat{\mathcal{L}}_N(\theta_1) - \hat{\mathcal{L}}_N(\theta_2)| \le
       |\hat{\mathcal{L}}_N(\theta_1) - \bar{\hat{\mathcal{L}}}_N(\theta_1)| + |\bar{\hat{\mathcal{L}}}_N(\theta_1) - \bar{\hat{\mathcal{L}}}_N(\theta_2)| + |\bar{\hat{\mathcal{L}}}_N(\theta_2) - \hat{\mathcal{L}}_N(\theta_2)| \\
       & |\mathcal{L}(\theta_1)-\mathcal{L}(\theta_2)| \le |\mathcal{L}(\theta_1)-\bar{\mathcal{L}}(\theta_1)| + |\bar{\mathcal{L}}(\theta_1)-\bar{\mathcal{L}}(\theta_2)| + |\bar{\mathcal{L}}(\theta_2)-\mathcal{L}(\theta_2)|
     \end{align*}
     it follows
     \begin{equation}
    \label{pac:lipschitz:loss:eq7}
      \begin{aligned}
      & \bP\left(\forall \theta_1, \theta_2 \in \Theta: \left|\hat{\mathcal{L}}_N(\theta_1) - \hat{\mathcal{L}}_N(\theta_1)\right| \le 2(C_{u,1}+C_{u,2}+\ln(\frac{1}{\delta})) \frac{16 L_{\ell}\ln(N)}{\sqrt{N}}) + \right. \\
      & \left. L_{\alpha}L_{\ell}\sqrt{n_\e}\frac{\ln(N)}{K\|\alpha_g\|_{\ell_1}} \|\theta_1 - \theta_2\|_2 \right) \ge 1-\delta \\
      & \forall \theta_1, \theta_2 \in \Theta:  \left|\mathcal{L}(\theta_1) - \mathcal{L}(\theta_2)\right| \\
      & \le  2(C_{u,1}+C_{u,2}+\ln(\frac{1}{\delta})) \frac{L \ln(N)}{\sqrt{N}}) +L_{\alpha}L_{\ell}\sqrt{n_\e}\frac{\ln(N)}{K\|\alpha_g\|_{\ell_1}} \|\theta_1 - \theta_2\|_2
   %\frac{1}{2\lambda} \left( \mathcal{C}(\|\bar{\alpha}\|_{\ell_1} \|\alpha_g\|_{\ell_1}, 2L_{\ell}\lambda,N)
      \end{aligned}
    \end{equation}
    Since $L_{\alpha}L_{\ell}\sqrt{n_\e}\frac{\ln(N)}{K\|\alpha_g\|_{\ell_1}}=L_{\Theta}L\ln(N)$ with the definition of $L_{\Theta}$
    as in Corollary \ref{lem:pac}, \eqref{pac:lipschitz:loss:eq7} can be rewritten as
\begin{equation}
    \label{pac:lipschitz:loss:eq8}
      \begin{aligned}
      & \bP\big(\forall \theta_1, \theta_2 \in \Theta: \left|\hat{\mathcal{L}}_N(\theta_1) - \hat{\mathcal{L}}_N(\theta_1)\right| \le 2(C_{u,1}+C_{u,2}+\ln(\frac{1}{\delta})) \frac{L\ln(N)}{\sqrt{N}}) +  \\
      & L_{\Theta}(\delta) L\ln(N)\|\theta_1 - \theta_2\|_2\big) \ge 1-\delta \\
      & \forall \theta_1, \theta_2 \in \Theta:  \left|\mathcal{L}(\theta_1) - \mathcal{L}(\theta_2)\right|  \\
      & \le  2(C_{u,1}+C_{u,2}+\ln(\frac{1}{\delta})) \frac{L\ln(N)}{\sqrt{N}}) +L_{\Theta}(\delta)L\ln(N) \|\theta_1 - \theta_2\|_2
   %\frac{1}{2\lambda} \left( \mathcal{C}(\|\bar{\alpha}\|_{\ell_1} \|\alpha_g\|_{\ell_1}, 2L_{\ell}\lambda,N)
      \end{aligned}
    \end{equation}
    hence \eqref{pac:lipschitz:loss:eq31}-\eqref{pac:lipschitz:loss:eq32}. 
    By using the union bound, 
    and that \eqref{catoni1} holds with probability $1-2\delta$, 
    \eqref{pac:lipschitz:loss:eq3} follows. 

%%     %If $\ell$ is $L_{\ell}$-Lipschitz and $\e_g$ is unbounded, then for $0 < \lambda < \frac{1}{4L_{\ell}}$ 
%%     fo
%%     \begin{align}
%%      & \bP(\forall \theta1,\theta_2 \in \Theta: \quad
%%          |\mathcal{L}_N(\theta_1) - \mathcal{L}_N(\theta_2)| \le L_{\Theta}(\delta) \ln(N) \|\theta_1 - \theta_2\|
%%           + \frac{1}{\lambda} \left(\mathcal{C}_{ub}(\lambda,L_{\ell},N) + \ln(\frac{3}{\delta}\right)
%%          ) \ge 1-\frac{\delta}{3} 
%%      \label{pac:lipsch:loss2:eq1}
%%          \\
%%      & \forall \theta_1,\theta_2: |\mathcal{L}(\theta_1)-\mathcal{L}(\theta_2)| \le L_{\Theta}(\delta) \ln(N) \|\theta_1 - \theta_2\|
%%       + \frac{1}{\lamda} \mathcal{C}_{ub}(\lambda,L_{\ell},N)
%%      \label{pac:lipsch:loss2:eq2}
%%     \end{align}

\textbf{Quadratic case}
Recall from \eqref{eq:initial:tildern1:quadratic0}, Subsection \ref{thm:main:quadratic} the definition of the loss
  function $\ell_C$ and the true and empirical errors
  ${}^C \mathcal{L}(\theta)$ and ${}^C \mathcal{L}_N(\theta)$.
  Recall that for $C=6G_e\|\e_g\|_{\psi_2}\sqrt{\ln(\frac{3N^2}{2\delta})}$. 
  %\eqref{thm:main:quadratic:eq21} and 
  %$\hat{\mathcal{L}}_N(\theta)(\omega) = {}^C\hat{\mathcal{L}}_N(\theta)(\omega)$

  Recall that  by Lemma \ref{lem:quadratic1} that 
  %\eqref{eq:initial:tildern1} %and hence \eqref{catoni1} for
  %the quadratic case are implied by the inequality
   \begin{equation}
  \label{quadratic:lipschitz:eq1}
  \begin{aligned}
   %& \mbox{\eqref{eq:initial:tildern1} holds for $\mathcal{L}(\theta)$ and $\mathcal{L}_N(\theta)$ replaced by ${}^C \mathcal{L}(\theta)$ and ${}^C \hat{\mathcal{L}}_N(\theta)$, and $\delta$ replaced by $\frac{2\delta}{3}$} \\
   & \sup_{t=0,\ldots,N-1, \theta \in \Theta} \max\{\|\y(t)-\y_{\theta}(t)\|_2, \|\y(t)-\y_{\theta}(t[0])\|_\infty\} \le C %6G_e (\|\e_g\|_{\psi_2})\sqrt{\ln\left(\frac{3N^2}{2\delta}\right)}
   \end{aligned}
   \end{equation}
   holds with probability at least $1-\frac{\delta}{3}$. 
   Also recall that $\ell_C$ is $L_{\ell}$-Lipschitz with   $L_{\ell} = 12G_e \|\e_g\|_{\psi_2} \sqrt{\ln\left(\frac{3N^2}{2\delta}\right)}$.

   By applying \eqref{pac:lipschitz:loss:eq8} to ${}^C \mathcal{L}_N(\theta)$ and ${}^C \mathcal{L}(\theta)$
   and replacing $\delta$ with $\frac{2\delta}{3}$ and using the union bound, it follows that
\begin{equation}
  \label{quadratic:lipschitz:eq3}
  \begin{aligned}
%   & \mbox{\eqref{catoni1} holds for $\mathcal{L}(\theta)$ and $\mathcal{L}_N(\theta)$ replaced by ${}^C \mathcal{L}(\theta)$ and ${}^C \hat{\mathcal{L}}_N(\theta)$, and $\delta$ replaced by $\frac{2\delta}{3}$} \\
   & \sup_{t=0,\ldots,N-1,\theta \in \Theta} \max\{\|\y(t)-\y_{\theta}(t)\|_2, \|\y(t)-\y_{\theta}(t[0])\|_\infty\} \le 6G_e (\|\e_g\|_{\psi_2})\sqrt{\ln\left(\frac{3N^2}{2\delta}\right)} \\
   %& \forall \theta1,\theta_2 \in \Theta: \\
    &      \sup_{\theta_1,\theta_2 \in \Theta}\left|{}^C \mathcal{L}_N(\theta_1) - {}^C \mathcal{L}_N(\theta_2)\right| \\
         & \le L_{\alpha}L_{\ell}\frac{\sqrt{n_\e}}{K\|\alpha_g\|_{\ell_1}} \ln(N) \|\theta_1 - \theta_2\|
           +  2\frac{16L_{\ell}\ln(N)}{\sqrt{N}}\left(C_{u,1} + C_{u,2}+\ln\left(\frac{2}{3\delta}\right)\right), 
  \end{aligned}
\end{equation}
holds with probability at least $1-\delta$, and
\begin{equation}
  \begin{aligned}
       %& \forall \theta1,\theta_2 \in \Theta: \quad
        &  \sup_{\theta_1,\theta_2 \in \Theta} \left|{}^C \mathcal{L}(\theta_1) - {}^C \mathcal{L}(\theta_2)\right| \\
        & \le L_{\alpha}L_{\ell}\frac{\sqrt{n_\e}}{K\|\alpha_g\|_{\ell_1}} \ln(N) \|\theta_1 - \theta_2\|
          + 2\frac{16L_{\ell}\ln(N)}{\sqrt{N}}\left(C_{u,1} + C_{u,2}\right)
   \end{aligned}
   \end{equation}
     where 
    $L_{\ell}=12G_e \|\e_g\|_{\psi_2} \sqrt{\ln\left(\frac{3N^2}{2\delta}\right)}$. 
    By using 
    \[ \sqrt{\ln\left(\frac{3N^2}{2\delta}\right)} \le \sqrt{2}(\ln(N))^{1/2} \sqrt{1+\ln\left(\frac{3}{2\delta}\right)} \]
    it follows that 
    \begin{align*}
        L_{\alpha}L_{\ell}\frac{\sqrt{n_\e}}{K} \ln(N)  \le
        12 L_{\alpha}G_e \|\e_g\|_{\psi_2} \sqrt{1+\ln\left(\frac{3}{2\delta}\right)}\sqrt{2}\frac{\sqrt{n_\e}}{K}
        (\ln(N))^{3/2}
    \end{align*}
    Recall from Subsection \ref{thm:main:quadratic}, proof of Theorem \ref{thm:main} that 
    if 
    \[
         \forall: \quad   t=0,1,\ldots,N-1: \quad  \max\{\|\y(t)-\y_{\theta}(t)\|_2, \|\y(t)-\y_{\theta}(t \mid 0)\|_\infty\} \le C
         %6G_e (\|\e_g\|_{\psi_2})\sqrt{\ln\left(\frac{3N^2}{2\delta}\right)}
    \]
   holds, then ${}^C \hat{\mathcal{L}}_N(\theta)=\hat{\mathcal{L}}_N(\theta)$.
   Moreover, recall \eqref{thm:main:quadratic:eq22}-\eqref{thm:main:quadratic:eq21} holds, and thus
   using \eqref{thm:bounded:alt:col:pf:eq6.0}, it follows that
   \begin{equation}
  \label{quadratic:lipschitz:eq04}
   \begin{aligned}
       \left|\mathcal{L}(\theta)-{}^C\mathcal{L}(\theta)\right| & \le 72G_e^2\|\e_g\|_{\psi_2}^2\frac{\ln\left(\frac{3N^2}{2\delta}\right)}{N} \le 72 \frac{L^2}{192^2} \frac{\ln(N)}{N} = \\
       & \frac{6}{16 \cdot 192} L^2 \frac{\ln(N)}{N} \le  \frac{L}{512} \frac{L(\ln(N))^{\frac{3}{2}}}{\sqrt{N}}
   \end{aligned}
  \end{equation}
  In then follows that
  \begin{equation}
  \label{quadratic:lipschitz:eq4}
   \begin{aligned}
       & \left|\mathcal{L}(\theta_1)-\mathcal{L}(\theta_2)\right| \le 
       \left|\mathcal{L}(\theta_1)-{}^C\mathcal{L}(\theta_1)\right| + \left|{}^C\mathcal{L}(\theta_1)-{}^C\mathcal{L}(\theta_2)\right| + \left|{}^C\mathcal{L}(\theta_2)-\mathcal{L}(\theta_2)\right| \\
       & \le \frac{\sqrt{n_\e}}{K} L\ln(N) \|\theta_1 - \theta_2\|
          + 2\frac{16L_{\ell}\ln(N)}{\sqrt{N}}\left(C_{u,1} + C_{u,2}\right) + \frac{2L}{512} \frac{L(\ln(N))^{\frac{3}{2}}}{\sqrt{N}}
       %&+L_{\alpha}L_{\ell}\frac{\sqrt{n_\e}}{K} L\ln(N) \|\theta_1 - \theta_2\|  
       %72G_e^2\|\e_g\|_{\psi_2}^2\frac{\sqrt{\ln\left(\frac{3N^2}{2\delta}\right)}}{N} \le \frac{6}{16*196} L^2 \frac{\ln(N)}{N} = 
   \end{aligned}
  \end{equation}
   By repeating the argument of the proof of Theorem \ref{thm:main} for the quadratic case,
   and using \eqref{quadratic:lipschitz:eq4} and 
   \begin{align*}
     & \frac{1}{L} 12 L_{\alpha}G_{e} \|\e_g\|_{\psi_2} \sqrt{1+\ln\left(\frac{3}{2\delta}\right)}\sqrt{2}\frac{\sqrt{n_\e}}{K}= \\
       & \frac{12 L_{\alpha}G_{e} \|\e_g\|_{\psi_2} \sqrt{1+\ln\left(\frac{3}{2\delta}\right)}\sqrt{2}\frac{\sqrt{n_\e}}{K}}{192 G_e \|\e_g\|_{\psi_2}\sqrt{2}\sqrt{1+\ln\left(\frac{3}{2\delta}\right)}}=
        L_{\alpha}\frac{\sqrt{n_\e}}{16K}
   \end{align*}
   it then follows that
   \eqref{quadratic:lipschitz:eq3} implies that 
   \begin{align}
   %& \mbox{\eqref{catoni1} holds with constants for the quadratic costs} \\
   %& \sup_{t=0,\ldots, N-1, \theta \in \Theta} \max\{\|\y(t)-\y_{\theta}(t)\|_2, \|\y(t)-\y_{\theta}(t \mid 0)\|_\infty\} \le C \\
   %6G_e (\|\e_g\|_{\psi_2})\sqrt{\ln\left(\frac{3N^2}{2\delta}\right)} \\
   & \mbox{\eqref{quadratic:lipschitz:eq1} holds} \\
   & \sup_{\theta_1,\theta_2 \in \Theta} |\mathcal{L}_N(\theta_1) - \mathcal{L}_N(\theta_2)| \\
   & \le L_{\Theta}(\delta) L(\ln(N))^{3/2} \|\theta_1 - \theta_2\|
           +  \frac{2L\ln(N)^{r}}{\sqrt{N}}\left(C_{u,1} + C_{u,2}+\ln\left (\frac{2}{3\delta}\right)\right)
  \label{quadratic:lipschitz:eq41}
   \end{align}
   which is implies \eqref{pac:lipschitz:loss:eq31}, and 
   \begin{align}
           \\
       & \sup_{\theta_1,\theta_2 \in \Theta}
          |\mathcal{L}(\theta_1) - \mathcal{L}(\theta_2)| \\
          &  \le L_{\Theta} L(\ln(N))^{3/2} \|\theta_1 - \theta_2\|
          + \frac{2L\ln(N)^{3/2}}{\sqrt{N}}\left(C_{u,1} + C_{u,2}\right) + \\
          & +  2\frac{L}{512} \frac{L (\ln(N))^{3/2}}{\sqrt{N}}
  \label{quadratic:lipschitz:eq42}
   \end{align}
  which is equivalent to
  \eqref{pac:lipschitz:loss:eq32}. 
   Since \eqref{quadratic:lipschitz:eq3} holds with probability at least $1-\delta$, it then follows that
   \eqref{pac:lipschitz:loss:eq31} also holds with probability at least $1-\delta$.

   Finally, from Corollary \ref{thm:bounded:alt:col} it follows that \eqref{catoni1} holds with probability at least $1-\delta$, and hence by the union bound \eqref{quadratic:lipschitz:eq3} follows.
   
   
  % from the proof of Theorem \ref{thm:main}, Subsection \eqref{thm:main:quadratic} it follows that 
  % that with probability at least $1-2\delta$, \eqref{eq:initial:tildern1} and  \eqref{quadratic:lipschitz:eq1} 
  % hold simultanously. Moreover, from the proof of Corollary \ref{thm:bounded:alt:col
  % In particular, 
  % as \eqref{catoni1} is implied by \eqref{eq:initial:tildern1} for
  % $\lambda=\frac{\sqrt{N}}{L\ln(N)^{3/2}}$, hence by using the union bound 
  % for \eqref{caton1},\eqref{quadratic:lipschitz:eq1} and \eqref{}
   
  % it follows that  
   %192 G_e\|\e_g\|_{\psi_2}\sqrt{\ln\left(\frac{3N^2}{2\delta}\right)}\ln(N)}$, it follows that
  % \begin{equation}
  %\label{quadratic:lipschitz:eq2}
  %\begin{aligned}
  % & \mbox{\eqref{catoni1} holds, and \eqref{quadratic:lipschitz:eq1} holds} 
   %for $\mathcal{L}(\theta)$ and $\mathcal{L}_N(\theta)$ replaced by ${}^C \mathcal{L}(\theta)$ and ${}^C \hat{\mathcal{L}}_N(\theta)$, and $\delta$ replaced by $\frac{2\delta}{3}$} \\
   %& \forall \quad   t=0,1,\ldots,N-1: \quad  \sup_{\theta \in \Theta}\max\{\|\y(t)-\y_{\theta}(t)\|_2, \|\y(t)-\y_{\theta}(t[0])\|_\infty\} \le C %6G_e (\|\e_g\|_{\psi_2})\sqrt{\ln\left(\frac{3N^2}{2\delta}\right)}
   %\end{aligned}
   %\end{equation}
   %holds with probability at least $1-2\delta$, and as \eqref{quadratic:lipschitz:eq3}

   


%empirical and true errors are Lipschitz even for quadratic loss function.
%  If $\ell$ is quadratic, then for any   $0 < \lambda < \frac{N}{12G_e^2\|\e_g\|_{\psi_2}^2\ln(\frac{3N^2}{2\delta})}$,
%    \begin{align}
%      & \bP( \ge 1-3\delta
%      \label{quadratic:lipschitz2:eq1}
%          \\
 %     & \forall \theta_1,\theta_2: |\mathcal{L}(\theta_1)-\mathcal{L}(\theta_2)| \le L_{\Theta}(\delta) \ln(N) \|\theta_1 - \theta_2\|
 %      + \frac{1}{\lamda} \mathcal{C}_{ub}(\lambda,L_{\ell},N)
 %     \label{quadratic:lipschitz2:eq2}
 %    \end{align}
\end{proof}
\begin{proof}[Proof of Corollary \ref{lem:pac}]
 Consider the inequality \eqref{pac:lipschitz:loss:eq3} which holds with probability at least $1-(2+r_{\mathrm{Lip,1}})\delta$.
If \eqref{pac:lipschitz:loss:eq3} holds, then for all $\theta \in \Theta$,
\[ \epsilon (\mathcal{L}(\theta)) - (\mathcal{L}(\tilde{\theta}))  \le |\epsilon (\mathcal{L}(\theta)-\mathcal{L}(\tilde{\theta}))|=|\mathcal{L}(\theta)-\mathcal{L}(\tilde{\theta})| \]
and hence 
\[
   \epsilon \mathcal{L}(\theta) \le \epsilon \mathcal{L}(\tilde{\theta}) 
      + L_{\Theta}(\delta) L(\ln(N))^r \|\theta - \tilde{\theta}\|_2
          + 2r_{\text{Lip},1}\frac{L\ln(N)^{r}}{\sqrt{N}}\left(C_{u,1} + C_{u,2}\right) + \\
            2r_{\text{Lip},2}  \frac{L}{512} \frac{L(\ln(N))^r}{\sqrt{N}}
\]
Then $E_{\tilde{\theta} \sim \rho_{\theta}} \|\tilde{\theta}-\theta\|^2_2 \le \frac{\sigma^2}{N}$ and hence
 \[ E_{\tilde{\theta} \sim \rho_{\theta}} \|\tilde{\theta}-\theta\|_2 \le  \sqrt{E_{\tilde{\theta} \sim \rho_{\theta}} \|\tilde{\theta}-\theta\|^2_2} \le \frac{\sigma}{\sqrt{N}} \]
Therefore
\begin{equation}
  \label{lem:pac:eq1}
\begin{aligned}
    & \mathcal{L}(\theta) \le E_{\tilde{\theta} \sim \rho_{\theta}}\mathcal{L}(\tilde{\theta}) + 
     L_{\Theta}(\delta) L(\ln(N))^r   \|\theta - \tilde{\theta}\|_2 \\
          & + 2r_{\text{Lip},1}\frac{L\ln(N)^{r}}{\sqrt{N}}\left(C_{u,1} + C_{u,2}\right) + 
          2r_{\text{Lip},2}  \frac{L}{512} \frac{L(\ln(N))^r}{\sqrt{N}} \\
          & \le \frac{\sigma L_{\Theta}(\delta)}{\sqrt{N}}+2r_{\text{Lip},1}\frac{L\ln(N)^{r}}{\sqrt{N}}\left(C_{u,1} + C_{u,2}\right) + 
          2r_{\text{Lip},2}  \frac{L}{512} \frac{L(\ln(N))^r}{\sqrt{N}}
\end{aligned}
\end{equation}
Similary, 
\begin{align*}
& \epsilon \hat{\mathcal{L}}_N(\tilde{\theta}) \le \epsilon \hat{\mathcal{L}}_N(\theta) + 
L_{\Theta}(\delta) L(\ln(N))^r \|\theta - \tilde{\theta}\|_2
           +  2r_{\text{Lip},1}\frac{L\ln(N)^{r}}{\sqrt{N}}\left(C_{u,1} + C_{u,2}+\ln\left(\frac{1}{\delta}\right)\right)
\end{align*}
and hence 
\begin{equation}
  \label{lem:pac:eq2}
\begin{aligned}
    & \epsilon E_{\tilde{\theta} \sim \rho_{\theta}}\hat{\mathcal{L}}_N(\tilde{\theta})  \le  \epsilon \hat{\mathcal{L}}_N(\theta)+
     L_{\Theta}(\delta) L(\ln(N))^r E_{\tilde{\theta} \sim \rho_{\theta}}\|\theta - \tilde{\theta}\|_2 \\
          & + 2r_{\text{Lip},1}\frac{L\ln(N)^{r}}{\sqrt{N}}\left(C_{u,1} + C_{u,2}+\ln\left(\frac{1}{\delta}\right)\right) \\
          & \le \frac{\sigma L_{\Theta}(\delta)}{\sqrt{N}}+2r_{\text{Lip},1}\frac{L\ln(N)^{r}}{\sqrt{N}}\left(C_{u,1} + C_{u,2}
          \ln\left(\frac{1}{\delta}\right)
          \right)
\end{aligned}
\end{equation}
In addition, since \eqref{catoni1} also holds, it follows that
\begin{equation}
  \label{lem:pac:eq3}
    \epsilon E_{\tilde{\theta} \sim \rho_{\theta}}\mathcal{L}(\tilde{\theta}) \le \epsilon E_{\tilde{\theta} \sim \rho_{\theta}}\hat{\mathcal{L}}_N(\tilde{\theta}) + \frac{L (\ln(N))^r}{\sqrt{N}} \left( \KL(\rho_{\theta} \| \pi) + \ln\left(\frac{C_{\delta}}{\delta}+C(1,\delta)\right) \right)
\end{equation}
By combining \eqref{lem:pac:eq1}, \eqref{lem:pac:eq2}, and \eqref{lem:pac:eq3}, it follows that
\begin{align*}
   \epsilon \mathcal{L}(\theta) \le & \epsilon E_{\tilde{\theta} \sim \rho_{\theta}}\mathcal{L}(\tilde{\theta}) + \frac{\sigma L_{\Theta}(\delta)}{\sqrt{N}}+2r_{\text{Lip},1}\frac{L\ln(N)^{r}}{\sqrt{N}}\left(C_{u,1} + C_{u,2}\right)  \\
          & + 2r_{\text{Lip},2}  \frac{L}{512} \frac{L(\ln(N))^r}{\sqrt{N}} \\
          & \le \epsilon E_{\tilde{\theta} \sim \rho_{\theta}}\hat{\mathcal{L}}_N(\tilde{\theta}) + \frac{L (\ln(N))^r}{\sqrt{N}} \left( \KL(\rho_{\theta} \| \pi) + \ln\left(\frac{C_{\delta}}{\delta}\right)+C(1,\delta)\right) + \\
          & + \frac{\sigma L_{\Theta}(\delta)}{\sqrt{N}}+2r_{\text{Lip},1}\frac{L\ln(N)^{r}}{\sqrt{N}}\left(C_{u,1} + C_{u,2}\right) + 2r_{\text{Lip},2}  \frac{L}{512} \frac{L(\ln(N))^r}{\sqrt{N}} \\
          & \le \epsilon \hat{\mathcal{L}}_N(\theta) + 
\frac{\sigma L_{\Theta}(\delta)}{\sqrt{N}}
           +  2r_{\text{Lip},1}\frac{L\ln(N)^{r}}{\sqrt{N}}\left(C_{u,1} + C_{u,2}+\ln\left(\frac{1}{\delta}\right)\right) \\
           & + \frac{L (\ln(N))^r}{\sqrt{N}} \left( \KL(\rho_{\theta} \| \pi) + \ln\left(\frac{C_{\delta}}{\delta}\right)+C(1,\delta) \right) + \\
           & + \frac{\sigma L_{\Theta}(\delta)}{\sqrt{N}}+2r_{\text{Lip},1}\frac{L\ln(N)^{r}}{\sqrt{N}}\left(C_{u,1} + C_{u,2}\right) + 2r_{\text{Lip},2}  \frac{L}{512} \frac{L(\ln(N))^r}{\sqrt{N}} \\
           & = \epsilon \hat{\mathcal{L}}_N(\theta)+\frac{2\sigma L_{\Theta}(\delta)}{\sqrt{N}} 
            + 4r_{\text{Lip},1}\frac{L\ln(N)^{r}}{\sqrt{N}}\left(C_{u,1} + C_{u,2}\right) + 2r_{\text{Lip},2}  \frac{L}{512} \frac{L(\ln(N))^r}{\sqrt{N}}  \\
           & + \frac{L(\ln(N))^r}{\sqrt{N}} \left(C(\theta) \ln(N) + \ln\left(\frac{C_{\delta}}{\delta^{2r_{\text{Lip},1}+1}}\right) + C(\epsilon,\delta) \right) \\
           & \le \epsilon \hat{\mathcal{L}}_N(\theta) + \frac{L(\ln(N))^{r+1}}{\sqrt{N}} \left(2\sigma L_{\Theta}(\delta)+4r_{\text{Lip},1}\left(C_{u,1} + C_{u,2}\right)  \right. \\
           & \left. + 2r_{\text{Lip},2}  \frac{L}{512}+ \left(C(\theta) + \ln\left(\frac{C_{\delta}}{\delta^{2r_{\text{Lip},1}+1}}\right) + C(\epsilon,\delta) \right)
           \right)
\end{align*}
where in the last step we used that $r \ge 0$ and $\ln(N) \ge 1$ for $N \ge 4$.
\end{proof}


\section{Proof of Corollary \ref{lem:bound:mean}}
\label{sec:proof_single_draw}

We start with the following remark explaining the notion of random models drawn
from a data-dependent posterior. 
\begin{remark}[Joint probability over data and parameters]
  \label{rem:jointProb}
Let $\rho^{\train}$ be a function on $\Omega \times \Theta$
which is measurable w.r.t. the direct product $\F \times \B$ of $\sigma$-algebras, and such that
for any $\omega \in \Omega$,
$\rho^{\train}(\omega): \Theta \ni \theta \mapsto \rho^{\train}(\omega,\theta)$ is a 
probability density on $\Theta$.
We refer to $\rho^{\train}$ as a \emph{random density
on $\Omega$}. Let $B$ be a subset of $\F \times \B$ and denote by
$P_{\theta \sim \rho^{\train}}(B)$ the random variable
$\omega \mapsto P_{\theta \sim \rho^{\train}(\omega)} (\{ \theta \mid (\omega, \theta) \in B\})$, where 
$P_{\theta \sim \rho^{\train}(\omega)}$ is the probability measure induced by the density $\rho^{\train}(\omega)$.
Define the probability measure
$\bP \times P_{\theta \sim \rho^{\train}}$
on $\F \times \B$ as follows:
\[
   \left(\bP \times P_{\theta \sim \rho^{\train}}\right)(B) \triangleq
   \bE[P_{\theta \sim \rho^{\train}}(B)]
   %\textcolor{red}{\bE[P_{\theta \sim \rho^{\train}}~|~B]}
\]
Similarly to the previous definition of relationship holding with probability over data, if $\diamond$ is a binary operator and 
$\{X_i,Y_i\}_{i \in I}$ is family of functions taking values in $\mathbb{R}$ and which are measurable w.r.t. $(\Omega \times \Theta,\F \times \B)$, 
then we say that $\forall i \in I: X_i \diamond  Y_i$, holds with probability  at least $1-\delta$ over data and over samples
from $\rho^{\train}$, if the set
$B=\{ (\omega,\theta) \in \Omega \times \Theta \mid \forall i \in I: X_i(\omega,\theta) \diamond  Y_i(\omega,\theta)\}$ is
measurable in $(\Omega \times \Theta,\F \times \B)$ and
$\left(\bP \times P_{\theta \sim \rho^{\train}}\right)(B) \geq 1-\delta$. 

Note that if a relationship holds with a probability at least $1-\delta$ over data, then it also holds with a probability at least $1-\delta$ over data and sampled from $\rho^{\train}$: if $F$ is an event from $\mathcal{F}$, then 
$
\left(\bP \times P_{\theta \sim \rho^{\train}}\right)(F) = \bP(F) \geq 1-\delta.
$

\end{remark}



\subsection{Proof of the \eqref{bound:emp} on empirical error}
%\label{sec:proof_single_draw}
First, we present a useful property of Gibbs posteriors.  
 \begin{lemma}[Property of Gibbs distribution]
  \label{lem:gibbs_property}
  \begin{align}
    \hat{\mathcal{L}}_N(\theta)+\frac{1}{\lambda}\ln\left(\frac{\rho^{\text{Gibbs}}(\theta)}{\pi(\theta)}\right) 
       = &  -\frac{1}{\lambda}\ln E_{\theta \sim \pi} e^{-\lambda \hat{\mathcal{L}}_N(\theta)} \nonumber \\
    & =\inf_{\rho \in \mathcal{M}_{\pi}} \left(\bE_{\theta \sim \rho} \hat{\mathcal{L}}_N(\theta) + \frac{1}{\lambda} \KL(\rho\|\pi)\right)
  \label{lem:gibbs_property:eq1} \\
   E_{\theta \sim \rho^{\text{Gibbs}}}\hat{\mathcal{L}}_N(\theta)+\frac{1}{\lambda} 
      \KL(\rho^{\text{Gibbs}}\|\pi) & =  -\frac{1}{\lambda} \ln   \bE_{\theta \sim \pi} e^{-\lambda \hat{\mathcal{L}}_N(\theta)} = \nonumber \\
    & =\inf_{\rho \in \mathcal{M}_{\pi}} \left(\bE_{\theta \sim \rho}(\hat{\mathcal{L}}_N(\theta)) + \frac{1}{\lambda} \KL(\rho\|\pi)\right)
  \label{lem:gibbs_property:eq2}
  \end{align}
 \end{lemma}
 \begin{proof}[Proof of Lemma \ref{lem:gibbs_property}]
  Using the definition of the Gibbs distribution $\rho^{\text{Gibbs}}(\theta)=\frac{\pi(\theta)e^{-\lambda \hat{\mathcal{L}}_N(\theta)}}{Z}$, $Z=E_{\theta \sim \pi} e^{-\lambda \hat{\mathcal{L}}_N(\theta)}$, it follows that
  \[ \ln\left(\frac{\rho^{\text{Gibbs}}(\theta)}{\pi(\theta)}\right) = -\lambda \hat{\mathcal{L}}_N(\theta) - \ln Z \]
  and hence 
  \[ \hat{\mathcal{L}}_N(\theta)+\frac{1}{\lambda}\ln\left(\frac{\rho^{\text{Gibbs}}(\theta)}{\pi(\theta)}\right) = \hat{\mathcal{L}}_N(\theta) - \hat{\mathcal{L}}_N(\theta) - \frac{1}{\lambda} \ln Z = - \frac{1}{\lambda} \ln Z \]
%Note that as $\ln$ is concave, $-\ln$ is convex, and hence by Jensens inequality
%Note that by the well-known property of
%Gibbs densities 
%-\frac{1}{\lambda} \ln Z = -\frac{1}{\lambda} \ln \bE_{\theta \sim \pi} e^{-\lambda \hat{\mathcal{L}}_N(\theta)} \le 
%     \inf_{\rho \in \mathcal{M}_{\pi}} \left(\bE_{\theta \sim \rho} \hat{\mathcal{L}}_N(\theta) + \frac{1}{\lambda} \KL(\rho\|\pi)\right)
%\]
This gives the first equality of \eqref{lem:gibbs_property:eq1}.
%gives the first desired equality.

As to  the first equality of \eqref{lem:gibbs_property:eq2}, it follows directly from the first equality of  \eqref{lem:gibbs_property:eq1} by taking the expectation with respect to $\theta \sim \rho^{\text{Gibbs}}$ and recalling that $\KL(\rho^{\text{Gibbs}}\|\pi) = \bE_{\theta \sim \rho^{\text{Gibbs}}} \left[\ln \frac{\rho^{\text{Gibbs}}(\theta)}{\pi(\theta)}\right]$.

The second equalities of \eqref{lem:gibbs_property:eq1}-\eqref{lem:gibbs_property:eq2} are 
%the first inequality of \eqref{lem:gibbs_property:eq2} and from the 
Donsker-Varadhan's variational formula \cite[Lemma 2.2]{alquier2021userfriendly}

%\[ E_{\theta \sim \rho^{\text{Gibbs}}}[\hat{\mathcal{L}}_N(\theta)] + \frac{1}{\lambda} \KL(\rho^{\text{Gibbs}}\|\pi) =  -\frac{1}{\lambda} \ln \frac{Z}=
%   \inf_{\rho \in \mathcal{M}_{\pi}} \left(\bE_{\theta \sim \rho}[\hat{\mathcal{L}}_N(\theta)] + \frac{1}{\lambda} \KL(\rho\|\pi)\right) =
%\]
%of the Gibbs distribution and the properties of the KL divergence.
 \end{proof}
 Next, we present a version of Theorem \ref{thm:main} where the inequality
holds not only  over data but over data and 
 any model randomly drawn from a data-dependent posterior.
 \begin{lemma}[Single draw bound]
\label{lem:single_draw}
 %That is, 
 Let $\rho^{\train}$ be a data-dependent posterior. 
 For any prior $\pi$,  any $\delta \in (0,0.5)$
 with a probability at least  $(1-2\delta)$ over the data and over all samples $\theta_{\star}$ drawn from $\rho_{\train}$,
 for any $\lambda$ satisfying the condition of Theorem \ref{thm:main}
  \begin{equation}
\label{RenyiBound:single_draw}
 \epsilon \mathcal{L}(\theta_{\star}) \le  \epsilon  \hat{\mathcal{L}}_N(\theta_{\star})+ 
 \frac{1}{\lambda} \left(\ln \frac{\rho^{\train}(\theta_{\star})}{\pi(\theta_{\star})} + \ln\dfrac{C_{\delta}}{\delta} + \frac{1}{2}\sum_{i=1}^{2} \ln \bE_{\theta \sim \pi} \exp(2\lambda \mathcal{C}_i(2\lambda,\theta_{\star},N,\delta,\epsilon))\right)
  \end{equation}
  In particular, with probability at least $(1-2\delta)$ over data and over all samples  $\theta_{\star}$ drawn from $\rho_{\train}$, 
  \begin{equation}
  \label{RenyiBound:single_draw1}
  \epsilon \mathcal{L}(\theta_{\star}) \le \epsilon \hat{\mathcal{L}}_N(\theta_{\star})+\frac{L(\ln(N))^{r}}{\sqrt{N}}(\ln\left(\frac{\rho^{\train}(\theta_{\star})}{\pi(\theta_{\star})}\right)+\ln\frac{C_{\delta}}{\delta} + 
  \mathcal{C}(\delta,\epsilon))
  \end{equation}
  for all $N \ge K_{\delta}$, where the constants $K_{\delta},C_{\delta},L,r$ are defined as in Corollary \ref{thm:bounded:alt:col}.
 Moreover, if  $\epsilon=1$ and 
 $\rho^{\train}=\rho_{\text{Gibbs}}$, then with probability at  least  $(1-2\delta)$
  over the data and over $\theta_{\star}$
  sampled from $\rho_{\text{Gibbs}}$,
  \begin{equation}
\label{RenyiBound:single_draw:giibs}
\begin{aligned}
 \mathcal{L}(\theta_{\star}) \le  
 %\frac{1}{\lambda} \left(-\ln(E_{\theta \sim \pi} e^{-\lambda \hat{\mathcal{L}}_N(\theta)})+ \ln\dfrac{C_{\delta}}{\delta} + \frac{1}{2}\sum_{i=1}^{2} \ln \bE_{\theta \sim \pi} \mathcal{C}_i(2\lambda,\theta,N,\delta,\epsilon))\right)  \\
 E_{\theta \sim \pi} \hat{\mathcal{L}}_N(\theta) +  \frac{1}{\lambda} \left(\ln\dfrac{C_{\delta}}{\delta} + \frac{1}{2}\sum_{i=1}^{2} \ln \bE_{\theta \sim \pi} \exp(\mathcal{C}_i(2\lambda,\theta,N,\delta,\epsilon))\right) \\
\end{aligned}
  \end{equation}
  In particular, if the Gibbs posterior $\rho^{\text{Gibbs}}$ is defined for 
  $\lambda=\lambda_N=\frac{\sqrt{N}}{L(\ln(N))^r}$ then for $N \ge K_{\delta}$
\begin{equation}
\label{RenyiBound:single_draw:giibs1}
\begin{aligned}
 \mathcal{L}(\theta_{\star}) \le  %&   \frac{1}{\lambda} \left(-\ln(E_{\theta \sim \pi} e^{-\lambda \hat{\mathcal{L}}_N(\theta)})+ \ln\dfrac{C_{\delta}}{\delta} + \frac{1}{2}\sum_{i=1}^{2} \ln \bE_{\theta \sim \pi} \exp(\mathcal{C}_i(2\lambda,\theta,N,\delta,\epsilon))\right)) \\
  -\frac{1}{\lambda} \ln E_{\theta \sim \pi} e^{-\lambda \hat{\mathcal{L}}_N(\theta)} + 
 \frac{L(\ln(N))^{r}}{\sqrt{N}}\left(\ln\frac{C_{\delta}}{\delta} +  \mathcal{C}(\delta,1)\right)
\end{aligned}
\end{equation}
% \frac{1}{\lambda}\left(\ln\dfrac{1}{\delta} + \Psi(\lambda,N) - Z \right)
%\frac{\bar{\mathcal{D}}_2(\rho^{\train}\|\pi) K}{
 % \sqrt{\delta  N}\delta_1}\left [ G_1 + \frac{4}{\sqrt{N}}G_2\right ]
%  \end{equation}
%  where $Z$ is the normalization constant from \eqref{eq:gibbs}.
  \end{lemma}
\begin{proof}[Proof of Lemma \ref{lem:single_draw}]
 Following the argument of the proof of \cite[Theorem 2.7]{alquier2021userfriendly}, it follows
 that for any measurable function
 $X$ on $\Omega \times \Theta$,
 $E_{\theta \sim \pi} e^{2\lambda X(\omega,\theta)} = E_{\theta \sim \rho^{\train}} e^{2\lambda X(\omega,\theta)-\ln \frac{\rho^{\train}(\omega,\theta)}{\pi(\theta)}}$ and hence 
 $\bE E_{\theta \sim \pi} e^{2\lambda X(\theta)}=\bE E_{\theta \sim \rho^{\train}}
 e^{2\lambda X(\theta)-\ln \frac{\rho^{\train}(\theta)}{\pi(\theta)}}$
 with the random variables $X(\theta): \omega \ni \Omega \mapsto X(\omega,\theta)$.
 %$\rho^{\train}(\theta): \omega \ni \Omega \mapsto \rho^{\train}(\omega,\theta)$.
 Notice that $\bE E_{\theta \sim \rho^{\train}}$ is the expectation operator
 corresponding to the probability measure
 $\bP \times P_{\theta \sim \rho^{\train}}$, and that
$\bE[E_{\theta \sim \pi} X]=E_{\theta \sim \pi} \bE[X]$ 
 , and hence, by Chernoff's bound:
 \begin{align*}
     & \bP \times P_{\theta \sim \rho^{\train}}\big (X < \frac{1}{2\lambda}\big(\ln \frac{1}{\delta} +  \ln \frac{\rho^{\train}(\theta)}{\pi(\theta)}+ \ln \bE E_{\theta \sim \pi} e^{2\lambda X}\big)\big) \\
     & = \bP \times P_{\theta \sim \rho^{\train}}\big ((X - \frac{1}{2\lambda} \ln \frac{\rho^{\train}(\theta)}{\pi(\theta)}) <  \frac{1}{2\lambda}\big(\ln \frac{1}{\delta} + \ln \bE E_{\theta \sim \pi} e^{2\lambda X}\big)\big)
     \le \frac{\bE E_{\theta \sim \rho^{\train}} e^{2\lambda X - \ln \frac{\rho^{\train}(\theta)}{\pi(\theta)}}}{e^{\ln \frac{1}{\delta} + \ln \bE E_{\theta \sim \pi} e^{2\lambda X}}} \\
     & =\frac{\bE E_{\theta \sim \pi} e^{2\lambda X}}{\frac{1}{\delta} \bE E_{\theta \sim \pi} e^{2\lambda X}}=\delta
 \end{align*}
 Hence, 
 \begin{equation}
  \label{lem:single_draw:quadratic:pf0}
   \begin{split}
    %& 1-\delta < \bP \times P_{\theta \sim \rho^{\train}}\big (X < \frac{1}{\lambda}\big(\ln \frac{1}{\delta} +  \\
    %& \ln \frac{\rho^{\train}(\theta)}{\pi(\theta)}+ \ln \bE E_{\theta \sim \rho^{\train}}
    %e^{2\lambda X(\theta)-\ln \frac{\rho^{\trai n}(\theta)}{\pi(\theta)}}\big)\big) \le \\
    %& \bP \times P_{\theta \sim \rho^{\train}}\big(
    X \le  \frac{1}{2\lambda}\big (\ln \frac{1}{\delta} +\ln \frac{\rho^{\train}(\theta)}{\pi(\theta)} + \ln E_{\theta \sim \pi} \bE e^{2\lambda X(\theta)}\big)
   \end{split} 
 \end{equation}
 with probability at least $1-\delta$ over data and samples from $\rho^{\train}$.
 Hence, by taking $X=\epsilon (\mathcal{L}(\theta)-V_N(\theta))$ and then $X=\epsilon (V_N(\theta))-\hat{\mathcal{L}}_N(\theta))$, and 
 by taking the union bound as in the proof of Theorem \ref{thm:initial}, it follows that
 \begin{equation}
\label{lem:single_draw:quadratic:pf1}
 \begin{split}
  \mathcal{L}(\theta) \le & \hat{\mathcal{L}}_N(\theta) + \frac{1}{\lambda} \Big(\ln \frac{1}{\delta} +\ln \frac{\rho^{\train}(\theta)}{\pi(\theta)}  \\
  &+\frac{1}{2} \Big (\ln E_{\theta\sim\pi}\bE[e^{2\epsilon\lambda(V_N(\theta)-\hat{\mathcal{L}}_N(\theta))}]
        +\ln E_{\theta\sim\pi}\bE[e^{2\epsilon\lambda (\mathcal{L}(\theta)-V_N(\theta))}] \Big )\Big)
 \end{split}
 \end{equation}

\textbf{Proof of \eqref{RenyiBound:single_draw} for bounded noise}
 Let us  apply the  inequality \eqref{lem:single_draw:quadratic:pf1} 
 %for $X(\omega,\theta)=\epsilon (\mathcal{L}(\theta)-V_N(\theta)(\omega))$ and 
 %$Y(\omega,\theta)=\epsilon (V_N(\theta)(\omega)-\hat{\mathcal{L}}_N(\theta)(\omega))$, 
 and use
 the bound on $\bE[e^{2\epsilon\lambda (V_N(\theta)-\hat{\mathcal{L}}_N(\theta))}]$ and $\bE[e^{2\epsilon\lambda (\mathcal{L}(\theta)-V_N(\theta))}]$ from
  Lemma \ref{lemma:bounded:mgf(L-V)}-\ref{lemma:bounded_mgf(V-hatL)}, from which \eqref{RenyiBound:single_draw} follows. 
  
  
\textbf{Proof of \eqref{RenyiBound:single_draw} for Lipschitz loss}
 Let us apply  \eqref{thm:main:pf:unbounded:eq1} from Subsection \ref{lipschitz:unbounded} to bound
 the bound on $\bE[e^{2\epsilon\lambda (V_N(\theta)-\hat{\mathcal{L}}_N(\theta))}]$ and $\bE[e^{2\epsilon\lambda (\mathcal{L}(\theta)-V_N(\theta))}]$, from which \eqref{RenyiBound:single_draw} follows.
  %Lemma \ref{{lemma:bounded:mgf(L-V)}, \ref{lemma:bounded_mgf(V-hatL)},\ref{basic_lemma:lipschitz3} 
   
\textbf{Proof of \eqref{RenyiBound:single_draw} quadratic loss}
Recall from Subsection \ref{thm:main:quadratic} the definition of the losses ${}^C \mathcal{L}(\theta)$, ${}^C \hat{\mathcal{L}}_N(\theta)$.
corresponding to the loss $\ell_C$. Note that by Lemma \ref{lem:quadratic2}  $\ell_C$ is $2C$-Lipschitz.
Notice that from the statement of the present lemma  proven for the case of Lipschitz losses
for $\lambda \le \frac{\sqrt{N}}{16 \cdot 2C\ln(N)}$,  with 
$C=6G_e\|\e_g\|_{\psi_2}\sqrt{\ln(\frac{3N^2}{2\delta})}$ probability $1-\frac{4}{3\delta}$ over data and parameters sampled
from $\rho^{\train}$, the following holds:
\begin{equation}
  \label{lem:single_draw:quadratic}
  \begin{aligned}
\epsilon {}^C\mathcal{L}(\theta)
         \leq  &  \epsilon {}^C \hat{\mathcal{L}}_N(\theta) +  
          \frac{1}{\lambda}\left( \ln\left(\frac{\rho(\theta)}{\pi(\theta)}\right) + \ln\dfrac{3}{2\delta} + \right. \\
          &  \left. \frac{1}{2} \sum_{i=1}^2 \ln \bE_{\theta \sim \pi} \exp\left(\bar{C}_i(2\lambda,\theta,N,\frac{2\delta}{3},\epsilon)\right)\right)
  \end{aligned}
\end{equation}
where $\bar{C}_i(2\lambda,\theta,N,\frac{2\delta}{3},\epsilon)$ satisfies \eqref{thm:main:quadratic:eq2}.
From Lemma \ref{lem:quadratic1} it follows that \eqref{lem:quadratic1:eq1} holds with probability at least
$1-\frac{2\delta}{3}$ 
over data and hence \eqref{lem:quadratic1} holds with probability at least $1-\frac{2\delta}{3}$ over data and parameters sampled from $\rho ^{\theta}$. Hence, by using the union bound, it follows that
\eqref{lem:single_draw:quadratic} and \eqref{lem:quadratic1:eq1} hold
simultaneously 
with probability at least  $1-2\delta$ over data and parameters sampled from $\rho^{\train}$.

Moreover, if 
\eqref{lem:quadratic1:eq2} holds, then %for $C= 6G_e (\|\e_g\|_{\psi_2})\sqrt{\ln\left(\frac{3N^2}{2\delta}\right)}$,
it holds that 
${}^C \hat{\mathcal{L}}_N(\theta)=\hat{\mathcal{L}}_N(\theta)$ 
and by and \eqref{thm:main:quadratic:eq21} -- \eqref{thm:main:quadratic:eq22} 
\[ 
    {}^C \mathcal{L}(\theta) \le \mathcal{L}(\theta) \le  {}^C \mathcal{L}(\theta)+ 72G_e^2\|\e_g\|_{\psi_2}^2 \frac{\ln\left(\frac{3N^2}{2\delta}\right)}{N}
\]
Hence, 
\begin{equation} 
  \label{lem:single_draw:quadratic1}
    \epsilon \mathcal{L}(\theta) \le \epsilon {}^C \mathcal{L}(\theta)+ 72G_e^2\|\e_g\|_{\psi_2}^2\frac{\ln\left(\frac{3N^2}{2\delta}\right)}{N} \max\{\epsilon,0\}
\end{equation}
 Hence, by using \eqref{lem:single_draw:quadratic} and ${}^C \hat{\mathcal{L}}_N(\theta)=\hat{\mathcal{L}}_N(\theta)$,
 it follows that if  \eqref{lem:single_draw:quadratic1} and \eqref{lem:quadratic1:eq1} both hold, then 
 \eqref{RenyiBound:single_draw} holds. Since \eqref{lem:single_draw:quadratic1} and \eqref{lem:quadratic1:eq1} holds
 with probability at least $1-2\delta$ over data and samples from $\rho^{\train}$, so does  \eqref{RenyiBound:single_draw}.

 \textbf{Proof of \eqref{RenyiBound:single_draw} implies \eqref{RenyiBound:single_draw:giibs}}
   It follows by applying Lemma \ref{lem:gibbs_property}. 

  \textbf{Proof of \eqref{RenyiBound:single_draw1}}
     From the proof of Corollary \ref{thm:bounded:alt:col} it follows that
     for $\lambda=\frac{\sqrt{N}}{L(\ln(N))^{r}}$  where
     $r$ is as in Corollary \ref{thm:bounded:alt:col},
     \[
     \frac{1}{2\lambda } \left(   \ln E_{\theta \sim \pi} \exp(\mathcal{C}_1(2\lambda,\theta,N,\delta,\epsilon)) +
           \ln E_{\theta \sim \pi} \exp(\mathcal{C}_2(2\lambda,\theta,N,\delta,\epsilon)) \right)
           \le \mathcal{C}(\delta,\epsilon)
     \]
     Using this inequality, it follows that \eqref{RenyiBound:single_draw} implies \eqref{RenyiBound:single_draw1}.
\end{proof}
\begin{lemma}
\label{lem:bound:mean1}
 If $\mathcal{L}(\theta)$ is convex, $\rho$ is a probability density over $\Theta$, and 
 $\theta_{\star}=E_{\theta \sim \rho} \theta$ is the mean 
 of the density $\rho$, then
 \begin{equation}
  \label{lem:bound:mean:eq1}
     \mathcal{L}(\theta_{\star}) \le E_{\theta \sim \rho}[\mathcal{L}(\theta)].
\end{equation}
\end{lemma}
\begin{proof}[Proof of Lemma \ref{lem:bound:mean1}]
 If $\mathcal{L}(\theta)$ is convex in $\theta$, then 
$\mathcal{L}(\theta_{\star})=\mathcal{L}(E_{\theta \sim \rho^{\train}} \theta)   \le E_{\theta \sim \rho^{\train}} \mathcal{L}(\theta)$ by Jensen's inequality.
\end{proof}
\begin{proof}[Proof of \eqref{bound:emp}]
  %\textbf{Single draw case}
    For the single draw case, from \eqref{RenyiBound:single_draw:giibs} of  Lemma \ref{lem:single_draw}, it follows
    that with probability at least $1-2\delta$ over data and samples from $\rho^{\text{Gibbs}}$,
    \begin{equation}
    \label{bound:emp:single_draw:pf:eq1}
        \mathcal{L}(\theta_{\star}) \le -\ln E_{\theta \sim \pi} e^{-\lambda \hat{\mathcal{L}}_N(\theta)} +
          \frac{L(\ln(N))^{r}}{\sqrt{N}}\left(\ln\frac{C_{\delta}}{\delta} +  \mathcal{C}(\delta,\epsilon)\right)
    \end{equation}
    for $\lambda=\sqrt{N}/(L(\ln(N))^{r})$
     For the MEP case,  using Lemma \ref{lem:bound:mean1} and Lemma \ref{lem:gibbs_property} and Corollary \ref{thm:bounded:alt:col}it follows that with probability at least $1-2\delta$ over data, 
     it holds that
     \begin{equation}
    \label{bound:emp:single_draw:pf:eq2}
    \begin{aligned}
        & \mathcal{L}(\theta_{\star}) \le E_{\theta \sim \rho^{\text{Gibbs}}} \mathcal{L}(\theta) \le E_{\theta \sim \rho^{\text{Gibbs}}} \hat{\mathcal{L}}_N(\theta)  + 
          \underbrace{\frac{L(\ln(N))^{r}}{\sqrt{N}}}_{1/\lambda}
          \left(\KL(\rho^{\text{Gibbs}} \mid \pi)+\ln\frac{C_{\delta}}{\delta} +  \mathcal{C}(\delta,\epsilon)\right) \\
       & = \left( E_{\theta \sim \rho^{\text{Gibbs}}} \hat{\mathcal{L}}_N(\theta)  + \frac{1}{\lambda} \KL(\rho^{\text{Gibbs}} \mid \pi)\right) + \frac{L(\ln(N))^{r}}{\sqrt{N}}\underbrace{\left(\ln\frac{C_{\delta}}{\delta} +  \mathcal{C}(\delta,\epsilon)\right)}_{\mathcal{C}_{em}(\delta,\sigma)} \\
       &  = -\ln E_{\theta \sim \pi} e^{-\lambda \hat{\mathcal{L}}_N(\theta)} +
         \frac{L(\ln(N))^{r}}{\sqrt{N}}\mathcal{C}_{em}(\delta,\sigma)
    \end{aligned}      
    \end{equation}
    with $\lambda=\sqrt{N}/(L(\ln(N))^{r})$.
    Using the second equality of \eqref{lem:gibbs_property:eq2}, it follows that
    for any posterior $\rho \in \mathcal{M}_{\pi}$,
    \begin{equation}
    \label{bound:emp:single_draw:pf:eq3}
          -\frac{1}{\lambda}\ln E_{\theta \sim \pi} e^{-\lambda \hat{\mathcal{L}}_N(\theta)} \le
            E_{\theta \sim \rho} \hat{\mathcal{L}}_N(\theta) + \frac{1}{\lambda}\KL(\rho \| \pi)
    \end{equation}
    Choose $\rho_{\theta_{\star}}$ as in Lemma \ref{pac:kl}, i.e., 
    \[ \rho_{\theta_{\star}}=\left\{ \begin{array}{rl}
                           \mathcal{N}(\theta_{\star},\frac{\sigma^2}{N}) & \mbox{ if \textbf{(Gauss)} holds} \\
                            \mathcal{U}(\theta_{\star},\frac{\sigma}{\sqrt{N}})  & \mbox{ if \textbf{(Uni)} holds}
                         \end{array} \right. 
   \]
      From Lemma \ref{pac:lipsch:loss} it follows that \eqref{pac:lipschitz:loss:eq31} holds with probability at least 
  $1-r_{\text{Lip},1}\delta$ over data, and hence over data and samples from $\rho^{\text{Gibbs}}$. 
    In particular, as \eqref{pac:lipschitz:loss:eq31} implies
   \begin{equation}
    \label{bound:emp:single_draw:pf:eq4}
    \begin{aligned}
    E_{\theta \sim \rho_{\theta_{\star}}} \hat{\mathcal{L}}_N(\theta) & \le \hat{\mathcal{L}}_N(\theta_{\star})+
     L_{\Theta}(\delta) L(\ln(N))^{r+1} \underbrace{E_{\theta \sim \rho_{\theta_{\star}}}\|\theta - \theta_{\star}\|_2}_{\le \frac{\sigma}{\sqrt{N}}}
           +  \\
           & 2r_{\text{Lip},1}\frac{L(\ln(N))^{r}}{\sqrt{N}}\left(C_{u,1} + C_{u,2}+\ln\left(\frac{1}{\delta}\right)\right), 
    \end{aligned}
    \end{equation}
   It then follows that \eqref{bound:emp:single_draw:pf:eq4}
   holds with probability at least $1-r_{\text{Lip},1}\delta$
   over data and samples from $\rho^{\text{Gibbs}}$.
   Then using the union bound on \eqref{bound:emp:single_draw:pf:eq1},\eqref{bound:emp:single_draw:pf:eq4},  and \eqref{bound:emp:single_draw:pf:eq3} and Lemma \ref{pac:kl} it follows that
   \begin{equation}
   \label{bound:emp:single_draw:pf:eq5}
   \begin{aligned}
         & \mathcal{L}(\theta_{\star}) \le  E_{\theta \sim \rho^{\text{Gibbs}}} \hat{\mathcal{L}}_N(\theta) 
          %\underbrace{\frac{L(\ln(N))^{r}}{\sqrt{N}}} \left(\KL(\rho^{\text{Gibbs}} \mid \pi)+\ln\frac{C_{\delta}}{\delta} +  \mathcal{C}(\delta,\epsilon)\right) \\
       %& = \left( E_{\theta \le \rho^{\text{Gibbs}}} \hat{\mathcal{L}}_N(\theta)  + \frac{1}{\lambda} \KL(\rho^{\text{Gibbs}} \mid \pi)\right) + \frac{L(\ln(N))^{r}}{\sqrt{N}}\underbrace{\left(\ln\frac{C_{\delta}}{\delta} +  \mathcal{C}(\delta,\epsilon)\right)}_{\mathcal{C}_{em}(\delta,\sigma)} \\
       \le - \ln E_{\theta \sim \pi} e^{-\lambda} \hat{\mathcal{L}}_N(\theta) + 
         \frac{L(\ln(N))^{r}}{\sqrt{N}}\left(\ln\frac{C_{\delta}}{\delta} +  \mathcal{C}(\delta,\epsilon)\right) \\
        & \le E_{\theta\sim \rho_{\theta_{\star}}}\hat{\mathcal{L}}_N(\theta) +
         \frac{1}{\lambda} \KL(\rho_{\theta_{\star}} \| \pi) + 
         \frac{L(\ln(N))^{r}}{\sqrt{N}}\left(\ln\frac{C_{\delta}}{\delta} +  \mathcal{C}(\delta,\epsilon)\right) \\
       & \le  \hat{\mathcal{L}}_N(\theta_{\star})+
     L_{\Theta}(\delta) L(\ln(N))^{r+1} \frac{\sigma}{\sqrt{N}}
           +  %\mathcal{C}(\delta,\epsilon) \\
           2r_{\text{Lip},1}\frac{L\ln(N)^{r}}{\sqrt{N}}\left(C_{u,1} + C_{u,2}+\ln\left(\frac{1}{\delta}\right)\right) \\
           & + \frac{L(\ln(N))^{r+1}}{\sqrt{N}}  C(\theta_{\star}) +  \frac{L(\ln(N))^{r}}{\sqrt{N}}\left(\ln\frac{C_{\delta}}{\delta} +  \mathcal{C}(\delta,\epsilon)\right) \\
           & \hat{\mathcal{L}}_N(\theta_{\star}) +
           \frac{L(\ln(N))^{r+1} \sigma}{\sqrt{N}}
           \big(\underbrace{L_{\Theta}(\delta)\sigma + 2r_{\text{Lip},1}(C_{u,1}+C_{u,2})}_{\mathcal{C}_{L,emp}} + C(\theta_{\star}) \\
           & +(\ln(\frac{C_{\delta}}{\delta})
           + 2r_{\text{Lin},1}\ln(\frac{1}{\sigma}))+C(\delta,\epsilon)\big) 
            \le \hat{\mathcal{L}}_N(\theta_{\star}) + \frac{L(\ln(N))^{r+1} \sigma}{\sqrt{N}} \mathcal{C}_{em}(\delta,\sigma)
    \end{aligned}   
  \end{equation}
  with probability at least $1-(2+r_{\text{Lip},1})\delta$ over data and $\theta_{\star}$ sampled from $\rho^{\text{Gibbs}}$.
  Similarly, by using  that \eqref{bound:emp:single_draw:pf:eq2} holds with probability at least $1-2\delta$, 
  the fact that \eqref{pac:lipschitz:loss:eq31} holds with 
  probability as least $1-r_{\text{Lip},1}\delta$ over data, and that \eqref{pac:lipschitz:loss:eq31} implies 
  \eqref{bound:emp:single_draw:pf:eq3}, it follows,
  using the union bound that for MEP,
  \begin{equation}
   \label{bound:emp:single_draw:pf:eq6}
   \begin{aligned}
         \mathcal{L}(\theta_{\star}) \le & E_{\theta \sim \rho^{\text{Gibbs}}} \hat{\mathcal{L}}_N(\theta) 
          %\underbrace{\frac{L(\ln(N))^{r}}{\sqrt{N}}} \left(\KL(\rho^{\text{Gibbs}} \mid \pi)+\ln\frac{C_{\delta}}{\delta} +  \mathcal{C}(\delta,\epsilon)\right) \\
       %& = \left( E_{\theta \le \rho^{\text{Gibbs}}} \hat{\mathcal{L}}_N(\theta)  + \frac{1}{\lambda} \KL(\rho^{\text{Gibbs}} \mid \pi)\right) + \frac{L(\ln(N))^{r}}{\sqrt{N}}\underbrace{\left(\ln\frac{C_{\delta}}{\delta} +  \mathcal{C}(\delta,\epsilon)\right)}_{\mathcal{C}_{em}(\delta,\sigma)} \\
          \le -\ln e^{-\lambda E_{\theta \sim \pi}\hat{\mathcal{L}}_N(\theta)}  + 
         \frac{L(\ln(N))^{r}}{\sqrt{N}}\left(\ln\frac{C_{\delta}}{\delta} +  \mathcal{C}(\delta,\epsilon)\right)  \\
        & \le E_{\theta\sim \rho_{\theta_{\star}}}\hat{\mathcal{L}}_N(\theta) +\frac{1}{\lambda} \KL(\rho_{\theta_{\star}} \| \pi) +  \frac{L(\ln(N))^{r}}{\sqrt{N}}\left(\ln\frac{C_{\delta}}{\delta} +  \mathcal{C}(\delta,\epsilon)\right) \\
        & \le \hat{\mathcal{L}}_N(\theta_{\star})+
     L_{\Theta}(\delta) L(\ln(N))^{r+1} \frac{\sigma}{\sqrt{N}} 
            + \frac{L(\ln(N))^{r+1}}{\sqrt{N}} \mathcal{C}(\theta_{\star})  \\
            & + 2r_{\text{Lip},1}\frac{L\ln(N)^{r}}{\sqrt{N}}\left(C_{u,1} + C_{u,2}+\ln\left(\frac{1}{\delta}\right)\right) \\
           & +  \frac{L(\ln(N))^{r}}{\sqrt{N}}\left(\ln\frac{C_{\delta}}{\delta} +  \mathcal{C}(\delta,\epsilon)\right) 
            \le \hat{\mathcal{L}}_N(\theta_{\star} + \frac{L(\ln(N))^{r+1} \sigma}{\sqrt{N}} \mathcal{C}_{em}(\delta,\sigma)
    \end{aligned}   
  \end{equation}
  holds with probability at least $1-(2+r_{\text{Lip},1})\delta$ over data.
   
    Finally, the case of MAP follows from Lemma \ref{lem:pac} by applying it to the random model $\theta_{\star}$:
    as  
    \[
    \begin{aligned}
     \forall \theta \in \Theta: 
    \mathcal{L}(\theta) \le 
    \hat{\mathcal{L}}_{N}(\theta) + &
     %\frac{C(\theta)}{\lambda}  \\
     \frac{L(\ln(N))^{r+1}}{\sqrt{N}} \left(
       \sigma L_{\Theta}(\delta) +  2r_{\text{Lip},1} \left(C_{u,1}+C_{u,2}\right) \right. \\
       & \left. +2r_{\text{Lip},2}\frac{L}{512} 
     +C(\theta)+\ln(\frac{C_{\delta}}{\delta})+ 2r_{\text{Lip},1}\ln(\frac{1}{\delta}) + \mathcal{C}(\delta,\epsilon)\right) 
     %\Psi(\lambda,\pi,N)}{\lambda}
    \end{aligned}
    \]
    holds with probability at least $1-(2+r_{\text{Lip},1})\delta$ over data, it then follows that 
    for $\theta=\theta_{\star}$, 
    \[
    \begin{aligned}
    \mathcal{L}(\theta_{\star}) \le 
    \hat{\mathcal{L}}_{N}(\theta_{\star}) + &
     \frac{L(\ln(N))^{r+1}}{\sqrt{N}} \left(
       \sigma L_{\Theta}(\delta) +  r_{\text{Lip},1} \left(C_{u,1}+C_{u,2}\right) \right. \\
       & \left. +2r_{\text{Lip},2}\frac{L}{512} 
     +C(\theta_{\star})+\ln\frac{C_{\delta}}{\delta^{2r_{\text{Lip},1}+1}} + \mathcal{C}(\delta,1)\right) \\
     & = \hat{\mathcal{L}}_{N}(\theta_{\star}) + 
     \frac{L(\ln(N))^{r+1}}{\sqrt{N}} \mathcal{C}_{em}(\delta,\sigma)
    \end{aligned}
    \]
% 
\end{proof}

\subsection{Proof of \eqref{bound:oracle}}
  From \eqref{catoni1} it follows that  each of the inequalities below hold with probability at least $1-2\delta$
  over data
  \begin{align}
      \forall \rho \in \mathcal{M}_\pi: E_{\theta \sim \rho} \mathcal{L}(\theta) \leq & E_{\theta\sim \rho} \hat{\mathcal{L}}_N(\theta)+  \frac{L(\ln(N))^{r}}{\sqrt{N}}(\KL(\rho ||\pi)+\ln\frac{C_{\delta}}{\delta} +  \mathcal{C}(\delta,1))
      \label{bound:oracle:eq1} \\
      \forall \rho \in \mathcal{M}_\pi: -E_{\theta \sim \rho} \mathcal{L}(\theta) \leq & -E_{\theta\sim \rho} \hat{\mathcal{L}}_N(\theta)+  \frac{L(\ln(N))^{r}}{\sqrt{N}}(\KL(\rho ||\pi)+\ln\frac{C_{\delta}}{\delta} + 
  \mathcal{C}(\delta,-1))
   \label{bound:oracle:eq2} 
  \end{align}
  In particular, by using the union bound it follows that \eqref{bound:oracle:eq1} and \eqref{bound:oracle:eq2} hold simultaneously with probability at least $1-4\delta$.  Using the fact that 
  for any posterior $\rho \in \mathcal{M}_{\pi}$, 
  \[ 
    E_{\theta\sim \rho^{\text{Gibbs}}} \hat{\mathcal{L}}_N(\theta)+  \frac{L(\ln(N))^{r}}{\sqrt{N}}\KL(\rho^{\text{Gibbs}} ||\pi) \le 
     E_{\theta \sim \rho} \hat{\mathcal{L}}_N(\theta) + \frac{L(\ln(N))^{r}}{\sqrt{N}}(\KL(\rho ||\pi)
  \]
 it  
  then follows that with probability at least $1-4\delta$
  over data %it follows that 
  \begin{align}
      \forall \rho \in \mathcal{M}_{\pi}: & E_{\theta \sim \rho^{\text{Gibbs}}} \mathcal{L}(\theta) \leq  E_{\theta\sim \rho^{\text{Gibbs}}} \hat{\mathcal{L}}_N(\theta)+  \frac{L(\ln(N))^{r}}{\sqrt{N}}(\KL(\rho^{\text{Gibbs}} ||\pi)+\ln\left(\frac{C_{\delta}}{\delta}\right) +  \mathcal{C}(\delta,1)) \nonumber \\
       & \le E_{\theta \sim \rho} \hat{\mathcal{L}}_N(\theta) + \frac{L(\ln(N))^{r}}{\sqrt{N}}\KL(\rho ||\pi)
       + \frac{L(\ln(N))^{r}}{\sqrt{N}}\left(\ln\left(\frac{C_{\delta}}{\delta}\right)+\mathcal{C}(\delta,1)\right) \nonumber \\
       & \le E_{\theta \sim \rho} \mathcal{L}(\theta) + \frac{L(\ln(N))^{r}}{\sqrt{N}}\left(\KL(\rho ||\pi)+\ln\frac{C_{\delta}}{\delta}+\mathcal{C}(\delta,-1)\right)+ \nonumber \\
       & + \frac{L(\ln(N))^{r}}{\sqrt{N}}\left(\KL(\rho ||\pi)+\ln\frac{C_{\delta}}{\delta}+\mathcal{C}(\delta,1)\right)
       \nonumber \\
       & =  E_{\theta \sim \rho} \mathcal{L}(\theta) + \frac{L(\ln(N))^{r}}{\sqrt{N}}\left(2\KL(\rho ||\pi)+2\ln\frac{C_{\delta}}{\delta}+\mathcal{C}(\delta,-1)+\mathcal{C}(\delta,1)\right)
      \label{bound:oracle:eq3}
  \end{align}
  For the posterior $\rho_{\theta_{true}}$ from Lemma \ref{pac:kl}, using
  \eqref{pac:lipschitz:loss:eq32} from Lemma \ref{pac:lipsch:loss}, it follows that
  \begin{align}
    & E_{\theta \sim \rho_{\theta_{true}}}  \mathcal{L}(\theta)  \le \mathcal{L}(\theta_{true})  \nonumber \\
    & +L_{\Theta}(\delta)L(\ln(N))^r E_{\theta \sim \rho_{\theta_{true}}}\|\theta_1 - \theta_2\|_2 \nonumber \\
          & + 2r_{\text{Lip},1}\frac{L\ln(N)^{r}}{\sqrt{N}}\left(C_{u,1} + C_{u,2}\right) +
            2r_{\text{Lip},2}  \frac{L}{512} \frac{L(\ln(N))^r}{\sqrt{N}} \nonumber \\
            & \le \mathcal{L}(\theta_{true})+ L_{\Theta}(\delta)\sigma \frac{L(\ln(N))^r}{\sqrt{N}}
          + 2r_{\text{Lip},1}\frac{L\ln(N)^{r}}{\sqrt{N}}\left(C_{u,1} + C_{u,2}\right) + 2r_{\text{Lip},2}  \frac{L}{512} \frac{L(\ln(N))^r}{\sqrt{N}}  \nonumber \\
        & \le \mathcal{L}(\theta_{true})+\frac{L\ln(N)^{r}}{\sqrt{N}} \mathcal{C}_{L,true}
      \label{bound:oracle:eq4}
  \end{align}
  By applying  \eqref{bound:oracle:eq3}  with $\rho=\rho_{\theta_{true}}$ and using that by Lemma \ref{pac:kl} 
  \[ \KL(\rho_{\theta_{true}} || \pi) \le \mathcal{C}(\theta_{true})\ln(N) \]
  and the fact that $\ln(N) \ge 1$ it follows that 
 \begin{align} 
    E_{\theta \sim \rho^{\text{Gibbs}}} \mathcal{L}(\theta) \le & \mathcal{L}(\theta_{true})+ \nonumber \\
      & \frac{L\ln(N)^{r+1}}{\sqrt{N}} \left( \mathcal{C}_{L,true} + 2\mathcal{C}(\theta_{true})
       +2\ln\frac{C_{\delta}}{\delta}+\mathcal{C}(\delta,-1)+\mathcal{C}(\delta,1) \right)
      \label{bound:oracle:eq41}
 \end{align}
 with probability at least $1-4\delta$.
 If $\theta_{\star}$ is MEP, then by Lemma \ref{lem:bound:mean1}
 $\mathcal{L}(\theta_{\star}) \le E_{\theta \sim \rho^{\text{Gibbs}}} \mathcal{L}(\theta)$. 
 Combining this with \eqref{bound:oracle:eq41}, \eqref{bound:oracle} for MEP follows.

 For single draw, from Lemma \ref{lem:single_draw}, \eqref{RenyiBound:single_draw:giibs1}, it follows that
 with probability at least $1-2\delta$ overa data and samples $\theta_{\star}$ from $\rho^{\text{Gibbs}}$
 it holds that 
 \begin{align}
  \mathcal{L}(\theta_{\star}) \le   -\frac{1}{\lambda} \ln E_{\theta \sim \pi} e^{-\lambda \hat{\mathcal{L}}_N(\theta)} + 
 \frac{L(\ln(N))^{r}}{\sqrt{N}}\left(\ln\frac{C_{\delta}}{\delta} +  \mathcal{C}(\delta,1)\right)
      \label{bound:oracle:eq5}
 \end{align}
 with $\lambda=\lambda_N=\frac{\sqrt{N}}{L(\ln{N})^r}$. 
 By using\eqref{lem:gibbs_property:eq1} from  Lemma \ref{lem:gibbs_property}, 
 it follows that for $\rho_{\theta_{true}}$ from Lemma \ref{pac:kl} 
 \begin{align}
      -\frac{1}{\lambda} \ln E_{\theta \sim \pi} e^{-\lambda \hat{\mathcal{L}}_N(\theta)} \le
       E_{\theta \sim \rho_{\theta_{true}}} \mathcal{\mathcal{L}}_N(\theta) + \frac{L(\ln(N))^{r}}{\sqrt{N}} \KL(\rho_{\theta_{true}} ||\pi) 
       %\frac{L(\ln(N))^{r+1}}{\sqrt{N}} \left( \mathcal{C}_{L,true} + \mathcal{C}(\theta_{true}) \right)
      \label{bound:oracle:e61}
\end{align}
Using that \ref{bound:oracle:eq2} holds with probability at least $1-2\delta$ over data, 
 and using \eqref{bound:oracle:eq4},  and $\ln(N) \ge 1$,
it follows that
\begin{align}
-\frac{1}{\lambda} & \ln E_{\theta \sim \pi} e^{-\lambda \hat{\mathcal{L}}_N(\theta)} \le 
       E_{\theta \sim \rho_{\theta_{true}}} \mathcal{\mathcal{L}}_N(\theta) + \frac{L(\ln(N))^{r}}{\sqrt{N}} \KL(\rho_{\theta_{true}} ||\pi)  \nonumber \\
       & \le E_{\theta \sim \rho_{\theta_{true}}} \mathcal{\mathcal{L}}(\theta) + \frac{L(\ln(N))^{r}}{\sqrt{N}} \KL(\rho_{\theta_{true}} ||\pi) \nonumber \\
       &+ \frac{L(\ln(N))^{r}}{\sqrt{N}}\left(\KL(\rho_{\theta_{true}} ||\pi)+\ln\frac{C_{\delta}}{\delta}+\mathcal{C}(\delta,-1)\right) \nonumber \\
       &  
        \le \mathcal{L}(\theta_{true})+\frac{L(\ln(N))^{r+1}}{\sqrt{N}} \left(\mathcal{C}_{L,true} + 2\mathcal{C}(\theta_{true}) + \ln\frac{C_{\delta}}{\delta} + \mathcal{C}(\delta,-1)\right)
       \label{bound:oracle:eq07}
\end{align}
with probability at least $1-2\delta$ over data and hence over data and samples from $\rho^{\text{Gibbs}}$.
By combining \eqref{bound:oracle:eq5} and \eqref{bound:oracle:eq07} using the union bounds, it follows that
with probability $1-4\delta$ over data and samples $\theta_{\star}$ from $\rho^{\text{Gibbs}}$,
\begin{align}
  \mathcal{L}(\theta_{\star}) \le & \mathcal{L}(\theta_{true})+\frac{L(\ln(N))^{r+1}}{\sqrt{N}} \left( \mathcal{C}_{L,true} + 2\mathcal{C}(\theta_{true}) + \ln\frac{C_{\delta}}{\delta} + \mathcal{C}(\delta,-1)) \right) +  \nonumber  \\
 &+  \frac{L(\ln(N))^{r}}{\sqrt{N}}\left(\ln\frac{C_{\delta}}{\delta} +  \mathcal{C}(\delta,1)\right) \nonumber \\
 & \le \mathcal{L}(\theta_{true})+\frac{L(\ln(N))^{r+1}}{\sqrt{N}} \left( \mathcal{C}_{L,true} + 2\mathcal{C}(\theta_{true}) + 2\ln\frac{C_{\delta}}{\delta} +  \mathcal{C}(\delta,1) + \mathcal{C}(\delta,-1)\right)
 \nonumber
\end{align}
from which the statement of \eqref{bound:oracle} for single draw follows, by replacing $\delta$ by $\frac{\delta}{2}$.

Finally, we consider the case of MAP. Since \eqref{bound:emp}, i.e.,
\begin{equation}
    \label{bound:oracle:eq7}
\begin{aligned}
  \mathcal{L}(\theta_{\star}) & \le  \hat{\mathcal{L}}_N(\theta_{\star}) + \\
   & \frac{2L(\ln(N))^{r+1}}{\sqrt{N}} 
  \Big( C(\theta_{\star})+\mathcal{C}_{L,true}+\mathcal{C}_{L,emp}+
     \ln\frac{C_{\delta}}{\delta^{2r_{\text{Lip},1}+1}} + \mathcal{C}(\delta,1) \Big)
\end{aligned}
\end{equation}
holds with probability at least $1-(2+r_{\text{Lip},1})\delta$.
From \eqref{lem:pac:eq1:uniform} applied to $\theta=\theta_{true}$ and $\epsilon=1$ it also follows that 
\begin{equation}
    \label{bound:oracle:eq8}
\begin{aligned}
  \hat{\mathcal{L}}_N(\theta_{true})  &\le \mathcal{L}(\theta_{true}) + \\
    & \frac{2L(\ln(N))^{r+1}}{\sqrt{N}} 
  \Big(
   C(\theta_{true})+ \mathcal{C}_{L,true}+\mathcal{C}_{L,emp}+
     \ln\frac{C_{\delta}}{\delta^{2r_{\text{Lip},1}+1}} + \mathcal{C}(\delta,-1) \Big)
\end{aligned}
\end{equation}
holds with probability at least $1-(2+\mathrm{r_{\text{Lip},1}})\delta$ over data.
If $\theta_{\star}$ is MAP, then $\theta_{\star}$
minimizes $\hat{\mathcal{L}}_N(\theta)-\frac{1}{\lambda} \ln \pi(\theta)$ and hence 
$\hat{\mathcal{L}}_N(\theta_{\star})-\frac{1}{\lambda} \ln \pi(\theta_{\star}) \le 
\hat{\mathcal{L}}_N(\theta_{true})-\frac{1}{\lambda} \ln \pi(\theta_{true})$.
Hence, 
\[
\hat{\mathcal{L}}_N(\theta_{\star}) \le 
\hat{\mathcal{L}}_N(\theta_{true})+\frac{1}{\lambda} \ln \left(\frac{\pi(\theta_{\star})}{\pi(\theta_{true})}\right)
\]
and then by using \eqref{bound:oracle:eq7}-\eqref{bound:oracle:eq8} and the fact that $\lambda=\frac{\sqrt{N}}{L(\ln(N))^r}$ and
hence $\frac{1}{\lambda} \le  \frac{L(\ln(N))^{r+1}}{\sqrt{N}}$,
it follow that with probability $1-(4+2r_{\text{Lip},1})\delta$ over data,
\begin{equation}
\label{pf:col1:eq6}
\begin{split}
 \hat{\mathcal{L}}_N(\theta_{\star})
        & \leq  
        \mathcal{L}(\theta_{true})  \\
        & +\frac{2L(\ln(N))^{r+1}}{\sqrt{N}} 
  \Big(
   2C(\theta_{true})+ 2\mathcal{C}_{L,true}+2\mathcal{C}_{L,emp}+
     2\ln\frac{C_{\delta}}{\delta} \\
     &+ \mathcal{C}(\delta,1)+\mathcal{C}(\delta,-1) \Big) 
        + \frac{2L(\ln(N))^{r+1}}{\sqrt{N}} \ln \left(\frac{\pi(\theta_{\star})}{\pi(\theta_{true})}\right)
\end{split}
\end{equation}
holds, hence, by replacing $\delta$ with $\delta/2$, it follows that  \eqref{bound:oracle} holds with probability
at least $1-(2+r_{\text{Lip},1})\delta$.

\subsection{Proof of \eqref{cor:param:est1:eq1}}
   If \textbf{(SConv)} holds, then by \cite[Theorem 2.1.11, Theorem 2.1.8]{NesterovBook}
   \[ m_{\Theta} \|\theta_{\star}-\theta_{true}\|^2 \le \mathcal{L}(\theta_{\star}) - \mathcal{L}(\theta_{true}) \]
   from which using \eqref{bound:oracle} the inequality \eqref{cor:param:est1:eq1} follows.

\subsection{Proof of \eqref{col1:eq1}}
From
\cite[Chapter 4, page 87]{LjungBook} 
it follows that 
$\sigma^2_{\e^s}=\mathcal{L}(\theta_{\mathrm{true}})=\inf_{\theta \in \Theta} \mathcal{L}(\theta)$, 
where $\sigma^2_{\e^s}=\bE[\|\e^s(t)\|^2_2]$, 
%%from \cite[page 683 eq. 17.31]{LindquistBook} for the former case, and it follows that
if $\ell$  is the quadratic loss function, 
i.e. the true system yields the best possible prediction in mean square sense.
%it then the following holds:
%for the latter case, it follows
\begin{equation}
\label{param_set:true_eq}
 \forall \theta \in \Theta:  \sigma^2_{\e^s}=\mathcal{L}(\theta_{true}) \le \mathcal{L}(\theta) 
\end{equation}
where $\sigma^2_{\e^s}$ is the variance of $\e^s(t)$.
That is, the parameter $\theta_{true}$ corresponding to the data generator results in the smallest possible generalisation loss for quadratic loss function.
Let $H_{\theta}$ be the transfer function of
the predictor parametrized by $\theta$

It can be shown (see \cite[Appendix, (82)]{eringisRenyi}) and \cite[Chapter 4, page 87]{LjungBook}) that
\begin{equation}
\label{rem:lqg:eq0}
\|H_{\theta_{\mathrm{true}}}-H_{\theta} \|_{H_2}^2 \le \frac{1}{m_{\w}} (\mathcal{L}(\theta) - \mathcal{L}(\theta_{true})), %~  \sigma^2_{\e^s}=\mathcal{L}(\theta_{true}) \le \mathcal{L}(\theta)
\end{equation}
Then \eqref{col1:eq1} follows from \eqref{bound:oracle} and \eqref{rem:lqg:eq0}.

%\subsubsection*{Discussion on the bound}
%The intuition behind the role of the noise level,
%system norms and stability for the bound above is the same as for Theorem \ref{thm:unbounded}.
%The bound above has the advantage that it converges to zero 
%zero at a rate $O(\frac{1}{\sqrt{N}})$, which is
%%\begin{corollary}[\textcolor{cyan}{convergence $O(1/\sqrt{N})$}]
%%\label{thm:bounded:alt:col}
%%If $\lambda_N=\sqrt{N}$,  and define the
%%constants 
%%\begin{color}{cyan}
%%\[
%%\begin{split}
%%%& \bar{G}_{f,i}=\sup_{\theta \in \Theta} \bar{G}_{f,i}(\theta), ~ i=1,2,  \\
%%% & G_{e}=\sup_{\theta \in \Theta} G_e(\theta), \quad G_{e,1}=\sup_{\theta \in \Theta} G_{e,1}(\theta), ~ \\
%% & \textcolor{red}{\mathcal{C}=\left(L G_{b,1} c_\rve+ 2L^2 G_{b,2} c_{\rve}^2\right)} \\
%%  %& \textcolor{red}{\left. 
%%  & G_{b,1}=G_{f,1}G_g, \quad  G_{b,2}=(G_e+2G_{e,1})\sqrt{n_{\rve}}
%%%& \mathcal{C} =4K\bar{G}_{f,2}+ \ln(\max{\bar{G}_{f,i}K,1})+8 c_\rve^4 n_\rve^2( G_e  + G_{e,1}^2G_e^2) \\%\end{split}
%%\]
%%where  $L=\sup_{\theta \in \Theta} L(\theta)$
%%%then 
%%%with  and 
%%%$\mathcal{C}_2=4K\bar{G}_{f,2}$, $\mathcal{C}=\mathcal{C}_1+\mathcal{C}_2$ 
%%%it follows that
%%For any prior $\pi$
%%\end{color}
%%\begin{equation}
%%\label{thm:bounded:alt:col:eq1}
%%     \textcolor{cyan}{\Phi(\lambda,\pi,N) \le \bar{\Phi}_{\pi}(\lambda,N) \le \mathcal{C}}
%%\end{equation}
%%In particular, for any prior $\pi$
%%and for any $\delta \in (0,0.5)$,
%%\begin{equation}
%%\label{catoni1}
%%\begin{split}
%%  & \forall \rho \in\mathcal{M}_\pi: E_{\theta \sim \rho } \mathcal{L}(\theta) \leq E_{\theta\sim \theta } \hat{\mathcal{L}}_N(\theta)+ \\
%%  & \frac{1}{\sqrt{N}}(\KL(\rho ||\pi)+\ln\frac{1}{\delta} + C)
%%\end{split}
%%\end{equation}
%%holds with probability at least $1-2\delta$.
%%\begin{color}{red}
%%    Moreover, for quadratic loss the constant $\mathcal{C}$ is of the form
%%    \[
%%      \begin{aligned}
%%        & \mathcal{C}=\left(2G_e G_{f,2}G_g(c_\rve\sqrt{n_\rve})^2+ \right. \\
%%        & \left. 2(2G_e)^2(G_e+2G_{e,1})(c_{\rve}\sqrt{n_{\rve}})^4 \right). 
%%      \end{aligned} 
%%    \]
%%\end{color}
%%\end{corollary} 
%%\textcolor{cyan}{
%%The bound \eqref{catoni1} is similar to classical PAC-Bayesian bounds for
%%i.i.d. data, and the constant $C$ describes the learning complexity of the hypothesis class and the data generator. 
%%As before, $\mathcal{C}$ increases with the amplitude of the noise and the maximum $\ell_1$ norm of the error system and the class of predictors. 
%%This corresponds to the usual intuition: it is harder to predict noisy data, and the generalization gap is higher for less stable predictors.
%%}
%%%for any given posterior $\hat{\rho}$,
%$\lim_{N \rightarrow \infty} \bar{\mathbb{B}}_N=0$
%with the rate of convergence $O(\frac{1}{\sqrt{N}})$.



\subsection{Remark on relationship with system identification and learning}
\label{app:indent}
%%\begin{remark}[Convergence, comparison with PAC bounds]
%%\label{bounded:convergence}
%%The bounds \eqref{catoni1},\eqref{pac:mean:bd}, \eqref{pac:mean:single} contain the term
%%$\KL(\rho^{\train} \mid \pi)/\sqrt{N}$,
%%$\ln\left(\frac{\rho^{\train}(\theta_{\star})}{\pi(\theta_{\star})}\right)/\sqrt{N}$, numerator of which
%%may depend on $N$ too: this is the case when
%%$\rho^{\train}$ minimizes \eqref{eq:initial:def:tildern}, e.g., it is the Gibbs posterior. Hence, in general, these terms need not converge to zero as $N \rightarrow \infty$, as the numerator may increase faster than $\sqrt{N}$. 
%%This is a general issue with PAC-Bayesian bounds. 
%%However, if constant posteriors are used, i.e., the PAC
%%bound\eqref{classic:pac} is used, then the right-hand side 
%%$\mathcal{L}(\theta)  \le \hat{\mathcal{L}}_N(\theta) + \frac{(\mathcal{C}+C(\theta)+2L+\frac{1}{\delta}}{\sqrt{N}}$
%%does converge to zero as $N \rightarrow \infty$. 
%%The reason for using both \eqref{catoni1},\eqref{pac:mean:bd}, \eqref{pac:mean:single} and 
%%\eqref{classic:pac}, is that while the latter is guaranteed to have the desired asymptotic behavior, for a fixed $N$
%%the bound \eqref{catoni1},\eqref{pac:mean:bd}, \eqref{pac:mean:single} may give tighter bounds for the generalization error than \eqref{classic:pac}.
%%To this end, priors and data dependent posterior have to be
%%carefully tuned. The latter is a separate research direction in PAC-Bayesian theory,e.g., \cite{Dziugaite2017},
%%with promising results. 
%%\end{remark}
%%\begin{remark}[Application to parameter estimation error]
%%\label{bounded:param_est}
%%The bound \eqref{col1:eq2} is
%%$O(\frac{1}{\sqrt{N}})$ 
%%for MEP and single draw, if  Gaussian priors are used.
%%The bound is 
%%$O(\frac{\ln(N)}{\sqrt{N}})$ for all cases if $\pi$ is the uniform distribution.
%%
%%For the case of
%%MAP and Gaussian prior $\pi \sim \mathcal{N}(\theta_m,P)$,
%%$C(\theta_{true},\theta_{\star})$,
%%the right-hand side of \eqref{col1:eq2},
%%the term $C(\theta_{true},\theta_{\star})$
%%contains the summand
%%$(\theta_{\star}-\theta_m)^TP^{-1}(\theta_{\star}-\theta_m)$, which measures the discrepancy
%%between MAP estimate and the mean of the prior. In principle, as $\theta_{\star}$
%%changes with $N$, we cannot guarantee
%%that %$\frac{C(\theta_{true},\theta_{\star}})
%%%{\sqrt{N}}$ 
%%right-hand side of \eqref{col1:eq2}
%%converges to zero as $N \rightarrow \infty$.
%%%If $\pi$ is uniform distribution, then 
%%%$C(\theta_{true},%\theta_{\star})=O(\ln(N))$
%%%and hence the right-hand side of %\eqref{col1:eq2} is $O(\frac{\ln(N)}{\sqrt{N}})$.
%%%$\bE[e^{\epsilon \lamnda (\m}]
%%
%%%\ref{lemma:bounded:mgf(L-V)}
%%
%%\end{remark}


\begin{remark}[Identification of FIR models]
    \label{rem:fir}
    Consider the SISO case ($n_y=n_u=1$), fix 
    an integer $T_0$ and let
    $\Theta=\mathbb{R}^{T_0}$ and 
    for any $\theta=(\theta_1,\ldots, \theta_{T_0})^T \in \Theta$, let
    $D_{\theta}=\theta_1$,
    \[ 
    A=\begin{bmatrix} 0 & 0 & \cdots & 0 & 0 \\ 
        1 & 0 & \cdots & 0 &0 \\
        \vdots & \vdots & \cdots & \vdots & \vdots \\
        0 & 0 & \cdots & 1 & 0
    \end{bmatrix}, ~ 
    B=\begin{bmatrix} 1 \\ 0 \\ \vdots \\ 0 \end{bmatrix}, ~
    C_\theta=\begin{bmatrix} \theta_2 \\ \vdots \\ \theta_{T_0} \\ 1 \end{bmatrix}^T
    \]
    Then $\hat{\rvy}_{\theta}(t)=u(t-T_0)+\sum_{i=0}^{T_0-1} \theta_{i+1} \rvu(t-i)$. Assume that there exists 
    $\theta_{true} \in \Theta$, such that 
    $\rvy(t)=\hat{\rvy}_{\theta_{true}}(t)+\rve(t)$,
    where $\rve(t)$ is i.i.d. with zero mean and variance $\sigma^2$, and independent of $\rvu(s)$, $s \in \mathbb{Z}$.
    Let $\hat{\rvu}(t)=\rvu(t)$ for $t > 0$ and
    let $\hat{\rvu}(t)=0$ for $t < 0$.
    Let $\Phi$ be the regressor matrix
    constructed from $\tilde{\rvu}$ as in
    \cite[(3.45), page 48]{pillonetto2022regularized} and
    let 
    $Y=\begin{bmatrix} \rvy(0) & \cdots &\rvy(N-1) \end{bmatrix}^T$.
    %$E=\begin{bmatrix} \rve(0) & \cdots &\rve(N-1) \end{bmatrix}^T$
    Then
    %$Y=\Phi \theta_{true} + E$ and 
    $\hat{\mathcal{L}}_N(\theta)=\frac{1}{N} \|Y-\Phi \theta\|_2^2$.
    Assume that the prior $\pi$ is the Gaussian density $\mathcal{N}(\tilde{\theta},P)$, $P=P^T > 0$. 
    Then a lengthy but straightforward calculation
    reveal that
    the Gibbs posterior is the  Gaussian density
    $\mathcal{N}(\tilde{\theta}+P\Phi^T(\Phi P \Phi^T + \frac{N}{2\lambda} I)^{-1}(Y-\Phi \tilde{\theta}), P-P\Phi^T(\Phi P\Phi^T + \frac{N}{2\lambda} I)^{-1}\Phi P)$, 
    which is the conditional density of
    $\theta$ w.r.t. $Y,\Phi$ presented in
    \cite[(4.6), page 101]{pillonetto2022regularized}, if
    $Y=\Phi\theta+E$ is assumed with 
    $E \sim \mathcal{N}(0,\frac{N}{2\lambda} I_{T_0})$.
    Then the MAP and MEP estimates are both
    $\theta_{\star}=\tilde{\theta}+P\Phi^T(\Phi P \Phi^T + \frac{N}{2\lambda} I)^{-1}(Y-\Phi \theta_{true})$
    and they solve  the regularized
    optimization problem
    $\argmin_{\theta \in \Theta} \lambda \hat{\mathcal{L}}_N(\theta) +  (\theta-\tilde{\theta})P^{-1}(\theta-\tilde{\theta})$. Moreover, they are also
    the maximum likelihood
    of the posterior $P(\theta \mid Y,\Phi)$. 

 Then Corollary \ref{lem:bound:mean} applies to $\theta_{\star}$,
  %inequality \eqref{catoni1}, and the average 
  %$E_{\theta \sim \rho_{\text{Gibbs}}}\hat{\mathcal{L}}_N(\theta)$
  %can readily be computed 
  %analytically, using the fact that 
  %$\hat{\mathcal{L}}_N(\theta)=\frac{\lambda}{N} \|Y-\Phi\theta\|_2^2$ is a quadratic function
  %of $\theta$ and $\rho_{\text{Gibbs}}$ is the
  %Gaussian distribution with mean and variance described above.    
  yielding  a 
  $O(\frac{(\ln(N))^r}{\sqrt{N}})$ bound on
  the generalization gap of $\theta_{\star}$
  and a $O(\frac{(\ln(N))^{r+1}}{\sqrt{N}})$ bound
  on the square parameter estimation error
  $\|\theta_{\star}-\theta_{true}\|_2^2$.
  \end{remark}
  %\begin{remark}[State-space models]
  % Consider a parametrization of LTIs of the form\eqref{eq:predictor}
  % such that $\hat{\rvy}_{\theta}(t \mid 0)=\Phi_t \theta$, where 
  % $\Phi$ is a suitable regressor matrix built from $\{\rvy(s),\rvu(s)\}_{s < t}$. For instance, 
  % 
 % 
 % 
 % \end{remark}
 \begin{remark}(\emph{Maximum likelihood estimates   of state-space models}). 
  \label{rem:ml}
  Consider the parametrization of LTI systems
  \begin{align}
        \hat{\x}(t+1)&=(\hat{A}_\theta + \hat{L}_{\theta} \hat{C}_{\theta}) \hat{\x}(t)
         +(\hat{B}_\theta+L_{\theta}D_{\theta})\rvu(t)+\hat{L}_\theta\e_{\theta}(t),  \nonumber \\
        \y_\theta(t)&=\hat{C}_\theta\hat{\x}(t)+\hat{D}_\theta\rvu(t) + \e_{\theta}(t)
  \label{eq:predictor_gen}
 \end{align}
 where %the matrices $A_{\theta},B_{\theta},L_{\theta},C_{\theta},D_{\theta}$ are as in \eqref{eq:predictor}, 
 $A_{\theta}$ and $A_{\theta}+L_{\theta}C_{\theta}$ are Schur, 
 $\e_{\theta}(t)$ is i.i.d., zero mean Gaussian with
 variance $\frac{1}{\lambda}$, $(\rvu,\e_{\theta})$ is jointly Gaussian, $\e_{\theta}(t)$ is independent of
 $\rvu(s)$ for $s < t$. 
 We can associate with \eqref{eq:predictor_gen} the predictor 
\begin{align}
        \hat{\x}(t+1)&=\hat{A}_\theta \hat{\x}(t)+\hat{B}_\theta \rvu(t)+\hat{L}_\theta\y(t),  \nonumber \\
        \hat{\y}_\theta(t \mid 0)&=\hat{C}_\theta\hat{\x}(t)+\hat{D}_\theta\rvu(t),
  \label{eq:predictor_gen1}
 \end{align}
 Assume that
 $\y$ is generated by an LTI system of the form \eqref{eq:predictor_gen} which satisfies  \cite[(1), INP B, page 427]{CainesBook}.

 In this case, by \cite[(2.7), page 422, Chapter 7]{CainesBook}
 the conditional likelihood function
 $p_{\theta}(\{y(t)\}_{t=-1}^{N-1} \mid \{u(t)\}_{t=0}^{N-1})$ 
 $\{\y_{\theta}(t)\}_{t=0}^{N-1}$ w.r.t
 $\{\rvu(t)\}_{t=0}^{N-1}$, evaluated at 
 $y(t)=\rvy(t)$ and $u(t)=\rvu(t)$, $t=-1,\ldots,N-1$
 equals 
 $Ze^{-\frac{N\lambda}{2} \hat{\mathcal{L}}_N(\theta)}$
 for a suitable constant $Z$, where the empirical error 
 $\hat{\mathcal{L}}_N(\theta)$ is defined for the quadratic loss. 
 That is, 
 i.e., up to an additive constant, its logarithm equals $\ln \rho^{\text{Gibbs}}(\theta)$.
 In particular, the maximum likelihood estimate 
 $\mathrm{argmax}_{\theta \in \Theta} p_{\theta}(\{y(t)\}_{t=0}^{N-1} \mid \{u(t)\}_{t=0}^{N-1})$, see
 \cite[Chapter 7,page 428]{CainesBook} 
 equals the MAP 
 $\mathrm{argmax}_{\theta \in \Theta} \rho^{\text{Gibbs}}(\theta)$ of the Gibbs posterior corresponding to the uniform distribution $\pi=\frac{1}{vol(\Theta)}$ as a prior. Here we assume that  $vol(\Theta)$  is finite.

 We can then use the bounds \eqref{bound:emp}, \eqref{bound:oracle}, \eqref{cor:param:est1:eq1}, \eqref{col1:eq1}, 
 for the generalization gap, the gap between the true error of the learned model and the best true error and the parameter estimation error respectively.
 %of the 
 %maximum likelihood estimate. 
 %If in \eqref{eq:predictor_gen} we assume that 
 %$\hat{x}$  stationary, jointly Gaussian with $\e_{\theta}$
 %and it satisfies \cite[INP B, page 427]{CainesBook},
 %then $p_{\theta}(\{y(t)\}_{t=0}^{N-1} \mid \{u(t)\}_{t=0}^{N-1})$ 
 \end{remark} 
 \begin{remark}[Bayesian methods for estimating LTI systems]
  \label{rem:bayesian}
    The Bayesian method for system identification of LTI systems \cite{NINNESS201040} corresponds to taking a single draw, MEP or MAP from the conditional distribution of the parameter given the past observations. If the parametrization is in innovation form, i.e., of the form \eqref{eq:predictor_gen} and all the signals are jointly Gaussian,  then by
    Remark \ref{rem:ml}
    the joint density of  $\{\y(t)\}_{t=-1}^{N-1}$ given 
    $\{\rvu(t)\}_{t=0}^{N-1}$ and 
    $\{\y(s),\rvu(s)\}_{s < 0}$ and a parameter 
    $\theta$ is $Ze^{-N\lambda \hat{\mathcal{L}}_N(\theta)}$ for a suitable constant, 
    and hence by \cite[Section 3, eq. (12)]{NINNESS201040}
    the conditional density of $\theta$ with respect to
    $\{\y(t),\rvu(t)\}_{t=0}^{N-1}$ is
    the Gibbs posterior 
    $\frac{e^{-N\lambda \hat{\mathcal{L}}_N(\theta)}p_{\theta}(\y(-1))\pi(\theta)}{E_{\theta \in \pi} e^{-\lambda \hat{\mathcal{L}}_N(\theta)}}$ for the quadratic loss and a suitable prior, 
    and hence 
    Gibbs posterior for the predictors \eqref{eq:predictor_gen1}, and hence taking a random draw, MEP or MAP from 
    this conditional density, as suggested by \cite{NINNESS201040}, corresponds to taking a random draw, MEP or MAP from
    the Gibbs posterior. In particular, the bounds  \eqref{bound:emp}, \eqref{bound:oracle}, \eqref{cor:param:est1:eq1}, \eqref{col1:eq1} apply.
 \end{remark}
  \begin{remark}[PEM methods]
    \label{rem:pem}
    To begin with if a PEM algorithm can be reduced to optimizing a cost function
    \( \hat{\mathcal{L}}_N(\theta)  + g(\theta),\)
    then the algorithm is equivalent to taking the MAP estimate of the Gibbs posterior with a prior $\pi$ such that 
    $-\frac{1}{\lambda}\ln \pi(\theta)=\frac{g(\theta)}{\lambda}$. In particular, \eqref{bound:emp}-\eqref{col1:eq1} from Corollary \ref{lem:bound:mean} applies in this case. For example, the quadratic regularization term $g(\theta)=
    (\theta - \theta_m)^T P^{-1} (\theta - \theta_m) $
    corresponds to $\pi \sim \mathcal{N}(\theta_m, P/\lambda)$. If $\pi$ can be interpreted as a Gaussian density, then MAP is MEP and the corresponding sharper
    bounds can be used. 

    Furthermore, many  PEM algorithms rely on choosing randomly the initial parameter estimate. The outcomes of such algorithms can be viewed as drawing a sample from a suitable posterior distribution. More precisely, consoider an algorithm
    $A_N$. For every initial value $\theta_0 \in \Theta$ and data $\mathcal{D}$, the algorithm
    generates an estimate $A_N(\theta_0,\mathcal{D})$.
    Most optimization algorithms used for PEM methods
    (Gradient, Gauss-Newton), etc. are of this form, if we fix the
    number of steps.
    %Assume that $A_N$ is smooth and smoothly invertable in 
    %$\theta_0$ and let us denote its inverse by $A_N^{-1}$, i.e. $A_N(A_N^{-1}(\theta,\mathcal{D}),\mathcal{D})=\theta$. 
    %Smooth invertability is often satisfied. For example,
    %if $A_N$ is the simple $1$-step gradient descend for minimizing the empirical error, 
    %i.e. $A_N(\theta,\mathcal{D})=\theta-\lambda \frac{d \mathcal{L}_N(\theta)}{d\theta}$ and 
    %the Hessian $\frac{d^2 \mathcal{L}_N(\theta)}{d\theta^2}$ is positive definite and has a non-zero minimal eigenvalue $\mu$,
    %Mthen the Jacobian  of $A_N$ w.r.t. $\theta$
    %is $I-\lambda \frac{d \mathcal{L}_N(\theta)}{d\theta} < I(1-\lambda\mu) < -I$ for large enough $\lambda$, i.e., the Jacobian in invertable for a suitable learning rate. Hence, $A_N$ is locally smoothly invertable.
    %Invertability of the Hessian of the empirical error holds if the empirical error is a convex function of the parameter.
    %The same analysis carries over to gradient descend executed for $K$-steps.
    Assume that the initial parameter values $\theta_{init}$
    are i.i.d., samples from a prior density $\pi$. Then the resulting parameter values
    $A_N(\theta_{init},\mathcal{D})$ are i.i.d. samples from
    a distribution $P_{A_N}(B)=\int_{A_N^{-1}(B,\mathcal{D})} \pi(\theta)dm(\theta)$ for any $B \in B_{\theta}$.
    Assume that the latter distribution has a density
    $\hat{\rho}$ which is in $\mathcal{M}_{\pi}$. 
    %$\hat{\rho}(\theta)=|\det \frac{dA_N^{-1}(\theta,\mathcal{D})}{d\theta}| \pi(A_N^{-1}(\theta,\mathcal{D})$
    %is a data dependent density of the outcome of applying 
    %$A_N$ to initial values $\theta_0$ sampled randomly
    %from $\pi$. 
    Then Lemma \ref{lem:single_draw} provides 
    an error bound for the true error of such an algorithm, when applied to an initial parameter value
    randomly sampled from $\pi$.
  \end{remark}
\begin{remark}[Strong convexity of the true error]
  \label{app:strong convexity}
 Strong convexity implies that there is a unique minimizer of the true error, which is
 the true system. It  can be seen as an \emph{identifiability and persistence of excitation assumption}, as it implies that there is no parameter value which gives the same
 input-output behavior as the true system (otherwise their true errores would be the same),
 and $\w$ is rich enough (the average response to $\w$ allows us to find the true system).
 In fact, strong convexity of the true error (at least locally), is a standard
 assumption  for studying the asympotic distribution of the parameter estimation error
 \cite[Theorem 9.1]{LjungBook}. 
\end{remark}
\begin{remark}[Bounds on the difference between  system matrices]
  \label{matrix:trans}
Assuming that the LTI systems $\Sigma(\theta)$ are all minimal,
%and the LTI system \eqref{eq:T0sys2} is minimal,
we can use \eqref{col1:eq1}
%(\ref{eq:RenyiBound})
and \cite[Theorem V.2]{oymak2021revisiting}
%\textcolor{black}{and,  \eqref{lqg:rem:eq1}--\eqref{lqg:rem:eq2} could be used 
to derive
an error bound (in high-probability) 
on the difference between 
the matrices of the learned and true system.
%depending on $\theta_{\star}$ 
%and  that of $\theta_{true}$, using 
%Theorem \ref{col1} or \eqref{rem:lqg:eq-1}.
Alternatively, if we consider parametrizations based on structure indices \cite[Chapter 3]{RalfPeeters} 
we could use the fact that $H_2$ norms can be used to approximate locally the Riemannian distance between parameter values  \cite[Chapter 4, page 129]{RalfPeeters}. 
\end{remark}
%\begin{remark}[Non-zero initial state \& existing finite-sample bounds]
%\label{rem:fin_samp:zero_init}
%
%
%
%Assuming that a deterministic initial state for the data generator is problematic for the following reason. Even if it is assumed that inputs before the start of the identification experiment was zero, as the presence of process noise would steer the system 
%to an initial state which depends on the past noises, i.e., which is a random variable. 
%In practice, one could pretend that learning the data generated by the system in
% the stationary regime due to past noise and input is equivalent to data points obtained
% by starting the system at zero initial state and discarding the first $T$ measurements for
% large $T$. However, if we are to integrate naively this approach to the learning algorithms
% \cite{oymak2021revisiting,lale2020logarithmic,Simchowitz_Foster_2020}, then this would
% lead to firing up time $T$. In turn, as the cited bounds are $O(T)$ and hold for $N >> T$, 
% so choosing $T$ large would make those bounds vacuous. 
% equivalent to discarding 
%tarting the identification experiment $T$ time steps earlie one could discard the first $T$ data points for learning,
%as it ignores
%the effect of the noise before the identification expreimen
%\end{remark}
\begin{remark}[Relationship with \cite{eringis2023pacbayesRNN,eringis2023pacbayesian}]
  \label{rem:deividas}
 For the case of bounded noise in the data generator, Theorem \ref{thm:main} provides a PAC-Bayesian bound on the
 generalization gap which is similar to \cite[Theorem 2, Corollary 1]{eringis2023pacbayesRNN} and \cite[Theorem 5.2]{eringis2023pacbayesian}. However, \cite{eringis2023pacbayesRNN} considers RNNs, and the corresponding 
 constants $L_{\rvv}(\theta),L_{g,\rvs}(\theta),L_{g,\rvv}$ correspond to $\|\hat{B}_{\theta}\|_2$, $\|\hat{C}_{\theta}\|_2$, and $\|\hat{D}_{\theta}\|_2$, respectively, and $\tau(\theta)$ corresponds to an upper bound on the spectrum of $\hat{A}_{\theta}$
 and $C(\theta)$ is such that $\|\hat{A}_{\theta}^k\|_{2} \le \tau(\theta)^k C(\theta)$.
   In our setting, $L_{\rvv}(\theta),L_{g,\rvs}(\theta),L_{g,\rvv}$ can be taken to be $C$ from Assumption \ref{as:parameterisation}, $C(\theta)$ can be taken to be $M$ and $\tau(\theta)$ can be taken to be $\gamma$ from Assumption \ref{as:parameterisation}.
    Then the constants $G_{\theta}$, $H_{\theta}$ from \cite{eringis2023pacbayesRNN}  
    are upper bounds on $\|\alpha_{\theta}\|_{\ell_1}$ and $\theta_{\infty}(\alpha_{\theta})$ respectively.
    Hence the term $\hat{\Psi}_{\pi}(\lambda,N)$ of 
    \cite[Theorem 2, Corollary 1]{eringis2023pacbayesRNN}  is an upper bound on the term
    $\frac{1}{2}\sum_{i=1}^2 \ln E_{\theta\sim\pi} \exp\left(\mathcal{C}_i(2\lambda,\theta,N,\delta,\epsilon)\right)$ from \eqref{eq:initial:tildern1}, i.e., the bound of \cite[Theorem 2, Corollary 1]{eringis2023pacbayesRNN} is
    more conservative than the one of this paper.

  As to \cite[Theorem 5.2]{eringis2023pacbayesian}, it uses $\|\alpha_{g,y}-c_{\theta}\|_{\ell_1} \le G_e(\theta)$ and $\theta_{\infty}(\alpha_{g,y}-c_{\theta}) \le G_{e,1}(\theta)$ in the bounds, where $c_{\theta}$ and $\alpha_{g,y}$ are defined in Lemma \ref{alpha:lemma2}. While this may appear less conservative, the use of $c_{\theta}$ makes the bound difficult to evaluate in practice, as $c_{\theta}$ depends on the parameters of the unknown data generator. In contrast, our bound uses only $\|\alpha_g\|_{\ell_1}$ and $\theta_{\infty}(\alpha_g)$, which, although they also depend on the data generator, are easier to estimate or upper bound in practice.

  Note that \cite{eringis2023pacbayesRNN,eringis2023pacbayesian} do not address the case of unbounded signals or parameter estimation errors. They also do not provide generalization bounds for single draw, MEP, or MAP estimators.
\end{remark}
\begin{remark}[Relationship with finite-sample bounds \cite{oymak2021revisiting,lale2020logarithmic,Simchowitz_Foster_2020,Ziemann1,Ziemann2}]
  \label{rem:finite-sample}
Recent works such as \cite{oymak2021revisiting,lale2020logarithmic,Simchowitz_Foster_2020,Ziemann1,Ziemann2} provide finite-sample parameter estimation bounds for specific algorithms (e.g., least-squares for ARX models, subspace methods), typically assuming the data generator starts from a deterministic initial state and often neglecting the stochastic component. However, assuming a deterministic initial state can be problematic, since process noise generally drives the system to a random initial state determined by past disturbances.

In contrast, our results: \textbf{(1)} apply to a broad class of learning algorithms (including PEM, maximum likelihood, Bayesian, and regularized FIR methods), \textbf{(2)} assume the data generator is in a stationary regime, \textbf{(3)} explicitly account for both deterministic and stochastic components, and \textbf{(4)} provide generalization gap bounds. The trade-off is that our bound on the square of the parameter estimation error is $O\left(\frac{(\ln N)^{r+1}}{\sqrt{N}}\right)$, which is more conservative than the $O\left(\frac{\ln N}{N}\right)$ rates in the cited works.
\end{remark}
%However, in contrast to the cited results, they apply to a wide variety of algorithms. 
%That is, the conservativeness of the PAC-Bayesian bound is the price to pay for its generality. 
%Moreover, PAC-Bayesian bounds allow us not only to bound the parameter estimation error, but
%also to bound the generalization loss.
% with a term of order $O(\frac{1}{\sqrt{N}})$.
%In turn, $O(\frac{1}{N})$ bounds on the $H_2$ norm of the difference between the true and estimated system do not translate automatically to 
%$O(\frac{1}{N})$ bounds on the generalization error in the presence of 
%noise. 
%That is, the presented bounds can be used to bound quantities not addressed by \citet{oymak2021revisiting,lale2020logarithmic,Simchowitz_Foster_2020}, and the presented bounds apply to a wider range of learning algorithms.
%The price to pay for this generality is that the bounds are more conservative when applied to the specific 
%setting of
%\citet{oymak2021revisiting,lale2020logarithmic,Simchowitz_Foster_2020}.
%From a practical point of view, the difference between %$O(\frac{1}{N})$ and 
%$O(\frac{1}{\sqrt{N}})$ need not always be critical, 
%as for relatively low $N$, higher constants can mask the advantages of the faster rates. A detailed numerical comparison remains topic of future research.


% \section{Numerical Experiments}
% \label{sec:numerical_examples}
% The aim of this section is to numerically evaluate the bounds in
% Theorem~\ref{thm:main} and Corollary~\ref{thm:bounded:alt:col}
% on a synthetic example.

% We use the stationary innovation model of Assumption~\ref{as:generator}
% with \(n_y=n_u=1\) and \(n_e=2\). The generator matrices are
% \begin{equation}
%     A_g=
%     \begin{bmatrix}
%         0.16 & -0.30\\
%         0 & -0.05
%     \end{bmatrix},
%     \qquad
%     K_g=
%     \begin{bmatrix}
%         0.33 & -0.75\\
%         0 & -0.09
%     \end{bmatrix},
%     \qquad
%     C_g=
%     \begin{bmatrix}
%         1 & 1\\
%         0 & 1
%     \end{bmatrix}.
% \end{equation}
% The innovation covariance is specified through
% \begin{equation}
%     L_Q=
%     \begin{bmatrix}
%         0.018 & 0.008\\
%         0.006 & 0.027
%     \end{bmatrix},
%     \qquad
%     Q_e=L_QL_Q^\top.
% \end{equation}

% We consider two innovation distributions. In the bounded case,
% \begin{equation}
%     \rve_g(t)=L_Q\tilde{\rve}_g(t),
%     \qquad
%     \tilde{\rve}_{g,i}(t)\sim\mathrm{Unif}[-\sqrt3,\sqrt3],
%     \quad i\in\{1,2\},
% \end{equation}
% independently over \(i\) and \(t\). In the unbounded case, we use the mixture
% \begin{equation}
%     \rve_g(t)=
%     \begin{cases}
%         L_Q\tilde{\rve}_g(t),
%         & \tilde{\rve}_g(t)\sim\mathcal N(0,I_2),
%         \quad \text{with probability } p,\\
%         L_Q\tilde{\rve}_g(t),
%         & \tilde{\rve}_{g,i}(t)\sim\mathrm{Unif}[-\sqrt3,\sqrt3],
%         \quad i\in\{1,2\},
%         \quad \text{with probability } 1-p,
%     \end{cases}
% \end{equation}
% with \(p=0.35\). Both cases satisfy
% \(\mathbb E[\rve_g(t)\rve_g(t)^\top]=Q_e\). For the unbounded case,
% \[
%     \|\rve_g\|_{\psi_2}\le 0.05.
% \]
% For the bounded case, the almost-sure innovation bound used in the constants is
% \begin{equation}
%     c_e
%     =
%     \sqrt3\,\|L_Q\|_{\infty\to\infty}
%     =
%     \sqrt3\max_i\sum_j |(L_Q)_{ij}|.
% \end{equation}

% The predictor is in the class of predictors from Assumption \ref{as:parameterisation} and uses only the input process, i.e., \(\rvw(t)=\rvu(t)\) with
% \(\hat D_\theta=0\). Specifically, we consider second-order SISO predictors parametrized by
% \begin{equation}
%     \theta=(r_1,r_2,\phi,s,\tilde b_1,\tilde b_2,\tilde c_1,\tilde c_2).
% \end{equation}
% Let
% \begin{align}
%     R_\theta
%     &=
%     Q_{\phi}
%     \begin{bmatrix}
%         r_1 & 0\\
%         0 & r_2
%     \end{bmatrix}
%     Q_{\phi}^\top,
%     &
%     Q_{\phi}
%     &=
%     \begin{bmatrix}
%         \cos\phi & -\sin\phi\\
%         \sin\phi & \cos\phi
%     \end{bmatrix},
%     \\
%     S_\theta
%     &=
%     \begin{bmatrix}
%         0 & s\\
%         -s & 0
%     \end{bmatrix},
%     &
%     A_{c,\theta}
%     &=
%     -\frac12 R_\theta(I+S_\theta).
% \end{align}
% The predictor matrices are chosen as
% \begin{equation}
%     \hat A_\theta
%     =
%     \left(I+A_{c,\theta}\right)
%     \left(I-A_{c,\theta}\right)^{-1},
%     \qquad
%     \hat B_\theta
%     =
%     R_\theta^{1/2}\tilde b,
%     \qquad
%     \hat C_\theta
%     =
%     \tilde c^\top R_\theta^{-1/2},
% \end{equation}
% where
% \[
%     \tilde b=(\tilde b_1,\tilde b_2)^\top,
%     \qquad
%     \tilde c=(\tilde c_1,\tilde c_2)^\top.
% \]
% The prior density \(\pi\) is uniform on the compact parameter set
% \(\Theta\) given by
% \begin{equation}
%     r_1,r_2\in[2,20],
%     \qquad
%     \phi\in[-\pi,\pi],
%     \qquad
%     s\in[-1,1],
% \end{equation}
% and
% \begin{equation}
%     \tilde b_1\in[-0.97,0.97],
%     \quad
%     \tilde b_2\in[-0.05,0.05],
%     \qquad
%     \tilde c_1\in[-0.095,0.095],
%     \quad
%     \tilde c_2\in[-0.02,0.02].
% \end{equation}
% This parametrization guarantees that \(\hat A_\theta\) is Schur for all
% \(\theta\in\Theta\).

% We consider the quadratic loss and the smooth absolute loss
% \begin{equation}
%     \ell_\tau(y,\hat y)
%     =
%     \sqrt{(y-\hat y)^2+\tau^2}-\tau,
%     \qquad
%     \tau=10^{-4}.
% \end{equation}
% Posterior expectations under the Gibbs posterior are approximated using
% fixed-step Hamiltonian Monte Carlo with \(200\) retained samples,
% \(200\) discarded burn-in samples, leapfrog step size \(0.015\), and \(4\)
% leapfrog steps per proposal.

% The true risk is approximated by a long stationary reference trajectory
% generated from the same innovation model after burn-in. We use
% \[
%     N\in\{3, 50, 10^2, 3\times10^2, 10^3, 3\times 10^3,
%     10^4, 10^5, 10^6, 10^7\},
% \]
% and report results over \(5\) independent data realizations.

% The quantities entering the bounds are computed as follows. Define
% \begin{equation}
%     \gamma
%     =
%     \sup_{\theta\in\Theta}
%     \left\|
%         R_\theta^{-1/2}\hat A_\theta R_\theta^{1/2}
%     \right\|_2 .
% \end{equation}
% This supremum is computed numerically over the prior support. Since
% \(\gamma<1\), the predictor impulse-response envelope is bounded by
% \begin{equation}
%     \|\bar\alpha\|_{\ell_1}
%     \le
%     \frac{M_BM_C}{1-\gamma},
%     \qquad
%     \theta_\infty(\bar\alpha)
%     \le
%     \frac{M_BM_C}{(1-\gamma)^2},
% \end{equation}
% where
% \[
%     M_B=\sup_{\theta\in\Theta}\|\tilde b\|_2,
%     \qquad
%     M_C=\sup_{\theta\in\Theta}\|\tilde c\|_2.
% \]
% The generator Markov parameters are computed from
% \begin{equation}
%     \alpha_g(0)=I_2,
%     \qquad
%     \alpha_g(k)=C_gA_g^{k-1}K_g,\quad k\ge1,
% \end{equation}
% with a geometric tail bound. The non-simplified
% Theorem~\ref{thm:main} bound is evaluated by Monte Carlo integration over
% the prior using \(4000\) deterministic Halton points, whereas the
% Corollary~\ref{thm:bounded:alt:col} bound uses the prior-support envelopes
% above.

% \begin{figure}[h!]
%     \centering
%     \input{semilogcor32}
%     \caption{PAC-Bayesian generalization gaps and bounds}
%     \label{fig:wide-minibatch-semilog}
% \end{figure}

% \begin{figure}[h!]
%     \centering
%     \input{loglogplotcor32}
%     \caption{PAC-Bayesian generalization gaps and bounds}
%     \label{fig:wide-minibatch-loglog}
% \end{figure}
  \newpage

%%  \section*{NeurIPS Paper Checklist}
%%
%%%%% BEGIN INSTRUCTIONS %%%
%%%%The checklist is designed to encourage best practices for responsible machine learning research, addressing issues of reproducibility, transparency, research ethics, and societal impact. Do not remove the checklist: {\bf The papers not including the checklist will be desk rejected.} The checklist should follow the references and follow the (optional) supplemental material.  The checklist does NOT count towards the page
%%%%limit. 
%%%%
%%%%Please read the checklist guidelines carefully for information on how to answer these questions. For each question in the checklist:
%%%%\begin{itemize}
%%%%    \item You should answer \answerYes{}, \answerNo{}, or \answerNA{}.
%%%%    \item \answerNA{} means either that the question is Not Applicable for that particular paper or the relevant information is Not Available.
%%%%    \item Please provide a short (1--2 sentence) justification right after your answer (even for \answerNA). 
%%%%   % \item {\bf The papers not including the checklist will be desk rejected.}
%%%%\end{itemize}
%%%%
%%%%{\bf The checklist answers are an integral part of your paper submission.} They are visible to the reviewers, area chairs, senior area chairs, and ethics reviewers. You will also be asked to include it (after eventual revisions) with the final version of your paper, and its final version will be published with the paper.
%%%%
%%%%The reviewers of your paper will be asked to use the checklist as one of the factors in their evaluation. While \answerYes{} is generally preferable to \answerNo{}, it is perfectly acceptable to answer \answerNo{} provided a proper justification is given (e.g., error bars are not reported because it would be too computationally expensive'' or ``we were unable to find the license for the dataset we used''). In general, answering \answerNo{} or \answerNA{} is not grounds for rejection. While the questions are phrased in a binary way, we acknowledge that the true answer is often more nuanced, so please just use your best judgment and write a justification to elaborate. All supporting evidence can appear either in the main paper or the supplemental material, provided in appendix. If you answer \answerYes{} to a question, in the justification please point to the section(s) where related material for the question can be found.
%%%%
%%%%IMPORTANT, please:
%%%%\begin{itemize}
%%%%    \item {\bf Delete this instruction block, but keep the section heading ``NeurIPS Paper Checklist"},
%%%%    \item  {\bf Keep the checklist subsection headings, questions/answers and guidelines below.}
%%%%    \item {\bf Do not modify the questions and only use the provided macros for your answers}.
%%%%\end{itemize} 
%% 
%%
%%%%% END INSTRUCTIONS %%%
%%
%%
%%\begin{enumerate}
%%
%%\item {\bf Claims}
%%    \item[] Question: Do the main claims made in the abstract and introduction accurately reflect the paper's contributions and scope?
%%    \item[] Answer: \answerYes{} % Replace by \answerYes{}, \answerNo{}, or \answerNA{}.
%%    \item[] Justification: the  main contribution of the paper is  a generalization bound for a class of dynamical systems, this contribution is clearly stated in the abstract and introduction. 
%%    \item[] Guidelines:
%%    \begin{itemize}
%%        \item The answer \answerNA{} means that the abstract and introduction do not include the claims made in the paper.
%%        \item The abstract and/or introduction should clearly state the claims made, including the contributions made in the paper and important assumptions and limitations. A \answerNo{} or \answerNA{} answer to this question will not be perceived well by the reviewers. 
%%        \item The claims made should match theoretical and experimental results, and reflect how much the results can be expected to generalize to other settings. 
%%        \item It is fine to include aspirational goals as motivation as long as it is clear that these goals are not attained by the paper. 
%%    \end{itemize}
%%
%%\item {\bf Limitations}
%%    \item[] Question: Does the paper discuss the limitations of the work performed by the authors?
%%    \item[] Answer: \answerYes{} % Replace by \answerYes{}, \answerNo{}, or \answerNA{}.
%%    \item[] Justification: the paper discusses the limitations of the proposed approach, including assumptions made and potential impacts on generalization, both in the main text and in the remarks at the end of the paper.
%%    \item[] Guidelines:
%%    \begin{itemize}
%%        \item The answer \answerNA{} means that the paper has no limitation while the answer \answerNo{} means that the paper has limitations, but those are not discussed in the paper. 
%%        \item The authors are encouraged to create a separate ``Limitations'' section in their paper.
%%        \item The paper should point out any strong assumptions and how robust the results are to violations of these assumptions (e.g., independence assumptions, noiseless settings, model well-specification, asymptotic approximations only holding locally). The authors should reflect on how these assumptions might be violated in practice and what the implications would be.
%%        \item The authors should reflect on the scope of the claims made, e.g., if the approach was only tested on a few datasets or with a few runs. In general, empirical results often depend on implicit assumptions, which should be articulated.
%%        \item The authors should reflect on the factors that influence the performance of the approach. For example, a facial recognition algorithm may perform poorly when image resolution is low or images are taken in low lighting. Or a speech-to-text system might not be used reliably to provide closed captions for online lectures because it fails to handle technical jargon.
%%        \item The authors should discuss the computational efficiency of the proposed algorithms and how they scale with dataset size.
%%        \item If applicable, the authors should discuss possible limitations of their approach to address problems of privacy and fairness.
%%        \item While the authors might fear that complete honesty about limitations might be used by reviewers as grounds for rejection, a worse outcome might be that reviewers discover limitations that aren't acknowledged in the paper. The authors should use their best judgment and recognize that individual actions in favor of transparency play an important role in developing norms that preserve the integrity of the community. Reviewers will be specifically instructed to not penalize honesty concerning limitations.
%%    \end{itemize}
%%
%%\item {\bf Theory assumptions and proofs}
%%    \item[] Question: For each theoretical result, does the paper provide the full set of assumptions and a complete (and correct) proof?
%%    \item[] Answer: \answerYes{} % Replace by \answerYes{}, \answerNo{}, or \answerNA{}.
%%    \item[] Justification: all the assumptions are formally stated, all the detailed proofs are in the supplementary material. 
%%    \item[] Guidelines:
%%    \begin{itemize}
%%        \item The answer \answerNA{} means that the paper does not include theoretical results. 
%%        \item All the theorems, formulas, and proofs in the paper should be numbered and cross-referenced.
%%        \item All assumptions should be clearly stated or referenced in the statement of any theorems.
%%        \item The proofs can either appear in the main paper or the supplemental material, but if they appear in the supplemental material, the authors are encouraged to provide a short proof sketch to provide intuition. 
%%        \item Inversely, any informal proof provided in the core of the paper should be complemented by formal proofs provided in appendix or supplemental material.
%%        \item Theorems and Lemmas that the proof relies upon should be properly referenced. 
%%    \end{itemize}
%%
%%    \item {\bf Experimental result reproducibility}
%%    \item[] Question: Does the paper fully disclose all the information needed to reproduce the main experimental results of the paper to the extent that it affects the main claims and/or conclusions of the paper (regardless of whether the code and data are provided or not)?
%%    \item[] Answer: \answerNA{} % Replace by \answerYes{}, \answerNo{}, or \answerNA{}.
%%    \item[] Justification: the paper contains no experiments.
%%    %the description of the experiments contains all the necessary details and the code is provided with the submission.
%%    \item[] Guidelines:
%%    \begin{itemize}
%%        \item The answer \answerNA{} means that the paper does not include experiments.
%%        \item If the paper includes experiments, a \answerNo{} answer to this question will not be perceived well by the reviewers: Making the paper reproducible is important, regardless of whether the code and data are provided or not.
%%        \item If the contribution is a dataset and\slash or model, the authors should describe the steps taken to make their results reproducible or verifiable. 
%%        \item Depending on the contribution, reproducibility can be accomplished in various ways. For example, if the contribution is a novel architecture, describing the architecture fully might suffice, or if the contribution is a specific model and empirical evaluation, it may be necessary to either make it possible for others to replicate the model with the same dataset, or provide access to the model. In general. releasing code and data is often one good way to accomplish this, but reproducibility can also be provided via detailed instructions for how to replicate the results, access to a hosted model (e.g., in the case of a large language model), releasing of a model checkpoint, or other means that are appropriate to the research performed.
%%        \item While NeurIPS does not require releasing code, the conference does require all submissions to provide some reasonable avenue for reproducibility, which may depend on the nature of the contribution. For example
%%        \begin{enumerate}
%%            \item If the contribution is primarily a new algorithm, the paper should make it clear how to reproduce that algorithm.
%%            \item If the contribution is primarily a new model architecture, the paper should describe the architecture clearly and fully.
%%            \item If the contribution is a new model (e.g., a large language model), then there should either be a way to access this model for reproducing the results or a way to reproduce the model (e.g., with an open-source dataset or instructions for how to construct the dataset).
%%            \item We recognize that reproducibility may be tricky in some cases, in which case authors are welcome to describe the particular way they provide for reproducibility. In the case of closed-source models, it may be that access to the model is limited in some way (e.g., to registered users), but it should be possible for other researchers to have some path to reproducing or verifying the results.
%%        \end{enumerate}
%%    \end{itemize}
%%
%%
%%\item {\bf Open access to data and code}
%%    \item[] Question: Does the paper provide open access to the data and code, with sufficient instructions to faithfully reproduce the main experimental results, as described in supplemental material?
%%    \item[] Answer: \answerNA{} % Replace by \answerYes{}, \answerNo{}, or \answerNA{}.
%%    \item[] Justification: the paper contains no experiments.
%%    %the paper provides open access to the data and code, with sufficient instructions to faithfully reproduce the main experimental results, as described in supplemental material.
%%    \item[] Guidelines:
%%    \begin{itemize}
%%        \item The answer \answerNA{} means that paper does not include experiments requiring code.
%%        \item Please see the NeurIPS code and data submission guidelines (\url{https://neurips.cc/public/guides/CodeSubmissionPolicy}) for more details.
%%        \item While we encourage the release of code and data, we understand that this might not be possible, so \answerNo{} is an acceptable answer. Papers cannot be rejected simply for not including code, unless this is central to the contribution (e.g., for a new open-source benchmark).
%%        \item The instructions should contain the exact command and environment needed to run to reproduce the results. See the NeurIPS code and data submission guidelines (\url{https://neurips.cc/public/guides/CodeSubmissionPolicy}) for more details.
%%        \item The authors should provide instructions on data access and preparation, including how to access the raw data, preprocessed data, intermediate data, and generated data, etc.
%%        \item The authors should provide scripts to reproduce all experimental results for the new proposed method and baselines. If only a subset of experiments are reproducible, they should state which ones are omitted from the script and why.
%%        \item At submission time, to preserve anonymity, the authors should release anonymized versions (if applicable).
%%        \item Providing as much information as possible in supplemental material (appended to the paper) is recommended, but including URLs to data and code is permitted.
%%    \end{itemize}
%%
%%
%%\item {\bf Experimental setting/details}
%%    \item[] Question: Does the paper specify all the training and test details (e.g., data splits, hyperparameters, how they were chosen, type of optimizer) necessary to understand the results?
%%    \item[] Answer: \answerNA{} % Replace by \answerYes{}, \answerNo{}, or \answerNA{}.
%%    \item[] Justification: the paper contains no experiments.
%%    %the paper provides all the necessary details for the experimental setup, including data splits, hyperparameters, and optimizer settings.
%%    \item[] Guidelines:
%%    \begin{itemize}
%%        \item The answer \answerNA{} means that the paper does not include experiments.
%%        \item The experimental setting should be presented in the core of the paper to a level of detail that is necessary to appreciate the results and make sense of them.
%%        \item The full details can be provided either with the code, in appendix, or as supplemental material.
%%    \end{itemize}
%%
%%\item {\bf Experiment statistical significance}
%%    \item[] Question: Does the paper report error bars suitably and correctly defined or other appropriate information about the statistical significance of the experiments?
%%    \item[] Answer: \answerNA{} % Replace by \answerYes{}, \answerNo{}, or \answerNA{}.
%%    \item[] Justification: the paper does not contain experiments. %the main contribution of the paper is a mathematical result, proof of which is included in the supplementary material. The experiments serve as an illustration, not as the justification of the 
%%    %main claim. Moreover, the nature of the experiments does not easily lend itself to 
%%    %computation of statistical significance. 
%%    %the paper provides error bars and statistical significance information for the experiments, including details on how they were calculated and the assumptions ma
%%    \item[] Guidelines:
%%    \begin{itemize}
%%        \item The answer \answerNA{} means that the paper does not include experiments.
%%        \item The authors should answer \answerYes{} if the results are accompanied by error bars, confidence intervals, or statistical significance tests, at least for the experiments that support the main claims of the paper.
%%        \item The factors of variability that the error bars are capturing should be clearly stated (for example, train/test split, initialization, random drawing of some parameter, or overall run with given experimental conditions).
%%        \item The method for calculating the error bars should be explained (closed form formula, call to a library function, bootstrap, etc.)
%%        \item The assumptions made should be given (e.g., Normally distributed errors).
%%        \item It should be clear whether the error bar is the standard deviation or the standard error of the mean.
%%        \item It is OK to report 1-sigma error bars, but one should state it. The authors should preferably report a 2-sigma error bar than state that they have a 96\% CI, if the hypothesis of Normality of errors is not verified.
%%        \item For asymmetric distributions, the authors should be careful not to show in tables or figures symmetric error bars that would yield results that are out of range (e.g., negative error rates).
%%        \item If error bars are reported in tables or plots, the authors should explain in the text how they were calculated and reference the corresponding figures or tables in the text.
%%    \end{itemize}
%%
%%\item {\bf Experiments compute resources}
%%    \item[] Question: For each experiment, does the paper provide sufficient information on the computer resources (type of compute workers, memory, time of execution) needed to reproduce the experiments?
%%    \item[] Answer: \answerNA{} % Replace by \answerYes{}, \answerNo{}, or \answerNA{}.
%%    \item[] Justification: no experiments in the paper.
%%    %the experiment was performed on a desktop computer, no special hardware or large-scale resources were required.
%%    \item[] Guidelines:
%%    \begin{itemize}
%%        \item The answer \answerNA{} means that the paper does not include experiments.
%%        \item The paper should indicate the type of compute workers CPU or GPU, internal cluster, or cloud provider, including relevant memory and storage.
%%        \item The paper should provide the amount of compute required for each of the individual experimental runs as well as estimate the total compute. 
%%        \item The paper should disclose whether the full research project required more compute than the experiments reported in the paper (e.g., preliminary or failed experiments that didn't make it into the paper). 
%%    \end{itemize}
%%    
%%\item {\bf Code of ethics}
%%    \item[] Question: Does the research conducted in the paper conform, in every respect, with the NeurIPS Code of Ethics \url{https://neurips.cc/public/EthicsGuidelines}?
%%    \item[] Answer: \answerYes{} % Replace by \answerYes{}, \answerNo{}, or \answerNA{}.
%%    \item[] Justification: the research conducted in the paper adheres to the NeurIPS Code of Ethics.
%%    \item[] Guidelines:
%%    \begin{itemize}
%%        \item The answer \answerNA{} means that the authors have not reviewed the NeurIPS Code of Ethics.
%%        \item If the authors answer \answerNo, they should explain the special circumstances that require a deviation from the Code of Ethics.
%%        \item The authors should make sure to preserve anonymity (e.g., if there is a special consideration due to laws or regulations in their jurisdiction).
%%    \end{itemize}
%%
%%
%%\item {\bf Broader impacts}
%%    \item[] Question: Does the paper discuss both potential positive societal impacts and negative societal impacts of the work performed?
%%    \item[] Answer: \answerNo{} % Replace by \answerYes{}, \answerNo{}, or \answerNA{}.
%%    \item[] Justification: the contribution of the paper is a theoretical result, no societal impacts are directly associated with the work.
%%    \item[] Guidelines:
%%    \begin{itemize}
%%        \item The answer \answerNA{} means that there is no societal impact of the work performed.
%%        \item If the authors answer \answerNA{} or \answerNo, they should explain why their work has no societal impact or why the paper does not address societal impact.
%%        \item Examples of negative societal impacts include potential malicious or unintended uses (e.g., disinformation, generating fake profiles, surveillance), fairness considerations (e.g., deployment of technologies that could make decisions that unfairly impact specific groups), privacy considerations, and security considerations.
%%        \item The conference expects that many papers will be foundational research and not tied to particular applications, let alone deployments. However, if there is a direct path to any negative applications, the authors should point it out. For example, it is legitimate to point out that an improvement in the quality of generative models could be used to generate Deepfakes for disinformation. On the other hand, it is not needed to point out that a generic algorithm for optimizing neural networks could enable people to train models that generate Deepfakes faster.
%%        \item The authors should consider possible harms that could arise when the technology is being used as intended and functioning correctly, harms that could arise when the technology is being used as intended but gives incorrect results, and harms following from (intentional or unintentional) misuse of the technology.
%%        \item If there are negative societal impacts, the authors could also discuss possible mitigation strategies (e.g., gated release of models, providing defenses in addition to attacks, mechanisms for monitoring misuse, mechanisms to monitor how a system learns from feedback over time, improving the efficiency and accessibility of ML).
%%    \end{itemize}
%%    
%%\item {\bf Safeguards}
%%    \item[] Question: Does the paper describe safeguards that have been put in place for responsible release of data or models that have a high risk for misuse (e.g., pre-trained language models, image generators, or scraped datasets)?
%%    \item[] Answer: \answerNA{} % Replace by \answerYes{}, \answerNo{}, or \answerNA{}.
%%    \item[] Justification: the paper poses no such risks.
%%    \item[] Guidelines:
%%    \begin{itemize}
%%        \item The answer \answerNA{} means that the paper poses no such risks.
%%        \item Released models that have a high risk for misuse or dual-use should be released with necessary safeguards to allow for controlled use of the model, for example by requiring that users adhere to usage guidelines or restrictions to access the model or implementing safety filters. 
%%        \item Datasets that have been scraped from the Internet could pose safety risks. The authors should describe how they avoided releasing unsafe images.
%%        \item We recognize that providing effective safeguards is challenging, and many papers do not require this, but we encourage authors to take this into account and make a best faith effort.
%%    \end{itemize}
%%
%%\item {\bf Licenses for existing assets}
%%    \item[] Question: Are the creators or original owners of assets (e.g., code, data, models), used in the paper, properly credited and are the license and terms of use explicitly mentioned and properly respected?
%%    \item[] Answer: \answerNA{} % Replace by \answerYes{}, \answerNo{}, or \answerNA{}.
%%    \item[] Justification: no assets are used. 
%%    %all assets used in the paper are properly credited, and their licenses and terms of use are explicitly mentioned and respected.
%%    \item[] Guidelines:
%%    \begin{itemize}
%%        \item The answer \answerNA{} means that the paper does not use existing assets.
%%        \item The authors should cite the original paper that produced the code package or dataset.
%%        \item The authors should state which version of the asset is used and, if possible, include a URL.
%%        \item The name of the license (e.g., CC-BY 4.0) should be included for each asset.
%%        \item For scraped data from a particular source (e.g., website), the copyright and terms of service of that source should be provided.
%%        \item If assets are released, the license, copyright information, and terms of use in the package should be provided. For popular datasets, \url{paperswithcode.com/datasets} has curated licenses for some datasets. Their licensing guide can help determine the license of a dataset.
%%        \item For existing datasets that are re-packaged, both the original license and the license of the derived asset (if it has changed) should be provided.
%%        \item If this information is not available online, the authors are encouraged to reach out to the asset's creators.
%%    \end{itemize}
%%
%%\item {\bf New assets}
%%    \item[] Question: Are new assets introduced in the paper well documented and is the documentation provided alongside the assets?
%%    \item[] Answer: \answerNA{} % Replace by \answerYes{}, \answerNo{}, or \answerNA{}.
%%    \item[] Justification: the paper does not release new assets.
%%    \item[] Guidelines:
%%    \begin{itemize}
%%        \item The answer \answerNA{} means that the paper does not release new assets.
%%        \item Researchers should communicate the details of the dataset\slash code\slash model as part of their submissions via structured templates. This includes details about training, license, limitations, etc. 
%%        \item The paper should discuss whether and how consent was obtained from people whose asset is used.
%%        \item At submission time, remember to anonymize your assets (if applicable). You can either create an anonymized URL or include an anonymized zip file.
%%    \end{itemize}
%%
%%\item {\bf Crowdsourcing and research with human subjects}
%%    \item[] Question: For crowdsourcing experiments and research with human subjects, does the paper include the full text of instructions given to participants and screenshots, if applicable, as well as details about compensation (if any)? 
%%    \item[] Answer: \answerNA{} % Replace by \answerYes{}, \answerNo{}, or \answerNA{}.
%%    \item[] Justification: the paper does not involve crowdsourcing nor research with human subjects.
%%    \item[] Guidelines:
%%    \begin{itemize}
%%        \item The answer \answerNA{} means that the paper does not involve crowdsourcing nor research with human subjects.
%%        \item Including this information in the supplemental material is fine, but if the main contribution of the paper involves human subjects, then as much detail as possible should be included in the main paper. 
%%        \item According to the NeurIPS Code of Ethics, workers involved in data collection, curation, or other labor should be paid at least the minimum wage in the country of the data collector. 
%%    \end{itemize}
%%
%%\item {\bf Institutional review board (IRB) approvals or equivalent for research with human subjects}
%%    \item[] Question: Does the paper describe potential risks incurred by study participants, whether such risks were disclosed to the subjects, and whether Institutional Review Board (IRB) approvals (or an equivalent approval/review based on the requirements of your country or institution) were obtained?
%%    \item[] Answer: \answerNA{} % Replace by \answerYes{}, \answerNo{}, or \answerNA{}.
%%    \item[] Justification: the paper does not involve crowdsourcing nor research with human subjects.
%%    \item[] Guidelines:
%%    \begin{itemize}
%%        \item The answer \answerNA{} means that the paper does not involve crowdsourcing nor research with human subjects.
%%        \item Depending on the country in which research is conducted, IRB approval (or equivalent) may be required for any human subjects research. If you obtained IRB approval, you should clearly state this in the paper. 
%%        \item We recognize that the procedures for this may vary significantly between institutions and locations, and we expect authors to adhere to the NeurIPS Code of Ethics and the guidelines for their institution. 
%%        \item For initial submissions, do not include any information that would break anonymity (if applicable), such as the institution conducting the review.
%%    \end{itemize}
%%
%%\item {\bf Declaration of LLM usage}
%%    \item[] Question: Does the paper describe the usage of LLMs if it is an important, original, or non-standard component of the core methods in this research? Note that if the LLM is used only for writing, editing, or formatting purposes and does \emph{not} impact the core methodology, scientific rigor, or originality of the research, declaration is not required.
%%    %this research? 
%%    \item[] Answer: \answerNA{} % Replace by \answerYes{}, \answerNo{}, or \answerNA{}.
%%    \item[] Justification: the paper does not involve LLMs as any important, original, or non-standard components.
%%    \item[] Guidelines:
%%    \begin{itemize}
%%        \item The answer \answerNA{} means that the core method development in this research does not involve LLMs as any important, original, or non-standard components.
%%        \item Please refer to our LLM policy in the NeurIPS handbook for what should or should not be described.
%%    \end{itemize}
%%
%%\end{enumerate}




%%  \newpage
%%
%%  \section*{NeurIPS Paper Checklist}
%%
%%%%% BEGIN INSTRUCTIONS %%%
%%%%The checklist is designed to encourage best practices for responsible machine learning research, addressing issues of reproducibility, transparency, research ethics, and societal impact. Do not remove the checklist: {\bf The papers not including the checklist will be desk rejected.} The checklist should follow the references and follow the (optional) supplemental material.  The checklist does NOT count towards the page
%%%%limit. 
%%%%
%%%%Please read the checklist guidelines carefully for information on how to answer these questions. For each question in the checklist:
%%%%\begin{itemize}
%%%%    \item You should answer \answerYes{}, \answerNo{}, or \answerNA{}.
%%%%    \item \answerNA{} means either that the question is Not Applicable for that particular paper or the relevant information is Not Available.
%%%%    \item Please provide a short (1--2 sentence) justification right after your answer (even for \answerNA). 
%%%%   % \item {\bf The papers not including the checklist will be desk rejected.}
%%%%\end{itemize}
%%%%
%%%%{\bf The checklist answers are an integral part of your paper submission.} They are visible to the reviewers, area chairs, senior area chairs, and ethics reviewers. You will also be asked to include it (after eventual revisions) with the final version of your paper, and its final version will be published with the paper.
%%%%
%%%%The reviewers of your paper will be asked to use the checklist as one of the factors in their evaluation. While \answerYes{} is generally preferable to \answerNo{}, it is perfectly acceptable to answer \answerNo{} provided a proper justification is given (e.g., error bars are not reported because it would be too computationally expensive'' or ``we were unable to find the license for the dataset we used''). In general, answering \answerNo{} or \answerNA{} is not grounds for rejection. While the questions are phrased in a binary way, we acknowledge that the true answer is often more nuanced, so please just use your best judgment and write a justification to elaborate. All supporting evidence can appear either in the main paper or the supplemental material, provided in appendix. If you answer \answerYes{} to a question, in the justification please point to the section(s) where related material for the question can be found.
%%%%
%%%%IMPORTANT, please:
%%%%\begin{itemize}
%%%%    \item {\bf Delete this instruction block, but keep the section heading ``NeurIPS Paper Checklist"},
%%%%    \item  {\bf Keep the checklist subsection headings, questions/answers and guidelines below.}
%%%%    \item {\bf Do not modify the questions and only use the provided macros for your answers}.
%%%%\end{itemize} 
%% 
%%
%%%%% END INSTRUCTIONS %%%
%%
%%
%%\begin{enumerate}
%%
%%\item {\bf Claims}
%%    \item[] Question: Do the main claims made in the abstract and introduction accurately reflect the paper's contributions and scope?
%%    \item[] Answer: \answerYes{} % Replace by \answerYes{}, \answerNo{}, or \answerNA{}.
%%    \item[] Justification: the  main contribution of the paper is  a generalization bound for a class of dynamical systems, this contribution is clearly stated in the abstract and introduction. 
%%    \item[] Guidelines:
%%    \begin{itemize}
%%        \item The answer \answerNA{} means that the abstract and introduction do not include the claims made in the paper.
%%        \item The abstract and/or introduction should clearly state the claims made, including the contributions made in the paper and important assumptions and limitations. A \answerNo{} or \answerNA{} answer to this question will not be perceived well by the reviewers. 
%%        \item The claims made should match theoretical and experimental results, and reflect how much the results can be expected to generalize to other settings. 
%%        \item It is fine to include aspirational goals as motivation as long as it is clear that these goals are not attained by the paper. 
%%    \end{itemize}
%%
%%\item {\bf Limitations}
%%    \item[] Question: Does the paper discuss the limitations of the work performed by the authors?
%%    \item[] Answer: \answerYes{} % Replace by \answerYes{}, \answerNo{}, or \answerNA{}.
%%    \item[] Justification: the paper discusses the limitations of the proposed approach, including assumptions made and potential impacts on generalization, both in the main text and in the remarks at the end of the paper.
%%    \item[] Guidelines:
%%    \begin{itemize}
%%        \item The answer \answerNA{} means that the paper has no limitation while the answer \answerNo{} means that the paper has limitations, but those are not discussed in the paper. 
%%        \item The authors are encouraged to create a separate ``Limitations'' section in their paper.
%%        \item The paper should point out any strong assumptions and how robust the results are to violations of these assumptions (e.g., independence assumptions, noiseless settings, model well-specification, asymptotic approximations only holding locally). The authors should reflect on how these assumptions might be violated in practice and what the implications would be.
%%        \item The authors should reflect on the scope of the claims made, e.g., if the approach was only tested on a few datasets or with a few runs. In general, empirical results often depend on implicit assumptions, which should be articulated.
%%        \item The authors should reflect on the factors that influence the performance of the approach. For example, a facial recognition algorithm may perform poorly when image resolution is low or images are taken in low lighting. Or a speech-to-text system might not be used reliably to provide closed captions for online lectures because it fails to handle technical jargon.
%%        \item The authors should discuss the computational efficiency of the proposed algorithms and how they scale with dataset size.
%%        \item If applicable, the authors should discuss possible limitations of their approach to address problems of privacy and fairness.
%%        \item While the authors might fear that complete honesty about limitations might be used by reviewers as grounds for rejection, a worse outcome might be that reviewers discover limitations that aren't acknowledged in the paper. The authors should use their best judgment and recognize that individual actions in favor of transparency play an important role in developing norms that preserve the integrity of the community. Reviewers will be specifically instructed to not penalize honesty concerning limitations.
%%    \end{itemize}
%%
%%\item {\bf Theory assumptions and proofs}
%%    \item[] Question: For each theoretical result, does the paper provide the full set of assumptions and a complete (and correct) proof?
%%    \item[] Answer: \answerYes{} % Replace by \answerYes{}, \answerNo{}, or \answerNA{}.
%%    \item[] Justification: all the assumptions are formally stated, all the detailed proofs are in the supplementary material. 
%%    \item[] Guidelines:
%%    \begin{itemize}
%%        \item The answer \answerNA{} means that the paper does not include theoretical results. 
%%        \item All the theorems, formulas, and proofs in the paper should be numbered and cross-referenced.
%%        \item All assumptions should be clearly stated or referenced in the statement of any theorems.
%%        \item The proofs can either appear in the main paper or the supplemental material, but if they appear in the supplemental material, the authors are encouraged to provide a short proof sketch to provide intuition. 
%%        \item Inversely, any informal proof provided in the core of the paper should be complemented by formal proofs provided in appendix or supplemental material.
%%        \item Theorems and Lemmas that the proof relies upon should be properly referenced. 
%%    \end{itemize}
%%
%%    \item {\bf Experimental result reproducibility}
%%    \item[] Question: Does the paper fully disclose all the information needed to reproduce the main experimental results of the paper to the extent that it affects the main claims and/or conclusions of the paper (regardless of whether the code and data are provided or not)?
%%    \item[] Answer: \answerYes{} % Replace by \answerYes{}, \answerNo{}, or \answerNA{}.
%%    \item[] Justification: the description of the experiments contains all the necessary details and the code is provided with the submission.
%%    \item[] Guidelines:
%%    \begin{itemize}
%%        \item The answer \answerNA{} means that the paper does not include experiments.
%%        \item If the paper includes experiments, a \answerNo{} answer to this question will not be perceived well by the reviewers: Making the paper reproducible is important, regardless of whether the code and data are provided or not.
%%        \item If the contribution is a dataset and\slash or model, the authors should describe the steps taken to make their results reproducible or verifiable. 
%%        \item Depending on the contribution, reproducibility can be accomplished in various ways. For example, if the contribution is a novel architecture, describing the architecture fully might suffice, or if the contribution is a specific model and empirical evaluation, it may be necessary to either make it possible for others to replicate the model with the same dataset, or provide access to the model. In general. releasing code and data is often one good way to accomplish this, but reproducibility can also be provided via detailed instructions for how to replicate the results, access to a hosted model (e.g., in the case of a large language model), releasing of a model checkpoint, or other means that are appropriate to the research performed.
%%        \item While NeurIPS does not require releasing code, the conference does require all submissions to provide some reasonable avenue for reproducibility, which may depend on the nature of the contribution. For example
%%        \begin{enumerate}
%%            \item If the contribution is primarily a new algorithm, the paper should make it clear how to reproduce that algorithm.
%%            \item If the contribution is primarily a new model architecture, the paper should describe the architecture clearly and fully.
%%            \item If the contribution is a new model (e.g., a large language model), then there should either be a way to access this model for reproducing the results or a way to reproduce the model (e.g., with an open-source dataset or instructions for how to construct the dataset).
%%            \item We recognize that reproducibility may be tricky in some cases, in which case authors are welcome to describe the particular way they provide for reproducibility. In the case of closed-source models, it may be that access to the model is limited in some way (e.g., to registered users), but it should be possible for other researchers to have some path to reproducing or verifying the results.
%%        \end{enumerate}
%%    \end{itemize}
%%
%%
%%\item {\bf Open access to data and code}
%%    \item[] Question: Does the paper provide open access to the data and code, with sufficient instructions to faithfully reproduce the main experimental results, as described in supplemental material?
%%    \item[] Answer: \answerYes{} % Replace by \answerYes{}, \answerNo{}, or \answerNA{}.
%%    \item[] Justification: the paper provides open access to the data and code, with sufficient instructions to faithfully reproduce the main experimental results, as described in supplemental material.
%%    \item[] Guidelines:
%%    \begin{itemize}
%%        \item The answer \answerNA{} means that paper does not include experiments requiring code.
%%        \item Please see the NeurIPS code and data submission guidelines (\url{https://neurips.cc/public/guides/CodeSubmissionPolicy}) for more details.
%%        \item While we encourage the release of code and data, we understand that this might not be possible, so \answerNo{} is an acceptable answer. Papers cannot be rejected simply for not including code, unless this is central to the contribution (e.g., for a new open-source benchmark).
%%        \item The instructions should contain the exact command and environment needed to run to reproduce the results. See the NeurIPS code and data submission guidelines (\url{https://neurips.cc/public/guides/CodeSubmissionPolicy}) for more details.
%%        \item The authors should provide instructions on data access and preparation, including how to access the raw data, preprocessed data, intermediate data, and generated data, etc.
%%        \item The authors should provide scripts to reproduce all experimental results for the new proposed method and baselines. If only a subset of experiments are reproducible, they should state which ones are omitted from the script and why.
%%        \item At submission time, to preserve anonymity, the authors should release anonymized versions (if applicable).
%%        \item Providing as much information as possible in supplemental material (appended to the paper) is recommended, but including URLs to data and code is permitted.
%%    \end{itemize}
%%
%%
%%\item {\bf Experimental setting/details}
%%    \item[] Question: Does the paper specify all the training and test details (e.g., data splits, hyperparameters, how they were chosen, type of optimizer) necessary to understand the results?
%%    \item[] Answer: \answerYes{} % Replace by \answerYes{}, \answerNo{}, or \answerNA{}.
%%    \item[] Justification: the paper provides all the necessary details for the experimental setup, including data splits, hyperparameters, and optimizer settings.
%%    \item[] Guidelines:
%%    \begin{itemize}
%%        \item The answer \answerNA{} means that the paper does not include experiments.
%%        \item The experimental setting should be presented in the core of the paper to a level of detail that is necessary to appreciate the results and make sense of them.
%%        \item The full details can be provided either with the code, in appendix, or as supplemental material.
%%    \end{itemize}
%%
%%\item {\bf Experiment statistical significance}
%%    \item[] Question: Does the paper report error bars suitably and correctly defined or other appropriate information about the statistical significance of the experiments?
%%    \item[] Answer: \answerNA{} % Replace by \answerYes{}, \answerNo{}, or \answerNA{}.
%%    \item[] Justification: the main contribution of the paper is a mathematical result, proof of which is included in the supplementary material. The experiments serve as an illustration, not as the justification of the 
%%    main claim. Moreover, the nature of the experiments does not easily lend itself to 
%%    computation of statistical significance. 
%%    %the paper provides error bars and statistical significance information for the experiments, including details on how they were calculated and the assumptions ma
%%    \item[] Guidelines:
%%    \begin{itemize}
%%        \item The answer \answerNA{} means that the paper does not include experiments.
%%        \item The authors should answer \answerYes{} if the results are accompanied by error bars, confidence intervals, or statistical significance tests, at least for the experiments that support the main claims of the paper.
%%        \item The factors of variability that the error bars are capturing should be clearly stated (for example, train/test split, initialization, random drawing of some parameter, or overall run with given experimental conditions).
%%        \item The method for calculating the error bars should be explained (closed form formula, call to a library function, bootstrap, etc.)
%%        \item The assumptions made should be given (e.g., Normally distributed errors).
%%        \item It should be clear whether the error bar is the standard deviation or the standard error of the mean.
%%        \item It is OK to report 1-sigma error bars, but one should state it. The authors should preferably report a 2-sigma error bar than state that they have a 96\% CI, if the hypothesis of Normality of errors is not verified.
%%        \item For asymmetric distributions, the authors should be careful not to show in tables or figures symmetric error bars that would yield results that are out of range (e.g., negative error rates).
%%        \item If error bars are reported in tables or plots, the authors should explain in the text how they were calculated and reference the corresponding figures or tables in the text.
%%    \end{itemize}
%%
%%\item {\bf Experiments compute resources}
%%    \item[] Question: For each experiment, does the paper provide sufficient information on the computer resources (type of compute workers, memory, time of execution) needed to reproduce the experiments?
%%    \item[] Answer: \answerYes{} % Replace by \answerYes{}, \answerNo{}, or \answerNA{}.
%%    \item[] Justification: the experiment was performed on a desktop computer, no special hardware or large-scale resources were required.
%%    \item[] Guidelines:
%%    \begin{itemize}
%%        \item The answer \answerNA{} means that the paper does not include experiments.
%%        \item The paper should indicate the type of compute workers CPU or GPU, internal cluster, or cloud provider, including relevant memory and storage.
%%        \item The paper should provide the amount of compute required for each of the individual experimental runs as well as estimate the total compute. 
%%        \item The paper should disclose whether the full research project required more compute than the experiments reported in the paper (e.g., preliminary or failed experiments that didn't make it into the paper). 
%%    \end{itemize}
%%    
%%\item {\bf Code of ethics}
%%    \item[] Question: Does the research conducted in the paper conform, in every respect, with the NeurIPS Code of Ethics \url{https://neurips.cc/public/EthicsGuidelines}?
%%    \item[] Answer: \answerYes{} % Replace by \answerYes{}, \answerNo{}, or \answerNA{}.
%%    \item[] Justification: the research conducted in the paper adheres to the NeurIPS Code of Ethics.
%%    \item[] Guidelines:
%%    \begin{itemize}
%%        \item The answer \answerNA{} means that the authors have not reviewed the NeurIPS Code of Ethics.
%%        \item If the authors answer \answerNo, they should explain the special circumstances that require a deviation from the Code of Ethics.
%%        \item The authors should make sure to preserve anonymity (e.g., if there is a special consideration due to laws or regulations in their jurisdiction).
%%    \end{itemize}
%%
%%
%%\item {\bf Broader impacts}
%%    \item[] Question: Does the paper discuss both potential positive societal impacts and negative societal impacts of the work performed?
%%    \item[] Answer: \answerNo{} % Replace by \answerYes{}, \answerNo{}, or \answerNA{}.
%%    \item[] Justification: the contribution of the paper is a theoretical result, no societal impacts are directly associated with the work.
%%    \item[] Guidelines:
%%    \begin{itemize}
%%        \item The answer \answerNA{} means that there is no societal impact of the work performed.
%%        \item If the authors answer \answerNA{} or \answerNo, they should explain why their work has no societal impact or why the paper does not address societal impact.
%%        \item Examples of negative societal impacts include potential malicious or unintended uses (e.g., disinformation, generating fake profiles, surveillance), fairness considerations (e.g., deployment of technologies that could make decisions that unfairly impact specific groups), privacy considerations, and security considerations.
%%        \item The conference expects that many papers will be foundational research and not tied to particular applications, let alone deployments. However, if there is a direct path to any negative applications, the authors should point it out. For example, it is legitimate to point out that an improvement in the quality of generative models could be used to generate Deepfakes for disinformation. On the other hand, it is not needed to point out that a generic algorithm for optimizing neural networks could enable people to train models that generate Deepfakes faster.
%%        \item The authors should consider possible harms that could arise when the technology is being used as intended and functioning correctly, harms that could arise when the technology is being used as intended but gives incorrect results, and harms following from (intentional or unintentional) misuse of the technology.
%%        \item If there are negative societal impacts, the authors could also discuss possible mitigation strategies (e.g., gated release of models, providing defenses in addition to attacks, mechanisms for monitoring misuse, mechanisms to monitor how a system learns from feedback over time, improving the efficiency and accessibility of ML).
%%    \end{itemize}
%%    
%%\item {\bf Safeguards}
%%    \item[] Question: Does the paper describe safeguards that have been put in place for responsible release of data or models that have a high risk for misuse (e.g., pre-trained language models, image generators, or scraped datasets)?
%%    \item[] Answer: \answerNA{} % Replace by \answerYes{}, \answerNo{}, or \answerNA{}.
%%    \item[] Justification: the paper poses no such risks.
%%    \item[] Guidelines:
%%    \begin{itemize}
%%        \item The answer \answerNA{} means that the paper poses no such risks.
%%        \item Released models that have a high risk for misuse or dual-use should be released with necessary safeguards to allow for controlled use of the model, for example by requiring that users adhere to usage guidelines or restrictions to access the model or implementing safety filters. 
%%        \item Datasets that have been scraped from the Internet could pose safety risks. The authors should describe how they avoided releasing unsafe images.
%%        \item We recognize that providing effective safeguards is challenging, and many papers do not require this, but we encourage authors to take this into account and make a best faith effort.
%%    \end{itemize}
%%
%%\item {\bf Licenses for existing assets}
%%    \item[] Question: Are the creators or original owners of assets (e.g., code, data, models), used in the paper, properly credited and are the license and terms of use explicitly mentioned and properly respected?
%%    \item[] Answer: \answerYes{} % Replace by \answerYes{}, \answerNo{}, or \answerNA{}.
%%    \item[] Justification: all assets used in the paper are properly credited, and their licenses and terms of use are explicitly mentioned and respected.
%%    \item[] Guidelines:
%%    \begin{itemize}
%%        \item The answer \answerNA{} means that the paper does not use existing assets.
%%        \item The authors should cite the original paper that produced the code package or dataset.
%%        \item The authors should state which version of the asset is used and, if possible, include a URL.
%%        \item The name of the license (e.g., CC-BY 4.0) should be included for each asset.
%%        \item For scraped data from a particular source (e.g., website), the copyright and terms of service of that source should be provided.
%%        \item If assets are released, the license, copyright information, and terms of use in the package should be provided. For popular datasets, \url{paperswithcode.com/datasets} has curated licenses for some datasets. Their licensing guide can help determine the license of a dataset.
%%        \item For existing datasets that are re-packaged, both the original license and the license of the derived asset (if it has changed) should be provided.
%%        \item If this information is not available online, the authors are encouraged to reach out to the asset's creators.
%%    \end{itemize}
%%
%%\item {\bf New assets}
%%    \item[] Question: Are new assets introduced in the paper well documented and is the documentation provided alongside the assets?
%%    \item[] Answer: \answerNA{} % Replace by \answerYes{}, \answerNo{}, or \answerNA{}.
%%    \item[] Justification: the paper does not release new assets.
%%    \item[] Guidelines:
%%    \begin{itemize}
%%        \item The answer \answerNA{} means that the paper does not release new assets.
%%        \item Researchers should communicate the details of the dataset\slash code\slash model as part of their submissions via structured templates. This includes details about training, license, limitations, etc. 
%%        \item The paper should discuss whether and how consent was obtained from people whose asset is used.
%%        \item At submission time, remember to anonymize your assets (if applicable). You can either create an anonymized URL or include an anonymized zip file.
%%    \end{itemize}
%%
%%\item {\bf Crowdsourcing and research with human subjects}
%%    \item[] Question: For crowdsourcing experiments and research with human subjects, does the paper include the full text of instructions given to participants and screenshots, if applicable, as well as details about compensation (if any)? 
%%    \item[] Answer: \answerNA{} % Replace by \answerYes{}, \answerNo{}, or \answerNA{}.
%%    \item[] Justification: the paper does not involve crowdsourcing nor research with human subjects.
%%    \item[] Guidelines:
%%    \begin{itemize}
%%        \item The answer \answerNA{} means that the paper does not involve crowdsourcing nor research with human subjects.
%%        \item Including this information in the supplemental material is fine, but if the main contribution of the paper involves human subjects, then as much detail as possible should be included in the main paper. 
%%        \item According to the NeurIPS Code of Ethics, workers involved in data collection, curation, or other labor should be paid at least the minimum wage in the country of the data collector. 
%%    \end{itemize}
%%
%%\item {\bf Institutional review board (IRB) approvals or equivalent for research with human subjects}
%%    \item[] Question: Does the paper describe potential risks incurred by study participants, whether such risks were disclosed to the subjects, and whether Institutional Review Board (IRB) approvals (or an equivalent approval/review based on the requirements of your country or institution) were obtained?
%%    \item[] Answer: \answerNA{} % Replace by \answerYes{}, \answerNo{}, or \answerNA{}.
%%    \item[] Justification: the paper does not involve crowdsourcing nor research with human subjects.
%%    \item[] Guidelines:
%%    \begin{itemize}
%%        \item The answer \answerNA{} means that the paper does not involve crowdsourcing nor research with human subjects.
%%        \item Depending on the country in which research is conducted, IRB approval (or equivalent) may be required for any human subjects research. If you obtained IRB approval, you should clearly state this in the paper. 
%%        \item We recognize that the procedures for this may vary significantly between institutions and locations, and we expect authors to adhere to the NeurIPS Code of Ethics and the guidelines for their institution. 
%%        \item For initial submissions, do not include any information that would break anonymity (if applicable), such as the institution conducting the review.
%%    \end{itemize}
%%
%%\item {\bf Declaration of LLM usage}
%%    \item[] Question: Does the paper describe the usage of LLMs if it is an important, original, or non-standard component of the core methods in this research? Note that if the LLM is used only for writing, editing, or formatting purposes and does \emph{not} impact the core methodology, scientific rigor, or originality of the research, declaration is not required.
%%    %this research? 
%%    \item[] Answer: \answerNA{} % Replace by \answerYes{}, \answerNo{}, or \answerNA{}.
%%    \item[] Justification: the paper does not involve LLMs as any important, original, or non-standard components.
%%    \item[] Guidelines:
%%    \begin{itemize}
%%        \item The answer \answerNA{} means that the core method development in this research does not involve LLMs as any important, original, or non-standard components.
%%        \item Please refer to our LLM policy in the NeurIPS handbook for what should or should not be described.
%%    \end{itemize}
%%
%%\end{enumerate}

 
\end{document}
